%% ============================================================
%%  PREAMBLE — Production-Grade Agentic AI: LangGraph + Python
%%  Extracted from main.tex.  Input ONCE from Tex/main.tex.
%% ============================================================

\documentclass[
  11pt,
  twoside,
  openright
]{book}

\usepackage{polyglossia}
\setdefaultlanguage{english}

\usepackage{fontspec}
\setmainfont{TeX Gyre Termes}[
  BoldFont   = TeX Gyre Termes Bold,
  ItalicFont = TeX Gyre Termes Italic
]
\setsansfont{TeX Gyre Heros}[
  BoldFont   = TeX Gyre Heros Bold,
  ItalicFont = TeX Gyre Heros Italic
]
\setmonofont{TeX Gyre Cursor}[
  Scale    = 0.88,
  BoldFont = TeX Gyre Cursor Bold
]

\usepackage{geometry}
\geometry{
  papersize  = {7in, 10in},
  top        = 1in,
  bottom     = 1.1in,
  inner      = 1in,
  outer      = 0.875in,
  headheight = 23.07pt,
  twoside
}

%% Fix fancyhdr \headheight warning (must be after geometry)
\setlength{\headheight}{23.07004pt}
\addtolength{\topmargin}{-0.00003pt}

\usepackage{xcolor}
\definecolor{BookBlue}       {HTML}{1A4A8A}
\definecolor{BookDark}       {HTML}{1C1C1E}
\definecolor{BookMid}        {HTML}{4A4A4A}
\definecolor{LearnBg}        {HTML}{EEF7EE}
\definecolor{LearnBorder}    {HTML}{27AE60}
\definecolor{TipBg}          {HTML}{EAF4FB}
\definecolor{TipBorder}      {HTML}{2980B9}
\definecolor{WarningBg}      {HTML}{FEF9E7}
\definecolor{WarningBorder}  {HTML}{F39C12}
\definecolor{NoteBg}         {HTML}{F4F4F4}
\definecolor{NoteBorder}     {HTML}{888888}
\definecolor{TakeawayBg}     {HTML}{F0F0FF}
\definecolor{TakeawayBorder} {HTML}{5B4FBE}
\definecolor{CodeBg}         {HTML}{F7F7F7}
\definecolor{CodeComment}    {HTML}{6A737D}
\definecolor{CodeKeyword}    {HTML}{D73A49}
\definecolor{CodeString}     {HTML}{032F62}
\definecolor{CodeConstant}   {HTML}{005CC5}
\definecolor{CodeAnnotation} {HTML}{6A737D}

\usepackage{microtype}
\usepackage{setspace}
\setstretch{1.18}

\usepackage{csquotes}
\usepackage{parskip}
\setlength{\parskip}{6pt}
\setlength{\parindent}{0pt}

\widowpenalty=10000
\clubpenalty=10000
\displaywidowpenalty=10000
\predisplaypenalty=10000
\raggedbottom
\emergencystretch=3em
\hbadness=3000
\vbadness=3000

\usepackage{needspace}
\usepackage{etoolbox}
\preto\section{\needspace{5\baselineskip}}
\preto\subsection{\needspace{4\baselineskip}}
\preto\subsubsection{\needspace{3\baselineskip}}

\usepackage{titlesec}
\renewcommand{\bottomtitlespace}{0.15\textheight}

\titleformat{\part}[display]
  {\centering\sffamily\Huge\bfseries\color{BookBlue}}
  {\sffamily\large\mdseries\partname\ \thepart}{20pt}{}

\titleformat{\chapter}[display]
  {\sffamily\huge\bfseries\color{BookBlue}}
  {\sffamily\normalsize\mdseries\chaptertitlename\ \thechapter}{12pt}{}
\titlespacing*{\chapter}{0pt}{0pt}{24pt}

\titleformat{\section}
  {\sffamily\Large\bfseries\color{BookBlue}}
  {\thesection}{1em}{}
\titlespacing*{\section}{0pt}{18pt}{6pt}

\titleformat{\subsection}
  {\sffamily\large\bfseries\color{BookDark}}
  {\thesubsection}{1em}{}
\titlespacing*{\subsection}{0pt}{14pt}{4pt}

\titleformat{\subsubsection}
  {\sffamily\normalsize\bfseries\color{BookDark}}
  {\thesubsubsection}{1em}{}
\titlespacing*{\subsubsection}{0pt}{10pt}{2pt}

\usepackage{fancyhdr}
\pagestyle{fancy}
\fancyhf{}
\fancyhead[LE]{\small\sffamily\color{BookMid}\leftmark}
\fancyhead[RO]{\small\sffamily\color{BookMid}\rightmark}
\fancyfoot[LE,RO]{\small\sffamily\color{BookMid}\thepage}
\renewcommand{\headrulewidth}{0.3pt}
\renewcommand{\footrulewidth}{0pt}

\fancypagestyle{plain}{
  \fancyhf{}
  \fancyfoot[LE,RO]{\small\sffamily\color{BookMid}\thepage}
  \renewcommand{\headrulewidth}{0pt}
}

\makeatletter
\renewcommand{\cleardoublepage}{%
  \clearpage
  \if@twoside
    \ifodd\c@page\else
      \thispagestyle{empty}\hbox{}\newpage
      \if@twocolumn\hbox{}\newpage\fi
    \fi
  \fi
}
\makeatother

\usepackage{rotating}
\usepackage{booktabs}
\usepackage{tabularx}
\usepackage{array}
\usepackage{longtable}
\renewcommand{\arraystretch}{1.35}

\newcolumntype{P}[1]{>{\RaggedRight\arraybackslash}p{#1}}

\usepackage{graphicx}
\usepackage{float}
\usepackage{caption}
\captionsetup{
  font      = {small,sf},
  labelfont = {bf,color=BookBlue},
  skip      = 6pt
}

\usepackage{listings}
\lstset{
  backgroundcolor   = \color{CodeBg},
  basicstyle        = \ttfamily\fontsize{9}{11}\selectfont,
  keywordstyle      = \color{CodeKeyword}\bfseries,
  stringstyle       = \color{CodeString}\itshape,
  commentstyle      = \color{CodeComment}\itshape,
  numbers           = left,
  numberstyle       = \tiny\color{BookMid},
  numbersep         = 6pt,
  breaklines        = true,
  frame             = single,
  framerule         = 0.4pt,
  rulecolor         = \color{NoteBorder},
  showstringspaces  = false,
  tabsize           = 2,
  captionpos        = b,
  breakatwhitespace = false,
  columns           = fullflexible,
  linewidth         = \linewidth,
  postbreak         = \mbox{\textcolor{gray}{$\hookrightarrow$}\space},
  xleftmargin       = 1.5em,
  framexleftmargin  = 1.5em,
  aboveskip         = 6pt,
  belowskip         = 4pt,
}

\lstdefinelanguage{Python}{
  morekeywords={
    False,None,True,and,as,assert,async,await,break,class,continue,def,
    del,elif,else,except,finally,for,from,global,if,import,in,is,lambda,
    nonlocal,not,or,pass,raise,return,try,while,with,yield
  },
  morekeywords=[2]{self,cls},
  sensitive=true,
  morecomment=[l]{\#},
  morestring=[b]",
  morestring=[b]',
  morestring=[s]{"""}{"""}, 
  morestring=[s]{'''}{'''},
  keywordstyle=\color{CodeKeyword}\bfseries,
  keywordstyle=[2]\color{CodeConstant},
  stringstyle=\color{CodeString},
  commentstyle=\color{CodeComment}\itshape,
  breaklines=true,
  breakatwhitespace=false,
  columns=fullflexible,
  keepspaces=true,
  showstringspaces=false,
  postbreak=\mbox{\textcolor{gray}{$\hookrightarrow$}\space},
}

\lstdefinelanguage{yaml}{
  morecomment=[l]{\#},
  morestring=[b]",
  morestring=[b]',
  literate=
    *{true}{{\color{CodeKeyword}true}}{4}
    {false}{{\color{CodeKeyword}false}}{5}
    {null}{{\color{CodeKeyword}null}}{4},
  sensitive=true,
  keywordstyle=\color{CodeKeyword}\bfseries,
  stringstyle=\color{CodeString},
  commentstyle=\color{CodeComment}\itshape,
  breaklines=true,
  columns=fullflexible,
  showstringspaces=false
}

\lstdefinelanguage{json}{
  morestring=[b]",
  morecomment=[l]//,
  morecomment=[s]{/*}{*/},
  sensitive=true,
  stringstyle=\color{CodeString},
  commentstyle=\color{CodeComment}\itshape,
}

\DeclareRobustCommand{\code}[1]{\texttt{#1}}
\providecommand{\term}[1]{\textit{#1}}
\providecommand{\chapterquote}[2]{%
  \vspace{1em}
  \begin{quote}\itshape#1\end{quote}
  \begin{flushright}--- #2\end{flushright}
  \vspace{1em}
}

\newcommand{\githubref}[1]{%
  \vspace{2pt}%
  {\small\sffamily\color{BookMid}Full source: \texttt{#1}}%
  \vspace{4pt}%
}

\usepackage{tcolorbox}
\tcbuselibrary{skins,breakable}

\newtcolorbox{learningbox}{
  colback=LearnBg, colframe=LearnBorder,
  before skip=12pt, fonttitle=\sffamily\bfseries\small,
  title={Learning Objective}, breakable, lines before break=3,
  left=8pt, right=8pt, top=6pt, bottom=6pt,
  boxrule=0.6pt, width=\linewidth
}
\newtcolorbox{tipbox}{
  colback=TipBg, colframe=TipBorder,
  before skip=12pt, fonttitle=\sffamily\bfseries\small,
  title={Tip}, breakable, lines before break=3,
  left=8pt, right=8pt, top=6pt, bottom=6pt,
  boxrule=0.6pt, width=\linewidth
}
\newtcolorbox{warningbox}{
  colback=WarningBg, colframe=WarningBorder,
  before skip=12pt, fonttitle=\sffamily\bfseries\small,
  title={Warning}, breakable, lines before break=3,
  left=8pt, right=8pt, top=6pt, bottom=6pt,
  boxrule=0.6pt, width=\linewidth
}
\newtcolorbox{notebox}{
  colback=NoteBg, colframe=NoteBorder,
  before skip=12pt, fonttitle=\sffamily\bfseries\small,
  title={Note}, breakable, lines before break=3,
  left=8pt, right=8pt, top=6pt, bottom=6pt,
  boxrule=0.6pt, width=\linewidth
}
\newtcolorbox{takeawaybox}[1][Key Takeaway]{
  colback=TakeawayBg, colframe=TakeawayBorder,
  before skip=12pt, fonttitle=\sffamily\bfseries\small,
  title={#1}, breakable, lines before break=3,
  left=8pt, right=8pt, top=6pt, bottom=6pt,
  boxrule=0.6pt, width=\linewidth
}

%% ── Source-Layout tree environment ─────────────────────────────────────
%% Usage:
%%   \begin{sourcelayout}
%%     \dirtree{%
%%       .1 src/.
%%       .2 agent/.
%%       .3 nodes/.
%%       .4 perceive\_node.py \DTcomment{§1.2}.
%%     }
%%   \end{sourcelayout}
%%
%% Requires: \usepackage{dirtree}  (loaded below)
%% Colours reuse NoteBg / NoteBorder (neutral grey — reference material,
%% not instructional content).
%% ────────────────────────────────────────────────────────────────────────
\usepackage{dirtree}

\newtcolorbox{sourcelayout}{%
  colback      = NoteBg,
  colframe     = NoteBorder,
  colbacktitle = NoteBorder,
  coltitle     = white,
  fonttitle    = \sffamily\bfseries\small,
  title        = {Source Layout},
  before skip  = 12pt,
  after skip   = 10pt,
  breakable,
  lines before break = 3,
  left   = 6pt,
  right  = 6pt,
  top    = 6pt,
  bottom = 6pt,
  boxrule = 0.6pt,
  width   = \linewidth,
}

%% dirtree global style: monospaced font, consistent with code listings
\renewcommand*\DTstylecomment{\rmfamily\color{CodeComment}\itshape}
\renewcommand*\DTstyle{\ttfamily\color{BookDark}}

\usepackage{amsmath}
\usepackage{amssymb}

\PassOptionsToPackage{hyphens}{url}
\usepackage{hyperref}
\usepackage{imakeidx}
\makeindex[columns=2, title=Index, intoc=true]
\urlstyle{tt}
\makeatletter
\g@addto@macro\UrlBreaks{\do\.\do\_\do\-}
\makeatother

\hypersetup{
  colorlinks  = true,
  linkcolor   = BookBlue,
  urlcolor    = BookBlue,
  citecolor   = BookBlue,
  pdftitle    = {Production-Grade Agentic AI: LangGraph + Python Frameworks},
  pdfauthor   = {Haythem Aber},
  pdfsubject  = {LangGraph, Python, Agentic AI, Production Systems},
  pdfkeywords = {LangGraph, Python, AI agents, production, observability, security}
}

\usepackage{tikz}
\usetikzlibrary{arrows.meta, positioning, shapes.geometric, shapes.misc, calc}
\usepackage{pdflscape}
\usepackage{ragged2e}

\newcommand{\component}[1]{\textsf{\textbf{#1}}}
\newcommand{\seechapter}[1]{\textit{(See Chapter~\ref{#1})}}
\newcommand{\seepart}[1]{\textit{(See Part~\ref{#1})}}

\newcommand{\partclose}[1]{%
  \cleardoublepage
  \vspace*{\fill}
  \begin{center}\sffamily\Large\color{BookBlue}#1\end{center}
  \vspace*{\fill}
  \cleardoublepage
}

%% Shared book-wide infrastructure layered on top of the author's preamble.
%% Keep this file loaded once by Tex/main.tex; chapter fragments must not load it.

\usepackage{calc} % Pandoc-generated longtable column dimensions.
\usepackage[
  backend=biber,
  style=verbose-note,
  autocite=footnote,
  sorting=nyt,
  maxbibnames=99
]{biblatex}
\addbibresource{Tex/bibliography/references.bib}

% Paper labels remain explicit and are not interchangeable.
\newcommand{\CoALACite}[1]{\footnote{\textsc{[CoALA]} See Sumers et al., \emph{Cognitive Architectures for Language Agents}, #1.}}
\newcommand{\FanLanCite}[1]{\footnote{\textsc{[Fan \& Lan]} See Fan and Lan, \emph{From Cognitive Architectures to Language Agents}, #1.}}

% The source paper is the theoretical basis in Parts I and II; examples and
% engineering decisions are explicitly identified in the prose.
\newcommand{\BookExampleLabel}{\textsc{[Book example]}}
\newcommand{\EngineeringLabel}{\textsc{[Engineering decision]}}

\usepackage[toc,nonumberlist]{glossaries}
\makeglossaries
\renewcommand*{\glossarypreamble}{\begin{multicols}{2}\raggedcolumns\small}
\renewcommand*{\glossarypostamble}{\end{multicols}}
\newglossaryentry{languageagent}{name={language agent},description={An agentic system that uses a language model as a core computational component and interacts with an environment}}
\newglossaryentry{cognitivelanguageagent}{name={cognitive language agent},description={A language agent that also uses internal processes to manage its information and behavior}}
\newglossaryentry{workingmemory}{name={working memory},description={Active information maintained for the current decision cycle}}
\newglossaryentry{longtermmemory}{name={long-term memory},description={Memory modules that retain information beyond the current decision cycle}}
\newglossaryentry{episodicmemory}{name={episodic memory},description={Memory of experience from earlier decision cycles}}
\newglossaryentry{semanticmemory}{name={semantic memory},description={Memory of knowledge about the world and the agent itself}}
\newglossaryentry{proceduralmemory}{name={procedural memory},description={Knowledge about how to perform actions and decision procedures, represented implicitly in model weights or explicitly in agent code}}
\newglossaryentry{grounding}{name={grounding},description={Procedures that execute external actions and process environmental feedback}}
\newglossaryentry{decisionprocedure}{name={decision procedure},description={The top-level procedure that coordinates the agent's repeated decision cycles}}


% Reusable review and recap blocks used throughout the book.
\newtcolorbox{chapterrecap}{
  colback=TakeawayBg, colframe=TakeawayBorder,
  before skip=10.5pt, fonttitle=\sffamily\bfseries\small,
  title={Chapter recap}, breakable, lines before break=3,
  left=8pt, right=8pt, top=5.25pt, bottom=5.25pt,
  boxrule=0.6pt, width=\linewidth
}
\newtcolorbox{reviewquestions}[1]{
  colback=NoteBg, colframe=NoteBorder,
  before skip=10.5pt, fonttitle=\sffamily\bfseries\small,
  title={#1}, breakable, lines before break=3,
  left=8pt, right=8pt, top=5.25pt, bottom=5.25pt,
  boxrule=0.6pt, width=\linewidth
}
\providecommand{\tightlist}{}

% Stable source references for implementation listings. The QR payload, the
% clickable source ID and "Open exact source" all consume the same URL argument.
% The visible source path/member and linked test path remain useful in print.
\newcommand{\BookSourceReferenceText}[9]{%
  {\sffamily\fontsize{8}{9}\selectfont\raggedright
    \setlength{\parskip}{0pt}%
    \href{#8}{\textbf{Source #1}}\hspace{0.55em}\textbar\hspace{0.55em}\texttt{#2}\par
    \textbf{Path:} \href{#8}{\path{#3}}\par
    \textbf{Member:} #4\hspace{0.5em}\textbf{Source lines:} #5\par
    \textbf{Test file:} \href{#9}{\path{#6}}\par
    \textbf{Tests:}\par \href{#9}{#7}\par
    \href{#8}{\textbf{Open exact source}}%
  }%
}

\newcommand{\BookSourceReferenceNoQR}[9]{%
  \par\smallskip\noindent
  \begin{minipage}{\linewidth}%
    \BookSourceReferenceText{#1}{#2}{#3}{#4}{#5}{#6}{#7}{#8}{#9}%
  \end{minipage}%
  \par\smallskip
}

\newcommand{\BookSourceReferenceQRContent}[9]{%
  \par\smallskip\noindent
  \begin{minipage}[c]{0.785\linewidth}%
    \BookSourceReferenceText{#1}{#2}{#3}{#4}{#5}{#6}{#7}{#8}{#9}%
  \end{minipage}%
  \hfill
  \begin{minipage}[c]{0.18\linewidth}%
    \centering
    % The source QR is a vector PDF asset generated with qrcode.sty.
    % Embedding it as a graphic lets hyperref create a real clickable area
    % around the full printed symbol (including its four-module quiet zone).
    \href{#8}{\mbox{{\setlength{\fboxsep}{0pt}\colorbox{white}{\includegraphics[height=\BookSourceQRHeight]{Tex/diagrams/source-qr/#1.pdf}}}}}\par
    {\scriptsize\sffamily Scan to inspect}%
  \end{minipage}%
  \par\smallskip
}

% Standard source references use a 19.5 mm QR including the quiet zone.
% Very dense commit-pinned URLs may opt into the 21 mm variant.
\newcommand{\BookSourceReferenceQR}[9]{%
  {\def\BookSourceQRHeight{19.5mm}\BookSourceReferenceQRContent{#1}{#2}{#3}{#4}{#5}{#6}{#7}{#8}{#9}}%
}
\newcommand{\BookSourceReferenceQRLong}[9]{%
  {\def\BookSourceQRHeight{21mm}\BookSourceReferenceQRContent{#1}{#2}{#3}{#4}{#5}{#6}{#7}{#8}{#9}}%
}


% Book-level metadata overrides the sample metadata retained in the supplied preamble.
\hypersetup{
  pdftitle={Reliable AI Agents: Architecture and Implementation},
  pdfauthor={Haythem Aber},
  pdfsubject={CoALA, cognitive architectures, agentic systems, Java production engineering},
  pdfkeywords={CoALA, Fan and Lan, cognitive architectures, agentic AI, Java, supply-chain disruption management, production systems}
}


%% Keep deep navigation in the PDF outline even though the printed contents
%% intentionally stops at first-level sections.
\hypersetup{bookmarksdepth=subsubsection}

%% Preserve the existing description-style glossary entries, but flow them
%% in two columns to reduce back-matter length without deleting definitions.
\newglossarystyle{twocolumnlist}{%
  \setglossarystyle{list}%
  \renewenvironment{theglossary}%
    {\begin{multicols}{2}\begin{description}}%
    {\end{description}\end{multicols}}%
}
\setglossarystyle{twocolumnlist}

%% Two-column printed table of contents. The standard chapter heading and
%% running marks are retained; only the layout and displayed depth change.
\makeatletter
\renewcommand{\tableofcontents}{%
  \chapter*{\contentsname
    \@mkboth{\MakeUppercase\contentsname}{\MakeUppercase\contentsname}}%
  \begin{multicols}{2}
    \setlength{\columnsep}{24pt}%
    \@starttoc{toc}%
  \end{multicols}%
}
\makeatother

\begin{document}
\frontmatter
\begin{titlepage}
\thispagestyle{empty}
\vspace*{1.35in}
\begin{center}
{\sffamily\Huge\bfseries\color{BookBlue} Reliable AI Agents: Architecture and Implementation}\par
\vspace{0.35in}
{\sffamily\Large Memory, Planning, and Recovery for Production Systems\\[0.08in]
Using Java, LangChain4j, and LangGraph4j}\par
\vspace{0.55in}
{\sffamily\large\bfseries Production AI Engineering}\par
\vspace{0.65in}
{\Large Haythem Aber}\par
\vfill
{\sffamily\large 2026}\par
\end{center}
\end{titlepage}

\thispagestyle{empty}
\vspace*{\fill}
\noindent Copyright \textcopyright{} 2026 Haythem Aber.\\
\noindent Book text: Creative Commons Attribution-NonCommercial-ShareAlike 4.0 International (CC BY-NC-SA 4.0).\\
\noindent Code examples and accompanying source code: Apache License 2.0.\\[0.5em]
\noindent The book distinguishes claims grounded in the cited research papers from explanatory examples and engineering decisions created for this work. Product X and its incident scenario are illustrative examples.\\[0.5em]
\noindent First edition, 2026.
\vspace*{\fill}
\clearpage

\chapter*{Dedication}
\addcontentsline{toc}{chapter}{Dedication}
\begin{center}\itshape Dedication to be supplied.\end{center}

\chapter*{Preface}
\addcontentsline{toc}{chapter}{Preface}
This book develops a source-disciplined path from cognitive architectures for language agents to the engineering of complex agentic systems in Java. The complete preface will be written after the chapter sequence is established.

\chapter*{How to Read This Book}
\addcontentsline{toc}{chapter}{How to Read This Book}
\begin{itemize}
  \item Part I develops the CoALA framework from its conceptual foundations to its decision process.
  \item Part II studies Fan and Lan's mechanism-level analysis as a distinct theoretical contribution.
  \item Part III applies the ideas in ten production-oriented Java case studies. Engineering choices are labelled as such and are not attributed to the papers.
\end{itemize}
The intended reader has basic programming knowledge. Each technical idea is introduced through a problem, followed by intuition, terminology, and a precise account.


\tableofcontents

\mainmatter
\part{CoALA: Understanding Language Agents as Cognitive Systems}
\label{part:coala}
\chapter[Why Do Language Agents Need Cognitive Architectures?]{Why Do Language Agents Need\\Cognitive Architectures?}
\label{ch:why-language-agents-need-cognitive-architectures}
\section*{Chapter at a glance}
\phantomsection
\addcontentsline{toc}{section}{Chapter at a glance}\label{chapter-at-a-glance}

A language model can produce an answer. A language agent can use a language model as part of a system that interacts with an environment. A cognitive language agent can additionally use the model in processes that manage the agent's own information and behavior. These are related, but they are not interchangeable descriptions.

This chapter explains why that distinction matters, where the idea of a cognitive architecture comes from, and why the CoALA paper uses that history to organize language agents. We will work toward the framework through a concrete problem rather than beginning with a list of definitions.

\begin{learningbox}
By the end of the chapter, you should be able to explain why an LLM alone is not a complete description of an agent system, why memory and control become important as agents grow more capable, and what CoALA contributes. You will also know what this chapter deliberately leaves for the chapters that follow.
\end{learningbox}

\section{A model call is not yet the whole system}\label{a-model-call-is-not-yet-the-whole-system}

Imagine that a payment service in a software company has started returning errors. An engineer asks a language model:

\begin{quote}
The payment service is returning HTTP 500 errors. What could be wrong?
\end{quote}

The model might return a useful list of possibilities: a recent deployment, a database connection problem, an unavailable dependency, or an unexpected input. That answer can help the engineer think. But notice what has happened---and what has not happened.

The model has received text and generated text. It has not inspected the service's logs. It has not checked whether a deployment occurred. It has not queried the database. It has not changed a configuration or restarted a process. Unless the relevant observations were included in the input, the model has no direct access to those facts in this interaction.

This is not a defect in the answer; it is a boundary of the interaction we have described. A \textbf{model call} is one invocation of a model with a particular input, producing an output. For a language model, the input is commonly a prompt represented as tokens and the output is a generated continuation. The surrounding software determines what information is supplied, what happens to the output, and whether another operation follows. A single call may answer, summarize, or draft effectively, but it does not by itself describe a system that repeatedly observes and acts in the world.

Now change the task. Instead of asking only for possible causes, the engineer asks a software system to investigate the incident. The system is expected to gather relevant evidence, decide which question to investigate next, perform permitted checks, interpret the results, and report what it recommends. It may need to repeat this process because the first observation does not settle the matter.

The moment we describe that larger task, the model call is no longer the whole story. We must account for the rest of the system.

\subsection{The environment, observations, and actions}\label{the-environment-observations-and-actions}

The \textbf{environment} is the part of the world with which a system interacts. In the incident example, it includes the running payment service and the surrounding operational systems from which information can be obtained or to which an operation can be directed. In a robot, the environment might include the physical room, objects, people, sensors, and actuators. In a web agent, it may include web pages and the site with which the agent interacts.

An \textbf{observation} is information the system receives from its environment. A service-health response, a log entry, a deployment record, or a human message could all be observations in our incident scenario. An observation does not automatically explain itself. It is information that becomes available for the system to use.

An \textbf{action} is an operation the system takes. An action might request a health check, retrieve a record, send a message, or issue a command to an external system. The action is not the same thing as text that merely describes the action. If a model generates the words \code{check the database}, nothing has been checked yet. Some component of the surrounding system must interpret the output and actually perform the operation.

That distinction matters because a language model can generate text that resembles an instruction without the instruction being valid, authorized, or executed. The model's output and the system's resulting behavior are separate parts of the overall process.

For now, keep three questions in mind:

\begin{enumerate}
\def\labelenumi{\arabic{enumi}.}
\tightlist
\item
  What information does the system receive and retain?
\item
  What operations can it perform?
\item
  How does it decide which operation should happen next?
\end{enumerate}

These questions will lead us to CoALA. We will not answer them fully yet. First, we need to understand how language-model systems expanded beyond isolated input-output interactions.

\section{How language-model systems move beyond a single response}\label{how-language-model-systems-move-beyond-a-single-response}

The CoALA paper distinguishes several ways of using large language models. The differences are easier to understand as a progression in the organization of the system, not as a ranking in which every later design is automatically better for every task \CoALACite{Fig. 1 and §§3.2--3.3}.

\subsection{From text input to text output}\label{from-text-input-to-text-output}

The simplest pattern is familiar: provide text to a language model and receive text in return.

\BookDiagram[0.42\textheight]
  {Tex/diagrams/generated/part1/chapter01/model-call.png}
  {A simple model call: supplied text is transformed into generated text. The surrounding application is not shown.}
  {fig:chapter01-model-call}

The application may display, store, or pass on the generated output, but those surrounding operations are outside the model call shown here. This pattern is sufficient for many tasks; it becomes an incomplete system description only when the intended behavior requires repeated interaction with an environment.

\subsection{Adding a sequence of model calls}\label{adding-a-sequence-of-model-calls}

A second step is to organize more than one model call into a sequence. The output of one call can become part of the input to another. For example, a first call might draft an answer, while a later call reviews that draft or uses it to produce a more focused result.

A \textbf{prompt} is the input provided to the model. It can contain a question, instructions, examples, observations, or text prepared by the surrounding program. \textbf{Prompt construction} is the work of assembling that input before the model is called.

In a few-shot prompt, examples are placed before the current question so that the model has examples of the desired pattern in its input. In retrieval-augmented generation, information retrieved from an external source is added to the prompt before the model generates an answer. In a chained interaction, one model call can produce an intermediate result that is supplied to a later call.

These techniques matter because they add structure around model generation. The system can prepare information, call the model, process the result, and call it again. The CoALA paper relates these forms of prompt construction and chaining to sequences of productions; we will develop that connection in Section 1.5.

However, a predefined sequence of calls is not the same as an agent that decides what to do based on new observations from an environment. A fixed sequence may always perform step A, then step B, then step C. An interactive agent can receive a result from the environment and use it to determine what to do next. The difference is whether the next step is simply prescribed in advance or is selected in response to the evolving situation.

\subsection{A language agent interacts with an environment}\label{a-language-agent-interacts-with-an-environment}

A \textbf{language agent}, in the sense used by CoALA, is a system that uses a large language model as a core computation unit to reason, plan, and act \CoALACite{§2.4}. The key distinction in the paper's Figure 1 is that the model is placed in a feedback loop with an external environment \CoALACite{Fig. 1B and §3.3}.

The system makes an observation available to the model, uses the model's output to determine an action, and invokes the relevant procedure. The environment may then provide further feedback.

\BookDiagram[0.62\textheight]
  {Tex/diagrams/generated/part1/chapter01/language-agent-loop.png}
  {A language agent places model use inside an interaction loop: software prepares model input, interprets the result, executes an action through an interface, and receives a subsequent observation.}
  {fig:chapter01-language-agent-loop}

The arrows form a loop because an action can change what the system observes next. The next observation may confirm an expectation, contradict it, or reveal that another investigation is necessary. The agent is therefore not merely producing an answer to the original prompt. It is participating in an ongoing interaction.

A feedback loop does not guarantee good decisions. It only describes the relationship between action and subsequent observation. The agent may still choose an unsuitable action, misinterpret evidence, or fail to account for relevant information. The architecture of the complete system matters precisely because the interaction has become more complex than one model call.

\begin{tipbox}
Keep the boundary between model output and system behavior explicit. A generated action description is still text until the surrounding agent software validates, authorizes, and executes it through an interface. The feedback loop begins with an actual interaction and the observation that follows it.
\end{tipbox}

\subsection{A cognitive language agent also manages internal information}\label{a-cognitive-language-agent-also-manages-internal-information}

The CoALA paper then describes a further development: \textbf{cognitive language agents}. These agents do not only use the model to select actions in the external environment. They also use it in processes that manage the agent's internal state, including reasoning and learning \CoALACite{Fig. 1C and §3.3}.

Consider the incident again. The system may need to compare the current symptoms with information retained from earlier incidents. It may reason over the evidence gathered so far. It may update information that could help it in a later interaction. Those operations are not identical to asking an external service for its current health or issuing a restart command. They concern the information maintained and processed within the agent system.

The distinction can be summarized like this:

\begin{longtable}[]{@{}
  >{\raggedright\arraybackslash}p{(\columnwidth - 2\tabcolsep) * \real{0.5000}}
  >{\raggedright\arraybackslash}p{(\columnwidth - 2\tabcolsep) * \real{0.5000}}@{}}
\toprule\noalign{}
\begin{minipage}[b]{\linewidth}\raggedright
Pattern
\end{minipage} & \begin{minipage}[b]{\linewidth}\raggedright
What the model is doing as part of the system
\end{minipage} \\
\midrule\noalign{}
\endhead
\bottomrule\noalign{}
\endlastfoot
Language model used for text input/output & Producing a response from supplied text \\
Language agent & Helping select actions in a feedback loop with an external environment \\
Cognitive language agent & Also participating in internal processes such as reasoning and learning that manage the agent's internal state \\
\end{longtable}

This table is a high-level reading aid, not a claim that real systems always fit into three perfectly separate boxes. A system may combine several patterns. The point is to notice what has been added: not just more calls, but more kinds of information processing and control.

The word \textbf{cognitive} can sound as though the system is being claimed to have human consciousness or subjective experience. That is not what the term means here. In this context, it refers to organizing processes associated with cognition---such as memory, reasoning, learning, and decision-making---within an artificial system. It is an architectural description, not a claim about consciousness.

The paper's Figure 1 makes this progression visible. Its first pattern shows an LLM taking text input and producing text output. Its second places the LLM in a direct feedback loop with an external environment. Its third adds processes through which the LLM helps manage the agent's internal state, including reasoning and learning. The progression gives us a useful starting point: as the system takes on more responsibilities, describing it only by saying that it ``uses an LLM'' tells us less and less about how it works.

\section{Why increasingly capable agents are difficult to describe consistently}\label{why-increasingly-capable-agents-are-difficult-to-describe-consistently}

Imagine reading documentation for two systems. Both are described as having memory, using tools, and performing reasoning. At first glance, they appear similar.

But in the first system, ``memory'' may mean a short history appended to each prompt. In the second, it may mean several organized stores that the agent can read from and write to over time. One system may call a tool whenever the model emits a particular command. Another may have a separate process that selects and executes actions according to its current state. The same feature labels do not tell us whether the underlying organization is the same.

This is the problem the CoALA paper identifies. Language-agent research had begun to use many different terms for related processes: tool use, grounding, actions, reasoning, and memory. Individual works often introduced their own terminology. That variety makes it harder to compare systems, understand how the field is evolving, or design new agents with clean and consistent abstractions \CoALACite{§1}.

A name such as ``memory'' is a useful beginning, but it is not a complete explanation. We need to ask what information is stored, how it is organized, when it is accessed, and how it affects what the agent does. Likewise, saying that a system uses ``tools'' does not tell us whether an operation reads from an external source, changes the outside world, or updates the system's own information.

This does not mean that every project must use the same internal implementation. Two systems can be designed differently and still be compared meaningfully. A conceptual framework provides common questions and distinctions; it does not require every system to be built from identical components.

CoALA's response is to organize language agents around three broad dimensions: information storage, action space, and decision-making procedure. Before we examine those dimensions, it helps to understand why the paper turns to the history of production systems and cognitive architectures.

\section{The historical bridge: production systems and cognitive architectures}\label{the-historical-bridge-production-systems-and-cognitive-architectures}

The CoALA paper does not present its framework as an arbitrary collection of software modules. It draws on ideas from the history of computing and artificial intelligence. The bridge begins with production systems.

\subsection{A production rule: a condition followed by an action}\label{a-production-rule-a-condition-followed-by-an-action}

A \textbf{production system} is built around rules that specify a precondition and an associated action. A precondition is a condition that must be satisfied before the action can be applied. When the condition is satisfied, the rule permits the associated action \CoALACite{§2.1}.

Consider a thermostat. It receives information about the temperature and whether the furnace is currently running. A simplified rule might say:

\begin{lstlisting}
IF the temperature is below the lower threshold
AND the furnace is off
THEN turn the furnace on.
\end{lstlisting}

Another rule might say:

\begin{lstlisting}
IF the temperature is above the upper threshold
AND the furnace is on
THEN turn the furnace off.
\end{lstlisting}

These rules connect a description of the current situation to a possible operation. The example is intentionally small, but it already shows two important ingredients: the system has information about the world, and the rules specify what may be done when that information meets certain conditions.

The word \textbf{production} comes from the idea that applying a rule can produce a new result. In early formal accounts discussed by CoALA, production systems were used to describe transformations of symbolic strings \CoALACite{§2.1}. A simple string-rewriting rule might replace one portion of a string with another:

\begin{lstlisting}
XYZ  ->  XWZ
\end{lstlisting}

Read this as: when the relevant pattern \code{XYZ} is present, it may be rewritten as \code{XWZ}. The rule describes a permitted transformation. It does not, by itself, tell us which rule should be selected when several rules could apply, or when the whole process should stop.

That leads to the next idea.

\subsection{Rules alone do not define the whole algorithm}\label{rules-alone-do-not-define-the-whole-algorithm}

Suppose a system has ten rules, and several of them can apply to the current situation. Which one should run first? Does the system always use a fixed priority order? Does it compare the possible results? Can it create a new subtask when it cannot determine what to do?

A set of production rules describes possible transformations. To obtain a particular algorithm or behavior, a system also needs \textbf{control flow}: the organization that determines how the rules are selected and applied over time \CoALACite{§2.2}.

A simple example is a priority order. If rules are examined in a fixed order, the ordering itself influences the result. A more sophisticated system may compare candidate operations, select one, apply it, update its state, and then evaluate what should happen next.

The lesson is broader than the thermostat. Knowing which operations are available is not enough to explain a system's behavior. We must also understand how the system chooses among them and how the process continues after an operation changes the situation.

\subsection{From production systems to cognitive architectures}\label{from-production-systems-to-cognitive-architectures}

Production systems became important in artificial intelligence as researchers used them to represent more complex problem-solving behavior. Their rules were no longer limited to rewriting strings. Preconditions could be checked against information about an agent's goals and its understanding of the world. Associated actions could modify internal information or interact with the environment.

As these systems grew, they needed more elaborate control. They could be connected to sensors, actuators, and knowledge bases. A system that receives observations, maintains information, and performs operations needs a way to organize all these processes.

A \textbf{cognitive architecture} is an organized account of the components and processes through which an artificial agent operates. In the historical line described by CoALA, such architectures explicitly represented processes associated with perception, memory, planning, and decision-making. They provided a way to reason about how these processes fit together and how control moves through the system \CoALACite{§2.3}.

The word ``architecture'' is important. It shifts attention from a single rule or operation to the organization of the complete system. It asks how information is represented, what processes can use it, how those processes are coordinated, and how behavior unfolds over time.

Again, ``cognitive'' does not mean conscious. The architecture is called cognitive because it organizes functions that are relevant to cognition. The term does not establish that the artificial system has human experiences, intentions, or understanding in the philosophical sense.

\subsection{A historical illustration: the Soar architecture}\label{a-historical-illustration-the-soar-architecture}

The CoALA paper uses Soar as a canonical example of a cognitive architecture built around production rules \CoALACite{§2.3}. We do not need to study Soar in detail here. It is useful because it shows how the pieces fit together.

In the account presented by CoALA, Soar has working memory for the agent's current circumstances and long-term memory divided into procedural, semantic, and episodic forms. Productions are stored in procedural memory. Their conditions are checked against working memory. The decision procedure uses productions to propose and evaluate possible actions, selects an action, and then applies it. If the system cannot choose because no action is valid or several actions are tied, Soar can create a subgoal to resolve the impasse. Learning can also update what is retained in long-term memory or modify behavior.

Do not mistake this short description for the complete CoALA framework. Soar is a historical example discussed in the paper, not the software framework that we will implement in this book. Its role here is to make a point: production rules become part of a much larger agent architecture when they are combined with memory, interaction with the environment, and a decision procedure.

\subsection{What this history contributes}\label{what-this-history-contributes}

The historical bridge gives us a way to understand why CoALA is concerned with more than the model's text output. Earlier cognitive architectures had to organize memory, possible actions, and control flow around production systems. The CoALA authors suggest that similar architectural principles can help organize systems built around language models.

That is the motivation. The next question is why language models make this connection interesting rather than simply requiring us to copy an older architecture unchanged.

\section{Why the production-system analogy is useful for language models}\label{why-the-production-system-analogy-is-useful-for-language-models}

A language model receives text as context and produces a continuation. There may be many possible continuations. The model assigns different probabilities to them, and generation selects a continuation according to the model's distribution and the decoding process.

The CoALA paper draws an analogy between this process and a production system \CoALACite{§3.1}. In a traditional production system, a rule specifies a possible way to transform a symbolic string. In a language model, the supplied text provides a context for possible continuations or modifications. From this perspective, the model defines a probability distribution over possible text-producing operations.

The analogy is useful, but it has limits.

\subsection{A probabilistic production system---not a literal rule engine}\label{a-probabilistic-production-systemnot-a-literal-rule-engine}

A traditional production rule can be written down as an explicit condition and action. A developer or researcher can inspect the rule directly. A modern language model, by contrast, has learned its behavior through pretraining and stores that behavior in a large set of parameters. The model's internal parameters do not expose a simple list of human-readable productions.

This creates an important contrast:

\begin{longtable}[]{@{}
  >{\raggedright\arraybackslash}p{(\columnwidth - 2\tabcolsep) * \real{0.5000}}
  >{\raggedright\arraybackslash}p{(\columnwidth - 2\tabcolsep) * \real{0.5000}}@{}}
\toprule\noalign{}
\begin{minipage}[b]{\linewidth}\raggedright
Traditional production system
\end{minipage} & \begin{minipage}[b]{\linewidth}\raggedright
Language model, viewed through the paper's analogy
\end{minipage} \\
\midrule\noalign{}
\endhead
\bottomrule\noalign{}
\endlastfoot
Rules are explicit symbolic operations. & Possible continuations are represented probabilistically. \\
A rule's precondition and action can be inspected directly. & The model's learned parameters are not a human-readable rule list. \\
Control may be defined by explicit rule-selection procedures. & The prompt and generation process influence which continuation is produced. \\
Behavior can be constrained by the specified rule set. & Pretraining gives broad flexibility, but the resulting behavior is less transparent and can be stochastic. \\
\end{longtable}

The analogy does not say that a language model literally contains a traditional production system. It gives us a way to compare two systems that transform or continue symbolic text, while keeping their differences visible.

The paper highlights both sides of the comparison. Large-scale pretraining gives language models a broad prior over text continuations, allowing them to handle many tasks without a developer writing an individual rule for every situation. But the model's learned parameters are opaque compared with explicit symbolic rules, and its probabilistic behavior makes analysis and control more challenging.

This tension is central to the motivation for architecture. The model brings flexible language-based computation. The surrounding agent system still needs an organized way to provide information, decide which operations to perform, and respond to what happens.

\subsection{Prompts can organize control}\label{prompts-can-organize-control}

A prompt is not merely a question. It can be constructed to influence what the model does next. The CoALA paper shows how several prompting methods can be understood as sequences of productions or as forms of control flow \CoALACite{§3.2}.

In a zero-shot interaction, the system supplies the current question and asks for an answer. In a few-shot interaction, examples are added before the current question. Those examples change the context from which the model generates its answer. In retrieval-augmented generation, the system retrieves external information and places it into the prompt so that the model can use it while responding.

A chained approach may also make more than one model call. For instance, the system can ask for a draft, then ask for a critique of that draft, then make another call that uses the critique. The calls form a sequence in which an intermediate result becomes input to the next step.

These are examples of how prompt construction can influence the process around an LLM. They should not be confused with a full account of an agent's decision-making procedure. Some sequences are predefined. An interactive agent, in contrast, can use new observations from the environment to determine what to do next.

The practical distinction is this:

\begin{itemize}
\tightlist
\item
  \textbf{Prompt construction} determines what information and instructions are presented to a model call.
\item
  \textbf{Prompt chaining} organizes multiple calls into a sequence.
\item
  \textbf{An agent feedback loop} connects model-supported decisions to actions and new observations from an environment.
\item
  \textbf{A cognitive language agent} may additionally use internal processes to manage information that influences future decisions.
\end{itemize}

These patterns can coexist. They describe different aspects of how a larger system is organized.

\subsection{Why bring cognitive architectures into the language-model era?}\label{why-bring-cognitive-architectures-into-the-language-model-era}

The CoALA authors argue that the history of cognitive architectures offers useful principles for designing LLM-based agents \CoALACite{§§2.3--3.3}. Earlier architectures had to organize memory, interactions with the world, learning, and decision-making around production systems. Language models change what can serve as the core computation: instead of requiring every behavior to be encoded as a hand-written rule, a model can generate flexible text and help perform reasoning tasks.

But the need for organization does not disappear. If anything, greater flexibility makes it more important to understand how the system's parts work together. A language model can propose text, but the architecture must account for what information is available to it, what operations the system can perform, and how a generated proposal becomes part of the system's behavior.

This is the bridge to CoALA. Before we introduce the framework formally, let's revisit the production incident and use it to make those questions concrete.

\section{A running example: investigating a production incident}\label{a-running-example-investigating-a-production-incident}

The following scenario is an illustrative example created for this book. It is not a case study from the CoALA paper.

A company operates Product X, whose payment service has begun returning HTTP 500 errors. The incident report says that some payment requests are failing. The agent is asked to investigate and help determine an appropriate response.

At first, the system has very little information. It knows that an incident has been reported, but it does not yet know whether the cause is a recent deployment, a database problem, a dependency failure, or something else. It should not treat any of these possibilities as established fact.

\subsection{Question 1: What information is available?}\label{question-1-what-information-is-available}

The initial report is one piece of information. More may become available from operational records, monitoring systems, prior incident records, service documentation, or people involved in the response.

The system must distinguish what it currently knows from what it could retrieve. It must also keep track of information produced during the investigation. For example, a check may show that the payment service is reachable but its database connection pool is exhausted. That observation changes the information available for the next decision.

We have not yet introduced CoALA's memory categories, and we should not pretend that this brief example defines them. It is enough to recognize the architectural question: where does the system keep the information that matters to the current investigation, and how can it obtain other information when needed?

\subsection{Question 2: What can the system do?}\label{question-2-what-can-the-system-do}

The system may be able to retrieve deployment history, inspect recent logs, request a service-health check, or ask an operator for clarification. These are possible operations. They are not all the same kind of operation: some obtain information from outside the agent, while others may involve internal processing of information already available to it.

The system might also be able to issue a remediation command. But the fact that an operation exists does not mean it should always be carried out. A command that changes production state is different from a read-only check. The agent's architecture must make it possible to reason about what operations are available and how decisions about them are made.

This book will later use the example to discuss engineering boundaries such as authorization and human approval. Those controls are important implementation decisions for a production system. They should not be misrepresented as requirements that CoALA itself prescribes.

\subsection{Question 3: How does the system decide what to do next?}\label{question-3-how-does-the-system-decide-what-to-do-next}

Suppose the first health check shows that the application is responding, but the database pool is exhausted. The agent might now retrieve database metrics, compare the observation with a recent deployment, or propose another diagnostic step. A new observation may change which possibilities deserve attention.

A fixed list of steps may be useful for a routine procedure, but an investigation often needs to respond to new information. The system must decide what candidate operation is appropriate in its current situation and what to do after the operation returns a result.

This is not just a question of which tool should be called. It is a question about how the system organizes its information, its possible actions, and the procedure by which it selects what happens next.

\subsection{The architectural questions exposed by the incident}\label{the-architectural-questions-exposed-by-the-incident}

The incident brings three questions into focus:

\begin{longtable}[]{@{}
  >{\raggedright\arraybackslash}p{(\columnwidth - 2\tabcolsep) * \real{0.5000}}
  >{\raggedright\arraybackslash}p{(\columnwidth - 2\tabcolsep) * \real{0.5000}}@{}}
\toprule\noalign{}
\begin{minipage}[b]{\linewidth}\raggedright
Question raised by Product X
\end{minipage} & \begin{minipage}[b]{\linewidth}\raggedright
Architectural dimension we will study
\end{minipage} \\
\midrule\noalign{}
\endhead
\bottomrule\noalign{}
\endlastfoot
What information does the agent have, and where is it maintained? & Information storage \\
What can the agent do internally or through its environment? & Action space \\
How does the agent choose and execute what happens next? & Decision-making procedure \\
\end{longtable}

This is the point of the running example. It gives us a practical reason to care about the distinctions that CoALA makes. It does not show that all production incident systems must be built in one particular way, and it does not make the details of this example part of the paper.

\section{CoALA: a framework for organizing language agents}\label{coala-a-framework-for-organizing-language-agents}

CoALA stands for \textbf{Cognitive Architectures for Language Agents}. The original paper proposes it as a conceptual framework for characterizing and designing general-purpose language agents \CoALACite{§4}. The framework positions the language model as the core component of a larger cognitive architecture, not as the complete agent by itself.

CoALA organizes agents along three dimensions. The table is a preview; Chapters 2--8 develop the distinctions and their relationships.

\BookDiagram[0.52\textheight]
  {Tex/diagrams/generated/part1/chapter01/coala-three-dimensions.png}
  {The three dimensions CoALA uses to organize a language agent. The detailed chapters examine each dimension and the relationships among them.}
  {fig:chapter01-coala-three-dimensions}

\begin{center}
\begin{tabularx}{\textwidth}{>{\bfseries}p{0.23\textwidth} X p{0.18\textwidth}}
\toprule
Dimension & Architectural question & Developed in \\
\midrule
Information storage & What information is active now, and what is retained for later use? & Chapter 3 \\
Action space & What internal operations and external interactions are available? & Chapters 4--5 \\
Decision-making procedure & How are candidate actions proposed, evaluated, selected, and executed over time? & Chapters 6--8 \\
\bottomrule
\end{tabularx}
\end{center}

The dimensions are distinct but interdependent: decisions draw on maintained information and available operations, while the results of actions can change what information is available in later cycles. CoALA supplies a shared vocabulary for describing those relationships without requiring one fixed implementation.

\section{What CoALA is---and what it is not}\label{what-coala-isand-what-it-is-not}

Before moving on, it is worth establishing a few boundaries. These boundaries will protect us from attributing more to the paper than it actually claims.

\subsection{What CoALA contributes}\label{what-coala-contributes}

The original paper makes several contributions relevant to this book \CoALACite{§§1 and 4--6}:

\begin{enumerate}
\def\labelenumi{\arabic{enumi}.}
\tightlist
\item
  It proposes a conceptual framework for organizing language agents using information storage, action space, and decision-making procedure.
\item
  It connects the development of language agents to the history of production systems and cognitive architectures.
\item
  It uses the framework to organize and compare a range of existing language-agent approaches.
\item
  It identifies directions for the development of more capable language agents.
\end{enumerate}

CoALA is therefore more than a diagram of memory boxes. The decision-making procedure is part of the framework. At the same time, CoALA is broader than one particular implementation or one fixed reasoning strategy.

\subsection{What CoALA does not prescribe}\label{what-coala-does-not-prescribe}

CoALA is a conceptual framework, not a ready-to-run library, runtime, or deployment architecture. It does not mandate a programming language, framework, database, model provider, or identical memory and action configuration for every agent. Those choices depend on the task and must be justified separately. Nor does mapping a program to CoALA prove that it is reliable, safe, or production-ready; requirements analysis, testing, security engineering, and operational evaluation remain necessary.

This separation will matter in Part III. When we build production-oriented Java systems, we will identify which ideas come from CoALA, which come from Fan \& Lan, and which are engineering decisions made for the example system. The papers will remain the theoretical foundations. Framework APIs, build tools, and implementation libraries will be consulted only to write correct code, not treated as additional cognitive theories.

\subsection{What this book will do with CoALA}\label{what-this-book-will-do-with-coala}

We will first understand the framework on its own terms. Part I follows the concepts in a deliberate order: architecture, memory, actions, planning, evaluation, selection, and the iterative decision process. We will avoid importing mechanisms from later research into our explanation of the original CoALA paper.

In Part II, we will study Fan \& Lan's mechanism-level analysis as a separate contribution. The purpose will not be to rewrite CoALA or claim that its original authors proposed every mechanism discussed there. We will preserve the distinction between the original framework and the later analysis.

In Part III, we will apply the concepts to one request-driven enterprise supply-chain disruption management system implemented in Java. Chapter 19 establishes its engineering foundations, and Chapter 20 develops the application. The design is framework-independent: its architecture, state and data contracts, control flow, persistence, safety boundaries, recovery procedures, and evaluation strategy are book-level engineering decisions. They illustrate one coherent set of choices and trade-offs; they are not implementations mandated by either paper.

\section{Where we go next}\label{where-we-go-next}

The Product X incident shows why a model's generated answer is only one part of an interactive system: the agent also needs organized information, available operations, and control over what happens next. The historical bridge from production systems to cognitive architectures explains why these concerns belong together. CoALA brings them into one framework through information storage, action space, and decision-making procedure. Chapter 2 develops that architectural map before the later chapters examine its components and processes.

\begin{chapterrecap}
A language model call describes text generation, not the full behavior of an interactive system. A language agent places model use in a feedback loop with an environment; a cognitive language agent also uses internal processes to manage information and behavior. CoALA draws on the history of cognitive architectures to organize language agents through information storage, action space, and decision-making procedure. Those dimensions give us a shared vocabulary without prescribing one implementation or proving that a system is production-ready.
\end{chapterrecap}

\begin{center}\rule{0.5\linewidth}{0.5pt}\end{center}

\section{Source and example note}\label{source-and-example-note}

The theoretical account in this chapter is grounded in \textbf{{[}CoALA{]}} only. The Product X payment-incident scenario and all of its specific details are illustrative material created for this book; they are not presented as a case study from the paper. The thermostat and string-rewriting examples are used to explain the historical discussion as it is presented in \textbf{{[}CoALA{]}}. The three diagrams are book-authored teaching diagrams maintained as Mermaid text sources and rendered to PNG during the build; they summarize concepts rather than reproduce the paper's artwork. No mechanism from Fan \& Lan or from another research framework is introduced in this chapter.


\chapter{The Architecture of CoALA}
\label{ch:coala-architecture}

\begin{learningbox}
By the end of this chapter, you should be able to read CoALA as one connected architecture rather than as a list of terms. You will be able to identify its three organizing dimensions, explain the relationship between the structural view and the decision-cycle view, distinguish the LLM from the larger agent, and trace how information, actions, and control fit together. The chapter deliberately previews memory and action categories without replacing the detailed treatment in Chapters 3--5.
\end{learningbox}

\section{From a collection of capabilities to an architecture}

Imagine an agent that can call a language model, retain information from earlier tasks, retrieve documents, invoke tools, and follow a program. We can list those capabilities, perhaps even identify the software components that implement them. But a list does not yet explain how they form one agent. Which information is active when a decision is made? How does the model use that information? Which operations can change internal information, and which can affect the environment? What decides which operation happens next?

Return to Product X, our book-created software incident scenario. The payment service is returning HTTP 500 errors. The agent may need the current incident report, recent deployment information, facts about the service, records of earlier incidents, and procedures for checking production systems. It may need to inspect logs, retrieve a runbook, reason about a possible cause, ask an operator a question, or run a health check. Those are not interchangeable activities. Some read retained information; some process information already available; some write new information; others interact with the running service or with a person.

The important question is not simply whether each capability exists. It is how each capability relates to the others and how the system coordinates their use. A component can be technically present but disconnected from the decision process. A tool can be registered but never selected. A memory store can contain useful records that are not retrieved at the right time. A model can produce a plausible instruction that no procedure actually executes. To understand the agent, we need to understand its organization.

In this chapter, we use the word \term{architecture} in a practical teaching sense: a map of the major elements of a system, the responsibilities assigned to them, and the relationships through which information and control move. This is an explanatory description for the reader, not a new formal definition claimed to appear verbatim in the CoALA paper. It is also not a deployment diagram. It does not tell us which cloud provider hosts a service, how many processes run, or which database product stores a record. Nor is it a class diagram that prescribes source-code types. It is a conceptual map that helps us ask what the agent contains and how its parts cooperate.

\begin{tipbox}
When reading an architecture diagram, do not stop at the boxes. Ask what each box is responsible for, what information enters and leaves it, and which part of the system determines when the interaction occurs.
\end{tipbox}

CoALA offers a common way to organize these questions. It describes a language agent along three dimensions: \term{information storage}, \term{action space}, and \term{decision-making procedure}. The LLM is central, but it is one component inside a larger cognitive architecture.\CoALACite{Section 4 and Figure 4}

\section{The three dimensions that organize CoALA}

The three dimensions can be introduced with three questions:

\begin{center}
\begin{tabularx}{\textwidth}{>{\bfseries}p{0.25\textwidth} X}
\toprule
Architectural dimension & Question it answers \\
\midrule
Information storage & What information does the agent maintain, and where is it available? \\
Action space & What can the agent do with its internal information, and how can it interact with the outside world? \\
Decision-making procedure & How does the agent choose and coordinate actions over time? \\
\bottomrule
\end{tabularx}
\end{center}

These are three views of one system, not three independent products that happen to be installed side by side. Information storage describes what the agent can draw on. The action space describes how the agent can use internal information and interact with its environment. The decision-making procedure coordinates the use of those capabilities. A description that omits one of the three dimensions leaves an important architectural question unanswered.

\subsection{Information storage: what information is maintained?}

A language model call receives an input and produces an output. In the simplest interaction, information supplied for one call is not automatically maintained as a structured state for the next call. A \gls{languageagent} that operates across multiple steps may need to preserve information between calls and decision cycles. CoALA therefore distinguishes \gls{workingmemory} from \gls{longtermmemory}. Long-term memory, in turn, includes episodic, semantic, and procedural forms.\CoALACite{Section 4.1}

For now, we need only the architectural contrast. Working memory is the active information for the current decision cycle. Long-term memory retains information that can be used beyond that immediate cycle. The long-term categories distinguish experience, knowledge, and how-to information. Chapter 3 will examine their meanings, boundaries, and examples in detail. Here, their purpose is to show that an architecture needs more than one undifferentiated store called “memory.”

\subsection{Action space: what can the agent do?}

An agent must do more than contain information. It must have operations for using it. CoALA distinguishes external actions from internal actions. External actions interact with an environment through grounding. Internal actions interact with the agent's own memory. The paper further describes three kinds of internal action: retrieval, reasoning, and learning.\CoALACite{Sections 4.2--4.5 and Figure 5}

At the level of this chapter, the distinctions are simple. Retrieval reads from long-term memory. Reasoning uses the LLM to process information in working memory. Learning writes information to long-term memory. Grounding executes external actions and processes environmental feedback. Chapters 4 and 5 will examine the operational details. For the architectural map, notice that “action” is broader than “call a tool”: some actions operate on internal information, while others affect or obtain information from the external environment.

\subsection{Decision-making procedure: how are actions coordinated?}

If an agent has information and operations, we still need to know how it determines what to do. CoALA describes decision-making as a repeated process involving planning and execution. During a decision cycle, the agent can use retrieval and reasoning to support planning. Planning proposes and evaluates candidate learning or grounding actions; a selected action is then executed. An observation may follow, and the next cycle begins.\CoALACite{Section 4.6 and Figure 4B}

This high-level description is enough to understand why the decision-making procedure is a separate architectural dimension. The action space tells us which operations are available. The decision procedure tells us how the agent coordinates them. We will separate planning, evaluation, selection, and execution in Chapters 6--8; this chapter establishes only the overall relationship.

\begin{takeawaybox}[The architecture in one sentence]
CoALA organizes a language agent by asking what information it maintains, what internal and external actions it can perform, and what decision procedure coordinates those actions in a repeated interaction with the world.
\end{takeawaybox}

\section{The structural map: components and relationships}

The three questions can be read as a set of responsibilities, but CoALA also shows how the elements relate. Its structural view places the LLM and agent code alongside working memory, long-term memory, the decision procedure, and interfaces to the external environment. The connections matter: the agent code constructs prompts and parses model outputs; memory supplies and receives information; the decision procedure invokes the action procedures; grounding carries actions and observations across the boundary between the agent and its environment.\CoALACite{Figure 4A and Section 4}

\begin{figure}[htbp]
\centering
\resizebox{\textwidth}{!}{%
\begin{tikzpicture}[
  box/.style={draw=BookBlue, rounded corners=2pt, thick, fill=NoteBg, minimum width=2.2cm, minimum height=0.85cm, align=center, font=\sffamily\small},
  memory/.style={draw=BookBlue, rounded corners=2pt, thick, fill=LearnBg, minimum width=2.35cm, minimum height=0.75cm, align=center, font=\sffamily\small},
  action/.style={draw=BookMid, rounded corners=2pt, fill=TipBg, minimum width=1.8cm, minimum height=0.68cm, align=center, font=\sffamily\small},
  flow/.style={-{Stealth[length=2.5mm]}, thick, draw=BookBlue},
  back/.style={-{Stealth[length=2.5mm]}, thick, draw=BookMid, dashed}
]
\draw[rounded corners=5pt, thick, draw=BookBlue] (-0.35,-3.45) rectangle (10.65,2.5);
\node[font=\sffamily\bfseries\color{BookBlue},anchor=west] at (-0.1,2.22) {Cognitive language agent};
\node[box] (llm) at (1.3,1.2) {LLM};
\node[box] (code) at (4.5,1.2) {Agent code};
\node[memory] (wm) at (7.8,1.2) {Working memory};
\node[box] (decision) at (1.7,-1.0) {Decision procedure};
\node[action] (retrieval) at (4.7,-0.35) {Retrieval};
\node[action] (reasoning) at (6.8,-0.35) {Reasoning};
\node[action] (learning) at (8.9,-0.35) {Learning};
\node[action] (grounding) at (9.0,-1.4) {Grounding};
\node[memory] (ltm) at (6.8,-2.35) {Long-term memory};
\node[box] (env) at (12.3,-2.35) {External environment};
\draw[flow] (code.west) -- (llm.east);
\node[font=\scriptsize,align=center] at (2.9,1.95) {prompt construction / output parsing};
\draw[flow] (wm.west) -- (code.east);
\node[font=\scriptsize] at (6.15,1.95) {active context};
\draw[flow] (decision) -- (code);
\draw[flow] (code) -- (retrieval);
\draw[flow] (code) -- (reasoning);
\draw[flow] (code) -- (learning);
\draw[flow] (retrieval) -- (ltm);
\draw[flow] (learning) -- (ltm);
\draw[back] (ltm) -- (wm);
\draw[flow] (code.south east) .. controls (6.0,-0.2) and (7.6,-1.4) .. (grounding.west);
\draw[flow] (grounding.east) -- (env.north west);
\draw[back] (env.west) .. controls (10.7,-3.05) and (9.1,-3.15) .. node[below,font=\scriptsize]{observations} (wm.south);
\end{tikzpicture}%
}
\caption{An original teaching redraw of the structural relationships in CoALA's Figure 4A. The drawing is schematic: arrows identify information or control relationships, not a required software topology.}
\label{fig:coala-structural-map}
\end{figure}

The figure is a teaching redraw, not a reproduction of the paper's artwork. Its purpose is to help us read the relationships. It does not prescribe that an implementation must use separate machines, processes, classes, or databases for every box. Nor should every arrow be mistaken for a network call. The diagram expresses roles and interactions at the conceptual level.

\subsection{The LLM and agent code}

The LLM is a core computational component of a \gls{cognitivelanguageagent}. It supplies language-based processing that can support reasoning and generate outputs used by the agent. But CoALA does not identify the LLM with the whole agent. The architecture also includes agent code: explicit procedures that construct model inputs, parse outputs, implement actions, and execute the decision-making process.\CoALACite{Section 4 and Figure 4A}

This distinction is easier to see if we separate two questions. First, what capabilities are available through the model? The LLM's learned parameters encode implicit knowledge that can influence the output it generates. Second, what procedures does the surrounding system explicitly execute? Agent code defines how prompts are built, how model outputs are interpreted, how memory is accessed, how grounding actions are performed, and how the decision process proceeds.

The paper also describes \gls{proceduralmemory} as having implicit and explicit forms: knowledge stored in the LLM weights and knowledge written in the agent's code. These are not identical storage mechanisms, even though both can contribute to how the agent behaves. Chapter 3 will examine this distinction as part of the wider memory architecture. For now, remember that an LLM call does not itself supply every operation needed to run the agent.

\subsection{Working memory as the active information hub}

Working memory occupies a central position because it maintains active information for the current decision cycle. CoALA describes it as a data structure that persists across LLM calls. The LLM input is synthesized from a subset of working memory, such as a prompt template populated with relevant variables. The model's output is parsed and stored back into variables in working memory, where it can be used to execute a corresponding action. Working memory also interacts with long-term memory and grounding interfaces.\CoALACite{Section 4.1}

That description gives us a useful correction to a common assumption: working memory is not simply another name for the current prompt. A prompt is an input representation prepared for one model invocation. Working memory is the larger active information structure from which that input can be assembled and into which results can be written. The current prompt may contain only the information relevant to a specific call.

Consider Product X. The active working information might include the incident goal, the latest observation, the current hypothesis, and the result of the most recent check. A particular model call may need only the latest error pattern and a request to compare two hypotheses. Another call may need a different subset, perhaps the result of a deployment check. The architecture can maintain more active information than is placed in any one prompt.

\begin{notebox}
The term “working memory” is used here in CoALA's architectural sense. Do not equate it automatically with a model vendor's context window, conversation transcript, cache, or a specific in-memory data structure. An implementation may connect these ideas, but the conceptual role is broader than any one mechanism.
\end{notebox}

\subsection{Long-term memory at the architectural level}

CoALA differentiates long-term memory into \gls{episodicmemory}, \gls{semanticmemory}, and procedural forms. The categories answer different questions: what happened in prior experiences, what the agent knows about the world and itself, and how the agent performs actions and decision procedures. The paper discusses both implicit procedural knowledge in model weights and explicit procedural knowledge in code.\CoALACite{Section 4.1}

This chapter does not attempt to define every category in full. It is enough to understand why the architecture distinguishes them. If the agent keeps an incident history, stores facts about the payment service, and contains executable procedures for checking the service, those are not the same kind of information. The design question is not simply “where is the memory database?” It is “what kind of information is this, what role does it play, and how does the agent access or update it?”

Chapter 3 will explore the memory categories, their relationships, and the implications of treating memory as organized information rather than as one transcript. Here the only point is architectural: long-term memory is connected to working memory through internal actions, and the decision procedure determines when those actions are used.

\subsection{The environment and the two-way connection}

The agent is not isolated inside its memory modules. \gls{grounding} procedures execute external actions and process environmental feedback into working memory. The CoALA paper describes physical environments, dialogue with humans or other agents, and digital environments as examples of external settings.\CoALACite{Section 4.2}

In Product X, a digital environment may expose service health, log records, deployment history, and database metrics. A dialogue environment may include an operator who can clarify whether a change was authorized. A physical environment is less central to this incident, but the same conceptual boundary applies to a robot receiving sensor data and issuing motor commands. The environment differs, while the architectural role of grounding remains recognizable.

The connection is two-way. Actions go from the agent toward the environment; observations come back. An action that is merely proposed in text has not yet crossed that boundary. A grounding procedure must execute the interaction and process the resulting feedback. That feedback can then become active information in working memory and influence a later decision.

\subsection{Reading the connections, not just the boxes}

An architecture diagram can be misleading when a reader treats it as an inventory. The value of the structural view lies in the relationships.

\begin{itemize}
  \item The decision procedure invokes explicit procedures in agent code.
  \item Agent code interacts with the LLM by preparing input and parsing output.
  \item Working memory provides active information for the current cycle and receives results.
  \item Retrieval and learning connect working memory with long-term memory through reads and writes.
  \item Grounding connects the agent with the external environment through actions and observations.
\end{itemize}

A single component may participate in several relationships. The LLM can help reason over working information, but it does not by itself decide how every operation is executed. Working memory may be the active information hub, but it is not the only memory store. Agent code implements procedures, but it is the decision procedure that coordinates how they are used over time. The environment contributes observations and receives actions, but it does not belong to the agent's internal memory architecture.

\section{Memory's place in the architecture}

It is useful to pause on memory because the structural diagram places it at the centre of several connections. A reader who thinks of memory as a single database may miss why CoALA separates working memory from long-term memory and then differentiates long-term memory into subtypes. The purpose of this section is not to teach those subtypes in full; it is to clarify their place in the architecture.

Working memory supports the current decision cycle. It holds relevant observations, active goals, knowledge retrieved from long-term memory, and results generated during reasoning. Long-term memory allows the agent to maintain information beyond the current cycle. Retrieval can bring relevant information into working memory; learning can write information from working memory into long-term memory. These connections make memory part of the agent's ongoing activity rather than a passive archive.

The categories also show that “remembering” is not one operation. Reading an earlier incident, retaining a fact about a service, and using an executable procedure are different roles. An architecture that labels all of them merely “memory” may hide important design differences. Conversely, the names of the categories do not require separate physical databases. The conceptual distinctions can be represented in different ways in software, provided the implementation makes their responsibilities clear.

\begin{warningbox}
Do not infer from the diagram that CoALA mandates a particular storage product, schema, vector database, file format, persistence layer, or retention policy. Those are implementation decisions. CoALA provides the conceptual memory distinctions; engineering choices must be justified separately.
\end{warningbox}

Chapter 3 will unpack working memory and the long-term memory categories. The current chapter only establishes the links: memory makes information available, actions read or write it, and the decision procedure determines when those operations are used.

\section{The action space: internal capabilities and external interaction}

The structural view also connects the agent to two broad kinds of action. External actions interact with an environment through grounding. Internal actions interact with internal memories. The internal category is further divided into retrieval, reasoning, and learning.\CoALACite{Sections 4.2--4.5 and Figure 5}

\subsection{Internal actions}

The three internal action categories are distinguished by what information they access and how they change it.

\begin{description}
  \item[Retrieval] reads information from long-term memory. Its role is to make relevant retained information available for the current decision process.
  \item[Reasoning] uses the LLM to process information in working memory and generate new information there. It reads from and writes to working memory.
  \item[Learning] writes information to long-term memory. In the CoALA account, learning is not limited to changing neural-network weights; it can include storing information in memory modules.
\end{description}

These definitions are deliberately brief. Chapter 4 will explain retrieval, reasoning, and learning as operational categories, including their directions of information movement and examples. This chapter only needs the distinction to make sense of the map.

\subsection{External action: grounding}

\gls{grounding} procedures execute external actions and process feedback from the environment into working memory. The paper groups external interaction into physical environments, dialogue with humans or other agents, and digital environments. A robot may receive visual and tactile observations and issue motor commands. A dialogue agent may ask a person a question or communicate with another agent. A digital agent may navigate a website, interact with an API, play a game, or execute code.\CoALACite{Section 4.2}

In each case, the action is more than a sentence describing an intention. It is an operation performed through an interface to the external world. The system must interpret the result as an observation that can inform later decisions. The agent's environment therefore forms part of the architecture's boundary: it is outside the agent's internal memory, but it is connected to the agent through grounding.

\subsection{Why planning is not the same as the entire action space}

Figure 5 in the CoALA paper provides a compact view of these categories. Reasoning and retrieval are internal actions that support planning. Learning is also an internal action, but it is not placed within the planning bracket in the figure. Grounding is external. This is an important distinction: planning uses certain internal actions to determine which learning or grounding action to select; it is not a synonym for every action available to the agent.\CoALACite{Figure 5}

\begin{figure}[htbp]
\centering
\begin{tikzpicture}[
  act/.style={draw=BookBlue, thick, fill=TipBg, minimum width=1.25cm, minimum height=0.8cm, align=center, font=\sffamily\small},
  group/.style={draw=BookBlue, thick, rounded corners=2pt}
]
\node[act] (reason) at (0,0) {Reasoning};
\node[act] (retrieval) at (1.55,0) {Retrieval};
\node[act] (learning) at (3.1,0) {Learning};
\node[act] (grounding) at (4.65,0) {Grounding};
% Grouping brackets are deliberately outside the action nodes so that the
% categories remain readable at print size.
\draw[thick,draw=BookBlue] (-0.68,0.72) -- (3.78,0.72);
\draw[thick,draw=BookBlue] (-0.68,0.72) -- (-0.68,0.58);
\draw[thick,draw=BookBlue] (3.78,0.72) -- (3.78,0.58);
\draw[thick,draw=BookBlue] (3.98,0.72) -- (5.32,0.72);
\draw[thick,draw=BookBlue] (3.98,0.72) -- (3.98,0.58);
\draw[thick,draw=BookBlue] (5.32,0.72) -- (5.32,0.58);
\node[font=\sffamily\bfseries\small,BookBlue] at (1.55,1.0) {Internal};
\node[font=\sffamily\bfseries\small,BookBlue] at (4.65,1.0) {External};
\draw[thick,draw=BookBlue] (-0.68,-0.72) -- (2.32,-0.72);
\draw[thick,draw=BookBlue] (-0.68,-0.72) -- (-0.68,-0.58);
\draw[thick,draw=BookBlue] (2.32,-0.72) -- (2.32,-0.58);
\node[font=\sffamily\small,BookBlue] at (0.82,-1.04) {Planning support};
\end{tikzpicture}
\caption{An original redraw of CoALA's action-space distinction. Reasoning and retrieval support planning; learning remains an internal action; grounding is external.}
\label{fig:coala-action-space}
\end{figure}

The distinction helps avoid a common overgeneralization. An agent's external tool calls are not its whole action space, because internal memory operations also matter. Conversely, not every internal operation is itself a planning step. Chapters 4 and 5 will examine the actions in detail; Chapters 6--8 will explain how planning works with them.

\section{The decision procedure: how the architecture operates over time}

A structural map tells us which elements exist and how they relate, but it does not fully show their order over time. For that, CoALA presents a \gls{decisionprocedure} view. The paper compares the decision procedure to a program's main procedure: a top-level routine that runs repeatedly, accepts new perceptual input, and calls action procedures in response.\CoALACite{Section 4.6}

\subsection{The top-level or main procedure}

Think of an ordinary software application. A program may contain functions for reading input, validating data, writing a file, and returning a result. The existence of those functions does not explain when they run. A top-level procedure determines the control flow: which functions are called, in what context, and what happens after they return.

CoALA's decision procedure plays a comparable organizing role for the agent. It coordinates action procedures and repeats the interaction with the environment. This comparison is a teaching analogy based on the paper's description; it does not mean that every implementation must contain a method literally named `main`, nor that the full decision process is a single fixed code listing.

\subsection{Planning and execution at a glance}

The paper's temporal view places planning before execution. During planning, retrieval and reasoning can help the agent propose and evaluate candidate learning or grounding actions. Selection chooses an action. Execution then carries out the selected action. An observation may be made, after which the decision cycle begins again.\CoALACite{Section 4.6 and Figure 4B}

\begin{figure}[htbp]
\centering
\begin{tikzpicture}[
  node distance=0.55cm,
  stage/.style={draw=BookBlue, thick, rounded corners=2pt, fill=TipBg, minimum width=2.2cm, minimum height=0.62cm, align=center, font=\sffamily\small},
  container/.style={draw=BookBlue, thick, rounded corners=4pt, inner sep=8pt},
  flow/.style={-{Stealth[length=2.4mm]}, thick, draw=BookBlue}
]
\node[stage] (obs) {Observation};
\node[stage,below=of obs] (proposal) {Proposal};
\node[stage,below=of proposal] (evaluation) {Evaluation};
\node[stage,below=of evaluation] (selection) {Selection};
\node[stage,below=0.8cm of selection] (execution) {Execution};
\draw[container] (-1.6,-3.3) rectangle (1.6,-0.4);
\node[font=\sffamily\bfseries\color{BookBlue}] at (0,-0.16) {Planning};
\draw[flow] (obs) -- (proposal);
\draw[flow] (proposal) -- (evaluation);
\draw[flow] (evaluation) -- (selection);
\draw[flow] (selection) -- (execution);
\draw[flow] (execution.east) -- ++(1.0,0) |- (obs.east) node[pos=0.28,right,font=\scriptsize,align=left]{next cycle / feedback};
\end{tikzpicture}
\caption{An original teaching redraw of CoALA's decision-cycle view. The stages are shown in the explanatory order used by the figure; the detailed chapters will examine interleaving and iteration.}
\label{fig:coala-decision-cycle}
\end{figure}

It is important not to read this high-level picture as a complete low-level algorithm. The figure helps us see the broad relation between planning and execution. It does not mean that the agent must perform one isolated reasoning call, then one isolated retrieval, then one proposal, then one evaluation, with no possibility of returning to an earlier subtask. The details of how planning's sub-stages relate, and how they may be revisited, are the subject of Chapters 6--8.

\subsection{A cycle is a useful view, not a rigid script}

The word \term{cycle} emphasizes repetition. It reminds us that an agent may receive new observations after an action and must decide what to do in light of them. The cycle also shows that the decision procedure does not merely select external tool calls. It can select a learning action that updates long-term memory as well as a grounding action that interacts with the environment.

At the same time, the diagram should not be mistaken for a universal, mandatory sequence that every CoALA agent must follow identically. CoALA is a conceptual framework. The paper describes planning in terms of proposal, evaluation, and selection, but the architecture is intended to organize varied agents and decision procedures. The detailed account in Chapters 6--8 will preserve the specific roles of these stages without imposing a rigid one-pass pipeline where the paper does not.

\section{A complete architectural reading of Product X}

We can now return to the payment incident and read it through the architecture. The following scenario is created for this book. Its purpose is to apply CoALA's distinctions, not to claim that the paper presents this incident or prescribes the specific operational steps.

The payment service is returning HTTP 500 errors. The agent's goal is to investigate the incident and provide a supported recommendation. It begins with an observation from the monitoring system: the error rate increased shortly after deployment 4.7. The agent does not yet know whether the deployment caused the failure.

\subsection{Information storage}

Working memory contains the active incident goal, the latest observations, relevant facts already in hand, and the current state of the investigation. Long-term memory may contain records of earlier incidents, facts about the payment service, and procedures for checking its health. The architecture distinguishes the role of active information from information retained for future use.

\subsection{Action space}

The agent may retrieve previous incidents or service facts from long-term memory. It may reason over the deployment timestamp and the error increase to formulate candidate explanations. It may learn by writing a new experience or fact to long-term memory when the procedure calls for it. It may ground by requesting a health check, reading logs through an external interface, or asking an operator for clarification.

These operations are different even when they are triggered by the same investigation. Retrieving a runbook does not itself inspect the service. Reasoning about a possible database failure does not query the database. A grounding procedure performs the external interaction and processes the result as an observation.

\subsection{Decision-making procedure}

The decision procedure coordinates what happens next. It can use retrieval and reasoning to support planning. It can propose candidate actions such as inspect database connectivity or compare the recent deployment with the error timeline. It can evaluate the candidates and select one for execution. Once the selected operation is carried out, an observation may update working memory. The next cycle can then take the new information into account.

The exact investigation strategy, authorization process, rollback policy, and operational safety controls are not defined by the CoALA architecture map itself. They would be engineering choices for a real incident-response system. CoALA helps us describe the cognitive organization of the agent; it does not automatically make the system safe to execute production changes.

\begin{takeawaybox}[What the example demonstrates]
A single incident can involve every architectural dimension. The agent maintains active and retained information, uses internal and external actions, and relies on a decision procedure to coordinate them. The example makes the architecture concrete without adding new CoALA mechanisms.
\end{takeawaybox}

\section{One framework, different agent configurations}

A framework is useful only if it can describe more than one narrowly constructed system. CoALA is intended to organize a range of language agents, not to require one fixed collection of components or one identical decision strategy. The paper applies the framework to different systems and uses it to compare their memory, actions, and decision-making characteristics.\CoALACite{Section 4 and Table 2}

This means that two agents can both be understood through CoALA while differing substantially in what they implement. One may have a simple memory arrangement and a small action space. Another may maintain several forms of long-term memory and use a more elaborate decision procedure. One agent may primarily interact through dialogue; another may act in a digital environment. The common framework provides the vocabulary for describing those differences.

The purpose of a common vocabulary is not to erase variation. It is to make variation easier to see. Instead of saying only that one system is “more agentic” or “uses more tools,” we can ask: What information does it maintain? What internal operations does it support? What external environment does it interact with? How does its decision procedure select and execute operations? Where are its capabilities limited?

This is also why a framework should not be confused with a product specification. CoALA gives us a way to describe and reason about language-agent organization. A particular application still needs requirements, domain constraints, interfaces, reliability properties, safety policies, and evaluation methods. Those are not supplied automatically by drawing the architecture.

\section{Common misreadings and architectural boundaries}

Several misunderstandings can arise when the architecture is reduced to a few memorable labels. It is worth correcting them before moving to the detailed chapters.

\begin{enumerate}
  \item \textbf{The LLM is the agent.} The LLM is central, but the architecture also includes agent code, memory, action procedures, and a decision procedure.
  \item \textbf{Working memory is just the prompt.} The prompt is synthesized from a relevant subset of working memory. Working memory is the active information structure that can persist across model calls.
  \item \textbf{All actions are tool calls.} CoALA includes internal memory actions as well as external grounding actions. Retrieval, reasoning, and learning are not interchangeable with external tool execution.
  \item \textbf{Planning is the whole action space.} Planning uses reasoning and retrieval to support candidate generation and evaluation. Learning is also an internal action, and grounding is an external action.
  \item \textbf{The diagram is a deployment plan.} The boxes are conceptual roles. They do not require one service, process, class, or database per box.
  \item \textbf{The decision cycle is a rigid one-pass script.} The cycle is a high-level view of repeated decision-making. Chapters 6--8 will examine the roles and revisiting of planning sub-stages carefully.
  \item \textbf{A CoALA description proves production readiness.} The framework is not a substitute for authorization, testing, security, reliability, observability, or domain validation.\index{architecture!cognitive} \index{memory!working} \index{memory!long-term} \index{grounding} \index{decision procedure}
\end{enumerate}

\begin{warningbox}
Do not import mechanisms from Fan and Lan into the explanation of the original CoALA framework. Part II will study that later mechanism-level analysis separately. Likewise, do not attribute implementation choices made in Part III to CoALA unless the original paper explicitly supports the claim.
\end{warningbox}

\section{How the remaining CoALA chapters build on this map}

This chapter has provided the map. The next chapters will zoom in on individual regions without losing the relationships between them.

\begin{itemize}
  \item \textbf{Chapter 3 --- Understanding Memory.} Working memory, episodic memory, semantic memory, procedural memory, and the boundaries between them.
  \item \textbf{Chapter 4 --- Internal Actions: Retrieval, Reasoning and Learning.} The internal operations, their inputs and outputs, and how they affect memory.
  \item \textbf{Chapter 5 --- External Actions and Grounding.} External environments, grounding procedures, observations, and action examples.
  \item \textbf{Chapter 6 --- Planning: From a Problem to Candidate Actions.} The proposal stage and the generation of candidate actions.
  \item \textbf{Chapter 7 --- Evaluation and Selection.} Candidate evaluation, values, selection methods, and the possibility of rejecting candidates.
  \item \textbf{Chapter 8 --- Interleaving, Iteration and the Complete CoALA Decision Process.} How the planning stages can be revisited and interleaved, and how the complete decision process operates.
\end{itemize}

The order matters pedagogically. We first need the whole map, then we can inspect its components and processes in more detail. The chapters will return to the same Product X example so that the reader can see how a single incident draws on memory, actions, and decision-making. The example remains a teaching aid, not a source claim.

\section{Chapter recap}

\begin{chapterrecap}
CoALA organizes language agents through three dimensions: information storage, action space, and decision-making procedure. The LLM is a core component within a larger architecture. Agent code implements procedures that connect the LLM to working memory, long-term memory, and the external environment. Internal actions include retrieval, reasoning, and learning; external actions use grounding. The decision procedure coordinates planning and execution in a repeated cycle. The structural diagram explains the components and their relationships, while the temporal diagram explains how the agent operates over time. Both views describe the same architecture from different angles.
\end{chapterrecap}

\begin{reviewquestions}{Review questions}
\begin{enumerate}
  \item Why is a list of capabilities not enough to explain an agent's architecture?
  \item What are the three dimensions CoALA uses to organize language agents?
  \item Why is working memory broader than the prompt supplied to one LLM call?
  \item How do internal actions differ from external grounding actions?
  \item Why are reasoning and retrieval associated with planning in the action-space diagram?
  \item What does the structural view show that the decision-cycle view does not, and vice versa?
  \item Why should the decision-cycle diagram not be read as a rigid one-pass script?
  \item Which decisions in the Product X example are illustrations rather than requirements prescribed by CoALA?
\end{enumerate}
\end{reviewquestions}

\section*{Source and example note}
The theoretical account in this chapter is grounded in \textsc{[CoALA]} only, principally Section 4 and Figures 4--5 of Sumers et al. The diagrams are original teaching redraws of the paper's conceptual relationships. Product X and all incident-specific details are examples created for this book. This chapter does not use Fan and Lan to reinterpret CoALA and does not prescribe a software implementation.

\chapter{Understanding Memory}
\begin{notebox}
Chapter placeholder. The detailed outline must be reviewed and approved before this chapter is written.
\end{notebox}

\chapter{Internal Actions: Retrieval, Reasoning and Learning}
\begin{notebox}
Chapter placeholder. The detailed outline must be reviewed and approved before this chapter is written.
\end{notebox}

\chapter{External Actions and Grounding}
\label{ch:external-actions-grounding}

\section*{Chapter at a glance}
\phantomsection
\addcontentsline{toc}{section}{Chapter at a glance}

An agent can retrieve information from its own long-term memory and reason about information already available to it. Those operations do not, by themselves, inspect a changing website, ask a person a question, move a robot, or query a service that exists outside the agent. To do those things, the agent must interact with an external environment.

CoALA calls this part of the action space \emph{grounding}. A grounding procedure executes an external action and processes environmental feedback so that it can become available to the agent, represented in working memory as text. This chapter explains what that means, distinguishes the environment from the procedure that interacts with it, and studies the physical, dialogue, and digital environments identified by the paper. We will continue to use Product X as a book-created illustration, while keeping the theoretical claims separate from the example.

\begin{learningbox}
By the end of this chapter, you should be able to explain why grounding is different from internal memory actions; distinguish an external environment, a grounding procedure, an action, and an observation; describe the three environment categories in CoALA; explain how different forms of environmental input can be made available to a language agent; and identify what the framework does not prescribe about tools, deployment, or safety.
\end{learningbox}

\section{From internal work to interaction with an external world}
\label{sec:external-world}

Return to Product X. The payment service is returning HTTP 500 errors. The agent has the incident report in working memory. It may also have retrieved an earlier incident and reasoned that a recent deployment is one possible explanation. Yet it still does not know whether the payment service is currently healthy, whether the deployment completed, or whether an operator observed a configuration change.

The missing information is not necessarily hidden in the agent's memory. It may exist in a monitoring system, a deployment dashboard, or the knowledge of a person who was present when the change occurred. The agent can reason about what it already knows, but reasoning cannot make a live external fact available merely by producing a plausible sentence. It must interact with the relevant environment.

This is the distinction between the internal actions studied in Chapter 4 and the external actions studied here. Retrieval reads from long-term memory into working memory. Reasoning processes information in working memory. Learning writes information to long-term memory. \gls{grounding} is different: it interacts with an \gls{environment} beyond the agent's internal memory.\index{grounding}\index{environment} The operation is external because the action is directed at something outside the agent's own information stores.\CoALACite{\S 4, Figure 5}

\subsection{What counts as an external environment?}

An \emph{environment} is the external setting with which an agent interacts. The word may suggest a physical location, but CoALA uses it more broadly. A robot's surroundings are an environment. A person or another agent communicating with the system can be part of an environment. A website, game, API, or code-execution setting can also be an environment.

The useful question is not simply whether the interaction uses a network connection or a tool-shaped function. Ask instead: is the agent interacting with a setting whose state or response is outside the agent's own internal memory? A database query can be external when it reads a changing external database. By contrast, retrieving a previously stored memory item is an internal action, even if the memory store happens to be implemented by a database server. The physical location of a component does not, by itself, determine its conceptual role.

This distinction will matter throughout the chapter. A tool wrapper may be the software through which an agent reaches an environment, but the wrapper and the environment are not the same thing. A function that sends a request is part of the agent's procedure. The service receiving the request is the external environment. The result of the interaction is information from that environment, not a memory item that the agent already possessed.

\section{Grounding: executing an action and processing feedback}
\label{sec:grounding-procedure}

CoALA describes grounding procedures as procedures that execute external actions and process environmental feedback into working memory as text. This definition joins two activities that are easy to overlook when we focus only on a tool call: something is done to an environment, and information about the environment's response is made available to the agent.\CoALACite{\S 4.2, ``Grounding actions,'' p. 10}

Consider a simple interaction. An agent sends a request to a digital service. The request is an external action. The service is the environment. The code that constructs and sends the request, handles the response, and makes the returned information available to the agent belongs to the grounding procedure. The response, if one is received, can be represented as an \gls{observation} in working memory.\index{observation}

These terms describe different roles:

\begin{itemize}
  \item \textbf{Environment:} the external setting with which the agent interacts.
  \item \textbf{Grounding procedure:} the procedure used to carry out an external interaction and process its feedback.
  \item \textbf{Action:} what the agent does to the environment, such as sending a query, communicating with a person, or issuing a physical command.
  \item \textbf{Observation:} information obtained from the environment and made available to the agent.
\end{itemize}

The distinction between an action and an observation is especially important. A request to inspect service health is not the same thing as the status returned by the service. The request is what the agent does. The returned status is information the agent receives. In a physical environment, the distinction may be between a robot's movement command and the sensory input that follows. In dialogue, it may be between a question sent to an engineer and the engineer's reply.

\BookDiagram[0.48\textheight]
  {Tex/diagrams/generated/part1/chapter05/grounding-interaction.png}
  {The grounding interaction: an external action is carried out in an environment, and any available environmental feedback is represented for the agent in working memory. This is a conceptual view, not a mandatory software pipeline.}
  {fig:grounding-interaction}

\subsection{Why CoALA uses a text-oriented abstraction}

The paper explains grounding as a way to simplify the agent's interaction with the outside world into something like a ``text game'': actions can be represented in language, and environmental feedback can be processed into textual observations. This makes the interaction legible to a language model that works primarily with text. It also provides a common conceptual interface across environments that are otherwise very different.\CoALACite{\S 4.2, p. 10}

The phrase does not mean that the external world is literally text, or that every sensor produces a textual result without additional processing. A camera captures an image, a microphone receives sound, and a service may return structured data. Those inputs may need a representation step before they can be used by a text-oriented reasoning process. In other systems, a multimodal model may receive some non-textual input directly. CoALA's point is the role of grounding in connecting an agent to its environment, not a claim that all environments share one identical conversion algorithm.

Nor does the definition promise that every action succeeds or produces a useful observation. A service may be unavailable. A person may not reply. A physical action may not produce the expected result. A request may return an error, a partial result, or no result at all. The procedure's responsibility is to execute the external interaction and process whatever feedback is available; the framework does not guarantee the outcome.

\begin{notebox}
Keep four things separate when reading a grounding example: the environment is where the interaction takes place; the procedure is how the agent interacts with it; the action is what the agent does; and the observation is what becomes known from the environment. An observation may be absent, delayed, incomplete, or unexpected. It is not the action itself.
\end{notebox}

\section{Physical environments: acting and perceiving in the physical world}
\label{sec:physical-environments}

Imagine an agent embodied in a robot. The robot cannot interact with a room by generating a sentence about the room. It needs a way to receive information from its surroundings and a way to affect those surroundings. The room, its objects, and the conditions in which the robot operates form a physical environment.

CoALA describes physical grounding as involving both perception and action. On the perceptual side, an agent may receive visual, audio, or tactile input. A system can process such inputs into textual observations, for example with a pretrained captioning model. On the action side, robotic planners can take language-based commands and translate them into actions that affect the physical world. The paper also notes that vision-language models can provide perceptual context to a language model.\CoALACite{\S 4.2, ``Physical environments,'' p. 10}

The key idea is the separation between what is sensed, how that information is represented, and what the agent does in response. Suppose a robot camera captures an image of a work area. The image is a perceptual input. A vision-language component may describe that a package is on a table and an obstacle is near the robot's path. That description can become a textual observation used by the language agent. A later command to move toward the package is an action. The resulting sensor information, if available, is another observation.

The representation is useful, but it is not identical to the physical scene. A caption may omit a detail, describe something ambiguously, or fail to capture a relevant condition. Likewise, a language-based command is not the movement itself. Some procedure must connect the command to the robot's physical capabilities. CoALA keeps these functions under the broad category of grounding without requiring the reader to adopt one robotics architecture.

\subsection{The two directions of physical interaction}

It helps to think of physical grounding as having two complementary directions. The environment can provide information to the agent through sensors. The agent can affect the environment through actuators or other physical mechanisms. These directions are connected, but they are not interchangeable.

For example, a robot may observe that an object is within reach. That observation does not move the object. A later action may attempt to grasp it. A subsequent observation may indicate whether the object was lifted, whether the grasp failed, or whether the situation changed in another way. The agent's next internal processing can then use the information that became available.

This chapter does not teach robot kinematics, sensor fusion, or actuator control. Those are substantial engineering topics in their own right. The conceptual point is simply that physical grounding includes receiving perceptual information and carrying out actions that affect a physical environment. The same external-action vocabulary can describe both.

\section{Dialogue with humans or other agents}
\label{sec:dialogue-environments}

Not every environment is a machine or a physical space. Another person can be part of the environment with which an agent interacts. A person may provide instructions, correct a misunderstanding, share information, or answer a question. An agent can also ask for help or clarification when it lacks information needed to continue.

CoALA treats dialogue with humans or other agents as one of the external grounding categories. The paper discusses accepting instructions and learning from people, as well as agents asking for help or clarification. It also notes work involving interaction among multiple agents, including social simulation, debate, and collaboration.\CoALACite{\S 4.2, ``Dialogue with humans or other agents,'' pp. 10--11}

Why call a message an external action rather than simply text generation? Because the relevant event is communication with another participant. The agent sends a message into a social or conversational environment and may receive a response from someone whose knowledge and behavior are not controlled by the agent. The output text is the content of the action, but the action is the act of communicating with another participant.

After the response arrives, the agent may reason about what it means. That later reasoning is an internal action. The external dialogue and the internal processing of the reply are different operations, even if both are represented as text. This is another example of why surface format alone does not determine an action's category.

\subsection{A clarification request in Product X}

\BookExampleLabel\ Suppose the Product X agent has observed the HTTP 500 errors and knows that the issue began around the time of a deployment. It does not know whether the deployment included a payment-service configuration change. The agent may send the on-call engineer a concise question asking whether such a change occurred.

The message is an external dialogue action. The engineer's reply, if one arrives, is an observation from the dialogue environment. The agent may then reason about the reply alongside the monitoring information already in working memory. The answer might clarify the timeline, add a new lead, or leave the uncertainty unresolved. The example does not assume that the engineer replies or that the reply establishes the cause.

\begin{tipbox}
A useful classification test is to ask whether the agent is communicating with someone outside itself or merely processing an existing message. Sending a clarification question is external dialogue. Interpreting the reply in working memory is internal reasoning. Saving a durable record of the interaction would be a separate learning action, as explained in Chapter 4.
\end{tipbox}

The chapter does not need to explain negotiation protocols, multi-agent coordination algorithms, or the design of conversational interfaces. Those subjects are broader than the CoALA grounding category being introduced here. The relevant distinction is that communication with a person or another agent is an external interaction that can provide instructions or feedback.

\section{Digital environments: websites, APIs, games, and code execution}
\label{sec:digital-environments}

A digital environment may look less like a ``world'' than a physical room does. It may be a website, an online game, an API, a database-backed application, or a code-execution setting. Yet it can still have state and behavior outside the agent's internal memory. A request can inspect that state, submit information, or cause a change. The interaction is external even though it occurs through software.

CoALA's digital category includes games, APIs, websites, and general code execution. The paper observes that digital grounding is often cheaper and faster than interaction with people or physical environments, making it a convenient testbed for language agents. It also identifies stateless digital APIs, such as search, calculator, or translator services, as special ``single-use'' digital environments.\CoALACite{\S 4.2, ``Digital environments,'' p. 11}

The qualification about stateless APIs matters. Some APIs provide a one-off response without maintaining a conversational or task state for the agent. These are the special single-use environments highlighted by the paper. It would be inaccurate to conclude that all digital environments are stateless. Websites, games, and stateful services can preserve information and change over time. A digital environment may have ongoing state even if a particular interface exposes it through a simple request-response operation.

\subsection{Why a tool-like interface can still be grounding}

Many agent systems describe external interactions as ``tool calls.'' That can be a convenient software term, but it should not obscure the conceptual role. A tool wrapper is an interface in the agent's code. The operation it performs may reach an external digital environment. For example, a search wrapper may send a query to a search service and return results. The wrapper is part of the grounding procedure; the search service and the information it returns are part of the external interaction.

By contrast, a function that retrieves an item from the agent's own long-term memory is an internal retrieval procedure. It may use similar technologies and even similar network protocols, but its conceptual role is different. It reads information from memory into working memory. An API call to a service outside the agent's internal memory interacts with a digital environment. The distinction depends on what the operation is doing in the architecture, not simply on whether the code uses HTTP, a database driver, or a tool interface.

\BookExampleLabel\ In Product X, querying a payment-service health endpoint is a digital grounding interaction because the agent is asking an external service for its current status. Reading a previously stored incident note is retrieval. The first may reveal a current service response; the second brings previously retained information into the current decision cycle. Both can inform the same investigation, but they are not the same action type.

\BookDiagram[0.52\textheight]
  {Tex/diagrams/generated/part1/chapter05/environment-categories.png}
  {The three external environment categories described by CoALA: physical environments, dialogue with humans or other agents, and digital environments. Examples are illustrative and the categories do not prescribe particular software stacks.}
  {fig:grounding-environment-categories}

\section{From different observations to working information}
\label{sec:grounding-observation-representation}

The three environments described so far produce very different kinds of information. A physical system may produce an image, sound, or tactile signal. A person may provide a sentence. A website may return a page, an API a structured response, and a code-execution environment a result or error message. The environment does not need to produce the same raw format in every case for the agent to use a common conceptual architecture.

CoALA emphasizes a textual representation of environmental feedback in working memory. In the examples discussed by the paper, perceptual data may be translated into text by a captioning or vision-language component. Digital results and dialogue responses may already be textual or may be converted into a form that the agent can process. The common point is that the grounding procedure makes environmental feedback available to the agent's internal process.\CoALACite{\S 4, Figure 4; \S 4.2}

The distinction between the external event and its representation is important. A text description of an image is not the image itself. A short status summary is not the full state of a service. A person's reply is not the agent's interpretation of that reply. Representations are useful because they let the agent work with information from outside, but they can be incomplete or lossy. The architecture should not confuse the description with the world that was described.

\BookDiagram[0.48\textheight]
  {Tex/diagrams/generated/part1/chapter05/observation-representation.png}
  {Different environmental inputs can be represented in forms available to working memory. The diagram is a teaching synthesis; it does not claim that physical, dialogue, and digital inputs use one identical conversion algorithm.}
  {fig:observation-representation}

Consider two Product X observations. A monitoring endpoint may return a structured status object such as a response code and a service-health field. An engineer may reply in natural language that a configuration change was made shortly before the errors began. These observations have different sources and formats. A grounding procedure can make both available in working memory in a form suitable for subsequent processing. The agent must still preserve the distinction between what the source reported and what it infers from that report.

\section{A Product X grounding walkthrough}
\label{sec:productx-grounding}

\BookExampleLabel\ We can now follow two external interactions in a deliberately limited scenario. The current working information says that the payment endpoint is returning HTTP 500 responses. The cause is not established.

First, the agent uses a digital grounding procedure to query a health or deployment-status endpoint. The action is the request. The external service is the environment. If the service returns a result, the grounding procedure processes it into an observation available in working memory. The result might confirm that the endpoint is responding with errors, or it might indicate a different condition. This chapter does not assume which result occurs.

Second, suppose the agent still lacks information about a recent operational change. It may communicate with the on-call engineer and ask whether a relevant deployment or configuration change occurred. That message is a dialogue action. If the engineer replies, the reply becomes another observation. The agent can now use both observations as input to later internal reasoning.

\BookDiagram[0.48\textheight]
  {Tex/diagrams/generated/part1/chapter05/productx-grounding.png}
  {A book-created Product X illustration showing a digital query and a separate dialogue request. Each interaction may provide an observation for working memory; neither result is assumed to be successful or conclusive.}
  {fig:productx-grounding}

The sequence is illustrative, not a workflow specified by CoALA. It does not claim that the agent must call a monitoring endpoint before asking a person, or that the two actions must always occur in this order. It also does not propose a diagnosis or remediation. The purpose is to make the grounding abstraction concrete: the agent acts outside its internal memory, and environmental feedback may then become available for further work.

\section{Grounding and the safety boundary of an action space}
\label{sec:grounding-safety}

An action being available to an agent does not make it harmless. An external action may change a digital system, affect a person through dialogue, or cause a physical consequence. CoALA explicitly notes that some grounding actions can cause external harm. Its examples include destructive commands in a shell, harmful speech in dialogue, and dangerous physical actions. The paper also notes that safety measures in existing systems are often task-specific, such as excluding particular operations, filtering dialogue, or limiting robots to controlled environments.\CoALACite{\S 6, ``Action space: thinking beyond external tools or actions.''}

This is a warning about the action space, not a complete safety architecture. CoALA does not specify a universal authorization model, approval workflow, audit-log format, rollback protocol, or production policy engine. Those are engineering concerns that will be addressed in the later Java case studies. Here the narrower lesson is that the external action space itself matters: what the agent can do influences the risks its behavior can create.

It is also useful to distinguish two questions. The first is conceptual: which operations belong to the agent's external action space? The second is operational: under what conditions should a particular operation be permitted, and what controls should surround it? CoALA helps organize the first question and highlights that the second cannot be ignored. It does not supply a complete answer to every production safety requirement.

\begin{warningbox}
Do not confuse an available grounding action with an authorized or safe action. The paper identifies the possibility of external harm and discusses task-specific safety measures; it does not prescribe a complete production authorization system. The concrete controls used in Part III will be identified as engineering decisions, not as mechanisms required by CoALA.
\end{warningbox}

\section{Common confusions and boundaries}
\label{sec:grounding-boundaries}

Several misunderstandings disappear once we keep the roles separate.

\begin{itemize}
  \item \textbf{Grounding is not retrieval.} Grounding interacts with an external environment. Retrieval reads from the agent's long-term memory into working memory.
  \item \textbf{An environment is not a tool wrapper.} A wrapper is part of the agent's procedure; the environment is the external setting reached through it.
  \item \textbf{An action is not an observation.} The action is what the agent does. The observation is information obtained from the environment, if feedback is available.
  \item \textbf{Dialogue is not always internal reasoning.} Sending a question to a person or another agent is external grounding. Interpreting the reply is internal reasoning.
  \item \textbf{Not every digital environment is stateless.} CoALA identifies stateless APIs as special single-use digital environments; games, websites, and other digital settings may have ongoing state.
  \item \textbf{Grounding does not decide what should be done.} It describes an external action and its feedback. The decision procedure that organizes planning and execution is treated in Chapters 6--8.
\end{itemize}

A reliable way to classify an operation is to ask what it interacts with and what role the result plays. If it reads an item from the agent's long-term memory, it is retrieval. If it processes working-memory contents to generate new active information, it is reasoning. If it writes information into long-term memory, it is learning. If it interacts with an external environment, it is grounding. A single larger task can contain several of these operations, but that does not make them interchangeable.

\section{Chapter recap}
\label{sec:grounding-recap}

\begin{chapterrecap}
CoALA distinguishes external grounding from internal memory actions. A grounding procedure executes an external action and processes environmental feedback so that it can become available in working memory as text. The external environment is the setting of the interaction; the grounding procedure is how the agent carries it out; the action is what the agent does; and an observation is information obtained from the environment. CoALA discusses physical environments, dialogue with humans or other agents, and digital environments. These categories share a conceptual grounding role, but the raw inputs and actions can differ. The framework highlights that external actions can cause harm, while leaving complete production safety controls to the system designer.
\end{chapterrecap}

\begin{center}
\begin{tabularx}{\linewidth}{@{}P{0.19\linewidth}P{0.34\linewidth}X@{}}
\toprule
\textbf{Environment} & \textbf{Examples in CoALA's discussion} & \textbf{What may be involved} \\
\midrule
Physical & Robots and embodied systems & Visual, audio, or tactile perception; physical actions through robotic procedures. \\
Dialogue & Humans and other agents & Instructions, questions, help, clarification, and other communication. \\
Digital & Games, websites, APIs, and code execution & Queries, website interaction, game actions, and execution in an external digital setting. \\
\bottomrule
\end{tabularx}
\end{center}

\begin{reviewquestions}{Review questions and short exercises}
\begin{enumerate}
  \item Explain why querying a live payment-service health endpoint is grounding, while retrieving a previously stored incident report is retrieval.
  \item In your own words, distinguish an environment from a grounding procedure.
  \item A robot camera captures an image, a vision-language component describes the scene, and a robot later moves toward an object. Identify the perceptual input, the representation, and the external action.
  \item An agent asks an engineer whether a deployment changed the payment configuration. Which part is the external dialogue action, and which part would be internal reasoning?
  \item Why is a stateless search API described as a special single-use digital environment? Why would it be inaccurate to call every digital environment stateless?
  \item A service returns an error rather than a health result. Has grounding still occurred? Explain the distinction between executing the action and receiving a useful observation.
  \item Give one example of an external action that could be harmful and explain why CoALA's discussion of action-space safety is not a complete production safety policy.
  \item \textbf{Classification exercise.} Classify each operation as grounding, retrieval, reasoning, or learning: (a) read a stored service manual into working memory; (b) query a live deployment endpoint; (c) interpret an engineer's reply; (d) preserve a completed incident record in episodic memory.
  \item Rewrite this sentence to distinguish action and observation: ``The agent's monitoring response checked the health endpoint.''
\end{enumerate}
\end{reviewquestions}

\section*{Source and example note}
The theoretical account is grounded primarily in \textsc{[CoALA]}, Section 4.2, with Section 4 and Figure 5 used to place grounding within the action space. Section 6 is used only for the paper's discussion of action-space safety. Product X and its incident details are book-created illustrations. The Mermaid figures are original teaching diagrams maintained as text sources and rendered to PNG during CI; they are not software architecture prescriptions. No Fan \& Lan mechanisms or Part III implementation decisions are used to reinterpret CoALA in this chapter.

\chapter{Planning: From a Problem to Candidate Actions}
\label{ch:planning-proposal}

\section*{Chapter at a glance}
\phantomsection
\addcontentsline{toc}{section}{Chapter at a glance}

An agent can have a current goal, useful information in working memory, access to long-term memory, and a set of possible actions. It still needs a way to move from the situation it faces to concrete possibilities for what it might do next. CoALA places that work in the \gls{planningstage} of its decision-making procedure. One of the planning sub-stages is called \gls{proposal}: it generates one or more \gls{candidateaction}.\index{planning stage}\index{proposal}\index{action candidate}

This chapter focuses on that proposal step. We will see how reasoning, sometimes supported by retrieval, can produce candidates; why a small action set may simply be enumerated; and how program structures or structured simulators can also generate possibilities. We will stop at the candidate set. Evaluating candidates, selecting one, and executing it are distinct responsibilities developed in Chapter 7, while the flexibility to interleave or iterate planning sub-stages is developed in Chapter 8.

\begin{learningbox}
By the end of this chapter, you should be able to explain where proposal belongs in CoALA's decision procedure; distinguish a candidate from an evaluated, selected, or executed action; describe how reasoning and optional retrieval can support candidate generation; recognize enumeration, code structures, and structured simulators as alternative proposal approaches; and explain what the proposal stage does not decide.
\end{learningbox}

\section{The problem before the first candidate}
\label{sec:problem-before-candidate}

Return to Product X. The payment endpoint is returning HTTP 500 errors. The incident goal is to restore normal service, but the cause is not yet known. The agent has a current observation and some service context. It may have retrieved an earlier incident or reasoned about possible explanations. What should it do next?

At this point, there may be several actions that are possible. It might inspect deployment metadata, query database connectivity, ask an operator about a recent configuration change, or examine another relevant service signal. Naming these possibilities is useful, but it is not yet a decision. The agent has not established which possibility is most valuable, which one should be chosen, or whether any proposed action should be carried out.

CoALA makes this distinction explicit by describing a planning stage and dividing it into sub-stages. The first of these is proposal. Its job is to generate one or more action candidates that can be considered by subsequent parts of the decision procedure.\CoALACite{\S 4.6, ``Planning'' and ``Proposal,'' pp. 12--13}

\section{Proposal's place in the decision cycle}
\label{sec:proposal-decision-cycle}

The decision-making procedure is the top-level or ``main'' program of a CoALA agent. It organizes repeated decision cycles. In the paper's account, the agent can use reasoning and retrieval during planning to develop alternatives. The decision procedure then moves through later planning work and executes the action that is selected. After execution, an observation may become available and the cycle can begin again.\CoALACite{\S 4.6, Figure 4B}

The planning stage is therefore not just a synonym for generating a sentence that names an action. It is a part of the larger program that helps the agent work out what it might do. Proposal is one sub-stage within planning. Evaluation and selection are other sub-stages, but they answer different questions. Execution is different again: it applies an action by invoking the relevant procedure in the agent's code.

\BookDiagram[0.46\textheight]
  {Tex/diagrams/generated/part1/chapter06/proposal-in-cycle.png}
  {Proposal within the broader decision cycle. The diagram locates proposal and marks evaluation, selection, and execution as later responsibilities; it is not a complete account of their mechanisms or a rigid one-pass ordering.}
  {fig:proposal-in-cycle}

\section{What is an action candidate?}
\label{sec:action-candidate}

Suppose the Product X agent produces the phrase ``inspect deployment metadata.'' Has the agent made its decision? Not yet. The phrase describes one action that may be possible. It is a candidate: a proposed possibility that can be considered by later parts of the decision procedure.

A candidate should be understood in relation to the action space introduced in Chapter 2. CoALA's action space includes internal actions and external grounding actions. For the proposal description in Section 4.6, the paper says that the usual approach is to use reasoning, and optionally retrieval, to sample one or more external grounding actions from the LLM. The same section explains the larger cycle as ultimately yielding an external grounding action or an internal learning action. This chapter follows the proposal paragraph's specific emphasis on external grounding candidates and does not generalize that paragraph into a complete account of every possible learning-action proposal mechanism.\CoALACite{\S 4.6, ``Proposal,'' p. 13}

It is helpful to distinguish four statuses:

\begin{itemize}
  \item A \textbf{candidate} is a possible action produced during proposal.
  \item An \textbf{evaluated candidate} is a candidate to which a value has been assigned during the later evaluation sub-stage, when multiple actions are proposed.
  \item A \textbf{selected action} is the candidate chosen by the selection sub-stage for what happens next.
  \item An \textbf{executed action} is an action that has actually been applied by invoking its procedure.
\end{itemize}

These distinctions describe different points in the decision process. A candidate can exist without receiving an evaluation value. An evaluated candidate is not automatically selected. A selected action has not necessarily been executed until the relevant procedure is invoked. Chapters 7 and 8 will examine the later decision stages; for now, keep the proposal boundary clear.

A candidate also differs from the procedure that implements it. ``Query deployment metadata'' is a description of a possible action. The agent's source code may contain a procedure that knows how to construct and send that query. The candidate is not itself the implementation. CoALA's architecture distinguishes the actions available to the agent from the code-based procedures through which those actions are carried out.

\subsection{A candidate must be understood within an action space}

A candidate is not an arbitrary sentence that happens to sound actionable. It is a possible action in the context of the agent's available action space. If the agent has no procedure for contacting a particular service and no available way to communicate with a person, simply generating the words ``contact the service'' or ``ask the engineer'' does not create that capability. The proposal description is about generating possibilities from the action space and procedures the agent has been given or can use, not about granting the agent a new ability merely by naming it.

This point follows from the architecture map established in Chapter 2. The action space identifies the kinds of internal and external actions available to the agent. The proposal stage works with that space. The concrete representation of a candidate can vary: it may be expressed in natural language, a structured action description, or another form used by the agent's code. CoALA does not prescribe one universal candidate schema. What matters for this chapter is the role the candidate plays: it is a possible action generated for later consideration.

\subsection{One candidate or several}

The paper explicitly allows proposal to generate one or more candidates. A narrow situation may lead to one plausible next action. A less constrained situation may produce several possibilities. Neither case changes the definition of proposal.

This is important because the number of candidates should not be confused with the sophistication of the overall decision procedure. A procedure that produces one candidate has still performed proposal. A procedure that produces several candidates has not necessarily performed evaluation or selection. The later stages may have more work to do when multiple alternatives are present, but that does not mean proposal itself has made a comparison.

\section{Using reasoning and retrieval to generate candidates}
\label{sec:reasoning-retrieval-proposal}

CoALA describes reasoning as the usual means of proposal. The LLM can process the current working information and produce one or more possible external grounding actions. Retrieval may also be used to bring relevant information from long-term memory into working memory, where it can provide additional context for reasoning. The paper's phrasing matters: reasoning and optionally retrieval. Retrieval can support proposal, but the text does not make it a mandatory precondition for every proposal.\CoALACite{\S 4.6, ``Proposal,'' p. 13; supporting action definitions in \S\S 4.3--4.4}

Consider the Product X situation. Working memory contains the HTTP 500 observation, the incident goal, and known service context. The agent can reason over that information to generate a possible next action. If earlier incident experience is relevant, retrieval may bring it into working memory before or as part of the proposal process. Reasoning can then use both the current observations and the retrieved context to generate candidates.

\BookDiagram[0.46\textheight]
  {Tex/diagrams/generated/part1/chapter06/reasoning-and-retrieval-proposal.png}
  {Reasoning can generate one or more action candidates from working information. Retrieval is shown as an optional source of additional context, not as a mandatory step in every proposal.}
  {fig:reasoning-retrieval-proposal}

The point is not that retrieval improves every proposal or that every action must be justified by a retrieved memory. Sometimes the current situation contains enough information to produce a possible action. Sometimes a prior episode or a stored fact can make a candidate more specific. The proposal description allows for this flexibility without defining one compulsory preparation sequence.

Likewise, one candidate may be enough in some cases, while other situations call for several possibilities. The number of candidates depends on the agent's action domain and proposal procedure. A single proposal is still a proposal; multiple proposals do not, by themselves, tell us which one should be chosen.

\subsection{Following the information path}

It can help to follow the information rather than memorize the labels. The current problem is already represented in working memory. If the agent uses retrieval, information is read from long-term memory into working memory. That retrieval operation does not itself generate the action candidate. It changes the information available to the next reasoning operation. The LLM can then process the active information and generate a candidate description. The candidate becomes part of the proposal result.

For example, suppose an earlier Product X incident involved a database connection-pool limit. Retrieval may make that episode available in working memory. Reasoning may then produce ``inspect current connection-pool configuration'' as one possible action. The earlier episode has informed the proposal, but it has not proved that the current incident has the same cause. A retrieved memory supplies context; it does not turn a possibility into a conclusion.

The reverse is also possible. If the current working information already identifies a useful next diagnostic action, the agent may propose it without retrieving an earlier episode. This is why the source wording ``reasoning (and optionally retrieval)'' should not be rewritten as a mandatory two-step recipe. The agent's procedure can use the information it needs for the task at hand.

The proposal mechanism is also not identical to the language model's raw output. A model may generate text that describes an action, but the surrounding agent procedure can parse or organize that output as an action candidate. This is consistent with the architecture's distinction between LLM reasoning and the explicit code that implements agent procedures. The chapter does not prescribe a particular parser or representation; it clarifies the function of the information produced.

\subsection{What the Product X candidates do and do not say}

\BookExampleLabel\ In Product X, the proposal procedure may generate candidates from several categories: inspecting deployment metadata, querying database connectivity, or asking the on-call engineer about a possible configuration change. The first two are digital grounding candidates; the third is a dialogue grounding candidate. The full unranked set is traced in Section~\ref{sec:productx-proposal}.

\section{When the action set is small: proposing all available actions}
\label{sec:enumerating-action-set}

It is tempting to assume that proposal always means asking the LLM to invent actions. CoALA explicitly gives a different possibility. In a simple domain with a limited set of available actions, the proposal stage may simply include all of them. The paper uses SayCan as an example: it discusses a fixed set of grounding skills in a robotic kitchen setting and notes that, for a simple domain with limited actions, proposal might include all actions.\CoALACite{\S 4.6, ``Proposal,'' p. 13; \S 5, SayCan case study}

Imagine a small environment in which only three actions are available: inspect an inventory item, move to a marked location, or report that the requested item is unavailable. If the domain is deliberately limited, the proposal procedure can enumerate those actions rather than asking the LLM to generate a new list. The proposal result is the set of available possibilities.

Enumeration is not evaluation. Listing all available actions does not assign a value to each one, compare their outcomes, or decide which action should be executed. It simply defines which possibilities are present for later consideration. This distinction is easy to lose when the action set is short: because all actions are visible, it may feel as though the decision has already been made. It has not.

Enumeration also depends on the action domain being sufficiently defined. The example does not imply that every agent has a small, fixed, or exhaustively enumerable set of actions. CoALA presents this as one possible proposal approach for simple domains, not a universal rule.


A small action set can make the boundary particularly easy to see. Imagine a service agent whose external action space for one narrow task contains exactly three diagnostic queries. A proposal procedure may put all three into the candidate set. At that moment the agent has not learned which query is most informative, which one is safest, or which one should be run first. The enumeration has answered only the availability question: what are the possible actions in this limited domain?

The scope of the action set is therefore important. If the available actions change, the enumerated candidate set changes. If the domain is not well-defined or the available possibilities are not exhaustively known, simple enumeration may not describe the proposal procedure. CoALA's example is intentionally modest: it shows that candidate generation can be straightforward when the domain is constrained. It does not claim that enumeration is appropriate for every agent.

\section{Program structures and structured simulators as proposal mechanisms}
\label{sec:program-simulator-proposal}

CoALA also describes more structured ways to generate candidates. The proposal process may use code structures such as \texttt{if--else} or \texttt{while--if} to produce possible actions. The paper further notes that agents deployed in well-defined domains may use structured simulators to generate plausible rollouts. These are examples of proposal mechanisms that are not limited to a fresh LLM sample.\CoALACite{\S 4.6, ``Proposal,'' p. 13}

A branch in program code can encode conditions under which different actions are proposed. For example, a domain-specific procedure may include one candidate when a relevant signal is present and another when it is absent. A loop structure may allow a proposal procedure to generate actions repeatedly while some condition remains true. The important point for this chapter is modest: code structure can participate in generating candidates. We are not defining a universal program structure, nor claiming that all agents should use the same branching or looping pattern.

A structured simulator offers another possibility in a domain where a model of the environment is available. It can generate plausible rollouts: possible sequences of states or events that may follow from a proposed course of action. In the proposal context, the simulator can help produce possibilities for consideration. The detailed use of simulations to evaluate outcomes or compare plans is outside this chapter. We will not turn this brief reference into a tutorial on search algorithms, planning languages, or simulator construction.


These examples also show why proposal need not be implemented entirely inside a model prompt. A program branch can narrow or expand the set of possible actions according to a domain condition. A loop can repeat a generation step under a condition. A simulator can produce possible sequences in a domain for which structured simulation is available. In each case, code or a simulator contributes to producing possibilities. The conceptual role is still proposal, even though the mechanisms differ.

The mention of structured simulators should not be mistaken for a claim that all candidate generation requires a complete world model. The paper identifies simulators as a possibility for well-defined domains. That is a qualification, not a universal requirement. Nor does a plausible rollout automatically tell us how valuable it is. A simulator-generated sequence may be included for later consideration, but the process of assigning values to alternatives and choosing what to do belongs to later parts of the decision procedure.

\BookDiagram[0.48\textheight]
  {Tex/diagrams/generated/part1/chapter06/candidate-generation-methods.png}
  {Four proposal approaches discussed in CoALA: LLM-based sampling, enumeration of a limited action set, code-structured generation, and structured-simulator rollouts. The figure compares ways candidates may be generated; it does not rank them.}
  {fig:candidate-generation-methods}

The four approaches answer the same narrow question: how might a set of possible actions be produced? They do not answer what value a candidate has, which candidate should be chosen, or whether the action should be executed. Keeping the question narrow prevents proposal from absorbing the work of the later planning stages.

\section{A Product X proposal walkthrough}
\label{sec:productx-proposal}

\BookExampleLabel\ We can now follow the path from a current problem to an unranked set of possibilities. The working information contains the HTTP 500 observation, the active goal of understanding the service failure, and the service context already available to the agent. The action space includes digital grounding procedures and dialogue interactions that the agent can use.

The proposal procedure reasons over the available working information. It may optionally retrieve a relevant prior incident or service fact if that information could help form a candidate. From this context, the agent might generate candidates such as inspecting deployment metadata, querying database connectivity, or asking an operator for clarification. The list illustrates the generation of possibilities; it does not establish which possibility is preferable or likely to solve the incident.

\BookDiagram[0.48\textheight]
  {Tex/diagrams/generated/part1/chapter06/productx-candidate-set.png}
  {A book-created Product X example showing current working information and available action types feeding proposal. The resulting candidates are explicitly unranked and unselected.}
  {fig:productx-candidate-set}

The chapter ends at this point. We have a candidate set. We have not assigned values to the candidates, compared their likely outcomes, selected one, or executed it. That stopping point is not an omission; it is the boundary that allows the next two chapters to explain the remaining parts without collapsing them into one vague idea of planning.


Notice how little the candidate set tells us about the underlying incident. The inclusion of deployment metadata does not prove that the deployment caused the errors. The inclusion of a database-connectivity query does not establish that the database is involved. The option to ask the operator does not imply that a change occurred. Proposal converts the current situation and available action space into possibilities. It does not convert those possibilities into facts about the world.

If the action space were narrower, the candidate set might look different. If it offered only a health query and an operator question, the proposal procedure could generate only those possibilities. If a relevant prior service fact were retrieved, reasoning might produce a more focused candidate. These variations alter the possibilities that are proposed; they do not change the boundary between proposal and later evaluation or selection.

\begin{warningbox}
A candidate is a possibility for later consideration, not an evaluated recommendation, a selected action, an approval, or an execution. Proposal does not establish that an action is good, safe, or likely to resolve the incident; those judgments belong to later decision stages and application-specific controls.
\end{warningbox}

\section{Common confusions and boundaries}
\label{sec:proposal-boundaries}

\begin{itemize}
  \item \textbf{Proposal is one sub-stage, not the whole planning stage.} Evaluation, selection, and execution remain separate responsibilities.
  \item \textbf{Retrieval is optional support.} CoALA describes reasoning and optionally retrieval for the usual LLM-based proposal approach.
  \item \textbf{Proposal is broader than open-ended LLM generation.} A limited action set may be enumerated; code structures and structured simulators may also generate possibilities.
  \item \textbf{No universal schedule is prescribed.} The flexibility to interleave or iterate planning sub-stages is treated in Chapter 8.
\end{itemize}

SayCan is mentioned only as CoALA's example of enumerating a limited action set, and structured simulators are retained at the level of the paper's proposal discussion rather than expanded into additional theories or algorithms.

\section{Chapter recap}
\label{sec:proposal-recap}

\begin{chapterrecap}
Proposal is the part of CoALA's planning stage that generates one or more candidate actions. A candidate is a possibility for later consideration, not an evaluated, selected, or executed action. CoALA describes reasoning as the usual way to sample external grounding candidates from the LLM, with retrieval as optional support. In simple domains with limited actions, the proposal stage may include all available actions. More structured agents may use program branches or loops, and well-defined domains may use structured simulators to generate plausible rollouts. The chapter stops once candidates exist; evaluation and selection belong to Chapter 7, and interleaving and iteration belong to Chapter 8.
\end{chapterrecap}

\begin{center}
\begin{tabularx}{\linewidth}{@{}P{0.24\linewidth}P{0.35\linewidth}X@{}}
\toprule
\textbf{Proposal approach} & \textbf{How candidates arise} & \textbf{What this approach does not decide} \\
\midrule
LLM-based sampling & Reasoning over working information, with optional retrieved context, produces one or more action descriptions. & Whether a candidate is valuable, selected, safe, or executed. \\
Limited-set enumeration & The proposal procedure includes all actions in a small, defined action set. & Which of the available actions should be preferred. \\
Program-structured generation & Branches or loops in agent code produce candidates under domain-specific conditions. & Whether the generated candidates should be executed. \\
Structured-simulator rollouts & A simulator generates plausible sequences or possibilities in a well-defined domain. & The detailed evaluation or ranking of those possibilities. \\
\bottomrule
\end{tabularx}
\end{center}

\begin{reviewquestions}{Review questions and short exercises}
\begin{enumerate}
  \item In CoALA, where does proposal sit in the decision procedure, and what does it produce?
  \item Explain the difference between a candidate action and the procedure that implements it.
  \item The Product X agent generates ``inspect deployment metadata'' and ``query database connectivity.'' Have these actions been evaluated or selected? Explain.
  \item Why does CoALA describe retrieval as optional support for the usual LLM-based proposal approach rather than as a mandatory precondition?
  \item When might the proposal stage simply include every available action? Why is that still different from evaluation?
  \item Give an example of how an \texttt{if--else} structure could contribute to candidate generation without deciding which candidate should be executed.
  \item What does a structured simulator contribute to proposal in a well-defined domain? What details are outside this chapter's scope?
  \item Compare LLM-based sampling and limited-set enumeration. What assumptions differ between the two approaches?
  \item \textbf{Classification exercise.} For each statement, identify whether it describes proposal, evaluation, selection, or execution: (a) ``inspect deployment metadata'' is added to a set of possibilities; (b) a value is assigned to each of three candidates; (c) one candidate is chosen for the next step; (d) the relevant procedure is invoked. Explain why only (a) belongs to this chapter.
  \item Rewrite the statement ``the agent selected the database check'' so that it describes only a proposal-stage event. A suitable answer should say that the proposal included the database check as one candidate.
\end{enumerate}
\end{reviewquestions}

\section*{Source and example note}
The theoretical account is grounded primarily in \textsc{[CoALA]}, Section 4.6, especially the description of the planning stage and the proposal sub-stage, with Figure 4B used only to locate proposal within the decision cycle. Sections 4.3 and 4.4 are recalled to explain retrieval and reasoning as supporting internal actions; Section 4.2 is recalled only to orient the reader to grounding actions as candidate types. SayCan is named only as the limited-action-set example discussed by CoALA. Product X remains a book-created teaching example. The Mermaid figures are original teaching diagrams maintained as text sources and rendered to PNG during CI. No Fan \& Lan mechanisms are used to reinterpret CoALA in this chapter.

\chapter{Evaluation and Selection}
\label{ch:evaluation-selection}

\section*{Chapter at a glance}
\phantomsection
\addcontentsline{toc}{section}{Chapter at a glance}

Chapter 6 ended with a set of possible actions. The agent had not yet decided which possibility should be carried forward. A candidate is only a possibility: it is not a finding about the world, a preferred action, or an action that has already been carried out. CoALA separates two further responsibilities. {\gls{evaluation}} assigns a value to proposed actions when there is more than one candidate. {\gls{selection}} uses the candidates and their values to decide what happens next. It may choose an action for execution, or reject the proposed set and return to proposal. These responsibilities are closely related, but they are not the same operation.

This chapter follows the distinction through a small Product X example. The examples of candidates, criteria, and numerical values are created for teaching; they are not incident findings or prescribed scoring rules from the paper.

\begin{learningbox}
By the end of this chapter, you should be able to explain why a candidate set is not yet a decision; describe evaluation as assigning values to proposed actions when multiple actions are considered; identify the value sources named by CoALA; distinguish evaluation from selection; explain the selection methods named in the paper; and locate rejection and return-to-proposal at the selection sub-stage.
\end{learningbox}

\section{From candidate actions to a decision}
\label{sec:evaluation-from-candidates}

Return to the Product X payment-service incident. The service is returning HTTP 500 errors. The proposal stage has produced three possible next actions:

\begin{enumerate}
  \item inspect the deployment metadata;
  \item query the database connectivity status;
  \item ask the on-call engineer whether a relevant configuration change occurred.
\end{enumerate}

At this point, the agent has a candidate set. The three candidates describe possible interactions with the environment. They do not establish that the deployment caused the errors, that the database is involved, or that a configuration change took place. Nor does their order mean that the first candidate is preferred. Proposal has generated possibilities; it has not yet assigned values or made a selection.

A reader might ask why the agent cannot simply take the first item. It could, if its decision procedure were designed to use a fixed ordering. But that would be one possible selection rule, not an evaluation of the alternatives. CoALA gives names to two different jobs that may be involved when several candidates are available. The first job assigns values. The second job uses the candidate set and those values to decide what happens next.

The distinction matters because a system can evaluate candidates and still need a separate method to decide what to do with the results. A value does not select itself. Equally, a selection method cannot compare values that have not been made available to it, unless the procedure uses some other basis for selection. CoALA's vocabulary helps us describe those responsibilities without treating them as one vague act of ``choosing.''

\CoALACite{\S 4.6, ``Evaluation'' and ``Selection,'' pp. 12--13}

\section{Evaluation: assigning values to proposed actions}\index{evaluation}
\label{sec:evaluation-values}

CoALA describes an evaluation sub-stage that assigns a value to each proposed action when multiple actions are proposed. The wording is conditional: evaluation is described for the case in which the proposal stage has produced more than one action. The paper does not require every agent to use this sub-stage in the same way, nor does it prescribe one universal evaluation formula.\CoALACite{\S 4.6, ``Evaluation,'' p. 12}

The word \emph{value} can sound abstract. In this chapter, an \gls{evaluationvalue} is the output assigned to a proposed action by an evaluation procedure. For now, think of it as information produced by an evaluation procedure to represent how a candidate fares under that procedure. A value may be numerical, but its meaning depends on how the particular agent defines and uses it. A value is not automatically a probability, a calibrated measure of confidence, a safety guarantee, or a prediction that the action will succeed.

Notice what is being evaluated. The object is a proposed action. If the agent is considering an inspection request, a database query, and a question to a human operator, the evaluation procedure may assign a value to each of those candidates. It is not assigning a universal quality score to the agent, and it is not declaring the current incident solved or understood.

Evaluation also differs from execution. Assigning a value to a proposed action does not perform that action. If a procedure estimates that inspecting deployment metadata is useful, the metadata has not thereby been inspected. The action remains a candidate until the selection and execution responsibilities are addressed.

\begin{notebox}
CoALA names the role of evaluation and gives examples of possible value sources, but it does not define a universal unit, range, normalization method, calibration guarantee, or domain-independent formula. In particular, the largest value in an illustrative example does not mean that an action is safest, most accurate, or certain to solve the incident; that interpretation depends on the invented criterion and values used in the example.
\end{notebox}

\section{Where evaluation values can come from}\index{evaluation value}\index{heuristic rules}\index{learned values}
\label{sec:evaluation-value-sources}

The CoALA paper names several possible sources for values: heuristic rules, LLM or perplexity values, learned values, LLM reasoning, and combinations of these. They are alternatives that illustrate the range of possible evaluation procedures. The paper does not require every agent to use all of them, and it does not specify a single weighting scheme that combines them.\CoALACite{\S 4.6, ``Evaluation,'' p. 12}

\subsection{Heuristic rules}

A heuristic is a rule or estimate that an agent designer specifies to help assess a candidate. It may encode a simple preference or a domain-informed judgment. For example, a book-created Product X procedure might assign a higher illustrative value to an inspection that is expected to clarify a recently changed subsystem than to an action that provides little information about the current question. The example explains what a heuristic could do; it is not a rule supplied by CoALA, and it does not establish that any subsystem actually changed.

A heuristic can be inexpensive and easy to understand. Its limitation is equally important: the rule reflects what its designer anticipated. A rule that is useful in one situation may be a poor guide in another. The CoALA paper names heuristic rules as a possible value source; it does not claim that heuristics are always reliable or that one particular heuristic should be used.

\subsection{LLM or perplexity values}

The paper also names LLM or perplexity values. The key idea is that a language-model-related measure can contribute to the value assigned to a candidate. This category should be kept narrow. Perplexity is a measure associated with how a language model assigns likelihood to text; it is not, by itself, a direct measurement of whether an external action is true, safe, useful, authorized, or likely to resolve a real-world incident.

For example, a model-derived value could help an evaluation procedure compare candidate descriptions under a particular representation. That does not mean the model has observed the outside environment or verified the consequences of the action. The value remains an output of an evaluation method, not environmental feedback.

\subsection{Learned values}

A value may also come from a learned estimator or value function. In this case, the procedure's ability to assign values is itself shaped by learning from data or experience. The important point for this chapter is not how to train such a function. It is that the evaluation source need not be a hand-written rule or a direct reasoning response from the LLM; it may be learned.

The phrase \emph{learned value} should not be read as a guarantee that the estimate is correct. A learned estimator can reflect the data and objective used to train it, and it may be uncertain when it encounters situations unlike those it has seen. CoALA lists learned values as one possible source without prescribing a training procedure or an evaluation guarantee.

\subsection{LLM reasoning}

The paper identifies LLM reasoning as another possible source. Here, the model can reason about a candidate and its possible consequences. CoALA gives the example of internally simulating possible grounding feedback as part of evaluating actions. This gives the agent a way to consider what might happen before it performs an external interaction.\CoALACite{\S 4.6, ``Evaluation,'' p. 12}

The word \emph{internally} is essential. A simulated consequence is generated as part of the agent's planning activity. It is not an observation returned by a website, API, person, robot sensor, or other environment. The agent may imagine that a query will reveal a connection error, but that imagined result does not establish that a connection error exists. Actual environmental feedback belongs to grounding after an external action is executed.

\begin{warningbox}
Do not turn a simulated outcome into a reported observation. A model can reason about what might happen if an action were carried out. Only an interaction with the relevant environment can provide feedback from that environment. These are different sources of information.
\end{warningbox}

\subsection{Combinations of value sources}

CoALA allows evaluation to use some combination of the sources it names. A procedure could, for example, use a hand-specified rule together with a learned estimate or LLM reasoning. The paper leaves room for different designs rather than prescribing one fixed combination.

A combination does not automatically make an evaluation better. Its behavior depends on what each source represents and how the procedure brings the results together. A chapter example may illustrate a combined procedure, but unless the combination is explicitly specified, it should not invent weights, thresholds, or a formula and then attribute that mechanism to CoALA.

\BookDiagram[0.58\textheight]
  {Tex/diagrams/generated/part1/chapter07/evaluation-value-sources.png}
  {Possible sources of values for proposed actions in CoALA. The branches are alternatives that an evaluation procedure may use; the figure does not imply that every source is mandatory or that all sources share one scale.}
  {fig:evaluation-value-sources}

The figure summarizes categories, not a required implementation. A particular agent might use one source, several sources, or a different value-estimation procedure. The common role is that the procedure assigns values to proposed actions when multiple candidates are being considered.

\section{Evaluation is not selection}
\label{sec:evaluation-not-selection}

This is the most important distinction in the chapter. Evaluation and selection are adjacent responsibilities, but one cannot be substituted for the other.

Evaluation produces values for candidate actions. Selection uses the actions and their values to decide what happens next. Suppose an evaluation procedure assigns values to three Product X candidates. Those values tell the decision procedure something about how the candidates fare under the chosen evaluation method. They do not, by themselves, tell us whether the procedure will choose the largest value, sample from a distribution, use a vote, or reject the proposed set.

This distinction becomes clearer if we hold the candidate set and its values constant. Imagine that two selection procedures receive exactly the same three candidate/value pairs. One may choose the candidate with the largest value. Another may use a value-sensitive distribution. A third may use a voting method. The selection outcomes need not be identical. The evaluation results are the same; the rule that acts on them differs.

\BookDiagram[0.48\textheight]
  {Tex/diagrams/generated/part1/chapter07/evaluation-versus-selection.png}
  {Evaluation produces candidate/value pairs; selection uses the candidates and values to decide what happens next. The boundary between the two roles is conceptual, not a requirement for separate software services.}
  {fig:evaluation-versus-selection}

The diagram also helps prevent a second confusion. Evaluation does not reject candidates in the CoALA description. The paper places the explicit option to reject the proposed actions and return to proposal under selection. Evaluation supplies values; selection makes the control decision.\CoALACite{\S 4.6, pp. 12--13}

\section{Selection: deciding what happens to the candidate set}\index{selection}
\label{sec:selection}

Given a set of proposed actions and their values, CoALA describes selection as the step that either selects one action to execute or rejects the proposed actions and loops back to proposal. Selection is therefore a control responsibility. It determines whether a candidate is carried forward to execution or whether the current set should not be used and the proposal sub-stage should be revisited.\CoALACite{\S 4.6, ``Selection,'' p. 13}

Selection is not execution. Selecting an action means that the decision procedure has identified which action is to be carried forward. Execution applies the action by invoking the relevant procedure from the agent's source code. If the selected action is an API call, the selection step does not itself send the request. If the selected action is a learning action, selection does not itself write the information to long-term memory. Those applications belong to the execution stage, which Chapter 8 examines as part of the complete decision process.

Selection is also not the same as human authorization, access control, or a production safety gate. A particular application may require human approval before a high-impact action is executed. That can be an important engineering requirement, but it is not what the CoALA passage means by the selection sub-stage. In this chapter, selection names the procedure's decision over the proposed candidate set and its values. It should not be used as a substitute for every other form of permission or oversight.

\section{Selection methods named by CoALA}\index{argmax}\index{softmax}\index{majority vote}
\label{sec:selection-methods}

The paper names argmax, softmax, majority vote, and other alternatives as possible selection methods, depending on the form of the action values. These examples show that selection is not tied to one universal rule. They should not be presented as mandatory stages that every agent must run.\CoALACite{\S 4.6, ``Selection,'' p. 13}

\subsection{Argmax}

\gls{argmax} selects the candidate associated with the largest represented value. If the values are comparable under the procedure's intended interpretation, the rule returns a candidate whose value is maximal among the alternatives. This is easy to explain because it makes the relationship between evaluation and selection visible: evaluation supplies the values; argmax uses them to identify a maximum.

The qualification about comparability matters. If two values are produced on incompatible scales, simply comparing their numerical sizes may not be meaningful. CoALA names argmax as an option; it does not claim that every evaluation method produces comparable values or that argmax is appropriate for every action space.

\subsection{Softmax}

\gls{softmax} is another value-dependent selection method named in the paper. At an intuitive level, instead of always treating the largest value as the only possible outcome, a softmax-based rule can use the values to determine a distribution over candidates. A higher value can receive greater selection weight, while other candidates may remain possible.

For a finite candidate set \(\mathcal C\), one standard mathematical illustration is
\begin{equation}
  P(a_i\mid\mathcal C)=
  \frac{\exp(v_i/\tau)}{\sum_{a_j\in\mathcal C}\exp(v_j/\tau)},
  \qquad \tau>0,
  \label{eq:ch7-softmax}
\end{equation}
where \(v_i\) is the value assigned to candidate \(a_i\), and \(\tau\) is a temperature parameter. A lower temperature concentrates more weight on higher-valued candidates; a higher temperature produces a flatter distribution. This equation makes the weighting behavior explicit, but it is an explanatory mathematical form used by this book, not a claim that CoALA mandates this exact parameterization.

The CoALA discussion does not prescribe a particular temperature, calibration method, or probability interpretation for every use of softmax. The chapter therefore uses the term to identify a family of value-sensitive selection behavior, not to promise that the resulting probabilities are calibrated estimates of real-world success.

\subsection{Majority vote}

The paper also names \gls{majorityvote} as an alternative. A voting procedure can determine a selection by aggregating votes over the candidate options. The outline of the method is enough for this chapter: selection can rely on a vote-based rule rather than only on a single maximum value.

CoALA does not specify a universal voting protocol in this passage. It does not say who or what must cast the votes, how many voters are required, how ties must be resolved, or what quorum should be used. Those details would belong to a particular implementation or a separate account of a voting method, not to the conceptual point being made here.

\subsection{Other alternatives}

The paper leaves room for another method suited to the form of the values. This flexibility is important. The named methods are examples, not a closed list of all possible selection procedures. The essential distinction remains the same: selection uses the proposed actions and their values to determine what happens next.

\BookDiagram[0.52\textheight]
  {Tex/diagrams/generated/part1/chapter07/selection-methods.png}
  {Different selection methods can operate on a proposed set and its values. The branches are alternatives, not successive steps in one required procedure.}
  {fig:selection-methods}

\section{A Product X example: from values to a selection decision}\index{Product X!evaluation and selection}
\label{sec:productx-evaluation-selection}

\BookExampleLabel\ The proposal stage has generated three candidates: inspect deployment metadata, query database connectivity, and ask the on-call engineer whether a relevant configuration change occurred. No candidate has been executed, and none of the candidates proves anything about the cause of the HTTP 500 errors.

For this walkthrough, imagine an evaluation procedure that uses one explicitly invented criterion: how useful the procedure estimates each candidate will be for reducing the agent's current uncertainty. To make the mechanics visible, suppose it assigns the following example values:

\begin{center}
\begin{tabularx}{\linewidth}{@{}P{0.46\linewidth}P{0.19\linewidth}X@{}}
\toprule
\textbf{Candidate} & \textbf{Illustrative value} & \textbf{What the value means here} \\
\midrule
Inspect deployment metadata & 0.80 & A high value under this invented uncertainty-reduction criterion. \\
Query database connectivity & 0.60 & A lower illustrative value under the same criterion. \\
Ask the on-call engineer & 0.45 & A still lower illustrative value under the same criterion. \\
\bottomrule
\end{tabularx}
\end{center}

These values are not from the CoALA paper, are not estimates of real incident outcomes, and are not probabilities that an action will succeed. They are numbers chosen for one book-created example so the reader can see how a selection method might consume values. A different criterion or evaluation procedure could assign different values.

To illustrate softmax without changing the meaning of those values, choose \(\tau=0.2\) for this example. Applying Equation~\ref{eq:ch7-softmax} to the values \(0.80\), \(0.60\), and \(0.45\) gives approximate selection weights of \(0.6486\), \(0.2386\), and \(0.1127\), respectively, for inspecting deployment metadata, querying database connectivity, and asking the on-call engineer. These are normalized selection weights under the chosen temperature. They are not calibrated probabilities that the actions will succeed, nor estimates of the incident's cause. The temperature is an illustrative parameter selected for the walkthrough, not a CoALA requirement.

If the selection method is argmax, it carries forward the candidate with the largest represented value: inspecting deployment metadata. This does not prove that the deployment caused the issue. It means only that, under the invented evaluation rule and the illustrated values, argmax returns that candidate. Another selection method could behave differently, and a procedure may instead reject the current set and return to proposal.

\BookDiagram[0.58\textheight]
  {Tex/diagrams/generated/part1/chapter07/productx-evaluation-selection.png}
  {Book-created Product X illustration: three candidates receive invented values under a stated example criterion, then a selection method acts on the candidate/value set. The numbers are not paper results or probabilities.}
  {fig:productx-evaluation-selection}

The example demonstrates the division of responsibility. Evaluation supplies values under a chosen method. Selection determines what to do with the proposed set. Execution has not yet occurred, and the incident's root cause remains unknown.

\section{Common confusions and chapter boundaries}
\label{sec:evaluation-boundaries}

Two boundaries deserve emphasis beyond the distinctions developed above:

\begin{itemize}
  \item \textbf{Simulation is not observation.} An internally imagined consequence is not feedback from the external environment; the grounding interaction must occur before such feedback is available.
  \item \textbf{Selection is not a complete safety architecture.} Human authorization, policy checks, access control, and operational safeguards may be required by an application, but they are not defined by CoALA's selection passage.
\end{itemize}

\section{Chapter recap}
\label{sec:evaluation-recap}

\begin{chapterrecap}
CoALA separates evaluation from selection. When multiple actions are proposed, evaluation assigns a value to each candidate. The paper names heuristic rules, LLM or perplexity values, learned values, LLM reasoning, and combinations as possible value sources. A value is an output of an evaluation procedure; the paper does not prescribe a universal scale or meaning. Selection uses the proposed actions and their values to select an action for execution or reject the proposed set and return to proposal. Argmax, softmax, majority vote, and other alternatives are named as possible selection methods. Selection is not execution, and neither value assignment nor selection should be confused with human authorization or an operational safety gate.
\end{chapterrecap}

\begin{reviewquestions}{Review questions and short exercises}
\begin{enumerate}
  \item What is the difference between a candidate action and an evaluated candidate action?
  \item Under what condition does CoALA describe evaluation as assigning values to proposed actions?
  \item Name the possible value sources listed in Section 4.6 and explain one in your own words.
  \item Why is an internally simulated consequence not an observation from the environment?
  \item Explain the difference between evaluation and selection without using the words ``score'' or ``choose.''
  \item What does argmax do, and what must be true of the represented values for the comparison to be meaningful?
  \item How does a softmax-based selection differ in spirit from always choosing only the largest value? What does the paper leave unspecified?
  \item What does it mean for selection to reject the proposed set and return to proposal?
  \item Classify these events: generating a database check, assigning values to candidates, selecting one candidate, and invoking its grounding procedure.
  \item Why is CoALA selection not the same thing as production authorization or human approval?
\end{enumerate}
\end{reviewquestions}

\section*{Source and example note}
The theoretical account is grounded in \textsc{[CoALA]}, Section 4.6, especially the Evaluation and Selection sub-stages, pp. 12--13. Product X is a book-created teaching example. All criteria and numerical values in its walkthrough are invented for explanation and are not results from the paper. No Fan \& Lan mechanisms are used to reinterpret CoALA.

\chapter{Interleaving, Iteration and the Complete CoALA Decision Process}
\begin{notebox}
Chapter placeholder. The detailed outline must be reviewed and approved before this chapter is written.
\end{notebox}


\part{Fan \& Lan: From Cognitive Architectures to Mechanism-Level Understanding}
\label{part:fan-lan}
\chapter{Why Feature-Level Descriptions Are Not Enough}
\begin{notebox}
Chapter placeholder. No prose has been drafted. This chapter will be outlined and approved before writing.
\end{notebox}

\chapter{The Seven Fields of a Mechanism}
\label{ch:seven-fields-mechanism}

\section*{Chapter at a glance}
\phantomsection
\addcontentsline{toc}{section}{Chapter at a glance}

Chapter 9 showed why a feature label cannot explain how an agent operates. We now need a consistent way to describe the operation itself. Fan and Lan represent each documented mechanism through seven fields: operative state substrate, control locus, transition trigger, persistence boundary, failure semantics, learning operator, and resource or governance policy. Together they form the tuple \(K = \langle S,C,T,P,F,L,R\rangle\). The fields are analytical questions, not seven required software components, a completeness score, or a maturity scale. We will learn them one at a time and then use the ACT-R record presented in the paper to see how they fit together.

\begin{learningbox}
By the end of this chapter, you should be able to read and write the seven-field tuple; distinguish state from control and a trigger from a vague invocation; describe persistence, failure, learning, and resource governance without reducing them to storage or logging; apply the fields to a source-grounded ACT-R example; and preserve unknown information instead of filling it by inference. You should also be able to explain why the tuple is separate from the paper's evidence-relation and migration-depth codes.
\end{learningbox}

\section{Why one structured representation helps}
\label{sec:seven-fields-why}

Return to the two descriptions from Chapter 9. Both systems are called memory systems. One is described as a searchable store. The other is described through activation-sensitive retrieval, latency, decay, competition, and utility. If we compare them only by their feature names, we can say that both provide access to past information. We cannot yet explain how they control retrieval or how their operation affects what happens later.

Writing a completely different narrative for every system makes comparisons difficult to inspect. A common set of questions helps expose the similarities and differences without forcing the systems into identical implementations. Fan and Lan introduce a seven-field mechanism representation for that purpose. For each mechanism, the reviewer records the operative state substrate (S), control locus (C), transition trigger (T), persistence boundary (P), failure semantics (F), learning operator (L), and resource or governance policy (R).\FanLanCite{§II-C}

The tuple is:
\[
  K = \langle S, C, T, P, F, L, R \rangle.
\]

This notation is compact, but it should not be mistaken for a software interface. It does not say that every agent must implement seven modules, expose seven APIs, or store seven corresponding objects. A single runtime component might participate in several fields. Conversely, one field may involve multiple components or state objects. The tuple describes the mechanism for analysis; it does not prescribe the system's implementation.

\BookDiagram[0.62\textheight]
  {Tex/diagrams/generated/part2/chapter10/seven-field-mechanism-tuple.png}
  {The seven-field mechanism representation. Each field is an analytical question about the same documented operation, not a required runtime component.}
  {fig:seven-field-mechanism-tuple}

The seven questions are related, but they are not seven consecutive steps. State and control are often tightly connected; a trigger may activate a control operation; persistence describes what survives beyond the immediate transition; failure and learning describe how outcomes are handled; and resource or governance policies constrain the operation. The relationships depend on the mechanism being examined.

\begin{notebox}
The tuple does not require every field to contain a fully specified value. A paper may describe the state and trigger clearly while saying little about learning or resource policy. An honest mechanism record preserves that gap. Missing documentation is not proof that a property is absent, and it is not permission to invent a plausible answer.
\end{notebox}

\section{State substrate (S): what the mechanism operates on}
\label{sec:field-state}

The first question is: what is the operative \gls{statesubstrate} that the mechanism consumes, maintains, or relies on? It is not simply the broad data category associated with a feature. It is the state involved in the operation being reconstructed.\FanLanCite{§II-C; Appendix A, Table VII}

Suppose a description says that an agent has memory. That description does not identify the memory state. A more useful account might identify chunks and buffers, beliefs and intentions, a workspace and candidate coalitions, or a trace of prior execution. These examples are not interchangeable. Each identifies a different kind of state with different possible relationships to control.

The reviewer should record only the state objects connected to the mechanism. If an architecture contains many data structures, the seven-field record is not an inventory of everything in the system. It focuses on the state consumed or maintained by the particular operation under study. A repository field or class name is not enough by itself; it must be linked to the operation and supported by the relevant source evidence.

\subsection{State is not the same as a feature name}

``Memory'', ``context'', ``plan'', and ``agent state'' are broad labels. A mechanism record should make them concrete. For example, the ACT-R entry in Fan and Lan's historical ledger identifies typed buffers, chunks, and procedural productions as state substrate. That is more informative than saying only that ACT-R has memory or rules.\FanLanCite{Appendix A, Table VII}

This does not mean that the reviewer must reverse-engineer every internal object. The task is bounded by the source being analyzed. If the paper names only some parts of the state, record those parts and leave the rest unspecified. Do not use familiarity with an architecture to fill gaps that the paper does not fill in its particular record.

\section{Control locus (C): who or what governs the transition}
\label{sec:field-control}

The second question is: where does the consequential control decision occur? The \gls{controllocus} identifies the process, component, or authority that governs the relevant choice or transition. This field is easy to confuse with state because a component may both hold information and participate in control. The reviewer must nevertheless ask which role the source actually assigns to it.\FanLanCite{§II-C}

A database may store an approval flag, but a separate controller may decide whether the flag is valid and whether a transition is allowed. A message queue may hold work items, but another process may decide which item is dispatched. A checkpoint may preserve a previous state, but it may not decide how recovery proceeds. These are \BookExampleLabel{} distinctions between where information is stored and where control is exercised.

The ACT-R summary identifies production matching and conflict resolution as the control locus. Its state includes typed buffers, chunks, and procedural productions, but the mechanism's control is not simply located in the storage objects. The matching and conflict-resolution process governs which applicable production is selected.\FanLanCite{Appendix A, Table VII}

\begin{takeawaybox}[State versus control]
State answers ``what information or structure is involved?'' Control answers ``which process or authority governs the relevant change?'' A component that stores a value does not automatically have authority to enforce the transition associated with that value.
\end{takeawaybox}

\section{Transition trigger (T): what activates the operation}
\label{sec:field-trigger}

The third question is: what identifiable event or condition activates the control operation? The \gls{transitiontrigger} is more precise than a statement such as ``the agent acts when needed.'' It names the occurrence that causes the mechanism to run or changes the conditions under which its control rule applies.\FanLanCite{§II-C}

A trigger can be an event, a state change, a match, a condition becoming true, or a completion that enables another operation. In the paper's historical ledger, examples include a buffer match or module event for ACT-R, selection or a typed impasse for Soar, and a blackboard change that enables a knowledge source in Hearsay-II. These are examples from the paper's catalog; the point is that each ties activation to an identifiable condition.\FanLanCite{Appendix A, Table VII}

A timestamp can be a trigger when the source explicitly connects that time condition to an operation. But the fact that a job runs on a schedule is not enough to establish a particular trigger semantics. A timer may wake a process; the process may then inspect other conditions before it exercises control. The reviewer should identify the event or condition that activates the actual transition being mapped, not merely the event that starts a surrounding program.

Similarly, a model call is not automatically a trigger. It may be one step in an operation, a request made after a trigger, or an event that has no autonomous control significance. The source must connect it to the mechanism's transition.

\section{Persistence boundary (P): what survives}
\label{sec:field-persistence}

The fourth question is: which parts of the state persist, for how long, and across which boundary? The \gls{persistenceboundary} describes the scope and lifetime of state. It may distinguish a short-lived buffer from information retained across multiple cycles, tasks, or sessions. Persistence is therefore not simply a yes/no property of the whole architecture. Different state objects in the same mechanism can have different lifetimes.\FanLanCite{§II-C; Appendix A, Table VII}

The ACT-R row makes this distinction explicit. It records that chunks and learned utilities persist, while buffers are transient. The state substrate and the persistence boundary are connected, but they answer different questions. The first identifies what state is involved; the second identifies how long the relevant state remains available.\FanLanCite{Appendix A, Table VII}

A durable database row or checkpoint does not, by itself, tell us what the mechanism treats as authoritative, when the state is released, or whether a later transition is controlled by it. The reviewer should record the boundaries described by the source and avoid assuming stronger persistence guarantees than the evidence supports.

Persistence also should not be confused with learning. Information may persist unchanged. A log may survive indefinitely without affecting future control. The learning field asks whether outcomes change how the mechanism behaves later; the persistence field asks what state remains available across time or boundaries.

\section{Failure semantics (F): what failure means to the mechanism}
\label{sec:field-failure}

The fifth question is: how does the mechanism represent, classify, or respond to failure? \Gls{failuresemantics} are more specific than the existence of an error message. An error may be recorded without changing control. A mechanism-level account asks whether a failure condition is represented in a way that affects what the system does next.\FanLanCite{§II-C; Appendix A, Table VIII}

The ACT-R entry lists retrieval failure, an unavailable buffer, and the absence of an applicable production. These conditions are useful because they identify situations in which the normal operation cannot proceed as expected. The table's purpose is not to claim that every failure has one universal recovery strategy. It records the failure conditions associated with the mechanism in the paper's reconstruction.\FanLanCite{Appendix A, Table VIII}

A failure log can help an operator understand what happened. But a log is not automatically a failure-control mechanism. To establish that connection, the source must show how the condition affects a transition: perhaps a process stops, an alternative is selected, a subproblem is created, or a method is reconsidered. The exact response must come from the documented mechanism.

\begin{warningbox}
Do not infer typed failure handling from the existence of exceptions, logs, retries, or error dashboards. Those may provide evidence for parts of a system, but the failure field asks how the mechanism represents failure and whether that representation participates in control.
\end{warningbox}

\section{Learning operator (L): how later behavior can change}
\label{sec:field-learning}

The sixth question is: how can experience or outcomes change later behavior or control? The \gls{learningoperator} records the mechanism by which experience updates something that affects future operation. Fan and Lan's field is broader than the act of writing information to storage. It asks how a result changes the behavior of the mechanism on a subsequent occasion.\FanLanCite{§II-C; Appendix A, Table VIII}

The ACT-R summary records activation and utility updates, as well as procedural learning. In this account, what happened earlier can affect later retrieval or production behavior. The learning field thus connects an outcome to a future control policy or operation.\FanLanCite{Appendix A, Table VIII}

This distinction prevents a common overclaim. Storing an execution trace is not automatically learning. Logging a success is not enough to show that later choices change. A memory write may preserve evidence; a learning operator uses experience to update later behavior, selection, or control. The record should identify the update relationship that the source supports.

Learning also need not mean model fine-tuning. A system may change a parameter, adjust a utility, update a policy, refine a rule, or adapt another part of its operation. The field is about the operator and its effect, not a predetermined learning technology.

\section{Resource or governance policy (R): what constrains operation}
\label{sec:field-resources}

The seventh question is: what resources, capacities, timing constraints, priorities, or governance rules constrain the mechanism? The \gls{resourcegovernancepolicy} makes these costs and constraints part of the mechanism description rather than treating them as unrelated implementation details.\FanLanCite{§II-C; Appendix A, Table VIII}

The ACT-R row records buffer capacity, module timing, and a central production bottleneck. These constraints matter because an operation does not run in an unlimited space. Its available state capacity, module timing, or centralized control can affect what is possible and when.\FanLanCite{Appendix A, Table VIII}

Resources are not limited to token budgets. Depending on the mechanism, relevant constraints may include memory capacity, latency, compute, timing, concurrency, queue capacity, or a limited workspace. Governance may include priorities, admission rules, or conditions on which actions can be taken. The reviewer must still stay close to the source: the field is not a license to attach every possible resource or policy concern to every mechanism.

\begin{notebox}
A field can be present in the analysis even when the source gives little detail. Record the actual constraint that is documented. If resource policy is not reported, say so; do not invent a token budget, scheduler, or governance policy to make the tuple look complete.
\end{notebox}

\section{A worked reading: ACT-R across all seven fields}
\label{sec:actr-worked-reading}

We can now read the ACT-R entry in Fan and Lan's historical ledger as one complete example. The account below uses the fields as the paper records them. It is not a full account of ACT-R, does not add details from external ACT-R sources, and should not be used as a substitute for the paper's original references.\FanLanCite{§III-A; Appendix A, Tables VII--VIII}

\begin{description}
  \item[State substrate (S).] Typed buffers, chunks, and procedural productions. These are the state structures the ledger identifies for the mechanism.
  \item[Control locus (C).] Production matching and conflict resolution. The control question is not answered by listing the buffers; it is answered by identifying the process that resolves which applicable production governs the next operation.
  \item[Transition trigger (T).] A buffer match or module event. The source identifies these as activation conditions for the mechanism.
  \item[Persistence boundary (P).] Chunks and learned utilities persist, while buffers are transient. One mechanism therefore contains state with different lifetimes.
  \item[Failure semantics (F).] Retrieval failure, an unavailable buffer, or no applicable production. These are the listed conditions under which the expected operation cannot proceed normally.
  \item[Learning operator (L).] Activation and utility updates, together with procedural learning. The ledger records updates that can affect later behavior.
  \item[Resource/governance policy (R).] Buffer capacity, module timing, and the central production bottleneck. These constraints shape the operational setting.
\end{description}

\BookDiagram[0.60\textheight]
  {Tex/diagrams/generated/part2/chapter10/actr-seven-field-reading.png}
  {The ACT-R record as summarized by Fan and Lan across the seven fields. The figure reproduces the analytical contents of the paper's ledger, not a deployment or implementation diagram.}
  {fig:actr-seven-field-reading}

The value of this reading is not that it gives ACT-R a score. It makes a mechanism inspectable. A reader can ask whether a different system has comparable state, which process controls production choice, what activates it, what persists, how failure appears, what is learned, and which capacity or timing constraints matter. Each answer may reveal a preserved property or a missing one.

The fields also constrain one another. A trigger should be connected to a control operation over identifiable state. A persistence statement should specify which state survives. A learning statement should identify what outcome changes future behavior. A resource statement should describe a real constraint on the operation. If one of those connections cannot be supported, it should remain an open question rather than be completed by inference.

\section{The fields are related, but they are not a universal sequence}
\label{sec:field-relationships}

It is tempting to read the tuple as a pipeline: state is loaded, a trigger fires, a controller acts, information persists, failure is handled, learning occurs, and resources are checked. That interpretation would be misleading. The order in which the fields are written is a convenient way to organize questions. It is not a universal temporal sequence and it does not imply that all seven are exercised as one chain during every execution.

For example, persistence may be a property of state that exists before a particular trigger occurs. Resource constraints may shape the control policy continuously rather than appearing only at the end of an operation. Failure semantics may be relevant only under certain conditions. Learning may occur after repeated outcomes or at a later update point. The relationships depend on the actual mechanism documented by the source.

\BookDiagram[0.50\textheight]
  {Tex/diagrams/generated/part2/chapter10/field-relationships.png}
  {Conceptual relationships among the seven fields. The connections indicate analytical dependencies, not a mandatory execution order or a universal runtime schedule.}
  {fig:field-relationships}

The diagram is therefore a relationship map. It reminds us that a field should be interpreted in relation to the other fields where the source provides that connection. It should not be read as seven software boxes connected in a required architecture.

\section{Unknown, absent, and not applicable are not interchangeable}
\label{sec:unknown-fields}

A structured representation can create pressure to fill every cell. If the reader sees a seven-column table, an empty cell may look unfinished. But filling it with an assumption would defeat the purpose of the method. The reviewer should distinguish at least three situations in working notes:

\begin{description}
  \item[Not reported in the available source.] The source does not provide enough information to make a positive claim about the field.
  \item[Explicitly absent or not present in the described mechanism.] The source itself supports a claim that the mechanism lacks a property or operation.
  \item[Not applicable to the mechanism as described.] The field may not apply in the same way to the operation being reconstructed, and that interpretation should be justified.
\end{description}

These are editorial ways to preserve uncertainty in the analysis. They are not additional codes introduced by Fan and Lan. The source paper's seven-field tuple remains exactly \(\langle S,C,T,P,F,L,R\rangle\). A blank or unknown field must not be silently converted into evidence of absence.\FanLanCite{§II-C; §VII-C}

\begin{warningbox}
Never fill an unknown field from a familiar architecture name, a framework feature, or a visually similar diagram. The representation exists to expose evidence gaps, not to conceal them behind plausible guesses.
\end{warningbox}

\section{From field extraction to a comparison claim}
\label{sec:fields-to-claim-packet}

Once the seven fields have been extracted for a historical mechanism and a modern case, a reviewer can compare the corresponding state and control relationships. Fan and Lan describe a claim packet that records the historical mechanism and primary source, the modern implementation and primary source, matching state, matching trigger, matching control or learning edge, preserved invariant, missing invariant, and then the separate evidence-relation and migration-depth codes.\FanLanCite{§II-C--§II-E}

The seven fields support the reconstruction. They do not replace the comparison tests. A field-by-field resemblance alone does not establish that two mechanisms are the same. The reviewer must still determine whether the nodes and edges correspond and whether at least one invariant is preserved. Equally, the tuple is not a substitute for implementation evidence when a claim depends on actual executable behavior.

\BookDiagram[0.55\textheight]
  {Tex/diagrams/generated/part2/chapter10/fields-to-claim-packet.png}
  {Seven-field extraction provides the mechanism description used in a later comparison packet. Evidence relation and migration depth remain independent classifications, not additional tuple fields.}
  {fig:fields-to-claim-packet}

The next chapter will apply this method across the ten historical architecture families and the modern system landscape reviewed by Fan and Lan. It will teach the detailed E1--E4 and D0--D4 coding rules, their evidence conditions, and the paper's historical-to-modern mapping ledger. This chapter has prepared the analytical language; it has not yet assigned codes to the full corpus.

\section{Common misreadings and the handoff to Chapter 11}
\label{sec:chapter10-common-misreadings}

The seven fields reconstruct one documented mechanism; they do not classify its historical relationship to a modern implementation. E1--E4 and D0--D4 are separate comparison axes applied only after the mechanism and its evidence have been reconstructed. The fields must therefore not be mistaken for additional codes or a maturity score.

Chapter 11 applies the representation across the historical architecture catalog and modern system landscape, then explains the evidence-relation and migration-depth rules.

\begin{chapterrecap}
Fan and Lan represent a documented mechanism with seven analytical fields: state substrate (S), control locus (C), transition trigger (T), persistence boundary (P), failure semantics (F), learning operator (L), and resource or governance policy (R). The tuple \(K=\langle S,C,T,P,F,L,R\rangle\) is not a software architecture, a score, or a mandatory component checklist. Each field asks a different question about one operation. The ACT-R example shows how the fields complement one another while preserving the source's limits. Unknown information must remain unknown rather than be inferred. The tuple supports mechanism reconstruction; the E1--E4 and D0--D4 axes separately classify evidence relation and migration depth. Chapter 11 will apply the method across the reviewed architectures and runtimes.
\end{chapterrecap}

\begin{reviewquestions}{Review questions and short exercises}
\begin{enumerate}
  \item Write the seven-field tuple in the correct order and explain what each symbol asks the reviewer to identify.
  \item Why is state substrate not simply another name for a feature label?
  \item A system stores an approval flag, while a separate controller validates it and authorizes the transition. Which field describes the stored state, and which describes the authority?
  \item Give an example of an identifiable trigger and explain why ``the agent acts when needed'' is not a sufficiently precise trigger description.
  \item How can one mechanism contain both transient and persistent state?
  \item Why is an error log not automatically a failure semantics description?
  \item Distinguish storing an outcome from learning from an outcome.
  \item Name three possible resource or governance constraints that are not token budgets.
  \item Reconstruct the ACT-R row from the paper across all seven fields. Which field records that buffers are transient?
  \item Why must an unknown field remain unknown rather than being inferred from the architecture name?
  \item Explain why the E1--E4 and D0--D4 axes are not fields in the seven-field tuple.
\end{enumerate}
\end{reviewquestions}

\section*{Source and example note}
\addcontentsline{toc}{section}{Source and example note}
This chapter is grounded in Fan and Lan (2026), primarily §II-C, with §II-A, §§II-D--II-E, §III-A, Appendix A (Tables VII--VIII), Appendix B, and §VII-C used for the mechanism definition, examples, comparison boundary, and validity limits. The ACT-R worked reading follows the paper's own ledger entries and does not attempt a complete account of ACT-R. The simplified approval, queue, and checkpoint illustrations are book-created examples. The seven-field tuple is attributed only to Fan and Lan; it is not a CoALA construct. The detailed historical-to-modern mapping and E/D coding rules are reserved for Chapter 11.\FanLanCite{\S II-A; \S II-C; \S II-D--II-E; \S III-A; Appendix A, Tables VII--VIII; Appendix B; \S VII-C}


\chapter{From Historical Cognitive Architectures to Modern Agents}
\begin{notebox}
Chapter placeholder. No prose has been drafted. This chapter will be outlined and approved before writing.
\end{notebox}

\chapter{Understanding Activation, Retrieval Latency, and Utility: B1}
\label{ch:activation-latency-utility}

\section*{Chapter at a glance}
\phantomsection
\addcontentsline{toc}{section}{Chapter at a glance}

An agent can remember far more than it can place in one prompt, and it can retrieve far more than it can safely rely on. A useful memory subsystem must therefore answer several different questions. Should a new observation be stored? Should an existing memory be revised, retained, or forgotten? Which valid memories should be fetched for this task? Which of those candidates deserve scarce context space, and in what order should they appear? Finally, did the selected information actually help the agent act well?

These questions are related, but they are not one decision. This chapter begins with the cognitive account of activation in ACT-R, then compares modern systems according to the lifecycle decision they implement. It uses that distinction to explain B1 in Fan and Lan's review: a proposed coupling of memory accessibility, predicted retrieval latency, downstream value, and maintenance cost at context assembly. B1 is a residual mechanism hypothesis, not a completed universal scoring algorithm. The chapter then turns the hypothesis into an auditable production design, using a running Product X support-agent example.

Each mechanism is treated through five questions: what the original source says; what its variables and assumptions mean; how related systems implement comparable ideas; how a production agent could represent and operate it; and what evidence would demonstrate that it works. This avoids presenting formulas without enough explanation to assess or implement them.

\begin{learningbox}
By the end of this chapter, you should be able to explain ACT-R activation and retrieval time; distinguish production utility from memory value; classify memory admission, retention, and retrieval mechanisms; compare selected modern systems without conflating their purposes; derive a bounded B1-oriented context selector; design event records that distinguish retrieval from actual use; and evaluate the resulting policy under quality, latency, cost, and reliability constraints.
\end{learningbox}

\begin{takeawaybox}[The chapter's organizing distinction]
A memory system makes at least three different lifecycle decisions: whether a memory enters persistent storage, whether it remains there or changes, and whether it is retrieved and inserted into a particular working context. B1 is primarily about the third decision and its feedback into future accessibility and utility. It should not be collapsed into a generic memory-admission or forgetting score.
\end{takeawaybox}

\section{What B1 claims—and what it leaves open}
\label{sec:ch12-b1}

Fan and Lan identify B1 as an activation--latency--utility coupling gap. Their proposed runtime view asks a memory item to carry four relevant quantities: its activation, its predicted \gls{retrievallatency}, its downstream \gls{memoryvalue}, and its maintenance cost. At \gls{contextassembly}, candidate items compete under a fixed resource budget. Selected items enter the working context; subsequent observations of use and outcome should update the accessibility and utility estimates used by later decisions. The review locates plausible implementation surfaces in adaptive memory systems and agent runtimes, but it does not present a single validated control law that solves the whole problem.\FanLanCite{\S VI-C, B1; Table V; Appendix C-A, R1; Appendix B}

That distinction matters. ACT-R already provides a model of declarative-memory activation and retrieval latency. Modern agent systems already implement different combinations of memory extraction, scoring, retention, retrieval, and context construction. The remaining question is how these signals should be coupled in a real agent: when one signal should dominate another, which costs belong to the current decision, how the selected set should be valued, and how observed downstream outcomes should change future estimates.

B1 also has a boundary with B3. B1 concerns the accessibility and selection of memory items for a task. B3 concerns competition and capacity in a shared workspace. A memory selector may decide which candidate facts enter a context, while a workspace controller may decide which partial result, subtask, or representation receives scarce shared capacity. Both involve competition, but their objects and triggers differ. Treating them as one problem would obscure the mechanism each chapter is meant to explain.\FanLanCite{Table V; \S VI-C, B1 and B3}

\subsection{Three lifecycle decisions}
\label{subsec:ch12-three-decisions}

The first task is to identify which decision a mechanism actually controls.

\begin{description}
  \item[Admission or encoding.] This is the \gls{memoryadmission} decision: should a newly observed candidate become persistent memory? The decision can include whether to store it as a new record, merge it with an existing record, reject it as unsupported, or encode it at a particular depth.
  \item[Retention, update, or forgetting.] This is the \gls{memoryretention} decision: what should happen to a memory already in storage? It may be strengthened, corrected, superseded, invalidated, consolidated, promoted to a durable tier, archived, or evicted.
  \item[Retrieval and context assembly.] Under \gls{contextassembly}, which currently valid memories should be fetched, selected, ordered, and inserted for a query, task, or possible action?
\end{description}

A single architecture can implement all three. The classification is about the operation, not the product name. An admission threshold answers whether a new record should persist. It does not automatically rank existing records for a particular question. An eviction policy chooses what loses storage priority. It does not necessarily identify the evidence needed for the next decision. A query-time ranker chooses candidates for a prompt, but it does not necessarily make a durable memory-admission decision.

The word \emph{admission} needs one further qualification. In this chapter, \emph{persistent-memory admission} means admitting a new record to long-term storage. Fan and Lan use memory admission in the B1 context-assembly setting for a different boundary: deciding which already-stored candidates are admitted into the active context for the current task. We call that \emph{context admission} when the distinction matters. A-MAC is discussed as a persistent-memory admission baseline; B1's central selection problem is context admission. The shared word does not make the two policies interchangeable.\FanLanCite{Table V; \S VI-C, B1; Appendix C-A, R1}

\begin{table}[tbp]
\centering
\small
\caption{Three distinct memory-lifecycle decisions and their triggers.}
\label{tab:ch12-lifecycle-decisions}
\begin{tabularx}{\linewidth}{@{}p{0.18\linewidth}p{0.29\linewidth}X@{}}
\toprule
\textbf{Decision} & \textbf{Typical trigger} & \textbf{Question being answered} \\
\midrule
Admission / encoding & A new observation or candidate memory is extracted & Is this information reliable, novel, durable, and useful enough to store? \\
Retention / update / forgetting & A stored memory is revisited, contradicted, aged, or competes for storage & Should this record be updated, kept, moved, invalidated, or removed? \\
Retrieval / context assembly & A task, query, or candidate action requires context & Which currently valid records will make this decision better, given context and latency limits? \\
\bottomrule
\end{tabularx}
\end{table}

\BookDiagram[0.54\textheight]{Tex/diagrams/generated/part2/chapter12/memory-lifecycle-decisions.png}{Admission, retention, and context-time retrieval are different decisions with different triggers and outcomes. The diagram distinguishes persistent-store lifecycle operations from query-specific context assembly.}{fig:ch12-memory-lifecycle}

\subsection{A running example: Product X}
\label{subsec:ch12-productx}

\BookExampleLabel{} Imagine a support agent for Product X. A customer reports that a payment retry failed. The agent may have access to several records: an error log from the current incident, an old note that a particular payment provider was disabled temporarily, a prior troubleshooting sequence, a database dependency map, and an outdated description of the retry policy.

The admission question is whether the newly observed incident facts are sufficiently supported and useful to preserve. The retention question is whether the old provider note remains valid, whether the policy record has been superseded, and whether a resolved incident should be summarized or archived. The retrieval question is which records should be placed into the context for the current diagnosis. A low-frequency fact about a recent configuration change may be more useful than a highly visited but generic troubleshooting note. Two individually weak records—a failed retry and a dependency relationship—may jointly identify the likely fault.

This example will recur throughout the chapter. It illustrates why memory age, frequency, semantic similarity, confidence, action value, and resource cost are separate signals rather than synonyms for one another.

\section{ACT-R: from declarative chunks to predicted retrieval}
\label{sec:ch12-actr}

\gls{actr} provides a foundational account of declarative memory.\index{ACT-R} Information is represented as chunks with typed slots. The current cognitive state and task provide retrieval cues; a retrieval request attempts to find a chunk matching those cues. The model assigns activation to candidate chunks. A sufficiently activated candidate can be retrieved, and activation is related to the predicted time required for retrieval. This is more specific than saying that an agent searches a vector database: the account defines a cognitive quantity and a relationship between that quantity and retrieval behavior \autocite{anderson2007actr,actr7manual}.

\subsection{Activation has several components}
\label{subsec:ch12-activation-components}

A useful decomposition of \gls{activation} is\index{activation}:
\begin{equation}
  A_i = B_i + S_i + P_i + \epsilon_i,
  \label{eq:ch12-activation}
\end{equation}
where \(A_i\) is the activation of candidate chunk \(i\), \(B_i\) is its base-level activation, \(S_i\) is activation spreading from the current context, \(P_i\) is the partial-matching contribution, and \(\epsilon_i\) represents activation noise. The terms explain different reasons that a chunk may become accessible:

\begin{itemize}
  \item \textbf{\gls{baselevelactivation}} captures the contribution of the chunk's presentation history. Repeatedly encountered information tends to have greater activation; older presentations contribute less as they recede in time.
  \item \textbf{\gls{spreadingactivation}} captures the influence of related information currently present in the model's context. A cue can increase the activation of associated chunks.
  \item \textbf{\gls{partialmatching}} allows imperfectly matching chunks to compete. Similarity can help, while mismatching slot values can impose a penalty.
  \item \textbf{Noise} captures stochastic variation rather than a deterministic rank that always returns the same winner.
\end{itemize}

These components should not be collapsed into a single embedding-similarity score. Cosine similarity can be one retrieval feature in an agent system, but it is not automatically ACT-R \gls{baselevelactivation}, \gls{spreadingactivation}, \gls{partialmatching}, or total activation. The ACT-R decomposition is a model of accessibility with defined sources of contribution, not a new name for every modern relevance score \autocite{actr7manual}.

\subsection{Base-level activation: frequency and recency}\index{base-level activation}
\label{subsec:ch12-base-level}

The base-level equation accumulates the influence of the chunk's presentation events:
\begin{equation}
  B_i = \ln\left(\sum_{j=1}^{n} t_j^{-d}\right)+\beta_i.
  \label{eq:ch12-base-level}
\end{equation}
Here, \(n\) is the number of relevant presentations, \(t_j\) is the elapsed time since presentation \(j\), \(d\) is the decay parameter, and \(\beta_i\) is a chunk-specific or model-specified base-level contribution. The logarithm compresses the sum. Each presentation adds a time-decayed contribution: many presentations can raise activation, while the influence of any one presentation declines as its elapsed time grows \autocite{actr7manual}.

The event semantics are important. In the ACT-R 7 reference manual consulted for this chapter, a presentation is associated with a chunk entering declarative memory or with a chunk being merged with an already existing declarative-memory chunk. It is therefore incorrect to state as a universal ACT-R rule that every database-style read or every unsuccessful retrieval attempt increments the presentation history. A model may define additional rehearsal or re-presentation behavior, but an implementation must state exactly which event changes the history \autocite{actr7manual}.

For an agent, at least six events should be distinguished: a memory is created or stored; retrieval is requested; retrieval succeeds; the memory is included in the prompt; the model actually uses it; and an outcome is observed. These events may happen close together, but they are not interchangeable. Refreshing recency on every read can create a positive feedback loop: once a memory is retrieved, it becomes more recent, which makes it more likely to be retrieved again even if it did not help.

ACT-R also describes an optimized approximation when the full presentation history is not retained:
\begin{equation}
  B_i \approx \ln\left(\frac{n}{1-d}\right)-d\ln(L)+\beta_i,
  \label{eq:ch12-base-level-approx}
\end{equation}
where \(L\) is the chunk's lifetime since creation. The approximation replaces the full sum over event timestamps with a compact expression based on presentation count and lifetime. It can reduce bookkeeping, but it loses information about the exact spacing of individual presentations. If an agent needs reproducible, event-level analysis, retaining timestamps is more informative than storing only a precomputed activation number. If memory volume makes that too expensive, the approximation should be recorded as the selected model rather than mistaken for the full-history calculation \autocite{actr7manual}.

\begin{notebox}
The base-level equation is a cognitive activation mechanism, not a universal retention rule. An agent may separately decide that a record must be deleted for privacy, invalidated because its facts changed, or retained because it is legally or operationally required. Those decisions are not resolved by substituting a different decay parameter into Equation~\ref{eq:ch12-base-level}.
\end{notebox}

\subsection{Spreading activation and association fan}
\label{subsec:ch12-spreading}

A current context can activate a chunk through associations. In the ACT-R account, the total spreading contribution can be written in the form
\begin{equation}
  S_i = \sum_k \sum_j W_{kj}S_{ji},
  \label{eq:ch12-spreading-total}
\end{equation}
where \(k\) indexes source buffers, \(j\) indexes source chunks in a buffer, \(W_{kj}\) is the amount of source activation allocated to a source chunk, and \(S_{ji}\) is the strength of association from source chunk \(j\) to candidate chunk \(i\). The default allocation distributes a buffer's source activation across its source chunks. When the source and candidate are associated, a fan-sensitive form is
\begin{equation}
  S_{ji}=S-\ln(\operatorname{fan}_{j}),
  \label{eq:ch12-spreading-fan}
\end{equation}
where \(S\) is an association-strength parameter and \(\operatorname{fan}_{j}\) counts the chunks associated with source chunk \(j\). A cue connected to many chunks distributes its influence more broadly; a more selective association can provide stronger support to a particular candidate. The exact treatment of unrelated chunks and fan is defined by the ACT-R model and its parameters \autocite{actr7manual}.

This mechanism offers a useful comparison with graph-based retrieval. A query may activate entities or nodes, and relevance may propagate through known relationships. That is conceptually related to contextual activation, but a graph algorithm such as Personalized PageRank is not ACT-R spreading activation. They have different mathematical update rules and should be compared as alternative mechanisms, not treated as equivalent implementations.

\subsection{Partial matching, noise, and retrieval threshold}
\label{subsec:ch12-partial-threshold}

Exact matching can be too brittle when a task cue differs slightly from how information was originally encoded. Partial matching allows a candidate to receive a graded contribution based on its slot similarities and mismatches. A chunk with a similar but not identical slot value may remain a candidate, while an important mismatch can reduce its activation. The similarity function and mismatch penalties are model parameters; the chapter does not assign a universal numeric penalty because the value depends on the model's slot semantics and task.

Noise means that retrieval is not necessarily a deterministic \texttt{top-1} operation. Repeated retrievals under nominally similar conditions may differ. A retrieval threshold \(\tau\) also matters: a candidate whose activation does not reach the threshold is not retrieved. The architecture can therefore fail to retrieve a chunk even when a close candidate exists. This is a useful reminder for production agents: retrieval should support an explicit no-result or abstention state rather than forcing the highest-ranked candidate into the prompt.

\subsection{Activation predicts cognitive retrieval time}\index{retrieval latency}
\label{subsec:ch12-retrieval-latency}

ACT-R relates a chunk's activation to retrieval latency:
\begin{equation}
  RT_i = F e^{-fA_i},
  \label{eq:ch12-actr-latency}
\end{equation}
where \(RT_i\) is predicted retrieval time, \(F\) is a latency scaling factor, and \(f\) controls the effect of activation. Higher activation predicts shorter retrieval time under the model's parameters. A failed retrieval has a corresponding latency determined by the threshold rather than by the activation of a successfully retrieved chunk \autocite{actr7manual}. Fan and Lan use this ACT-R memory structure as the historical reference point for their adaptive-memory mapping, while noting that modern memory policies can change the decision semantics from an accessibility law to a policy action.\FanLanCite{§ III-A; § V-A, ACT-R to adaptive memory control; § V-B, Memory changes from accessibility law to policy action; Appendix B}

A production agent must distinguish this \emph{cognitive prediction} from \emph{infrastructure latency}. The latter is the wall-clock time consumed by embedding lookup, database access, graph traversal, network I/O, reranking, serialization, or prompt construction. One is predicted from a cognitive model; the other is measured from software execution. They can both matter, but neither should be silently substituted for the other.

There is also a double-counting trap. If a selector rewards high ACT-R activation and separately rewards low \(RT_i\) from Equation~\ref{eq:ch12-actr-latency}, it may reward the same activation advantage twice. One defensible design uses activation as the cognitive accessibility signal and treats ACT-R latency as a derived prediction. Another uses activation together with independently measured infrastructure latency. An evaluation can compare these choices through ablation; a selector should not include both as independent bonuses without explaining the intended interpretation.

\begin{takeawaybox}[What ACT-R contributes to B1]
ACT-R supplies a principled account of why a memory becomes accessible and how activation relates to retrieval time. It does not, by itself, tell an agent how valuable that memory will be for a particular downstream action, how much infrastructure cost its retrieval will incur, or whether it deserves space in the next prompt.
\end{takeawaybox}

\subsection{Sidebar: Rational Analysis and future demand}
\label{subsec:ch12-rational-analysis}

Rational Analysis offers a useful interpretation of memory accessibility. A memory that has occurred often, or has occurred recently, may be more likely to be needed again given the statistics of the environment. From this perspective, an activation signal can approximate the likelihood that information will be relevant to future processing. The environmental basis of memory provides a broader account of why frequency, repetition, and spacing matter for what remains accessible \autocite{anderson2022environmental}.

But future demand is not downstream action value. A frequently retrieved generic troubleshooting note may have a high probability of being requested again while contributing little to the current incident. A rarely accessed configuration change may be decisive for one particular diagnosis. Rational Analysis helps explain why accessibility is useful; B1 asks how accessibility should interact with task-specific value, latency, and cost.

\section{Production utility is not memory value}
\label{sec:ch12-production-utility}

ACT-R also associates utility with productions, which are action or control rules. A familiar expected-value expression is
\begin{equation}
  U_i = P_iG-C_i,
  \label{eq:ch12-production-expected-utility}
\end{equation}
where \(P_i\) is the estimated probability that production \(i\) will achieve its intended result, \(G\) is the value of that result, and \(C_i\) is the expected cost of carrying out the production. This expression is useful as an expected-value intuition: an action can be attractive because it is likely to succeed and valuable if successful, but unattractive when its execution cost is high.

The ACT-R \gls{productionutility} learning update is a separate statement:\index{production utility}
\begin{equation}
  U_i(n)=U_i(n-1)+\alpha\left[R_i(n)-U_i(n-1)\right],
  \label{eq:ch12-production-learning}
\end{equation}
where \(\alpha\) is a learning rate and \(R_i(n)\) is the reward associated with the production's \(n\)-th use, with time-related cost accounted for in the reward definition. The update moves the previous utility estimate toward the newly observed reward. It is a learning rule for a production, not a formula for the value of a declarative-memory chunk \autocite{anderson2007actr,actr7manual}.

The distinction should be reflected in the data model. A production has an identity, an applicability condition, and an action outcome. A memory has content, provenance, a temporal validity interval, and a context-dependent contribution to some task or action. A field named simply \texttt{utility} is ambiguous unless its owner and update process are explicit.

\begin{table}[tbp]
\centering
\small
\caption{Production utility and memory value describe different objects.}
\label{tab:ch12-production-vs-memory}
\begin{tabularx}{\linewidth}{@{}p{0.22\linewidth}XX@{}}
\toprule
 & \textbf{Production / action utility} & \textbf{Memory value} \\
\midrule
Object & An action or production rule & A memory item or selected memory set \\
Question & Is this action worth executing? & Will this information improve this task or action? \\
Conditioning & Goal, expected success, execution cost & Query, task, candidate action, other selected memories \\
Feedback & Reward or outcome of action execution & Observed use, task outcome, attribution uncertainty \\
Lifecycle & Learned policy state for a production & Context-dependent estimate plus memory lifecycle metadata \\
\bottomrule
\end{tabularx}
\end{table}

\section{Modern memory systems: compare the decision they actually make}
\label{sec:ch12-modern-systems}

Modern agent memory research contains many useful formulas, but not every formula answers the same question. The following systems are best understood by locating their primary operation in the memory lifecycle.

\subsection{Generative Agents: a query-time relevance score}\index{Generative Agents}
\label{subsec:ch12-generative-agents}

Generative Agents rank candidate memories using recency, importance, and relevance. Recency decays exponentially according to the elapsed simulated game time since a memory was last retrieved. Importance is assigned by an LLM when the memory is created, and relevance is based on the semantic similarity between the memory and the current query. The three signals are normalized and combined:
\begin{equation}
  \operatorname{Score}_i = \alpha_r R_i + \alpha_I I_i + \alpha_s S_i,
  \label{eq:ch12-generative-score}
\end{equation}
where \(R_i\), \(I_i\), and \(S_i\) denote recency, importance, and semantic relevance. In the paper's reported configuration, the three weights are equal and the highest-scoring memories that fit the context budget are selected. The recency decay factor is 0.995 per simulated game hour; under that per-hour interpretation, the half-life is approximately 138 hours. This is a simulation-time parameter, not a claim about 138 hours of real-world time \autocite{park2023generative}.

The design is attractive because the signals are understandable. It also illustrates a feedback risk: retrieval updates the memory's recency, and recency influences the next score. A memory that is retrieved frequently can remain prominent even when no outcome evidence shows that it was useful. A production implementation should therefore record retrieval, prompt inclusion, actual use, and task outcome separately. One possible policy is to refresh an accessibility history on a well-defined rehearsal/use event rather than on every mechanical read.

\subsection{MemoryBank: a simplified retention curve}\index{MemoryBank}
\label{subsec:ch12-memorybank}

MemoryBank models retention with an exponential strength-and-time function:
\begin{equation}
  R=e^{-t/S},
  \label{eq:ch12-memorybank}
\end{equation}
where \(t\) is elapsed time and \(S\) is memory strength. In the reported simplified mechanism, a new memory starts with an initial strength, recall increases that strength, and the elapsed time is reset. The authors present the model as a simplification \autocite{zhong2024memorybank}.

The equation and ACT-R base-level activation both involve time and repeated exposure, but they are not the same model. ACT-R accumulates contributions from multiple presentation events using a power-law decay. MemoryBank uses a strength/time retention representation and a recall-driven strengthening rule. It is therefore most useful here as a retention comparison, not as a replacement for ACT-R activation or a general context-selection law.

\subsection{MemoryOS: tiered lifecycle management}\index{MemoryOS}
\label{subsec:ch12-memoryos}

MemoryOS explicitly organizes conversational memory into short-term, mid-term, and long-term personal stores, providing a concrete \gls{memoryretention} and tier-management mechanism.\index{memory retention} Short-term dialogue pages move into mid-term storage through a dialogue-chain FIFO process. Mid-term storage groups pages into topic segments, and its update mechanism uses a heat score to govern deletion and promotion. Retrieval first identifies relevant segments and then retrieves pages within those segments; the retrieved material is combined with short-term context and long-term persona information \autocite{kang2025memoryos}.

The segment heat function is
\begin{equation}
  H = \alpha N_{\text{visit}} + \beta L_{\text{interaction}} + \gamma R_{\text{recency}},
  \label{eq:ch12-memoryos-heat}
\end{equation}
with
\begin{equation}
  R_{\text{recency}}=\exp\left(-\frac{\Delta t}{\mu}\right),
  \label{eq:ch12-memoryos-recency}
\end{equation}
where \(N_{\text{visit}}\) is the number of segment retrievals, \(L_{\text{interaction}}\) is the total number of dialogue pages in the segment, and \(\Delta t\) is elapsed time since the last retrieval, measured in seconds. The paper gives a configurable time constant \(\mu\), with \(10^7\) seconds as its example. When segment length exceeds capacity, the lowest-heat segments are evicted; segments above a threshold can be transferred to long-term personal memory \autocite{kang2025memoryos}.

This is a lifecycle mechanism. It captures visit frequency, interaction volume, and recency to manage storage tiers. The heat score should not be relabeled as downstream task value: a large, frequently accessed segment may be expensive or redundant for the next prompt, and it may still contain outdated information. MemoryOS is useful as a concrete example of promotion and eviction, while B1 focuses on context-time competition and feedback.

\subsection{A-MAC: explicit memory admission}\index{A-MAC}
\label{subsec:ch12-amac}

Adaptive Memory Admission Control (A-MAC) treats entry into long-term memory as a structured decision. It evaluates five factors: future utility, factual confidence, semantic novelty, temporal recency, and a content-type prior. Its score is
\begin{equation}
  \mathcal{S}(m)=w_1\mathcal{U}(m)+w_2\mathcal{C}(m)+w_3\mathcal{N}(m)
   +w_4\mathcal{R}(m)+w_5\mathcal{T}(m),
  \label{eq:ch12-amac}
\end{equation}
where the feature functions map a candidate to \([0,1]\), weights are non-negative and sum to one, and a candidate is admitted when \(\mathcal{S}(m)\geq\theta\). The authors combine rule-based feature computation with one LLM-assisted utility assessment and fit weights and threshold through cross-validation. Their reported evaluation is on LoCoMo; the result is a research-paper claim under that experimental setup, not an independently established production guarantee \autocite{zhang2026amac}.

The five features have distinct roles. Utility estimates potential usefulness in future interactions. Confidence asks whether the candidate is supported by evidence. Novelty discourages redundant records. Recency represents temporal decay. Type prior gives stable categories—such as durable preferences—a different prior from transient states. The separation is valuable because a fact can be useful but unsupported, novel but short-lived, or reliable but redundant.

A-MAC primarily decides whether a candidate should be admitted, updated, or rejected as persistent memory. This is an example of \gls{memoryadmission}.\index{memory admission} That decision can reduce later memory bloat and hallucination propagation, but it does not answer which of the already-admitted records should enter a particular prompt. Its score is an important baseline for the admission layer, not a universal B1 context selector.

\subsection{The wider B1 baseline set: adaptive memory and value--cost governance}
\label{subsec:ch12-b1-baselines}

Fan and Lan's B1 comparison identifies AgeMem, Memory-R1, MemCon, DeltaMem, A-MAC, and CURATOR as executable adaptive baselines; DAM is conceptual. These are distinct systems, not one interchangeable algorithm. The table records the review's classification and a bounded role for each comparator without attributing undocumented internals.\FanLanCite{§ IV-B; § V-A; § V-B; § VI-C, B1; Table VI; Appendix C-A, R1; Appendix F-B}

\begin{table}[H]
\centering
\footnotesize
\caption{Memory-focused baselines identified in Fan and Lan's B1 analysis.}
\label{tab:ch12-b1-baselines}
\begin{tabularx}{\linewidth}{@{}p{0.18\linewidth}p{0.14\linewidth}X@{}}
\toprule
\textbf{System} & \textbf{Status} & \textbf{Use in this chapter} \\
\midrule
AgeMem\index{AgeMem} & Executable & Use as an adaptive-policy comparator; do not infer the full B1 coupling. \\
Memory-R1\index{Memory-R1} & Executable & Compare its adaptive-control policy on its own terms. \\
MemCon\index{MemCon} & Executable & Keep distinct from persistent admission and context ranking. \\
DeltaMem\index{DeltaMem} & Executable & Compare directly; test mechanism-specific equivalence rather than assume it. \\
A-MAC & Executable & Compare its explicit persistent-memory admission policy. \\
CURATOR\index{CURATOR} & Executable & Closest value--cost baseline; its reported signals are explained below. \\
DAM\index{DAM} & Conceptual & Conceptual comparator only; not an executable B1 baseline. \\
\bottomrule
\end{tabularx}
\end{table}

Fan and Lan describe CURATOR as fitting helpfulness online and using retrieval propensity, value, harm, and byte cost to govern keep/share/trust decisions. It is therefore the closest value--cost comparator in this set. This is the review's mechanism-level summary; the chapter neither reconstructs CURATOR's full implementation nor invents an equation.\FanLanCite{Appendix C-A, R1; Appendix B}

The paper is explicit about the novelty boundary. Adaptive admission and value--cost governance already have implementations; CURATOR is the closest value--cost comparator. B1 therefore does not claim novelty merely for assigning values to memories, learning admission policies, or considering resource cost. Its residual hypothesis is the unified coupling of accessibility/activation, predicted retrieval latency, downstream action utility, and maintenance cost at context assembly, with observed use feeding back into future accessibility and utility. That claim must be tested against the closest executable baselines, not against a straw-man top-$k$ retriever alone.\FanLanCite{§ VI-C, B1; Table V; Table VI; Appendix D, B1 experiment template}

\subsection{Chen and Cheng: a multi-factor memory-value model}\index{Chen and Cheng}
\label{subsec:ch12-chen-cheng}

Chen and Cheng's 2026 preprint proposes a multi-factor estimate of memory value:
\begin{equation}
  V(m)=\sum_i w_i f_i(m).
  \label{eq:ch12-chen-cheng}
\end{equation}
Its factors cover emotional intensity, goal relevance, value alignment, self/user relevance, task utility, reliability, and usage history. The authors describe learning factor weights against a downstream objective and using the resulting value estimate to inform encoding depth, forgetting risk, and retrieval ranking \autocite{chencheng2026remember}.

This is close to B1 in its attention to interpretable memory value, but the comparison must remain precise. A learned multi-factor value estimate does not by itself supply ACT-R activation, the activation-derived latency relationship, independent retrieval infrastructure latency, set-dependent downstream contribution, and all resource-maintenance costs in one validated runtime selector. The source is an arXiv preprint, not a peer-reviewed production standard. It is appropriate to treat it as relevant, recent prior art whose claims require the normal scrutiny of a preprint and whose exact operating assumptions must be preserved.

\subsection{Zep and Graphiti: temporal facts and hybrid retrieval}\index{Zep and Graphiti}
\label{subsec:ch12-zep-graphiti}

A memory can be well retrieved and still be wrong for the current time. The \gls{contextassembly} stage\index{context assembly} must therefore apply validity checks before selection. Temporal knowledge systems therefore need to distinguish when the system learned a fact from when that fact was true in the world. Graphiti's temporal model distinguishes fields such as \texttt{created\_at} (when the system learned the fact), \texttt{valid\_at} (when the fact became true), \texttt{invalid\_at} (when it ceased to be true), and \texttt{expired\_at} (when the system learned that it had expired). The exact fields depend on the version and integration, but the conceptual distinction is essential: system knowledge time and world-validity time are not the same \autocite{zepgraphitidocs}.

Graphiti also supports hybrid retrieval, combining semantic/vector search with lexical retrieval such as BM25 and graph-linked search. A common way to combine ranked candidate lists is Reciprocal Rank Fusion:
\begin{equation}
  \operatorname{RRF}(d)=\sum_{r\in\mathcal{R}}\frac{1}{k+\operatorname{rank}_r(d)},
  \label{eq:ch12-rrf}
\end{equation}
where \(d\) is a candidate document, \(\mathcal{R}\) is the set of ranked result lists, \(\operatorname{rank}_r(d)\) is the rank assigned by list \(r\), and \(k\) is a smoothing constant. RRF combines ranks rather than assuming that scores from different retrieval systems share the same scale. The chapter does not attribute a particular value of \(k\) to a vendor implementation unless that value is verified for the version in use \autocite{cormack2009rrf,zepgraphitidocs}.

Temporal validity and hybrid candidate generation belong before final context assembly. A stale fact should not become eligible merely because its embedding is similar to the query. A robust selector first applies provenance, validity, privacy, and conflict checks, then scores or optimizes the remaining candidates. Retrieval quality begins with constructing a defensible candidate set.

\subsection{HippoRAG: graph propagation from query entities}\index{HippoRAG}
\label{subsec:ch12-hipporag}

HippoRAG combines language-model extraction with a knowledge graph and query-seeded Personalized PageRank to support multi-hop retrieval. A query is connected to entities; relevance then propagates through the graph so that related facts can be surfaced even when they are not direct lexical matches. This is useful when a diagnosis or question depends on a chain of relationships rather than a single nearest-neighbour match \autocite{gutierrez2024hipporag}.

The analogy to ACT-R spreading activation is conceptual: both allow contextual cues to influence associated information. The algorithms are not the same. ACT-R's spreading activation uses its buffer-source and association-strength formulation. HippoRAG uses graph structure and a PageRank-style propagation process. A system designer should choose between them based on the retrieval problem and the required evidence, not call one a direct implementation of the other.

\subsection{MMR: relevance with redundancy control}
\label{subsec:ch12-mmr}

Selecting the top \(k\) individually relevant memories can fill a prompt with near-duplicates. Maximal Marginal Relevance (MMR) selects a candidate by balancing its relevance to the query against its similarity to items already selected:
\begin{equation}
  m^* = \arg\max_{m\notin S}
  \left[\lambda\,\operatorname{Sim}(m,q)
  -(1-\lambda)\max_{s\in S}\operatorname{Sim}(m,s)\right],
  \label{eq:ch12-mmr}
\end{equation}
where \(S\) is the selected set, \(q\) is the query, and \(\lambda\in[0,1]\) controls the relevance--diversity trade-off. MMR is a practical diversity heuristic. It is not a direct estimate of downstream action value \autocite{carbonell1998mmr}.

The distinction becomes clear if we define the value of a selected set as
\begin{equation}
  V(S\mid q,a),
  \label{eq:ch12-set-value}
\end{equation}
where \(q\) is the query and \(a\) is a possible action or decision. The marginal contribution of a candidate memory is
\begin{equation}
  \Delta V(m\mid S,q,a)=V(S\cup\{m\}\mid q,a)-V(S\mid q,a).
  \label{eq:ch12-marginal-value}
\end{equation}
This formulation captures two interactions. Redundancy occurs when the new memory adds little because the selected set already contains equivalent evidence. Complementarity occurs when two memories together have more value than either one alone. For Product X, an incident log and a dependency relation may jointly isolate a fault even if neither is decisive by itself.

A true estimate of \(V(S\mid q,a)\) requires downstream evidence or a defensible model. MMR can approximate one part of the problem by reducing semantic redundancy, but it cannot discover every complementarity relation or prove that a diversified set improves task success. The implementation should distinguish the heuristic it uses from the outcome it hopes to improve.

\section{Constructing the B1 integration}
\label{sec:ch12-b1-integration}

The mechanisms above provide useful pieces: ACT-R contributes accessibility and predicted cognitive latency; Generative Agents combine understandable query-time signals; MemoryBank and MemoryOS show retention and tier management; A-MAC provides an admission baseline; Chen and Cheng offer a value model; graph retrieval and MMR help construct a candidate set. B1 asks how the relevant dimensions should interact at the moment of context assembly and how outcomes should update the policy. This section provides an engineering formulation to make that question testable. It is a book-level synthesis, not an equation published by Fan and Lan or by ACT-R.

\subsection{Four signals and five cost categories}
\label{subsec:ch12-four-signals-costs}

A B1-oriented selector should preserve the distinction among four kinds of information:
\begin{enumerate}
  \item \textbf{Activation/accessibility:} how accessible is the memory under the current task context?
  \item \textbf{Predicted retrieval latency:} how long does the selected retrieval model predict that retrieval will take?
  \item \textbf{Downstream value:} how much is this memory expected to contribute to the current task or candidate action, considering what is already selected?
  \item \textbf{Cost:} what incremental resources are consumed by storing, retrieving, including, validating, refreshing, or maintaining it?
\end{enumerate}

The fourth item is not a single timeless cost, and should be distinguished from \gls{memoryvalue}. At least five categories should be recorded:

\begin{description}
  \item[Persistent storage and indexing.] Space, embedding creation, index insertion, replication, and any background consolidation.
  \item[Retrieval and reranking.] Query execution, vector search, BM25, graph expansion, filtering, and reranker computation.
  \item[Context cost.] Tokens consumed by the memory and the associated inference/input-token cost or lost opportunity to include other context.
  \item[Latency cost.] Wall-clock time, including serial dependencies, queueing, and tail latency.
  \item[Maintenance cost.] Validation, temporal refresh, deduplication, source re-checking, deletion, and privacy-compliance work.
\end{description}

A persistent storage cost is usually incurred when the memory is admitted or maintained. It is not automatically a fresh marginal cost every time the memory is considered for a prompt. Subtracting the full historical embedding or storage cost from every prompt would repeatedly charge a sunk cost. It can enter a prompt-time objective if the system explicitly amortizes it over expected uses, assigns a marginal refresh cost, or converts limited storage into a shadow price. The accounting rule must be stated.

\subsection{Latency depends on execution structure}
\label{subsec:ch12-latency-structure}

Candidate retrievals may run serially, in parallel, or with partial dependencies. If retrievals are genuinely serial, their elapsed retrieval time is approximately
\begin{equation}
  T_{\mathrm{serial}}(S)=\sum_{i\in S}T_i.
\end{equation}
If they are independent and fully parallel, the elapsed time is approximately
\begin{equation}
  T_{\mathrm{parallel}}(S)=\max_{i\in S}T_i.
\end{equation}
Real systems often sit between these extremes. A graph expansion may depend on a first-stage entity retrieval; a reranker may wait until vector and lexical results are ready; a concurrency limit may force some tasks to queue. The general solution is to measure end-to-end wall-clock latency or model the critical path of the actual dependency graph. Do not sum every per-item latency when the implementation overlaps work, and do not assume maximum latency when dependencies are serial.

A practical selector can impose a deadline or latency budget:
\begin{equation}
  T_{\mathrm{observed}}(S)\leq D,
  \label{eq:ch12-latency-budget}
\end{equation}
where \(D\) is a task-specific deadline. The predicted cognitive latency from ACT-R may inform accessibility, but infrastructure latency should be measured separately and should include relevant tail behavior such as p95, not only average time.

\subsection{A set-aware, budget-constrained context objective}
\label{subsec:ch12-context-objective}

For a query \(q\), candidate action \(a\), and selected memory set \(S\), an illustrative objective is
\begin{equation}
  \max_{S\subseteq\mathcal{M}_{q}}
  \quad \widehat{V}(S\mid q,a)
  -\lambda_r C_{\mathrm{retrieval}}(S)
  -\lambda_c C_{\mathrm{context}}(S)
  -\lambda_l C_{\mathrm{latency}}(S),
  \label{eq:ch12-b1-objective}
\end{equation}
subject to
\begin{align}
  \operatorname{Tokens}(S)&\leq B,\\
  T_{\mathrm{observed}}(S)&\leq D,\\
  m\in S&\Rightarrow \operatorname{Valid}(m)\land\operatorname{Allowed}(m),
  \label{eq:ch12-b1-constraints}
\end{align}
where \(\mathcal{M}_q\) is the candidate set after query-time retrieval and hard filtering, \(\widehat{V}\) estimates set-level downstream value, \(B\) is the context-token budget, and \(D\) is the latency deadline. The \(\lambda\) terms convert the relevant resource costs into comparable units or shadow prices. Their values must be calibrated to the deployment objective, not chosen as universal constants.

The objective intentionally does not subtract total historical storage cost by default. Persistent costs belong in the admission and maintenance budgets unless the design assigns them an explicit amortized marginal cost. If an application has a strict memory-store capacity, storage can be managed by a separate lifecycle policy or represented as a long-horizon constraint rather than repeatedly charged as a per-prompt penalty.

A selector may use activation to estimate accessibility, with ACT-R's latency model as a derived prediction. Or it may combine activation with separately measured infrastructure latency. In either case, the accounting should avoid counting the same activation signal twice. The formulation also values the selected set, not just a sum of independent item values. A simple implementation can approximate set value using MMR or another redundancy penalty; a stronger implementation can estimate marginal outcome contribution from controlled experiments.

\subsection{A bounded selection procedure}
\label{subsec:ch12-selection-procedure}

\EngineeringLabel{} The following procedure is a practical baseline for Product X. It does not claim to solve the research problem. It makes assumptions visible and produces an auditable trace:

\begin{enumerate}
  \item \textbf{Build candidates.} Run the appropriate semantic, lexical, and graph retrieval paths. Preserve candidate provenance and the retrieval method that produced each candidate.
  \item \textbf{Apply hard gates.} Reject or quarantine records outside the authorized privacy scope, records with invalid temporal status, and unsupported or superseded records that cannot be reconciled. Hard policy constraints should not be traded away for a high relevance score.
  \item \textbf{Compute query-time features.} Estimate activation components, relevance, confidence, expected downstream contribution, token size, and infrastructure cost. Record the model and parameter versions.
  \item \textbf{Select a feasible set.} Add candidates according to estimated marginal value per resource, subject to token and critical-path latency budgets. Use a redundancy heuristic where direct set-value estimates are unavailable.
  \item \textbf{Order and serialize.} Order records according to the prompt template and evidence requirements. Keep source identifiers and temporal qualifiers with the content so the agent can distinguish facts from hypotheses.
  \item \textbf{Observe outcomes.} Log retrieval, inclusion, use evidence, and task outcome as separate events. Update only those estimates justified by the observed evidence.
  \item \textbf{Degrade safely.} When evidence is weak, contradictory, too expensive, or unavailable before the deadline, retrieve fewer items, ask a targeted clarification, use a safe fallback, or abstain.
\end{enumerate}

The policy should preserve a no-result state. A system that always inserts the highest-scoring record will eventually insert a misleading one when every candidate is stale, unsupported, or irrelevant. Selective retrieval includes the ability to say that no candidate passes the required threshold.

\BookDiagram[0.58\textheight]{Tex/diagrams/generated/part2/chapter12/b1-context-assembly-feedback.png}{A B1-oriented context assembly path: hard eligibility gates precede value and cost comparison; the selected set is constrained by token and elapsed-time budgets, while outcomes update future estimates.}{fig:ch12-b1-control-loop}

\section{Product X production design: state, events, and contracts}
\label{sec:ch12-production-design}

A production system must be able to reconstruct why a memory was retrieved and why it entered a prompt. A single mutable row with \texttt{last\_accessed\_at}, \texttt{activation}, and \texttt{utility} is insufficient. Activation depends on the current query, context cues, model parameters, and access history. Value depends on the task and candidate action. Prompt position varies from one prompt to the next. Costs depend on the actual path taken. These are event or estimate properties, not immutable characteristics of the memory record.

\subsection{Stable memory record}
\label{subsec:ch12-memory-record}

A \texttt{MemoryRecord} stores the stable identity and content metadata:

\begin{itemize}
  \item \texttt{memory\_id}, \texttt{memory\_type}, \texttt{content\_ref}, and \texttt{version};
  \item \texttt{source\_id}, \texttt{source\_type}, and provenance or supporting evidence references;
  \item \texttt{stored\_at} and \texttt{source\_observed\_at};
  \item \texttt{valid\_from}, \texttt{valid\_to}, \texttt{invalidated\_at}, and \texttt{superseded\_by};
  \item confidence, reliability, sensitivity, and privacy scope; and
  \item estimated token size and physical storage size.
\end{itemize}

The record should not contain one supposedly universal activation value. A stored baseline may be useful for diagnostics, but the active score must be recomputed for the current query or recorded as a query-time snapshot. Temporal fields must allow the agent to distinguish the time a statement was recorded from the period during which it was true.

\subsection{Access and prompt-inclusion events}
\label{subsec:ch12-access-inclusion}

A \texttt{MemoryAccessEvent} records an attempted retrieval or access operation. It should include the query/task/action context identifier, timestamp, retrieval operation and index version, candidate identifiers and match results, activation-component snapshot, activation algorithm and parameter version, predicted retrieval latency, observed infrastructure latency, and whether retrieval succeeded, missed, or failed a threshold.

A \texttt{PromptInclusionEvent} records what happened during one \gls{contextassembly}. It should include the prompt identifier, memory identifier and version, prompt template and model version, included/rejected status and reason, order index, token start/end, token count, and context-position zone. Position belongs to this event because the same memory can appear near the beginning of one prompt and in the middle of another. Research on long-context behavior shows that model use of relevant information can depend on its position in the input; the pattern is an empirical effect to test for each model and task, not a universal law \autocite{liu2024lostmiddle}.

\subsection{Outcome, attribution, and cost events}
\label{subsec:ch12-outcome-events}

A \texttt{MemoryOutcomeEvent} records the result of the task or action: success, failure, correction, escalation, user confirmation, or another domain-specific outcome. An \texttt{OutcomeMemoryAttribution} links an outcome to zero or more memory identifiers and records the asserted contribution direction, confidence, attribution method, and whether actual model use was observable. One outcome can depend on several memories; one memory can contribute to many outcomes. Attribution should not be treated as ground truth simply because the memory was present in the prompt.

A \texttt{MemoryCostSnapshot} records the cost components relevant to a decision: storage and indexing, embedding or refresh, retrieval and reranking, token/context consumption, measured retrieval and end-to-end latency, and maintenance or privacy cleanup. A \texttt{MemoryValueEstimate} records a query/task/action-conditioned estimate of \gls{memoryvalue} and downstream contribution, uncertainty, estimator version, and calibration state. A \texttt{ProductionUtilityState} stores learned utility for an action or production separately from memory-value estimates. A \texttt{PolicySnapshot} stores the activation and selection algorithm versions, parameters, weights, thresholds, token and latency budgets, concurrency assumptions, policy result, and decision trace.

\begin{table}[tbp]
\centering
\small
\caption{Minimum auditable records for a B1-oriented memory subsystem.}
\label{tab:ch12-production-schema}
\begin{tabularx}{\linewidth}{@{}>{\raggedright\arraybackslash}p{0.29\linewidth}>{\raggedright\arraybackslash}X@{}}
\toprule
\textbf{Record} & \textbf{Purpose and essential fields} \\
\midrule
\texttt{Memory\allowbreak Record} & Stable identity, type, content reference, version, source/provenance, stored/observed times, temporal validity, confidence, privacy, token and storage size. \\
\texttt{Memory\allowbreak Access\allowbreak Event} & Query/task context, retrieval operation, candidates, activation components and version, predicted cognitive latency, observed infrastructure latency, success/miss. \\
\texttt{Prompt\allowbreak Inclusion\allowbreak Event} & Prompt/model/template, memory version, inclusion decision/reason, order, token offsets/count, context position. \\
\texttt{Memory\allowbreak Outcome\allowbreak Event} & Task/action result, evaluation status, observed result, timestamp. \\
\texttt{Outcome\allowbreak Memory\allowbreak Attribution} & Many-to-many links between outcomes and memory records, contribution direction, confidence, attribution method. \\
\texttt{Memory\allowbreak Cost\allowbreak Snapshot} & Storage, indexing, embedding, refresh, retrieval, reranking, token, latency, maintenance, and deletion costs. \\
\texttt{Memory\allowbreak Value\allowbreak Estimate} & Query/task/action-conditioned value, uncertainty, estimator version, calibration state. \\
\texttt{Production\allowbreak Utility\allowbreak State} & Production/action identity, learned utility, reward history, update rate, algorithm version. \\
\texttt{Policy\allowbreak Snapshot} & Selector/activation versions, parameters, thresholds, budgets, concurrency assumptions, selected set, rejection reasons, and trace. \\
\bottomrule
\end{tabularx}
\end{table}

\BookDiagram[0.56\textheight]{Tex/diagrams/generated/part2/chapter12/memory-event-evidence-chain.png}{The auditable event chain separates retrieval, prompt inclusion, observed use, and outcome attribution. Presence in a prompt is not sufficient evidence of use or causal contribution.}{fig:ch12-event-evidence-chain}

The event design supports four statements that must not be conflated:
\begin{enumerate}
  \item The system \emph{retrieved} the memory.
  \item The system \emph{included} the memory in a particular prompt.
  \item The model \emph{used} the memory in its reasoning or answer; this may require evidence such as a citation, trace, or controlled intervention.
  \item The memory \emph{improved the outcome}; this is a causal or counterfactual claim that requires stronger evaluation than co-occurrence.
\end{enumerate}

\subsection{Lifecycle contract and operational safeguards}
\label{subsec:ch12-lifecycle-contract}

A clear subsystem contract can be expressed as a sequence of operations:

\begin{enumerate}
  \item \texttt{extractCandidates(observation)} returns atomic candidate records with source spans and confidence evidence.
  \item \texttt{admitOrUpdate(candidate, store)} applies the admission policy and records whether the candidate was inserted, merged, rejected, or sent for review.
  \item \texttt{maintain(store, now)} checks temporal validity, contradictions, source freshness, privacy scope, and storage policy; it records invalidation, supersession, consolidation, promotion, or eviction.
  \item \texttt{retrieve(query, task)} generates candidates through the configured retrieval paths and records an access event.
  \item \texttt{assembleContext(candidates, budgets)} applies hard gates and set selection, then records an inclusion event for every considered memory.
  \item \texttt{observeOutcome(task, result)} records the outcome and any supported memory-attribution evidence.
  \item \texttt{updateEstimates(events)} updates the appropriate accessibility, value, or production-utility state without conflating their owners or learning signals.
\end{enumerate}

\EngineeringLabel{} Operational safeguards should include idempotent event identifiers, immutable or append-only audit history, versioned policy parameters, bounded retries, deterministic replay of selection where possible, secure deletion for privacy-sensitive content, and a tested path for temporal corrections. A correction should supersede or invalidate an earlier fact rather than merely add a contradictory sentence that future retrieval may select at random. The subsystem must also respect deletion obligations across content, embeddings, graph edges, caches, derived summaries, and backups according to the application's retention policy.

\section{Failure modes and evaluation}
\label{sec:ch12-evaluation}

Memory quality cannot be evaluated by retrieval frequency alone. A memory can be retrieved often because a feedback loop makes it prominent. It can be included in a prompt but ignored by the model. It can be used and still be wrong. It can improve one answer while causing a future failure through stale or misaligned replay. Evaluation must therefore separate the stages and measure downstream task performance.

\subsection{Failure modes to test}
\label{subsec:ch12-failure-modes}

\begin{description}
  \item[Retrieval feedback loops.] Retrieval refreshes recency or heat, which makes the same memory more likely to be retrieved again. Frequency can become self-reinforcing without verified usefulness.
  \item[Error propagation.] An unsupported or mistaken memory is replayed as prior experience and influences later tasks. Xiong et al. study how memory-management choices affect experience-following behavior and identify error propagation and misaligned experience replay as consequential risks \autocite{xiong2026memorymanagement}.
  \item[Temporal invalidity.] A once-correct fact remains semantically similar to the query but no longer describes the world. Validity-time filtering must precede ranking.
  \item[Conflicting or superseded records.] Two versions disagree; without provenance, versioning, and explicit supersession, the selector may include both or choose the older one.
  \item[Low-frequency decisive evidence.] A memory with few visits may be essential for a rare but high-impact task. A policy based only on frequency, heat, or recency can forget it.
  \item[Overconfidence and unsupported content.] A high LLM-assigned importance score cannot replace evidence confidence and provenance.
  \item[Context-position effects.] A relevant memory may be less effectively used at some prompt positions. The Lost in the Middle study demonstrates position sensitivity in its tested long-context tasks and models; a production agent should measure the effect for its own model, template, and task \autocite{liu2024lostmiddle}.
  \item[Attribution error.] Presence is mistaken for use, and use is mistaken for causal benefit. Outcome attribution should carry confidence and method metadata.
  \item[Privacy and lifecycle leakage.] Deleted or restricted content remains in embeddings, summaries, caches, or a derived knowledge graph unless lifecycle controls cover all copies.
\end{description}

\subsection{Benchmarks and what they measure}
\label{subsec:ch12-benchmarks}

LoCoMo evaluates very long-term conversational memory, including multi-session questions and event understanding across extended conversations. It is useful for testing whether an agent preserves and retrieves information over time, but a benchmark score alone does not prove that a selector has reliable causal memory attribution or production resource behavior \autocite{maharana2024locomo}.

LongMemEval evaluates several long-term memory abilities: information extraction, multi-session reasoning, temporal reasoning, knowledge updates, and abstention. Its explicit attention to updates and abstention is especially relevant to a subsystem that must refuse stale or insufficient evidence rather than always return a candidate \autocite{wu2025longmemeval}.

\subsection{Metrics beyond answer quality}\index{memory evaluation}
\label{subsec:ch12-metrics}

A useful evaluation dashboard should include:

\begin{itemize}
  \item retrieval precision and recall, measured against evidence labels;
  \item correct abstention and false-abstention rates;
  \item downstream task success, error rate, escalation rate, and recovery quality;
  \item stale-memory, contradiction, and unsupported-memory rates;
  \item error propagation and misaligned replay rates across later tasks;
  \item admission and retention quality, including redundant storage and loss of decisive memories;
  \item prompt tokens, retrieval time, end-to-end latency p50 and p95, and cost per completed task; and
  \item robustness when facts change, source quality varies, candidate volume increases, or the context budget tightens.
\end{itemize}

Metrics should be reported for the full system and for each stage. For example, high retrieval recall can coexist with low prompt-inclusion recall if the selector is too aggressive, and high inclusion can coexist with low downstream use. Without stage-specific denominators, the system may appear to improve while the actual failure simply moves from retrieval to selection or from selection to generation.

\subsection{Ablations and counterfactual tests}\index{ablation studies}
\label{subsec:ch12-ablations}

Fan and Lan's Appendix D specifies a B1 experiment template: the trace includes activation, predicted latency, memory value, selected memories, and downstream use; primary outcomes include task success, relevant recall, and use precision; resource outcomes include token cost, stale-memory rate, and selector latency. The rejection condition includes a simpler selector matching performance or a fall in relevant recall. The experiments below operationalize and extend that template; they are proposed tests, not completed results.\FanLanCite{Appendix D, B1; § VI-C, B1; Table VI}

A B1 evaluation should test the contribution of each mechanism rather than comparing only one final score. A reasonable sequence of ablations is:

\begin{enumerate}
  \item activation-only selection;
  \item task-value-only selection;
  \item activation plus task value;
  \item add measured infrastructure latency;
  \item add explicit retrieval and context costs;
  \item add set-aware redundancy/complementarity handling;
  \item add position-aware ordering; and
  \item compare alternative feedback rules for retrieval, prompt inclusion, and verified use.
\end{enumerate}

The tests should include single-memory and pairwise cases. In a single-memory test, remove one candidate and measure whether the answer or action changes. In a pairwise test, compare two individually weak memories that may be complementary against two redundant memories. Include wrong-memory contamination tests, temporal-update tests, deletion tests, and missing-evidence cases that should trigger abstention. Vary candidate count, token budgets, concurrency, and deadlines so that the selector is not evaluated only under a convenient small-context condition.

\begin{warningbox}
A memory appearing in a successful prompt does not prove that it caused success. Causal claims require controlled interventions or counterfactual comparisons, and any attribution model should disclose its uncertainty. Outcome feedback that treats every included memory as helpful can reinforce irrelevant or incorrect content.
\end{warningbox}

\section{Boundaries, advanced comparisons, and handoff}
\label{sec:ch12-boundaries}

Several adjacent ideas are useful but should not take over the main argument. Classic cache policies such as Belady's optimal replacement and GreedyDual-Size can clarify resource-aware eviction, but they are cache-management analogies rather than cognitive activation theories. Inverse propensity scoring and doubly robust estimators may help evaluate policies from logged decisions, but they require careful treatment of action propensities, support, and confounding; they belong in an advanced evaluation treatment. FSRS-style scheduling and MaRS-style forgetting taxonomies can inform retention policy, but neither should be presented as a general answer to query-time B1 context assembly.

Neural parametric memory is another boundary. A system such as Titans updates a learned neural memory module rather than selecting only explicit declarative records for a prompt. That is an important memory architecture, but it changes the model's parameterized state and requires a different mechanism analysis. This chapter therefore distinguishes parametric learning from explicit memory-item activation and does not import neural-memory update equations into the B1 main argument \autocite{behrouz2025titans}.

Honda and colleagues' ACT-R-inspired agent-memory work is relevant as an implementation bridge between cognitive activation ideas and LLM agents. The accessible abstract describes a vector-based activation mechanism incorporating temporal decay, semantic similarity, usage frequency, and probabilistic noise. Detailed claims about its exact timestamp semantics or complete equations should be made only against the full paper text \autocite{honda2026actragentmemory}.

\begin{chapterrecap}
\begin{enumerate}
  \item ACT-R separates base-level history, contextual spreading activation, partial matching, and noise. Activation predicts retrieval time, but does not itself provide downstream action value.
  \item Production utility belongs to actions or productions. Memory value belongs to an item or a set of items under a query and possible action.
  \item Admission, retention/forgetting, and query-time context assembly are separate lifecycle decisions. A-MAC, MemoryBank, MemoryOS, Generative Agents, and modern graph retrieval should be classified at the level of the operation they implement.
  \item B1 asks how activation/accessibility, predicted retrieval latency, downstream value, and costs should interact when a constrained context is assembled—and how outcomes should update future estimates.
  \item A credible implementation separates cognitive predictions from infrastructure measurements, sunk costs from marginal costs, individual values from set values, and retrieval from inclusion, use, and causal benefit.
  \item Evaluation must include downstream outcomes, abstention, stale-memory and error-propagation rates, prompt position, token use, latency, and resource cost—not retrieval frequency alone.
\end{enumerate}
\end{chapterrecap}

\begin{reviewquestions}{Review and design questions}
\begin{enumerate}
  \item Why is an admission score not a substitute for query-time context selection? Give an example where the two decisions would choose differently.
  \item What are the four components in the ACT-R activation decomposition, and what does each contribute?
  \item Why is it incorrect to assume that every database read should increment ACT-R base-level presentation history?
  \item Explain the difference between the ACT-R retrieval-latency prediction and infrastructure wall-clock latency.
  \item Why can a selector double-count activation if it rewards both activation and the latency derived from activation?
  \item How does production utility differ from memory value? Which event updates each estimate?
  \item What is the role of the A-MAC threshold, and why is it not a prompt-time relevance threshold?
  \item When would MMR improve a context, and what kind of memory complementarity can it fail to detect?
  \item Which cost categories should be charged at admission, maintenance, retrieval, or prompt assembly?
  \item Design one experiment that distinguishes retrieval, prompt inclusion, actual use, and causal contribution.
\end{enumerate}
\end{reviewquestions}

\section*{Source and example note}
\addcontentsline{toc}{section}{Source and example note}

The historical cognitive mechanism in this chapter is grounded in the ACT-R literature and the ACT-R 7.30+ reference manual. The B1 residual hypothesis, its distinction from B3, the selected modern memory baselines, and the proposed experimental boundary are taken from Fan and Lan's mechanism-level review, especially the ACT-R mapping, the translation analysis of memory from accessibility law to policy action, § VI-C (B1), Tables V--VI, Appendix B, Appendix C-A (R1), Appendix D, and Appendix F-B. The reviewed paper is a bounded mapping review; this chapter does not treat its proposed B1 coupling as an experimentally validated algorithm.\\FanLanCite{§ V-A; § V-B; § VI-C, B1; Tables V--VI; Appendices B, C-A, D, and F-B}

Product X, its support incident, example records, event names, schema, and selector procedure are teaching and engineering examples created for this book. They are not implementations reported by Fan and Lan or by ACT-R. The set-value objective and operational contracts are the book's proposed synthesis and must be evaluated against the named baselines before any effectiveness claim is made.

\section*{Looking ahead}
\addcontentsline{toc}{section}{Looking ahead}

B1 is concerned with which explicit memories become accessible and enter a particular context, and how their estimated value and cost should shape that competition. The next mechanism study turns from memory accessibility to typed failure and recovery: Chapter 13 (B2) examines how an agent identifies an impasse, protects parent state, creates a bounded recovery substate, returns a typed resolution, and learns reusable recovery procedures. B1 remains distinct from B3, which will later address content competition and shared-workspace broadcast.

\chapter{Understanding Typed Impasses, Substates and Compilation: B2}
\label{ch:typed-impasses-substates-compilation}

\section*{Chapter at a glance}
\phantomsection
\addcontentsline{toc}{section}{Chapter at a glance}

An agent can encounter a failure and still have no safe idea what should happen next. A tool call may time out after the remote service has already acted. Two candidate operations may be equally plausible. A workflow may have contradictory preconditions. A recovery attempt may discover useful evidence while the original task state changes underneath it. If the runtime stores only a string such as \code{failed}, it cannot express why progress stopped, what kind of recovery is appropriate, which state must remain authoritative, or what evidence would make a result safe to use.

This chapter studies B2 in Fan and Lan's review: the proposed coupling of typed failure diagnosis\index{typed failure diagnosis}, a protected recovery substate\index{recovery substate}, typed result return, and resolution compilation. Its historical anchor is \gls{soar}\index{Soar}. In Soar, operators\index{operator} are proposed for a state and preferences\index{operator preferences} are evaluated by a decision procedure. When the procedure cannot select a valid next operator, a typed \gls{impasse}\index{impasse} can create a child substate\index{substate} whose problem-solving is related to a parent superstate\index{superstate}. Results created in the child can contribute to the parent; dependencies help keep subgoal reasoning consistent with its context; and chunking\index{chunking} can derive reusable procedural knowledge from the result and its dependencies. The chapter distinguishes this native mechanism from Fan and Lan's proposed modern runtime analogue.

The practical question is not whether a language agent should retry more intelligently. It is how to build a recovery subsystem that records the failure, preserves the parent task, limits what a child process may do, validates what returns, and learns only when a reusable rule has defensible applicability conditions. We develop a concrete reference design with component boundaries, typed contracts, a state machine, persistence and concurrency rules, external-side-effect safeguards, a candidate-rule lifecycle, implementation pseudocode, and tests. Product X\index{Product X} provides a running payment-timeout example. All modern schemas and controller interfaces are book-level engineering proposals, not claims about the exact internal designs of Soar or the modern systems reviewed by Fan and Lan.

\begin{learningbox}
By the end of this chapter, you should be able to explain Soar's operator-selection impasses and subgoal processing; distinguish tie, conflict, constraint-failure, and no-change impasses; describe superstates, substates, results, the Goal Dependency Set\index{Goal Dependency Set}, and chunking; reconstruct B2's four residual invariants; distinguish partial modern baselines from the full proposed coupling; design a typed recovery subsystem with explicit parent/child authority and safe result return; and evaluate the mechanism through state traces, safety metrics, controlled baselines, and ablations.
\end{learningbox}

\section{Failure is not yet a recovery policy}
\label{sec:ch13-failure-policy}

Consider an agent that submits a payment request and receives a timeout. The error message tells us that the client did not receive a response in time. It does not establish whether the provider received the request, whether the provider settled the payment, whether the request was rejected, or whether the response was lost after the side effect occurred. A blind retry can create a duplicate charge. Stopping immediately may leave a task unresolved. A reflection note saying ``check the payment status next time'' may help later, but it does not protect the current task state or constrain what the next recovery action is allowed to do.

A production agent needs to separate at least six operations:

\begin{enumerate}
  \item \textbf{Observe.} Capture what the runtime and tool actually reported, including timestamps, request identifiers, and any partial output.
  \item \textbf{Classify.} Decide which failure condition is supported by the evidence. A raw symptom and an attributed cause are not the same record.
  \item \textbf{Scope.} Determine whether recovery is permitted, which parent state must be preserved, and which operations a recovery process may perform.
  \item \textbf{Recover.} Investigate or repair the problem within the permitted scope.
  \item \textbf{Return and validate.} Produce an explicit result and check that it still applies to the parent task and its current dependencies.
  \item \textbf{Learn selectively.} Determine whether the episode supports a reusable rule, and whether that rule has been validated sufficiently to be used again.
\end{enumerate}

An ordinary retry loop typically combines several of these steps inside one control path. It observes a failed call, perhaps chooses a different parameter, and invokes the tool again. That may be entirely appropriate for a known transient failure with an idempotent operation. It is not enough for an uncertain external outcome, an inconsistent task state, or a recovery that could change what the parent task means. The point is not to replace all retries with a sophisticated subgoal architecture. It is to identify when a retry is merely an action and when the system needs a more explicit recovery control process.

\subsection{The parent task must remain meaningful}

A recovery process is opened because the parent task has encountered a problem. Its authority is therefore derivative: it exists to help the parent proceed, not to silently redefine the parent goal. The parent may contain a customer request, verified facts, constraints, completed work, pending operations, and external side effects already initiated. Recovery can add evidence, but speculative findings should not overwrite verified parent state.

This introduces a distinction that will recur throughout the chapter. A checkpoint is a representation of state that may allow a workflow to resume. A B2-like recovery branch needs more: a typed reason for its creation, a bounded child scope, a relation to an authoritative parent, an explicit result contract, and a control rule that validates the result before the parent resumes. Fan and Lan specifically caution that a checkpoint alone does not establish a Soar-like typed impasse and substate transition.\FanLanCite{Appendix E-B, Checkpoint versus typed impasse; \S VI-C, B2}

\begin{warningbox}
An internal snapshot cannot undo an external action. Restoring the parent state after a timeout does not reverse a payment, sent email, deleted object, or submitted form. A recovery design must track external effects separately and use service-supported idempotency or verification where available. Compensation is a new external action with its own authorization and failure modes, not an automatic consequence of rolling back an internal branch.
\end{warningbox}

\subsection{The B2 question}

The residual question is whether one runtime can couple four properties that current systems, in Fan and Lan's review, mostly implement separately: typed failure diagnosis; a protected child recovery state; a typed result returned to the parent; and compilation of the conditions that made the resolution valid. The failure type should help define the recovery scope. The return contract should determine whether the result can be accepted and what dependencies are relevant to later reuse. The compiled rule should not be a generic reminder extracted from a successful trajectory; it should be guarded by the conditions that explain why the resolution worked.

\BookDiagram[0.52\textheight]{Tex/diagrams/generated/part2/chapter13/b2-four-invariants.png}{The four B2 invariants and their dependency chain. Typed diagnosis scopes recovery; protected state and a bounded child isolate it; typed return exposes evidence; and dependency-aware compilation is considered only after validation.}{fig:ch13-b2-invariants}

\section{Soar: the historical mechanism behind B2}
\label{sec:ch13-soar}

Soar is organized around problem solving as the selection and application of operators to a state. A state represents the current problem-solving situation; an \gls{operator} represents a possible transformation; and the architecture's decision procedure evaluates \glspl{preference} to select an operator. Working memory holds the current state and its relevant working-memory elements. Production rules contribute knowledge and preferences used during the decision cycle. This is enough of Soar to understand the mechanism studied here; the chapter does not attempt to reproduce the full architecture.\autocite{soar2025architecture}

\subsection{Operators and preferences}

Suppose a problem-solving state proposes three operators. The proposal process may establish that each is acceptable, that one is better than another, that one is required, or that another is rejected or prohibited. Preferences are not simply numeric scores. They have defined semantics that the architecture's fixed decision procedure evaluates. If the available preferences distinguish one winning operator, the system can select it and apply it. If they do not, or if they are inconsistent in a way that prevents selection, the decision procedure can report an impasse.

The important point is that the impasse appears at a specific control point: operator selection\index{operator selection}. It is not a general-purpose label attached after any failure in the environment. A tool timeout in a language agent may be a reason for a proposed recovery policy to open a child scope, but it is not literally a Soar operator-selection impasse.

\subsection{Four typed impasses}
\label{subsec:ch13-impasse-types}

Soar's documentation distinguishes four categories. The type tells us what is wrong with the current selection problem; it does not merely name the fact that the system is unhappy.\autocite{soar2025architecture}

\begin{description}
  \item[Tie impasse.]\index{impasse!tie} Two or more acceptable operators remain tied because the current preferences do not distinguish them. The problem is insufficient discrimination among candidates.
  \item[Conflict impasse.]\index{impasse!conflict} Better/worse preferences conflict, for example when one preference says that A is better than B while another says B is better than A, and neither candidate is rejected, prohibited, or required. The problem is inconsistency in the preference relations.
  \item[Constraint-failure impasse.]\index{impasse!constraint failure} The constraints themselves are incompatible: for example, more than one value is required, or a value is both required and prohibited. Additional preferences do not automatically solve a contradiction in the constraints; the conflicting preference conditions must change or retract.
  \item[No-change impasse.]\index{impasse!no-change} The decision procedure selects no new operator. The manual distinguishes a \emph{state no-change} case, where no acceptable or required operator is available (or acceptable candidates have been rejected), from an \emph{operator no-change} case, where an operator may have been selected but no further productions match during application, or no operator is selected in the next decision phase.
\end{description}

These distinctions matter because different causes need different resolution work. A tie may call for further evaluation that makes one option preferred, rejects alternatives, or establishes indifference. A conflict requires removing or superseding incompatible preference relations. Constraint failure requires changing the conflicting constraints, not simply adding another arbitrary preference. A state no-change problem may require proposing a usable operator; an operator no-change problem may require applying the current operator, changing the state so its proposal no longer applies, or proposing and preferring another operator. These are source-specific descriptions of Soar's decision procedure, not a generic taxonomy for all software exceptions.\autocite{soar2025architecture}

\begin{table}[tbp]
\centering
\small
\caption{Soar's typed impasses identify different operator-selection problems.}
\label{tab:ch13-impasse-types}
\begin{tabularx}{\linewidth}{@{}p{0.20\linewidth}p{0.37\linewidth}X@{}}
\toprule
\textbf{Type} & \textbf{What is blocked} & \textbf{Resolution direction} \\
\midrule
Tie & Several acceptable candidates are not distinguished & Add discriminating preference knowledge, reject alternatives, or establish indifference \\
Conflict & Preference relations contradict one another & Remove or resolve the inconsistent preference relations \\
Constraint failure & Required/prohibited constraints are mutually incompatible & Retract or change the conflicting constraints \\
No-change & No new operator can be selected or applied productively & Propose a valid operator, apply the current one, or change the state so selection can proceed \\
\bottomrule
\end{tabularx}
\end{table}

\subsection{A substate is a child problem with a parent relation}
\label{subsec:ch13-substates}

When an impasse is detected, Soar creates a new state to represent a subgoal: the problem of resolving the impasse. The new child contains a description of the impasse and the state in which it arose. From the child's perspective, that higher-level state is its \gls{superstate}. The child is a \gls{substate} because its problem-solving is nested under the parent state. It can propose and apply its own operators. Those operators may be similar to the parent operators, or they may be quite different. An additional impasse can create another subgoal, resulting in a stack of related states.\autocite{soar2025architecture}

This is a more structured relation than a separate process receiving a copied prompt. The child exists to resolve a specific parent-level problem. The parent-child relationship is explicit in the architecture's working-memory representation. It is not accurate, however, to translate the Soar superstate into an immutable database snapshot or an operating-system process. Those may be useful engineering analogies for a modern runtime, but they are not the native Soar semantics.

\subsection{Results: what is allowed to cross the boundary?}
\label{subsec:ch13-results}

A Soar subgoal result is a working-memory element or preference created in the substate and linked directly or indirectly to a superstate. The architecture can identify these elements as results because their creation is related to higher-level working-memory context. When the substate is removed, child-local contents that are not results disappear, while results linked to a higher state remain available to the parent. The result is therefore not any arbitrary child message. It is an element with a defined relationship to the higher-level state.\autocite{soar2025architecture}

An impasse is resolved when processing in the subgoal creates results that let the parent state make progress—for example, when the parent can now select an operator. The child can then be removed, including its substructure and non-result elements. The parent resumes with the result that matters to its own problem-solving.\autocite{soar2025architecture}

\subsection{Resolved, eliminated, and regenerated are different outcomes}
\label{subsec:ch13-impasse-lifecycle}

An impasse does not always disappear because the child successfully solved it. Soar distinguishes three important cases:

\begin{description}
  \item[Resolved.] The subgoal creates results that enable progress in the state where the impasse arose. The parent can make a new decision. This is the result-producing path associated with procedural learning.
  \item[Eliminated.] A higher-level change or external input makes the old subgoal irrelevant. The problem the child was solving no longer matters. The impasse is removed, but this is not the same as resolving it through a result; the manual says Soar does not learn from this elimination path.
  \item[Regenerated.] The child has relied on information in a higher-level context that changes. The old substate is removed and a new one is created so problem-solving can restart against the updated context. The old child reasoning is not treated as a valid resolution of the changed problem.
\end{description}

The \gls{goaldependencyset} (GDS) supports this consistency mechanism. For each subgoal/substate, Soar records relevant higher-level working-memory elements that were tested in the production dependencies leading to persistent, o-supported child structures linked to the substate. If a dependency changes, the substate can be regenerated. Not every fact that changes in a superstate causes regeneration; the dependency set is not the set of all parent facts, but the subset relevant to the supported structures under the manual's rules. The manual also distinguishes i-supported structures and o-supported results in its dependency accounting.\autocite{soar2025architecture}

This is not simply a general-purpose version counter. A version number can tell a modern runtime that some parent field changed. A dependency model asks whether the changed field was actually part of the basis for the child result. A coarse parent hash is safer in some settings but may invalidate more branches than necessary. A precise dependency set is more selective, but it requires the system to record dependencies faithfully.

\subsection{Chunking: deriving procedural knowledge from results}
\label{subsec:ch13-chunking}

Soar's \gls{chunking} is a mechanism for learning procedural knowledge from subgoal processing. A newly created result in a higher-level state can be summarized by a learned production. The conditions of the learned production represent the relevant state elements and dependency chain that supported the result; its action creates the result. Variables allow the learned production to apply in a similar future state, potentially avoiding the need to perform the same deliberation again. This is different from storing a transcript of the child, retaining a natural-language reflection, or caching the final action without its applicability conditions.\autocite{soar2025architecture,soar2025procedural}

There is a timing nuance in the official manual that matters for source fidelity. The architecture chapter describes chunking as a learning opportunity in the context of resolving an operator impasse. The separate procedural-learning chapter says chunks are created when a subgoal creates a result in a superstate, and that chunk generation occurs as soon as the result is created rather than waiting for the substate to be finally removed. These descriptions should not be flattened into a universal rule that learning only starts after the child closes. The relevant mechanism is the result and its dependency analysis during subgoal processing. The same manual distinguishes chunks from justifications and discusses Explanation-based Behavior Summarization (EBBS) as a later treatment of generality and correctness issues; this chapter does not equate the proposed modern rule compiler with all those Soar mechanisms.\autocite{soar2025architecture,soar2025procedural}

The central lesson is that the system does not simply store every successful trajectory. It derives a production from the conditions that support a result. If the learned rule is too broad, it may fire in situations where its action is no longer valid. Dependency analysis and generalization are therefore part of the learning mechanism, not optional annotations added afterward.

\begin{table}[tbp]
\centering
\small
\caption{Soar result handling and chunking: related operations with distinct roles.}
\label{tab:ch13-result-chunking}
\begin{tabularx}{\linewidth}{@{}p{0.23\linewidth}XX@{}}
\toprule
\textbf{Operation} & \textbf{Question answered} & \textbf{What it contributes} \\
\midrule
Subgoal result & What child-produced information is linked to the parent? & A working-memory element or preference that remains relevant after child removal \\
Dependency tracking & Which higher-level facts supported persistent child structures? & A basis for detecting context changes that make the subgoal inconsistent \\
Chunking & Can the result and its supporting conditions become reusable procedural knowledge? & A generalized production that may create a similar result in a future matching state \\
\shortstack[l]{Elimination/\\regeneration} & Is the old subgoal still relevant and consistent? & Removal or recreation of the child without treating stale reasoning as a resolved result \\
\bottomrule
\end{tabularx}
\end{table}

\subsection{A small Soar cycle}

\BookExampleLabel{} Imagine a parent state with two acceptable operators, \code{routeA} and \code{routeB}. The preference structure does not distinguish them, so the decision procedure detects a tie impasse. A child state is created to resolve the choice. Within that child, the system retrieves a current constraint about a route closure and uses it to generate a result preference that favors route B. The result is linked to the parent state. The parent can now select route B, and child-local structures that are not results can be removed. Chunking may derive a production whose conditions include the relevant parent facts that supported the preference, and whose action generates that result.

Now add a counterexample: while the child reasons, the route closure information changes. If the child result depended on the old closure fact, the child reasoning is no longer consistent with the parent context. The system must not treat the old result as if it were a valid answer to the new state. The dependency mechanism allows Soar to regenerate the substate. This is not an unsuccessful retry; it is a response to a changed problem context. The example is a teaching illustration of the documented semantics, not a quoted Soar tutorial scenario.

\BookDiagram[0.60\textheight]{Tex/diagrams/generated/part2/chapter13/soar-impasse-substate-cycle.png}{Soar's operator decision may create a typed impasse and a child substate. Results can support the parent; dependency changes can instead regenerate the child, while an irrelevant impasse can be eliminated.}{fig:ch13-soar-cycle}

\section{Translating the mechanism into a language-agent runtime}
\label{sec:ch13-translation}

Fan and Lan propose an execution-ledger design in which a failure is classified, the relevant parent state is frozen, a bounded recovery subgraph is opened, a typed resolution is returned, and a reusable recovery rule is compiled. They identify this coupling as B2 and keep it separate from B5 because B2 concerns recovery inside an episode, while B5 concerns general skill invocation and governance across tasks.\FanLanCite{\S VI-B, Table V; \S VI-C, B2}

The translation is useful precisely because it is not a literal port of Soar. Soar's trigger is an operator-selection impasse. A modern runtime may open recovery after a tool error, a failed postcondition, or an uncertain external outcome. Soar has a native superstate/substate relationship and working-memory dependency semantics. A modern agent must choose an engineering representation for the parent version, child scope, evidence, permissions, and return validation. Soar results are working-memory elements/preferences linked to higher states; a modern runtime may use a typed resolution object. Soar chunking is a procedural learning mechanism; a modern recovery-rule compiler needs an explicit representation and validation policy.

The four invariants proposed by Fan and Lan can be stated as follows:

\begin{enumerate}
  \item \textbf{Typed diagnosis.} The runtime records why the current execution is blocked and uses that type to define the recovery problem.
  \item \textbf{Protected substate.} Recovery proceeds in a bounded child scope while the parent state remains authoritative.
  \item \textbf{Typed return.} The child returns a structured resolution whose evidence and applicability can be checked against the parent.
  \item \textbf{Resolution compilation.} The system can derive a reusable rule from the conditions that made a validated resolution correct.
\end{enumerate}

These invariants depend on one another. A branch without a typed trigger is only a generic fork. A typed diagnosis without protected state can still corrupt the task it is trying to repair. A protected child without a typed return leaves the parent to interpret arbitrary output. A reusable rule without dependency conditions can replay a repair outside the context where it worked. Fan and Lan's proposal is a residual hypothesis about the coupling, not a report that an existing system has already completed it.\FanLanCite{\S V-A; \S V-B; \S VI-B, Table V; \S VI-C, B2}

\section{Modern systems: the closest pieces are distributed}
\label{sec:ch13-modern-baselines}

The modern comparison is not a leaderboard. Fan and Lan identify partial mechanisms that preserve different edges of the historical control chain. PALADIN and AgentDebugX contribute failure diagnosis and recovery selection. Shepherd provides a reversible, forkable execution trace. AgentHER and Reflexion change later behavior using failed experience. The review's conclusion is that these mechanisms cover diagnosis, protected branching, and experience-based learning in separate systems; the full typed-impasse, protected-substate, typed-return, condition-safe-compilation coupling remains the residual.\FanLanCite{\S V-A, Soar to typed recovery and experience learning; Appendix C-B--C-D}

\subsection{Diagnosis and recovery}

PALADIN classifies tool failures and retrieves recovery actions. AgentDebugX organizes a Detect--Attribute--Recover--Rerun loop. These approaches make the failure type or attributed cause relevant to the recovery choice, which preserves an important control edge: different failures can lead to different recovery actions. In Fan and Lan's ledger, the corresponding typed-recovery record is provisionally coded E3/D3. That code applies to the specific mapped edge, not to the entire system or B2 as a whole.\FanLanCite{Appendix C-B, R2a; Appendix F-B}

\subsection{Reversible execution and later learning}

Shepherd makes execution traces first-class and supports reversible operations such as forking and replaying. Fan and Lan map that mechanism to the protected-branch edge and provisionally code the correspondence E2/D3. Its missing properties in the review are typed impasse creation, typed result return, and condition-safe compilation. AgentHER and Reflexion make failed experience available to change later behavior. Their mapped learning edge is provisionally E3/D4, but it is not the same as compiling the necessary dependencies of a scoped Soar subgoal result.\FanLanCite{Appendix C-C, R2b; Appendix C-D, R3; Appendix F-B}

\subsection{Runtime boundaries are insertion points, not proof}

LangGraph, OpenHands, Magentic-One, and Shepherd provide state, orchestration, execution, or branching boundaries at which a recovery design could be hosted. A checkpoint or graph state can preserve workflow information. A reversible trace can provide a branch substrate. An orchestration loop can coordinate agents or actions. None of these facts alone establishes typed failure diagnosis, protected child-state semantics, typed return, and safe compilation. The chapter therefore treats the runtime as the place where the control law could be implemented, while the specialized systems are partial baselines for particular edges.\FanLanCite{\S IV-A--IV-C; \S V-A; Appendix E-B}

\begin{table}[tbp]
\centering
\small
\caption{Modern B2-related systems preserve different parts of the control chain in Fan and Lan's review.}
\label{tab:ch13-modern-baselines}
\begin{tabularx}{\linewidth}{@{}p{0.19\linewidth}p{0.30\linewidth}p{0.20\linewidth}X@{}}
\toprule
\textbf{Case} & \textbf{Preserved edge} & \textbf{Provisional record} & \textbf{Not established as a complete B2 chain} \\
\midrule
PALADIN & Failure classification selects a recovery action & E3/D3, typed recovery & Protected child state, typed return, condition-safe compilation \\
AgentDebugX & Detection and attribution inform recovery and rerun & E3/D3, typed recovery & Parent isolation, typed result return, safe compilation \\
Shepherd & Reversible trace, fork, and replay & E2/D3, protected branch & Typed impasse, parent-linked result, dependency-safe rule compilation \\
AgentHER / Reflexion & Failed experience changes later behavior & E3/D4, learning edge & Compilation from the dependency structure of a scoped resolution \\
\bottomrule
\end{tabularx}
\end{table}

The E and D codes are interpreted as in Chapter 11: evidence relation and migration depth answer different questions. They are neither additive nor one combined maturity score. The codes above are copied as case-specific provisional records from Fan and Lan; this chapter does not independently recode the source corpus.\FanLanCite{\S II-D; Appendix C-B--C-D}

\BookDiagram[0.48\textheight]{Tex/diagrams/generated/part2/chapter13/b2-modern-baseline-map.png}{Reviewed modern cases preserve different B2-related edges: diagnosis/recovery, reversible branching, or later behavior change.}{fig:ch13-modern-baselines}

\section{B2 as a composed runtime hypothesis}
\label{sec:ch13-b2-hypothesis}

Fan and Lan state B2 as five linked operations: classify the failure, preserve the relevant parent state, open a bounded recovery subgraph, return a typed resolution, and compile only the conditions that made the resolution valid. The chain is a source-level residual hypothesis, not a prescribed software pipeline. The key boundaries are that diagnosis does not grant action authority, a child does not own the parent goal, a returned result remains subject to validation, and a successful episode does not automatically justify a reusable rule. The detailed responsibilities, records, state transitions, persistence rules, and validation contracts follow in the implementation design below.\FanLanCite{\S VI-C, B2}

\section{From B2 to an implementable recovery subsystem}
\label{sec:ch13-implementation}

The theory-to-practice bridge is where the four invariants become software responsibilities. The design below is an \EngineeringLabel{} reference pattern, not an implementation specification contained in Fan and Lan's review. It is intentionally language-neutral; the interfaces can be expressed in Java records and services, Kotlin data classes, or another typed application language.

\subsection{Component architecture and authority}
\label{subsec:ch13-components}

A minimal architecture separates observation, policy, local recovery, parent commit, and learning. That separation matters because the component that proposes a diagnosis should not also be the only authority that decides whether a state-changing recovery is safe.

\begin{description}
  \item[Execution adapter.] Wraps a tool or action boundary and emits a normalized execution event. It records request identity, raw outcome, timestamp, tool version, and any provider-side operation identifier.
  \item[Failure classifier.] Converts an event and available evidence into a typed finding. The classifier can use deterministic rules, a model, or both. A model's suggested label is evidence for the governor, not the authority to open a branch.
  \item[Parent-state store.] Maintains the authoritative parent task state and creates a versioned view for a recovery episode. It also references verified facts, goal invariants, pending actions, dependency versions, and the external side-effect ledger.
  \item[Recovery governor.] Checks eligibility, selects a recovery policy, creates the child scope, and authorizes its capability allowlist. It owns state transitions.
  \item[Recovery runner.] Executes the child subgraph in a separate namespace. The runner can inspect permitted evidence and return a proposed resolution; it cannot directly mutate the parent's state.
  \item[Resolution validator.] Checks result type, evidence, expected parent version, dependency freshness, postconditions, and side-effect claims.
  \item[Parent committer.] Applies a validated state delta only if the expected parent version is still current. It rejects stale results rather than silently merging incompatible state.
  \item[Recovery ledger.] Stores an append-only trace connecting event, classification, parent snapshot, branch operations, child result, validation, commit, and any rule candidate.
  \item[Rule compiler and registry.] Derive candidate guards/actions from validated traces, test them, and maintain quarantine, shadow, promotion, invalidation, and rollback states.
\end{description}

The model can help interpret events, propose a diagnosis, or suggest a recovery action. The governor controls branch creation. The validator controls result acceptance. The committer controls parent mutation. The rule registry controls promotion. These authorities should be explicit in code and tests, rather than implied by prompt instructions.

% The source-level BookDiagram call is the validator's authoritative link
% between the accessible Mermaid source and this chapter's generated PNG.
% Keep the figure in reading order before the typed-contract listings.
\BookDiagram[0.46\textheight]{Tex/diagrams/generated/part2/chapter13/b2-reference-implementation.png}{A proposed implementation separates model suggestions from runtime-controlled branch creation, result validation, parent commit, and recovery-rule promotion.}{fig:ch13-reference-architecture}

\subsection{Typed contracts}
\label{subsec:ch13-contracts}

The following records are \EngineeringLabel{} proposals. They are not native Soar objects and are not schemas supplied by Fan and Lan. Their purpose is to make each handoff explicit enough to implement, validate, and audit.

\paragraph{ExecutionFailure.}

This record captures what the action boundary observed before diagnosis begins:

\begin{lstlisting}[language=Java,caption={Raw action-boundary event (illustrative Java record).},label={lst:ch13-execution-failure}]
record ExecutionFailure(
    UUID eventId,
    UUID parentId,
    UUID operationId,
    Instant observedAt,
    String rawError,
    OutcomeCertainty outcomeCertainty,
    String toolName,
    String toolVersion,
    List<EvidenceRef> evidenceRefs,
    String externalOperationRef
) {}
\end{lstlisting}

The key field is \code{outcomeCertainty}. A timeout can mean that the caller lacks a result, not that the remote operation did not happen. The raw event should be immutable; later attribution must not rewrite what the adapter actually observed.

\paragraph{FailureFinding.}

A typed finding adds interpretation: a failure category, execution phase, evidence, attribution status, uncertainty, classifier version, retryability, and side-effect risk. A category such as \code{payment\_outcome\_unknown} is useful only if its evidence and semantics are defined. The category should not be a free-form model string that silently changes meaning across model versions.

\paragraph{ParentSnapshot.}

The parent view should contain a stable parent ID, version, state digest, goal reference, invariant references, verified facts, pending operation, dependency versions, side-effect ledger reference, and capture time. A state digest detects changes to the captured representation; a dependency list supports more selective invalidation. Neither should be confused with a complete transaction or rollback of the external world.

\paragraph{RecoveryScope.}

The child scope carries a recovery ID, parent ID/version, failure ID, status, allowed and denied capabilities, local attempt budget, deadline, local state reference, dependency set, and opening time. The scope should explicitly record what the child can read, what it can write locally, and what actions require a separate authorization. A read-only recovery branch is the safest initial vertical slice.

\paragraph{Resolution.}

A resolution envelope should contain an ID, kind, parent version, evidence references, dependency versions, proposed delta, postconditions, side-effect summary, and creation time. Outcome kinds may include \code{RESOLVED}, \code{STILL\_UNKNOWN}, \code{BLOCKED}, \code{PARENT\_STALE}, and \code{ESCALATE}. The result is a typed claim supported by evidence; it is not a natural-language report that the parent must interpret without validation.

\paragraph{RecoveryRuleCandidate.}

A candidate rule records its guard, action/plan reference, dependency set, source trace references, tool/API/model versions, validation suite, promotion status, expiry/invalidation conditions, and rollback target. A candidate must be distinguishable from a rule already authorized for live execution.

\begin{table}[tbp]
\centering
\small
\caption{Proposed B2 records and their ownership boundaries.}
\label{tab:ch13-contracts}
\begin{tabularx}{\linewidth}{@{}p{0.22\linewidth}p{0.31\linewidth}X@{}}
\toprule
\textbf{Record} & \textbf{Created by} & \textbf{What it must not be confused with} \\
\midrule
ExecutionFailure & Action/tool adapter & Attributed root cause \\
FailureFinding & Classifier plus evidence & Raw error text or authorization to act \\
ParentSnapshot & Parent-state store & A rollback of external side effects \\
RecoveryScope & Recovery governor & A generic checkpoint or unbounded agent fork \\
Resolution & Recovery runner, then checked by validator & Arbitrary child text or automatic parent mutation \\
\shortstack[l]{RecoveryRule\\Candidate} & Compiler over a validated trace & A promoted/live rule \\
\bottomrule
\end{tabularx}
\end{table}

\subsection{Runtime state machine}
\label{subsec:ch13-state-machine}

A modern recovery runtime needs legal states and guarded transitions\index{recovery state machine}. The names below are a proposed implementation vocabulary, not Soar terminology:

\begin{description}
  \item[RUNNING.] The parent task is executing normally.
  \item[FAILURE\_CAPTURED.] An action boundary emitted a failure or ambiguous outcome event.
  \item[CLASSIFIED.] The classifier produced a typed finding with evidence and uncertainty.
  \item[PARENT\_PROTECTED.] The parent version, invariants, pending operation and relevant dependencies have been captured.
  \item[RECOVERY\_OPEN.] The governor created a bounded child scope and granted capabilities.
  \item[RECOVERING.] Child-local investigation or permitted recovery actions are running.
  \item[RESULT\_PROPOSED.] The child returned a typed resolution envelope.
  \item[RESULT\_VALIDATED.] Evidence, postconditions, parent version and dependencies passed validation.
  \item[PARENT\_COMMITTED.] A permitted state delta was applied against the expected parent version.
  \item[ESCALATED / ABORTED.] Local recovery could not safely continue or policy rejected it.
  \item[ELIMINATED / REGENERATED.] The recovery problem became irrelevant or a dependency changed so the child must be discarded and recomputed. These are engineering analogues informed by Soar's distinctions, not an assertion that every runtime must use these exact labels.
\end{description}

No transition may jump directly from \code{RESULT\_PROPOSED} to \code{PARENT\_COMMITTED}. A typed result is only a proposal until the validator has checked the evidence, state version and postconditions. Compilation is a separate learning path after a result and its outcome evidence have been recorded; it is not a mandatory terminal action of every recovery episode.

\begin{table}[tbp]
\centering
\small
\caption{Guarded transitions in the proposed recovery controller.}
\label{tab:ch13-state-transitions}
\begin{tabularx}{\linewidth}{@{}>{\raggedright\arraybackslash}p{0.28\linewidth}>{\raggedright\arraybackslash}p{0.34\linewidth}X@{}}
\toprule
\textbf{Transition} & \textbf{Required guard} & \textbf{Resulting effect} \\
\midrule
\shortstack[l]{RUNNING to\\FAILURE\_CAPTURED} & Adapter observes a defined failure or ambiguous outcome & Append immutable execution event \\
\shortstack[l]{FAILURE\_CAPTURED to\\CLASSIFIED} & Finding type and supporting evidence recorded; uncertainty not hidden & Recovery eligibility can be evaluated \\
\shortstack[l]{CLASSIFIED to\\PARENT\_PROTECTED} & Parent ID/version and goal invariants available & Capture authoritative parent view and dependencies \\
\shortstack[l]{PARENT\_PROTECTED to\\RECOVERY\_OPEN} & Governor approves type, scope and capabilities & Create bounded child namespace \\
\shortstack[l]{RECOVERING to\\RESULT\_PROPOSED} & Child emits a recognized result kind and evidence envelope & No parent mutation yet \\
\shortstack[l]{RESULT\_PROPOSED to\\RESULT\_VALIDATED} & Evidence, postconditions, parent version and dependencies pass & Resolution may be considered for commit \\
\shortstack[l]{RESULT\_VALIDATED to\\PARENT\_COMMITTED} & Compare-and-swap succeeds and delta remains authorized & Parent state advances exactly once for that accepted version \\
\shortstack[l]{Any active child\\to ELIMINATED /\\REGENERATED} & External context makes task irrelevant, or a material dependency changes & Discard stale branch; restart or stop without learning from stale result \\
\shortstack[l]{Any unsafe or\\unsupported path to\\ESCALATED /\\ABORTED} & Budget exhausted, permissions absent, evidence insufficient, or invariant threatened & Preserve trace; do not force a parent update \\
\bottomrule
\end{tabularx}
\end{table}

\subsection{Persistence, concurrency, and external effects}
\label{subsec:ch13-persistence}

A robust implementation needs two persistence domains. The branch-local state can change while the child reasons. The parent remains authoritative and advances only through an accepted commit. An append-only ledger records transitions and evidence; it should not be the same object as a mutable prompt transcript. A branch record can reference a parent version and store local findings, proposed deltas, and tool observations without changing the parent's verified facts.

At return, the parent committer should verify an expected version. If the current parent version is still the version against which the child reasoned, and the dependencies it used remain valid, a compare-and-swap\index{compare-and-swap} update can apply the result. If the version differs, the system should inspect whether the changed fields intersect the child's dependency set. A whole-state hash is a safe, coarse first implementation: any change rejects the child result. More precise dependency versions reduce unnecessary invalidation, but they require a reliable dependency model. The trade-off should be measured rather than assumed.

Concurrency makes this more important. Two recovery children may be opened from the same parent version. Both may return plausible results. The runtime cannot commit both if they make incompatible changes. One result may win the compare-and-swap; the other must be revalidated against the new parent or regenerated. The commit path should serialize transitions per parent version or provide an explicit arbitration policy. A child branch must never silently merge a stale patch merely because its JSON is syntactically valid.

External effects\index{external side effects} require a separate ledger. Suppose a payment call times out. Retrying the same request may be safe only if the provider's idempotency contract guarantees that repeated requests with the same idempotency key represent the same logical operation. If that contract is absent or unknown, the recovery child should query the original operation's status before initiating another side effect. A local rollback cannot reverse the provider's action. A compensating transaction is another operation, which may fail and requires its own authority.

Event processing also needs idempotency\index{idempotency}. A timeout callback may be delivered more than once; an adapter may emit a failure event and later receive a delayed success response. Stable event IDs and operation IDs allow the ledger to deduplicate observations. They do not create universal exactly-once semantics for external APIs. The system should state what guarantee it actually has—at-most-once dispatch, at-least-once dispatch with provider idempotency, or uncertain delivery—and test the behavior under duplicated and delayed events.

\subsection{Validating a result before returning it}
\label{subsec:ch13-result-validation}

A result validator should check at least the following:

\begin{enumerate}
  \item the result kind is recognized and allowed for this failure category;
  \item each decisive claim points to evidence with provenance and a source/tool version;
  \item the parent ID and expected version match the branch's captured parent;
  \item dependencies used by the child remain current;
  \item the postconditions are independently checked where possible;
  \item any proposed state delta is within the child capability boundary;
  \item external side effects are accounted for and were authorized; and
  \item uncertainty remains explicit, with a defined reject/escalate outcome.
\end{enumerate}

The last point prevents a dangerous collapse: a child may successfully diagnose that the settlement state is unknown. That is a useful resolution of the diagnostic question, but it is not permission to declare the payment settled or to retry it. A typed result can resolve one subproblem while returning an unresolved status for the parent's larger task.

\subsection{Dependency-aware compilation and rule governance}
\label{subsec:ch13-rule-compiler}

A modern analogue of chunking should start from the validated result and its dependencies\index{dependency-aware compilation}. A practical compiler pipeline can:

\begin{enumerate}
  \item collect the branch trace and the validated resolution;
  \item identify the smallest set of parent facts, tool contracts, and environmental conditions needed for the result;
  \item express those conditions as explicit guards;
  \item bind a bounded action or plan reference to the guard;
  \item validate the candidate on held-out traces, counterexamples and current tool contracts;
  \item retain it in quarantine or shadow mode before live use; and
  \item promote it only under a stated policy, retaining versions and a rollback path.
\end{enumerate}

The compiler should not turn an LLM-generated narrative into executable authority. It should produce a reviewable intermediate representation: predicates, version constraints, an allowed plan reference, postconditions, and a maximum action scope. A candidate can be rejected when the trace's causal dependencies are unclear, when the outcome was uncertain, or when the resolution depended on human intervention. The absence of a safe rule is a valid result.

Overgeneralization is the central learning risk. A rule learned from one payment provider may not apply to a second provider; a status endpoint can change semantics after an API version upgrade; an idempotency guarantee may expire; and a recovery that was safe for one customer/account state may be unsafe for another. A compiled candidate must therefore include its environment/tool version and invalidation conditions. The registry should support disablement and rollback when a counterexample appears\index{recovery rule!promotion and rollback}. This is a book-level governance proposal inspired by Soar's dependency-based learning, not a claim that Fan and Lan prescribe this exact compiler design.

\subsection{A practical vertical-slice implementation plan}
\label{subsec:ch13-vertical-slice}

An engineering team should not implement all four invariants at once. The following sequence introduces observability and authority before learning:

\begin{enumerate}
  \item instrument action boundaries and persist immutable execution events;
  \item introduce a small typed failure taxonomy, including explicit unknown/ambiguous outcomes;
  \item capture parent versions and goal invariants before state-changing actions;
  \item open isolated child scopes with a read-only capability allowlist;
  \item return a typed resolution envelope and validate it before the parent resumes;
  \item add version-checked parent commits and stale-result rejection;
  \item record dependency references and test regeneration/invalidation;
  \item store recovery-rule candidates, initially disabled;
  \item evaluate candidates in replay or shadow mode, then promote only under a documented policy; and
  \item measure the full B2 trace and ablations before claiming that the coupled mechanism adds value.
\end{enumerate}

A reasonable project layout might place action adapters and normalized failures in an execution package, parent/child state contracts in a recovery domain package, the governor and validator in an application layer, persistence adapters in an infrastructure layer, and the candidate-rule pipeline behind a separate learning boundary. This prevents a model wrapper from becoming the implicit owner of state transitions. The exact package names are project decisions; the important requirement is that the commit and promotion authorities remain explicit and testable.

\subsection{Reference pseudocode}
\label{subsec:ch13-pseudocode}

The following Java-like pseudocode shows the intended authority boundaries. It is illustrative rather than a drop-in implementation. In particular, the model can suggest a failure type or recovery plan, but only the governor can open a child and only the committer can advance the parent.

\begin{lstlisting}[language=Java,caption={B2-style recovery orchestration (illustrative pseudocode).},label={lst:ch13-recovery-orchestration}]
RecoveryOutcome onFailure(ParentId parentId,
                          ExecutionFailure event) {
    ledger.append(event);

    FailureFinding finding = classifier.classify(event);
    ledger.append(finding);

    RecoveryDecision decision = governor.evaluate(finding);
    if (!decision.allowed()) {
        return escalate(parentId, finding, decision.reason());
    }

    ParentSnapshot parent = parentStore.capture(parentId);
    RecoveryScope scope = governor.open(parent, finding, decision);
    ledger.append(scope);

    Resolution proposed = recoveryRunner.run(scope);
    ledger.append(proposed);

    ValidationResult check = validator.validate(parent, scope, proposed);
    ledger.append(check);
    if (!check.accepted()) {
        return escalate(parentId, finding, check.reason());
    }

    CommitResult committed = committer.applyIfVersion(
        parent.id(), parent.version(), check.approvedDelta());
    ledger.append(committed);

    if (committed.accepted()) {
        ruleCompiler.observeValidatedTrace(scope, check, committed);
    }
    return RecoveryOutcome.from(committed);
}
\end{lstlisting}

The call to \code{observeValidatedTrace} records a candidate-learning opportunity; it does not directly promote a rule. Promotion requires a distinct validation and registry decision. The code also omits retries around transient storage failures, transaction-outbox implementation details, and error-specific policies; those belong in a production implementation, not in a short illustration.

\subsection{Implementation tests that expose the missing edges}

A B2 implementation should include tests at every contract boundary:

\begin{itemize}
  \item \textbf{Stale parent.} Open a child at parent version 7, advance the parent to version 8, and ensure the child cannot commit a patch based on version 7.
  \item \textbf{Dependency-only change.} Change an unrelated parent field and compare coarse hash invalidation with a dependency-aware validator; confirm that only dependencies declared by the child govern selective reuse.
  \item \textbf{Duplicate failure event.} Deliver the same event ID twice and verify that the ledger deduplicates it and does not open two independent recovery branches.
  \item \textbf{Unknown external outcome.} Simulate a payment timeout after the provider has committed the charge. The child must query status or return \code{STILL\_UNKNOWN}; it must not create a second charge.
  \item \textbf{Capability violation.} Have the child request a forbidden write or payment operation and verify that the governor/adapter rejects it regardless of model output.
  \item \textbf{Nested branch loop.} Make a child encounter a repeated failure and confirm that local depth/attempt limits lead to escalation rather than infinite recursion.
  \item \textbf{Rule overgeneralization.} Change a tool contract or API version and verify that a candidate's guard invalidates or requires revalidation before reuse.
  \item \textbf{False diagnosis.} Supply evidence consistent with two failure classes and ensure uncertainty is retained or escalated rather than silently promoted to a single cause.
\end{itemize}

\section{Product X: a worked implementation from timeout to guarded rule}
\label{sec:ch13-productx}

\BookExampleLabel{} Reuse the Product X support agent from Chapter 12. A customer has asked the system to complete a payment. The agent submits the original payment operation with a provider-supported idempotency reference. The request times out. The incident details in this section are illustrative; they are not evidence about a real payment provider.

The parent task records the customer request, the original operation ID, the invariant ``do not create a second charge'', verified account information, the pending payment action, and the current parent version. The execution adapter emits an \code{ExecutionFailure} event: the request timed out, the provider operation reference is known, and the outcome certainty is unknown. The classifier returns \code{payment\_outcome\_unknown}, not \code{payment\_failed}. The distinction prevents the governor from treating a missing response as evidence that the provider did nothing.

The governor opens a child recovery scope. Its capability allowlist contains read-only status lookup for the original operation and retrieval of the provider's operation metadata. It explicitly excludes new payment creation, cancellation, refund, and any mutation to customer account state. The child can assemble evidence and return a typed resolution. It cannot change the parent or issue a second payment.

Three outcomes are possible:

\begin{description}
  \item[SETTLED.] The provider returns authoritative evidence that the original operation settled. The child returns the evidence reference and provider operation version. The parent validator checks that it is the same operation and that the parent version/dependencies are current. The parent records completion and does not submit another payment.
  \item[NOT\_SETTLED.] The provider returns authoritative evidence that the original operation did not settle. The parent can now consider the permitted retry policy. The new action remains a parent-level operation subject to its ordinary approval and idempotency rules; it is not performed merely because the child found a result.
  \item[STILL\_UNKNOWN.] The status endpoint is unavailable or returns insufficient evidence. The child returns uncertainty with the failed query details. The parent escalates or waits; it does not guess that the payment failed.
\end{description}

\BookDiagram[0.62\textheight]{Tex/diagrams/generated/part2/chapter13/productx-b2-recovery-example.png}{An illustrative Product X payment-timeout recovery. A read-only child checks the original operation, returns a typed status, and leaves the parent to validate and decide.}{fig:ch13-productx}

The result is not yet a compiled rule. Suppose several independent incidents later exhibit the same provider-specific timeout behavior and the status query reliably distinguishes settled from not-settled outcomes. A compiler may propose a guarded recovery rule: when the provider name and API version match, the operation carries the relevant idempotency reference, the raw error is the known timeout class, and the provider's status API confirms non-settlement, return a \code{NOT\_SETTLED} resolution to the parent. The rule should not match an unknown provider, missing operation reference, changed API contract, or ambiguous status response.

A counterexample makes the guard's importance clear. Suppose the status API changes its semantics and now returns a cached pre-settlement state. A rule that ignores API version or status freshness may authorize a duplicate attempt. The candidate must be disabled or invalidated when the provider contract changes. The episode may still be useful as a failure trace, but it is no longer evidence that the old rule remains correct.

This example demonstrates the bridge from mechanism to practice. The failure is typed; the parent and its invariant are preserved; the child scope is bounded and read-only; the returned result is typed and evidence-backed; the parent validates before it acts; and a reusable rule is only considered when the conditions that made the resolution correct can be stated and tested.

\section{Represent B2 through the seven fields}
\label{sec:ch13-seven-fields}

Chapter 10 introduced the fixed tuple
\[
  K = \langle S,C,T,P,F,L,R\rangle,
\]
where state, control, transition, persistence, failure, learning, and resources/governance are analytical fields. We use the tuple exactly as defined there; it is not a claim that each field maps to one software module.

\begin{table}[tbp]
\centering
\small
\caption{B2 represented through the book's seven-field mechanism vocabulary.}
\label{tab:ch13-seven-fields}
\begin{tabularx}{\linewidth}{@{}p{0.15\linewidth}p{0.39\linewidth}X@{}}
\toprule
\textbf{Field} & \textbf{Soar mechanism} & \textbf{B2 runtime hypothesis / engineering design} \\
\midrule
State (S) & Working-memory state, operators, preferences, substates/superstates, results and dependency information & Execution ledger, parent version, typed failure, child recovery trace, typed resolution and guarded rule candidate \\
Control (C) & Decision procedure and production/operator-selection machinery & Recovery governor authorizes classification-driven branch creation, result acceptance and candidate promotion \\
Transition (T) & Operator selection or typed impasse; result/dependency/context change & Failure capture, classification, branch opening, validation, parent commit, invalidation and later rule matching \\
Persistence (P) & Child-local structures are removed; results linked to higher states remain; learned chunks persist & Child-local namespace, append-only trace, authoritative parent, versioned rule registry \\
Failure (F) & Typed impasse, resolution, elimination and regeneration & Misdiagnosis, unresolved child, stale parent/dependency, invalid result, forbidden side effect or unsafe rule reuse \\
Learning (L) & Chunking derives a production from a result and its supporting dependencies & Proposed guarded recovery-rule extraction and outcome validation; distinct from cross-episode reflection alone \\
\shortstack[l]{Resources/\\governance (R)} & Report only documented Soar behavior; do not invent a resource policy & Bound local recovery operations; wider budgets, uncertainty and interruption policy remain B6's subject \\
\bottomrule
\end{tabularx}
\end{table}

The tuple makes an important point visible. A system may expose checkpoints or branch operations under state and persistence, while leaving the transition trigger and control authority unspecified. It may classify a tool error and select a retry, while leaving parent isolation and result return unspecified. It may learn from a failed trajectory but not capture the conditions that made a specific child result valid. These are not minor implementation details; they are the differences between the component edges and the proposed B2 coupling.

\section{Failure modes, safety and authority}
\label{sec:ch13-safety}

\subsection{Misdiagnosis and uncertainty collapse}

A false classification can open an unsafe recovery path. Treating an unknown payment outcome as a confirmed failure is a concrete example. A classifier should preserve the original event and evidence, record its model/rule version, and return uncertainty when multiple explanations remain possible. Evaluation should measure per-type precision and recall, but it should also measure the cost of a wrong type because not all errors have equal consequences.

\subsection{Parent corruption and stale dependencies}

The child may accidentally mutate shared state, return a result tied to an old parent version, or ignore a dependency that changed during recovery. Namespace isolation, immutable parent views, version checks, dependency tracking and compare-and-swap commits reduce these risks. A coarse whole-state hash is easier to implement but can trigger unnecessary regeneration. A fine-grained dependency map may reduce wasted work but can itself be incomplete. Both designs must be tested against injected changes.

\subsection{External side effects and duplicate execution}

An internal branch cannot undo the outside world. Tools can commit effects even when the agent does not receive a response. A system must record whether an operation is read-only, idempotent, compensatable, irreversible, or of unknown semantics. For state-changing operations, idempotency keys and status verification are preferred where the external service supports them. A compensating action must be separately authorized and logged. The recovery child should default to the smallest capability set that can resolve the diagnostic question.

\subsection{Unsafe compilation and rule proliferation}

An incorrect diagnosis can be compiled into a persistent shortcut. A missing dependency can make the guard too broad. A tool/API change can invalidate a previously correct rule. An environment-specific repair can fail in another account or tenant. Rule candidates therefore need provenance, version constraints, tests, shadow evaluation, promotion criteria and rollback. The system also needs a way to decline compilation when evidence is insufficient. More stored rules do not necessarily mean better recovery.

\subsection{Unbounded recursion and authority confusion}

A child that opens a child that opens another child can recurse indefinitely. Local attempt, depth and operation limits are necessary, with escalation as a terminal outcome. These are B2-local safeguards; Chapter 16 examines wider resource allocation and interruption. Authority must also be explicit: a language model may propose what to do, but the runtime governor and policy checks must decide whether an action is allowed, and a separate committer must decide whether a returned state change is valid.

\section{Evaluation: test the full coupling}
\label{sec:ch13-evaluation}

Fan and Lan specify the required B2 trace fields as failure type, parent hash, substate trace, returned resolution and compiled rule. Their primary outcomes include recovery by failure type, repeated-failure rate and transfer to later similar cases. Safety and efficiency measures include diagnosis latency, parent-state corruption and unsafe rule reuse. The rejection condition is that a simple retry baseline matches the recovery outcome or compilation causes harm. The appendix also requires matched model, prompt, tools, environment and task budget, preregistration, and ablations that remove the new state object, authoritative control edge, or learning update.\FanLanCite{Appendix D, Table X, B2; \S VI-C, B2}

\subsection{Required state trace}

Every evaluation episode should record the raw execution event, typed finding, parent version/hash, child scope, branch trace, capability use, dependency set, returned resolution, validation result, commit outcome, compiled candidate and any subsequent rule invocation. The fields required by the paper should remain easy to extract; the additional details are engineering instrumentation needed to explain why a result was accepted or rejected.

\subsection{Baselines and comparisons}

A controlled comparison should include:

\begin{enumerate}
  \item a simple retry policy;
  \item typed diagnosis and recovery selection without a protected child state;
  \item reversible branch/replay without typed return validation;
  \item cross-episode reflection or hindsight learning without dependency-safe compilation;
  \item the complete proposed B2 composition; and
  \item human-supervised escalation for cases where the risk makes automatic recovery inappropriate.
\end{enumerate}

Each baseline isolates a control edge. It should not be falsely represented as a single named product unless the evidence supports that attribution.

\subsection{Metrics}

Report recovery success by failure type, repeated-failure rate, transfer to a later similar case, diagnosis latency, total recovery latency, number of recovery operations, cost, escalation precision, parent-state corruption, invariant violations, harmful rule reuse and compile candidate rejection. Define each denominator. For example, repeated-failure rate should count only cases in which a genuinely comparable later failure occurs; transfer should measure whether the compiled rule applies correctly to a held-out similar situation. A system that reduces retries by abandoning more tasks is not necessarily improving recovery.

\subsection{Ablations and rejection criteria}

The required ablations\index{ablation studies} separate the contribution of state, control and learning:

\begin{itemize}
  \item Remove the new parent/child state object but retain the rest of the logic. This tests whether explicit substate representation contributes beyond the controller code.
  \item Remove the authoritative transition while retaining logs. This tests whether the runtime control edge, rather than instrumentation alone, changes behavior.
  \item Remove the learning update while retaining immediate recovery. This separates in-episode recovery from later reuse and transfer.
\end{itemize}

Add counterfactual tests for stale parent versions, missing dependencies, duplicated events, uncertain side effects, forbidden capabilities and changed tool contracts. Pre-register the task population, experimental unit, base model and dependency versions, stochastic replications, sample size, metric formulas, equivalence/non-inferiority margin, tests or uncertainty intervals, exclusions and stopping rule. A simpler baseline matches only if it is non-inferior under the pre-registered primary margin while using no more resources. The source paper supplies this protocol as a future experiment, not completed experimental results.\FanLanCite{Appendix D; \S II-E; \S VI-C, B2}

\begin{warningbox}
Do not claim B2 works because the recovery branch returned a plausible answer or because a later task happened to succeed. The evaluation must show that the proposed state and control coupling caused the improvement, while measuring state integrity and harmful rule reuse. If the retry baseline is non-inferior, or if compilation increases harmful reuse, the B2 claim should be rejected or narrowed.
\end{warningbox}

\section{Boundaries with the neighboring bundles}
\label{sec:ch13-boundaries}

B2 depends on neighboring mechanisms without absorbing them:

\begin{description}
  \item[B1, Chapter 12.] Memory accessibility and context selection may supply facts or prior recovery evidence to a child. B1 does not define the failure trigger, protected child scope or parent return protocol.
  \item[B3, Chapter 14.] Workspace competition may decide which recovery findings or candidate outputs receive scarce shared capacity. It is not the protocol that opens and validates a recovery substate.
  \item[B4, Chapter 15.] Persistent intentions and reconsideration authority define the goal that recovery must preserve. Fan and Lan's dependency figure identifies B4 as a scope for recovery; this chapter does not teach the BDI mechanism in detail.
  \item[B5, Chapter 17.] B2 may produce a reusable recovery rule, but it remains an in-episode recovery transition. B5 concerns general skill invocation and governance across tasks. Fan and Lan close B5 as a residual candidate after closest-baseline screening against GraSP.
  \item[B6, Chapter 16.] Local recovery needs a bound and escalation path. B6 studies broader uncertainty propagation, resource allocation, interruption and stopping authority.
\end{description}

The final synthesis in Chapter 18 will examine how B1, B2, B3, B4 and B6 compose, while retaining B5 as the closed GraSP case. Fan and Lan's proposed dependencies suggest that the mechanisms should be tested compositionally rather than assumed to add value merely because their diagrams can be connected.\FanLanCite{\S VI-D; Figure 8; \S VI-E}

\begin{chapterrecap}
Soar represents unresolved operator selection as a typed impasse and opens a substate to solve the corresponding subgoal. The child has a relation to its superstate; results are linked to higher-level context; dependency tracking can regenerate a subgoal when its relevant context changes; and chunking derives procedural knowledge from results and their dependencies. Fan and Lan translate this family of ideas into B2: typed diagnosis, protected recovery, typed return and resolution compilation. A production implementation needs explicit records, authority boundaries, a guarded state machine, version/dependency checks, side-effect handling, rule validation and a measurable fallback. The complete coupling remains a falsifiable research hypothesis, not a validated standard.
\end{chapterrecap}

\begin{reviewquestions}{Review and application questions}
\begin{enumerate}
  \item Why is a Soar impasse not the same as a generic tool exception?
  \item Distinguish tie, conflict, constraint-failure, and the two forms of no-change impasse.
  \item What makes a substate different from an independent agent process or a checkpoint?
  \item What qualifies as a Soar subgoal result, and why is it not arbitrary child text?
  \item Explain the difference between resolving, eliminating and regenerating an impasse.
  \item What does the Goal Dependency Set record, and why can a dependency change require regeneration?
  \item How should the procedural-learning manual's chunk-creation timing be explained alongside the architecture chapter's resolved-impasse discussion?
  \item Which edge does PALADIN or AgentDebugX preserve? Which edge does Shepherd preserve? What remains missing?
  \item Design a typed resolution for a payment timeout where the settlement outcome is unknown.
  \item Why must an accepted child result be checked against a parent version and dependency set?
  \item What can an internal rollback undo, and what can it not undo about external side effects?
  \item Describe a safe path from a validated recovery trace to a promoted reusable rule.
  \item Which tests would reveal an overgeneralized recovery rule after an API version change?
  \item Which B2 ablations separate the value of state representation, runtime control, and learning?
  \item Why does compilation of an in-episode recovery rule not make B2 identical to B5?
\end{enumerate}
\end{reviewquestions}

\section*{Source and example note}
\phantomsection
\addcontentsline{toc}{section}{Source and example note}

The historical mechanism is presented from Fan and Lan's Soar reconstruction and the Soar Architecture documentation they cite. The official Procedural Knowledge Learning manual is also consulted to state the timing of chunk creation and distinguish chunking from EBBS. The B2 residual bundle, modern baseline mappings, case-specific E/D records and experiment template are attributed to Fan and Lan. Product X, the payment-timeout sequence, the component architecture, schemas, state-machine labels, Java-like pseudocode, database/concurrency recommendations, idempotency safeguards and rule-promotion pipeline are book-created teaching or engineering proposals. They should not be mistaken for implemented functionality in Soar, PALADIN, AgentDebugX, Shepherd, AgentHER, Reflexion, or the Fan and Lan paper.

\FanLanCite{\S III-A; \S V-A--V-B; \S VI-B--VI-E; Appendices C-B--C-D, D, E-B, and F-B--F-C}
\autocite{soar2025architecture,soar2025procedural}

\section*{Looking ahead}
\phantomsection
\addcontentsline{toc}{section}{Looking ahead}

Chapter 14 turns from recovery of a blocked parent decision to a different competition problem: how typed content from memory, agents, tools, monitors and plans competes for limited workspace capacity, and how selected content is broadcast to downstream processing. A recovery result from B2 may become a candidate for that workspace, but workspace admission is not the same mechanism as opening, validating or compiling a recovery substate.

\chapter{Content Competition, Workspace and Broadcast: B3}
\begin{notebox}
Chapter placeholder. No prose has been drafted. This chapter will be outlined and approved before writing.
\end{notebox}

\chapter{Persistent Intentions, Reconsideration and Method Authority: B4}
\begin{notebox}
Chapter placeholder. No prose has been drafted. This chapter will be outlined and approved before writing.
\end{notebox}

\chapter{Uncertainty, Resources, Interruption and Stopping: B6}
\label{ch:uncertainty-resources-b6}

\section*{Chapter at a glance}
\phantomsection
\addcontentsline{toc}{section}{Chapter at a glance}

This chapter connects calibrated uncertainty to resource allocation, action authorization, interruption, and stopping. It reconstructs the complementary historical contributions of EPIC and Sigma, then develops a testable runtime design with evidence-linked uncertainty, resource vectors, guarded scheduling, persistence, and evaluation.

\begin{learningbox}
By the end of this chapter, you will be able to explain why a confidence value attached to one model or tool response is not the same as uncertainty that remains available to the whole runtime; reconstruct the complementary contributions of EPIC and Sigma; distinguish resource scheduling from uncertainty-conditioned action selection; interpret B6's E3/D3 record without treating it as a score for an entire agent; and design a testable control path that carries evidence-linked uncertainty into resource allocation, tool authorization, interruption, and stopping. You will also be able to distinguish a pause from a cancellation, a cancellation request from a confirmed cancellation, and budget exhaustion from successful completion.
\end{learningbox}

\section{The investigation that cannot afford to do everything}
\label{sec:ch16-opening}

Imagine that Product X's payment service begins returning intermittent HTTP 500 responses. The first alert is incomplete: the cause could be a recent deployment, a database connection problem, a downstream-provider delay, or malformed requests. The agent can inspect deployment metadata, query a read-only health endpoint, retrieve a bounded log window, compare error rates, or run a more expensive diagnostic.

Each investigation path consumes a different combination of tokens, model calls, tool quotas, compute, wall-clock time, and verification capacity. Shared bottlenecks make these costs non-interchangeable: a database may impose a strict query limit, an operator may be available only briefly, and one diagnostic may occupy the only permitted execution slot. A cheap token operation can still be expensive in time or external quota.

The agent must also distinguish investigation authority from operational authority. Read-only diagnostics may be permitted while restarting a service, rolling back a deployment, or issuing a payment mutation requires separate approval. When a second alert raises urgency, the runtime must decide whether lower-priority work can be interrupted, whether it has reached a safe checkpoint, and what state and remote effects must be checked before resumption. The later Product X trace works through these decisions in sequence.

The common control question is how evidence-linked uncertainty should change investigation, resource use, authorization, interruption, and the decision to continue or stop. An estimate that disappears after planning cannot govern later stages; an estimate treated as permission confuses probability with authority; and interruption that loses execution state can make the incident harder to resolve rather than safer.

\section{B6 in the mechanism-level review}
\label{sec:ch16-b6-review}

Fan and Lan identify B6 as the residual coupling between uncertainty and resource control. The term \gls{uncertaintyestimate} will refer to a proposition- or outcome-specific representation rather than a global confidence score.\index{uncertainty estimate} They propose that one calibrated uncertainty object accompany each consequential proposal and propagate through memory admission, planning, tool authorization, execution monitoring, and stopping. They connect that state to resource allocation and to the ability to interrupt work at a cognitive decision boundary. Their closest modern cases include systems that manage agent resources and systems that use uncertainty to select tools, request help, verify, explore, or stop. The review does not claim that those controls are absent. It argues that the reviewed systems do not yet supply the complete calibrated cross-stage propagation and cognitive preemption/interruption chain as one mechanism.\FanLanCite{\S IV-A--IV-B; \S V-A, EPIC/Sigma mapping; \S VI-C, B6; Appendix C-Q}

The distinction matters. A runtime may expose a scheduler, task queue, interrupt event, cancellation method, or approval callback. A specialized agent may estimate uncertainty or ask a human for help. A search algorithm may reduce exploration when its remaining budget shrinks. Each is a substantive control capability. But the presence of several capabilities in one software stack does not prove that one calibrated uncertainty state travels between them and changes their decisions. The question is whether the causal edge is present and inspectable.

\subsection{The proposed coupling}
\label{sec:ch16-b6-proposed-coupling}

The source hypothesis can be read as a chain. A consequential proposal is associated with an uncertainty estimate. The estimate has a defined target and \gls{calibration} context.\index{calibration} It remains traceable as the proposal moves through planning and memory selection. The scheduler considers it alongside the demand for scarce resources. The authorization gate considers it alongside consequence, reversibility, scope, and policy. Execution monitoring incorporates new observations without silently erasing the estimate's history. The runtime can then continue, verify, request help, yield, interrupt, or stop. A verified outcome remains linked to the estimate and decisions that preceded it.

This is not a demand that every stage use the same numerical threshold. A retrieval decision, a production restart, and a human-escalation decision have different consequences and different error costs. Nor does ``one uncertainty object'' mean one scalar that describes the entire agent. It means that the uncertainty associated with a particular consequential proposition or proposal has a durable, versioned lineage that later stages can interpret without losing its identity.

\begin{tipbox}
\textbf{Read the word ``propagate'' carefully.} It does not mean copy the same number into every component and pretend it remains current. New observations may produce a new estimate revision. Propagation means preserving the target, provenance, calibration context, and derivation relationship so that downstream decisions can identify exactly which uncertainty information they consumed.
\end{tipbox}

\subsection{What the chapter will and will not claim}

B6 is a residual hypothesis produced by a bounded mapping review. It is not a theorem that no system anywhere has combined uncertainty and resources in this way. Fan and Lan's corpus and coding rules define the scope of the result; their own validity discussion warns against treating an unobserved complete coupling as proof of an empty research area. Accordingly, the claim is limited: the reviewed corpus does not establish the complete coupling. That is materially different from saying that modern agents cannot do this.\FanLanCite{\S II-A--II-E; \S VII; Appendix F-C}

The chapter also distinguishes a source-level statement from an engineering choice. Fan and Lan specify the control problem, the proposed insertion point, the closest baselines, and the falsifying outcomes. They do not prescribe the Java records, database schema, resource units, scheduler objective, calibration method, threshold values, event names, or transaction protocol developed below. Those are proposals made to turn the hypothesis into a system that can be implemented and rejected by experiments.

\section{EPIC: why time and resource contention belong inside the model}
\label{sec:ch16-epic}

A common way to describe an agent is to list the components it has: a perception module, a reasoning module, a tool executor, and perhaps several parallel workers. That list says little about how the components compete for time. Two operations may both be available but still be unable to run at the same time because they require the same limited processor or channel. A third operation may be independent and able to proceed while the first two wait. Without timing and contention, the architecture does not explain which operation finishes first or why another one is delayed.

Fan and Lan use \gls{epic} as a historical resource-control reference.\index{EPIC} EPIC models perceptual, cognitive, and motor processors with explicit processing times, parallel pathways, queues, and bottlenecks. A perceptual event can initiate processor activity, but shared bottlenecks and motor timing constrain which operations overlap and which must wait. The contribution is not merely that EPIC has several processors. It is that the organization of those processors and their timing predicts latency and interference.\FanLanCite{\S III-D, EPIC; Fig. 4; Appendix F-A}\autocite{kieras1997epic}

\BookDiagram[0.52\textheight]{Tex/diagrams/generated/part2/chapter16/epic-resource-bottlenecks.png}{Historical teaching reconstruction of the EPIC resource idea: separate processor paths can overlap, while shared bottlenecks and channel capacity constrain timing.}{fig:ch16-epic}

\subsection{Parallel work does not mean unlimited work}
\label{sec:ch16-epic-parallel}

Consider three abstract operations in an incident investigation:

\begin{itemize}
  \item a log query that occupies a database connection;
  \item a local analysis that consumes CPU but no database connection;
  \item a deployment-history query that uses a separate read-only service.
\end{itemize}

The local analysis may overlap with either query if compute is available. The two database-dependent operations may not overlap if the only permitted connection is already occupied. A fourth task may be ready to run but blocked by the same bottleneck. Counting four available tasks does not tell us that four tasks can execute simultaneously.

The architecture must represent demand and capacity, including \gls{resourcecontention} when multiple activities require the same limited resource.\index{resource contention} If a resource has capacity one, admitting two operations that each require the whole capacity is not parallelism; it is an oversubscription decision. If a process can use a shared service only at a restricted rate, the scheduler must account for the service's limit rather than assuming that local concurrency is free. The system's temporal behavior follows from these constraints.

A useful way to picture this is a timeline. An operation begins when its required resources can be reserved. An independent operation may start in parallel. An operation that needs an occupied bottleneck waits. If a higher-priority activity arrives, the runtime may decide to defer or preempt lower-priority work, but it must also account for whether that work can be paused safely. A processor model that predicts duration is not yet a full interruption policy; it supplies the resource and timing facts on which a policy can act.

\subsection{From processors to modern runtime resources}
\label{sec:ch16-epic-modern}

A modern language-agent runtime has a different physical organization from EPIC. Its resources may include model inference capacity, tokens, tool quotas, browser sessions, database connections, containers, memory, wall-clock time, or human review. Some are continuously divisible; some are discrete; some are shared; some are leased; some cannot be recovered after a call has been made. A model call and a database query cannot be treated as interchangeable units simply because both consume budget.

The useful historical lesson is narrower: resource demand, concurrency, queues, durations, and bottlenecks should be explicit enough to influence scheduling. A thread pool or semaphore can provide an implementation primitive. It does not automatically become a cognitive resource policy. The mechanism appears when the runtime's task state and resource demand determine what can start, what waits, which operation gets priority, and whether work should yield. Fan and Lan make this connection when they argue that the modern analogue must expose scheduling, preemption, and parallelism as consequences of task state and resource demand.\FanLanCite{\S III-D, EPIC; \S V-A; \S V-B}

\subsection{Boundary of the EPIC comparison}

EPIC does not provide the entire B6 mechanism. Its historical contribution is explicit timing and resource contention. It does not, by itself, establish a calibrated uncertainty record that follows a language-agent proposal from memory admission through tool authorization and execution monitoring. That complementary relationship is the reason Fan and Lan pair EPIC with Sigma rather than treating scheduling and confidence estimation as unrelated features.

\section{Sigma: why uncertainty must remain part of the state}
\label{sec:ch16-sigma}

Suppose a model estimates that a deployment change probably caused an incident. The agent retrieves a deployment record, asks a diagnostic tool for evidence, receives a new observation, and then chooses between a further query and a report. If the original uncertainty value is attached only to the first model response, it may no longer describe the proposition used by the later decision. The tool result might support the hypothesis, weaken it, or reveal that the original question was underspecified. A later component needs to know what was estimated, from which evidence, under what calibration context, and how the estimate changed.

Fan and Lan use \gls{sigma} as the historical reference for uncertainty propagation.\index{Sigma} Sigma uses factor graphs as a shared substrate for symbolic, probabilistic, and decision computations. Evidence can trigger message passing; updated beliefs support action selection; and learning can update parameters or structures in the same representation. The point for B6 is not that Sigma supplies a ready-made LLM confidence object. Its contribution is that uncertainty remains part of the state as evidence moves through inference and decision.\FanLanCite{\S III-D, Sigma; Fig. 4; Appendix F-A}\autocite{rosenbloom2016sigma}

\BookDiagram[0.54\textheight]{Tex/diagrams/generated/part2/chapter16/sigma-uncertainty-propagation.png}{Historical teaching reconstruction of Sigma's factor-graph idea: evidence updates belief distributions that inform expected utility, action selection, and learning.}{fig:ch16-sigma}

\subsection{A local score and a propagated state are different}
\label{sec:ch16-local-vs-propagated}

A local confidence score answers a narrow question at one point in a process. For example, a model may produce a probability for whether a log excerpt contains a database timeout. A tool may provide a calibrated probability that a robot action will succeed. A classifier may return a confidence for a particular label. These estimates can be useful. The limitation is not that local uncertainty values are always invalid; it is that they do not become a runtime-wide control state merely because they are available.

Consider the difference between these two sequences:

\begin{enumerate}
  \item A model emits ``0.8 confidence.'' The planner copies that number into a prompt. The tool gate, scheduler, monitor, and stopping policy each make separate decisions without knowing the estimate's target, calibration set, freshness, or source evidence.
  \item A proposal carries a versioned estimate tied to a defined outcome. A later observation creates a linked revision. The scheduler records which revision it used to allocate resources. The authorization gate records whether the same estimate was still applicable to the action being proposed. The monitor records the updated evidence and whether the estimate became stale or out of domain.
\end{enumerate}

The second sequence is not simply more logging. The identity and \gls{uncertaintylineage} of the estimate are available at the control boundaries where they can change execution.\index{uncertainty lineage} That is the bridge to B6.

\subsection{What the factor graph analogy does not authorize}

The proposed runtime record is not a copy of Sigma's factor graph. A JSON object containing a probability and evidence identifiers does not reproduce Sigma's inference semantics. Nor does Fan and Lan specify a calibration algorithm for every type of prediction. Sigma is used here as a historical model of a shared uncertainty-bearing representation, while B6 asks whether a comparable state can influence control across modern agent stages.

A modern system may need several uncertainty records at once. The probability that a deployment caused an HTTP error is not identical to the probability that restarting a service will restore availability without data loss. A diagnostic estimate and an action-outcome estimate have different targets. Combining them requires a model of how the propositions relate. If the relationship is not known, the runtime must preserve the estimates separately rather than inventing independence or multiplying their values.

\subsection{The two historical lessons together}
\label{sec:ch16-epic-sigma-together}

EPIC and Sigma contribute complementary control ideas. EPIC makes demand, duration, overlap, and bottlenecks explicit. Sigma makes uncertainty-bearing state persist through inference and decision. B6 asks whether a modern runtime can connect them: a consequential proposal carries a calibrated, evidence-linked uncertainty lineage; resource decisions use it alongside resource demand; tool authorization and execution monitoring consume it at explicit boundaries; and interruption or stopping can occur when the resulting risk and resource state justify a control transition.\FanLanCite{\S V-A, EPIC/Sigma mapping; \S V-B; \S VI-C, B6}

\section{Modern systems already control parts of B6}
\label{sec:ch16-modern-baselines}

A residual hypothesis is meaningful only after the closest modern baselines have been considered. Fan and Lan separate general runtime families from mechanism-focused systems. The distinction prevents a framework interface from being mistaken for an implemented uncertainty policy, and it prevents a research prototype from being mistaken for a portable production invariant. Their resource and uncertainty examples are substantive, but the review treats them as partial controls rather than as proof that the full B6 coupling is already present.\FanLanCite{\S IV-A--IV-C; \S V-A--V-B; Appendix E-A; Appendix F-B}

\subsection{Resource-side controls}
\label{sec:ch16-resource-baselines}

\paragraph{AIOS.} AIOS proposes an operating-system layer for LLM agents, with scheduling and managers for context, memory, storage, tools, and access. It gives resource governance an explicit runtime control surface rather than leaving every resource decision inside model-generated text. Fan and Lan use it as a resource-governance precedent. The boundary is that a scheduler and its managers do not, by themselves, prove that one calibrated uncertainty state propagates through every downstream decision.\autocite{meiaiOS2024}\FanLanCite{\S IV-B; \S V-A; Appendix F-B, AIOS}

\paragraph{AgentRM.} AgentRM treats agent-system resources as operating-system-like concerns. Fan and Lan describe its feedback-driven admission, compaction, hibernation, and scheduling controls as evidence that adaptive resource management has been operationalized. It is therefore a stronger comparator than a static token limit. Its scope remains resource management: it does not establish the complete calibrated uncertainty, tool-authorization, preemption, interruption, and stopping chain required by B6.\autocite{she2026agentrm}\FanLanCite{\S IV-B; \S V-A; Appendix F-B, AgentRM}

\paragraph{Budget-aware value search.} Budget-aware value search changes which search-tree node is selected as the remaining budget changes. This demonstrates that a budget can influence the agent's reasoning path rather than merely being checked after the path has been generated. Its mapped control is narrower than an uncertainty lineage that follows memory, planning, authorization, execution monitoring, and stopping.\autocite{li2026bavt}\FanLanCite{\S IV-B; \S V-A; Appendix F-B, budget-aware value search}

\subsection{Uncertainty-conditioned action controls}
\label{sec:ch16-uncertainty-baselines}

\paragraph{UALA.} The Uncertainty-Aware Language Agent uses uncertainty quantification to moderate interaction between an agent and the external world. Fan and Lan include it as an uncertainty-conditioned tool-interaction baseline. That is an important control edge: uncertainty can influence whether the agent reaches outside its current internal state. It is not evidence that the same calibrated object governs resource allocation, every later tool, interruption, and terminal stopping.\autocite{han2024uala}\FanLanCite{\S IV-B; \S V-A; Appendix F-B, UALA}

\paragraph{KnowNo.} KnowNo uses conformal prediction to align uncertainty with help seeking in multi-step robotic planning. It is a concrete example of uncertainty causing an agent to request human help rather than continuing as if every prediction were reliable. The boundary is equally important: a help trigger does not, on its own, implement a general multi-resource allocator or a persistent interruption authority across an entire language-agent runtime.\autocite{ren2023knowno}\FanLanCite{\S IV-B; \S V-A; Appendix F-B, KnowNo}

\paragraph{Calibrate-Then-Act.} Calibrate-Then-Act makes cost--uncertainty trade-offs explicit in environment exploration and stopping. It demonstrates that the expected value of further information can affect whether an agent continues to investigate or commits to an answer. Fan and Lan use it as evidence for an uncertainty-conditioned action/stopping edge, not as proof that a common uncertainty record has propagated through all the stages in B6.\autocite{ding2026calibratethenact}\FanLanCite{\S IV-B; \S V-A; Appendix F-B, Calibrate-Then-Act}

\paragraph{Utility-Guided Agent Orchestration.} This is the closest cost-sensitive action baseline named in the B6 discussion. It selects among responding, retrieving, calling a tool, verifying, and stopping by considering estimated gain, step cost, heuristic uncertainty, and redundancy. That is real orchestration control; it would be inaccurate to describe it as having no resource or action policy. Fan and Lan's specific boundary is that the uncertainty signal is heuristic rather than calibrated, and the system does not provide cross-stage propagation or interruption authority. The residual therefore survives a cost-aware action selector.\autocite{liu2026utilityorchestration}\FanLanCite{\S VI-C, B6; Appendix C-Q; Appendix F-B, Utility-Guided Agent Orchestration}

\begin{table}[tbp]
\centering
\small
\caption{The modern controls contribute important parts of B6, but not the complete coupling as characterized by Fan and Lan.}
\label{tab:ch16-modern-baselines}
\begin{tabularx}{\linewidth}{P{0.22\linewidth}XX}
\toprule
\textbf{Control family} & \textbf{What the reviewed cases establish} & \textbf{What remains open for B6} \\
\midrule
Resource governance & Scheduling, admission, compaction, hibernation, or budget-sensitive search can change execution. & A calibrated uncertainty lineage is not shown to propagate into all resource and risk decisions. \\
Uncertainty-conditioned action & Uncertainty can influence tool interaction, help seeking, exploration, verification, or stopping. & These controls do not by themselves establish multi-resource allocation, persistent cross-stage state, or interruption authority. \\
General runtime boundaries & State, tools, approvals, checkpoints, cancellation, interruption, and resumption expose insertion points. & An interface boundary is not itself a calibrated control law or evidence that a system uses uncertainty to govern it. \\
\bottomrule
\end{tabularx}
\end{table}

\subsection{Runtime interfaces are insertion points, not B6 implementations}
\label{sec:ch16-runtime-insertion-points}

General-purpose runtimes make practical boundaries visible. LangGraph exposes stateful graph execution with checkpoint, interrupt, and resume concepts. The OpenAI Agents SDK interface examined by Fan and Lan exposes run state, tools, handoffs, guardrails, approval, cancellation, and resume-related behavior. These are useful boundaries for an implementation discussion. They should not be promoted into evidence that the framework supplies calibrated uncertainty propagation or a B6 scheduler. The interface evidence is version-bounded in the review, and the scientific mechanism claim remains grounded in the paper's source analysis.\autocite{langgraph2026repo,openaiagents2026repo}\FanLanCite{\S IV-A; Appendix E-A, Table XI}

The distinction is easy to miss because the same word can name different levels of control. A cancellation method is an API operation. A safe interruption policy must decide when cancellation is appropriate, which work can stop, what state is preserved, and what can be said about remote side effects. A checkpoint stores enough state to resume. A B6 controller must also decide whether the estimate, budget, permission, and evidence on which the checkpoint depended remain valid. A scheduler assigns work. A B6 controller needs to connect that assignment to calibrated uncertainty and action risk.

\subsection{The comparison result}

The modern landscape is not empty. Resource managers demonstrate that admission and scheduling can respond to system state. Uncertainty-aware systems demonstrate that confidence, calibration, and cost can influence tools, help, exploration, and stopping. Runtime frameworks expose the places where state and control could be connected. The remaining hypothesis is about the coupling and authority of these mechanisms: calibrated uncertainty must accompany consequential proposals, change resource allocation and tool governance across stages, and support cognitive preemption and interruption at a defined decision boundary.\FanLanCite{\S V-A--V-B; \S VI-C, B6; Appendix C-Q}

\BookDiagram[0.60\textheight]{Tex/diagrams/generated/part2/chapter16/b6-modern-baseline-map.png}{Modern resource-governance and uncertainty-conditioned systems contribute separate control edges; the B6 hypothesis concerns their calibrated cross-stage coupling.}{fig:ch16-modern-baselines}

\section{What E3/D3 says about B6}
\label{sec:ch16-ed-coding}

Chapter 11 introduced two independent coding axes. Evidence relation describes how a modern mechanism corresponds to a historical one. Migration depth describes how far the mapped mechanism has been implemented. These are separate questions, not two parts of a single maturity score. E3 identifies a functionally comparable control solution whose internal organization differs materially from the historical mechanism. D3 identifies a runtime control law in which an explicit trigger executes the mapped action and changes state; it does not require experience to update future policy. D4 is the level associated with outcome-dependent changes to future control policy.\FanLanCite{\S II-D; Table II; Appendix C-Q}

Fan and Lan's B6 mapping is E3/D3. The evidence relation is a functional approximation: the resource and uncertainty systems solve comparable control problems without establishing a direct historical transfer from EPIC or Sigma. The depth code reflects explicit modern control transitions such as budget-sensitive scheduling, help seeking, tool arbitration, exploration, verification, or stopping. It does not say that the full B6 chain is present, because the residual is precisely the missing calibrated cross-stage propagation plus cognitive preemption and interruption.\FanLanCite{\S V-A, EPIC/Sigma mapping; \S V-B; \S VI-C, B6; Appendix C-Q}

\begin{warningbox}
\textbf{Do not read E3/D3 as ``almost complete.''} E3 and D3 answer different questions. E3 is not a lower quality grade than E2, and D3 is not a percentage of completion. The pair describes one bounded mapping claim in the paper's frozen corpus. It is not a score for AIOS, AgentRM, UALA, KnowNo, Utility-Guided Orchestration, or the quality of modern agents as a whole.
\end{warningbox}

\subsection{Why the resource and uncertainty parts form one bundle}
\label{sec:ch16-one-bundle}

The merge rule in Fan and Lan's review retains atomic residuals together when they share a runtime insertion boundary and one supplies state or authority needed by the other. For B6, a scheduler without uncertainty can allocate resources, but it cannot use the relevant calibrated risk state to decide when verification, help, or interruption is worth its cost. An uncertainty estimator without resource and action control can report a number, but the number does not determine what the runtime does next. The shared insertion points are the scheduler and tool-authorization boundary. The coupling can be tested by removing the state object or its authoritative control edge while leaving other behavior in place.\FanLanCite{\S II-D; \S VI-B, Table V; \S VI-C, B6; Table VI}

\subsection{The source's four central failure conditions}

Fan and Lan name four ways the proposed bundle can fail. First, \emph{sequential miscalibration}: an estimate that appears useful at one stage may become misapplied as the proposal moves into a different target, domain, or context. Second, \emph{excessive abstention}: uncertainty-aware control may request help or decline action so often that it prevents useful completion. Third, \emph{verification cost without risk reduction}: more checks consume resources without reducing the relevant risk. Fourth, \emph{interruption that prevents correct completion}: a controller may interrupt or preempt work at the wrong time and make the task fail. These are specific falsifiers, not generic warnings to ``be careful with confidence.''\FanLanCite{\S VI-C, B6; Table VI; Appendix D, Table X}

\section{Reconstructing B6 with the seven mechanism fields}
\label{sec:ch16-seven-fields}

The seven-field representation from Chapter 10 is a way to reconstruct a mechanism, not a mandatory software schema. It asks what state the mechanism consumes, who controls it, what triggers a transition, what persists, how failure is represented, how experience changes future behavior, and how resources or governance constrain it. B6 is a useful test of the representation because the source identifies some fields clearly and leaves others deliberately open.

\begin{table}[tbp]
\centering
\small
\caption{B6 reconstructed through the seven fields. Source-supported facts are separated from implementation details that remain unspecified.}
\label{tab:ch16-seven-fields}
\begin{tabularx}{\linewidth}{P{0.15\linewidth}XX}
\toprule
\textbf{Field} & \textbf{B6 source record} & \textbf{What remains unspecified} \\
\midrule
State (S) & Budget, queue, context, uncertainty, and step cost appear in the resource/uncertainty record. & A Java object model, database schema, versioning strategy, and calibration payload are engineering proposals. \\
Control (C) & Scheduler and tool-authorization boundaries influence resource/action decisions. & The paper does not prescribe one process boundary or universal scheduler implementation. \\
Transition (T) & Scheduling, help, retrieval, tool use, verification, and stopping choices are identified; new evidence can affect control. & Exact event taxonomy, thresholds, tie-breaking, and priority rules are not specified. \\
Persistence (P) & Uncertainty is intended to accompany consequential proposals across stages; common runtime state includes uncertainty and budget. & Durable schema, retention, expiry, crash recovery, and cross-session persistence are not specified for B6. \\
Failure (F) & Sequential miscalibration, excessive abstention, verification cost without risk reduction, and harmful interruption are explicit failures. & Concrete exception types and recovery transitions are implementation proposals. \\
Learning (L) & Sigma's historical mechanism can update parameters or structures; the synthesis discusses outcome-linked learning. & No B6-specific online calibration algorithm is prescribed, and the modern cases do not all implement one. \\
Resources (R) & Token, time, tool, verification, concurrency, risk, preemption, interruption, and stopping are relevant resource/governance concerns. & The paper gives no single objective function, universal conversion rate, or global threshold. \\
\bottomrule
\end{tabularx}
\end{table}

The table is intentionally conservative. In particular, a reader should not infer a complete persistence design from the word ``propagate.'' Propagation establishes a control requirement: downstream stages must be able to access an estimate with its identity and context. The exact retention period, database representation, or recovery protocol is not specified. The same restraint applies to learning. Sigma's historical parameter/structure updates do not prove that a modern B6 system has a calibrated online update rule. A book-level implementation may propose one, but it must remain labelled as a proposal and evaluated independently.\FanLanCite{\S II-C; \S III-D, Sigma; \S VI-A; Appendix C-Q; Appendix E-A}

\section{Designing a runtime that can act on uncertainty}
\label{sec:ch16-runtime-design}

We can now turn the residual hypothesis into an implementable design. The design below is deliberately explicit about authority. The model proposes explanations and actions. It may estimate resource demand or uncertainty. It does not reserve authoritative budget, grant its own tool permission, override an interruption, or declare that an external side effect has been undone. Those are runtime decisions owned by components whose state and transitions are durable and testable.

\subsection{Start with authoritative state outside the model call}
\label{sec:ch16-authoritative-state}

A runtime cannot enforce a policy whose only representation is a paragraph in a prompt. The authoritative state should live in versioned records that can be read and changed by the components responsible for control. For a B6-capable system, the relevant state includes:

\begin{itemize}
  \item the active goal and scope owned by the B4 intention boundary;
  \item the evidence and admitted context available through B1/B3 mechanisms;
  \item the uncertainty lineage attached to each consequential proposition or proposal;
  \item the resource budget, reservations, consumption, and deadlines;
  \item the active work, tool-call state, and any in-flight external side effect;
  \item recovery state when a typed failure is being handled through B2; and
  \item the currently selected skill or workflow and its verification status, including any migrated skill governor used as the B5 baseline.
\end{itemize}

The B6 controller does not own all of these objects. It consumes the relevant state and governs risk/resource decisions across the boundaries. The intention governor remains the authority for changing or abandoning the goal. A recovery controller owns the typed recovery branch. A skill governor owns its graph and verification semantics. B6 supplies the risk and resource policy that can approve further bounded work, reserve verification capacity, ask for help, interrupt, or stop. This separation keeps a risk decision from silently becoming a goal decision.

\BookDiagram[0.63\textheight]{Tex/diagrams/generated/part2/chapter16/b6-cross-stage-control-loop.png}{Proposed B6 control loop: uncertainty is linked to a consequential proposal and consumed at memory, planning, scheduling, authorization, monitoring, and stopping boundaries.}{fig:ch16-b6-loop}

\subsection{The uncertainty lineage record}
\label{sec:ch16-uncertainty-record}

\EngineeringLabel{} The runtime should attach a versioned uncertainty record to each consequential proposition or proposal lineage. It should not use one global confidence value to describe an entire task. A payment-status proposition, a deployment-cause hypothesis, and a proposed production restart are different prediction targets. They may depend on overlapping evidence, but they do not automatically share one probability or one calibration guarantee.

A proposed record can contain the following fields:

\begin{itemize}
  \item \texttt{\detokenize{estimateId}} — stable identity for this estimate revision;
  \item \texttt{\detokenize{subjectRef}} — the proposition, event, or action outcome being estimated;
  \item \texttt{\detokenize{targetDefinition}} — the exact event or outcome whose probability or prediction set is represented;
  \item \texttt{\detokenize{prediction}} — a probability, interval, prediction set, or other typed output suitable for the target;
  \item \texttt{\detokenize{modelVersion}} and \texttt{\detokenize{calibratorVersion}} — the estimator and calibration artifact that produced the value;
  \item \texttt{\detokenize{calibrationDomain}} — the task, action class, environment, or population for which calibration evidence applies;
  \item \texttt{\detokenize{calibrationEvidenceRef}} — the versioned evaluation record supporting the calibration claim;
  \item \texttt{\detokenize{evidenceRefs}} and \texttt{\detokenize{parentEstimateId}} — provenance and derivation links;
  \item \texttt{\detokenize{observedAt}}, \texttt{\detokenize{expiresAt}}, and \texttt{\detokenize{stageVersion}} — time, freshness, and process-stage metadata; and
  \item \texttt{\detokenize{status}} — for example \texttt{\detokenize{CALIBRATED}}, \texttt{\detokenize{OUT_OF_DOMAIN}}, \texttt{\detokenize{STALE}}, \texttt{\detokenize{UNVALIDATED}}, or \texttt{\detokenize{UNKNOWN}}.
\end{itemize}

This schema is not in Fan and Lan. It is a proposed contract that makes the propagation requirement operational. The target definition is essential. An estimate that a log entry contains a timeout is not interchangeable with an estimate that the timeout caused a payment failure. An estimate about diagnosis is not an estimate about the safety of a restart. If the target changes, the runtime must create a new estimate or explicitly derive one under a documented model. It must not silently carry the old number forward.

A lineage can have several revisions. A new observation should create a successor estimate with a link to the previous estimate and the evidence that changed it. The old revision remains available for audit. Downstream decisions record the exact revision they consumed. This permits a reviewer to reconstruct why the scheduler reserved resources or why the tool gate denied a proposal even if later evidence changed the estimate.

\subsection{Calibration is target-specific evidence}
\label{sec:ch16-calibration}

Calibration is often described too casually. A model's verbal statement that it is ``90 percent confident'' is not, by itself, evidence that events assigned that probability occur 90 percent of the time. Calibration is evaluated against outcomes for a defined prediction task and a relevant population. Its meaning depends on how the target is defined, how the examples are sampled, and whether the calibration population resembles the situation where the estimate is used.

A calibrated estimate for a read-only diagnostic classification is not automatically calibrated for a high-impact write action. A model version change, a new tool, a different environment, or a shift in the severity distribution may make an earlier calibration record less applicable. The runtime therefore needs explicit status for stale, out-of-domain, or unvalidated estimates. In those cases the system may choose a conservative policy, require verification, or ask for help. The exact response is an engineering decision; the important point is that an unvalidated score must not masquerade as a calibrated one.

\begin{warningbox}
\textbf{Calibration is not authorization.} A well-calibrated prediction describes expected outcomes under a stated target and population. It does not grant permission to restart a service, write to a payment system, disclose data, or take another consequential action. Authorization depends on the principal's policy, the action's scope and consequence, reversibility, and any required human approval.
\end{warningbox}

The chapter does not prescribe one calibration method. The estimator might use a calibrated classifier, a conformal prediction set, or another validated technique appropriate to the task. These are not interchangeable outputs; the contract must represent the method's semantics rather than forcing every result into a scalar probability. KnowNo is a concrete example of conformal prediction used for help-seeking in a robotic planner. Its role in the chapter is as a comparator identified by Fan and Lan, not as the universal calibration algorithm for every B6 domain.\autocite{ren2023knowno}\FanLanCite{\S IV-B; \S VI-C, B6; Appendix F-B}

\subsection{Propagation without inventing independence}
\label{sec:ch16-uncertainty-propagation}

The same uncertainty estimate may be consumed by several downstream components, but new observations do not necessarily provide independent evidence. If two log queries retrieve overlapping windows, treating them as independent confirmations can exaggerate the strength of the conclusion. If a model summarizes an observation and a later model call summarizes that summary, the second text is not a new independent sensor. If the target changes from the proposition ``a database timeout is present'' to ``a database timeout caused the incident,'' the old estimate does not automatically transfer.

The proposed runtime should therefore distinguish reuse from update. A downstream component may reuse an existing estimate if its target, context, and freshness remain applicable. A new observation may produce a linked revision. A derived proposition requires an explicit derivation step. Aggregation requires assumptions about the relationship among evidence sources and should be tested rather than improvised. Where dependence is unknown, retain the sources and uncertainty separately or use a conservative decision rule. Do not multiply probabilities merely because several records exist.

This design is a practical expression of the Sigma-inspired idea that uncertainty remains part of the state during inference and decision. It is not a claim that a list of estimate records reproduces Sigma's factor graph or message-passing algorithm. The runtime contract preserves lineage and makes the control decisions inspectable; the inference method that creates a calibrated estimate remains a separately evaluated component.

\subsection{A resource vector, not one fungible number}
\label{sec:ch16-resource-vector}

\EngineeringLabel{} The scheduler should represent resource demand as a \gls{resourcevector} with typed units.\index{resource vector} A proposal might consume model tokens, one model round trip, two external tool calls, several seconds of wall-clock time, a database connection, and one verification slot. It is misleading to collapse those quantities into a single ``cost'' unless the system has a validated domain-specific conversion model. A spare token budget does not compensate for an expired deadline, and unused local CPU does not create permission to call a rate-limited production endpoint.

A proposed resource vector can include:

\begin{itemize}
  \item inference tokens or another model-capacity unit;
  \item model calls or round trips;
  \item external tool calls, separated by tool and quota domain;
  \item wall-clock time and absolute deadline;
  \item compute capacity or execution slots;
  \item context/memory capacity; and
  \item verification or human-review capacity when it is explicitly budgeted.
\end{itemize}

Each dimension has a unit, a capacity or limit, a reservation rule, a consumption record, and a reconciliation rule. Some limits are global to a run; others are shared across runs or enforced by an external service. Some resources are reserved before dispatch and settled afterward. Others, such as wall-clock deadline, are not refundable. A design must say which semantics apply to each dimension.

\begin{table}[tbp]
\centering
\small
\caption{An illustrative resource vector for the Product X investigation. The values are hypothetical teaching values, not measurements or recommendations.}
\label{tab:ch16-resource-vector}
\begin{tabularx}{\linewidth}{P{0.20\linewidth}P{0.16\linewidth}XX}
\toprule
\textbf{Resource} & \textbf{Unit} & \textbf{Why it is separate} & \textbf{Possible control rule} \\
\midrule
Inference capacity & tokens & Model usage is not equivalent to tool or elapsed-time use. & Reserve a maximum for each stage and retain a global cap. \\
Model calls & calls & Each round trip has latency and provider limits. & Admit only when a call reservation is available. \\
External tools & calls by tool & A database query and a deployment API call have different quotas and risks. & Maintain separate quotas and action-specific approval. \\
Wall-clock time & \shortstack[l]{seconds /\\deadline} & Waiting can be costly even when tokens are cheap. & Use absolute deadlines and reserve time for verification. \\
Compute & slots or CPU-time & Parallel work can contend for shared capacity. & Schedule against available concurrency and safe preemption points. \\
Verification & slots or review minutes & Verification may be the scarce resource that prevents unsafe completion. & Protect a reserve; do not spend it on low-value exploration. \\
\bottomrule
\end{tabularx}
\end{table}

The scheduler can still optimize a task-level objective, but the resource constraints remain explicit. A constrained selection problem may maximize expected task progress or expected risk reduction subject to the capacity of each resource. If the designer chooses to combine risk reduction and task progress into one utility, the conversion weights are policy assumptions that must be exposed and evaluated. They are not supplied by Fan and Lan. The runtime must also retain hard constraints: an unauthorized action cannot become eligible merely because its score is high, and a resource vector cannot be overdrawn because an optimizer thinks the benefit is large.

\BookDiagram[0.59\textheight]{Tex/diagrams/generated/part2/chapter16/b6-resource-vector-scheduler.png}{Proposed resource-vector scheduler: candidates request typed resources, reservations enforce separate limits, and new evidence can change priority without treating unlike resources as fungible.}{fig:ch16-resource-vector}

\subsection{Budget allocation and scheduling}
\label{sec:ch16-budget-scheduling}

A useful scheduler considers at least five things: expected task progress, the uncertainty relevant to the next decision, the consequence of being wrong, the resources required by the candidate action, and the hard constraints that govern whether the action may occur. It should decide among several responses, not simply compute a confidence threshold. The options may include a cheaper diagnostic, another verification, a request for human help, a deferral, a safe preemption, or a stop with an explicit explanation.

An implementation should reserve capacity before dispatch. Otherwise two concurrent workers can each observe that one tool call remains and both start it. Reservation creates an authoritative claim on a resource. When the operation finishes, the runtime reconciles the reservation against actual consumption. A reservation must be released or settled exactly once, even when duplicate events arrive or a worker retries after a crash.

The scheduler should also avoid consuming all capacity on exploration. A verification reserve can protect a later check that is required before reporting success. A deadline reserve can leave time for an operator handoff. A per-stage cap can prevent a single uncertain branch from monopolizing the run. These policies are proposed engineering controls. They must be tested against fixed-budget and simpler budget-aware baselines rather than assumed to be beneficial.

Concurrency introduces further requirements. A shared resource may be temporarily unavailable. A high-priority task may preempt work that is safe to pause. A low-priority task can starve if urgent tasks arrive continuously. Priority inversion can occur when a high-priority operation waits for a resource held by low-priority work. The implementation needs an explicit policy for fairness, priority aging, emergency reserve, and resource release. It must also distinguish resources that can be preempted from external actions that cannot be undone once dispatched.

\subsection{Tool authorization is a separate authority}
\label{sec:ch16-tool-authorization}

Uncertainty is one input to authorization, not its substitute. The policy should consider the requested operation, the principal's authority, the exact target and scope, the action's consequence, reversibility, possible data exposure, required verification, and any human approval condition. A read-only query can be allowed under one policy while a write operation is denied. The same probability estimate may inform both decisions, but it does not make them equivalent.

A proposed authorization record should include the subject and principal, the exact tool and operation, target scope, the goal/intention reference, the uncertainty estimate revision consumed, consequence and reversibility classification, policy version, required verifier or human approver, decision, reason, and expiry. The decision should be bound to the action as requested. If the model changes the target, payload, or operation, the old authorization must not silently apply to the new request.

This distinction blocks a dangerous shortcut: ``the model is confident, therefore the action is allowed.'' Even a well-calibrated estimate that a restart will fix the problem does not prove that the agent has permission to restart production. Conversely, a low estimated risk does not waive a policy that requires human approval for an irreversible action. The authorization component owns that decision; the model supplies a proposal and supporting evidence.

\BookDiagram[0.62\textheight]{Tex/diagrams/generated/part2/chapter16/b6-authorization-and-interruption.png}{Proposed authority boundary: the model proposes, uncertainty informs the decision, policy authorizes the exact action, and a separate monitor can request safe interruption.}{fig:ch16-auth-interrupt}

\subsection{Scheduling, preemption, interruption, cancellation and stopping}
\label{sec:ch16-interrupt-stop}

These terms are often used as if they all mean ``stop the task.'' They do not. A runtime needs distinct states and semantics because pausing local reasoning is not the same as cancelling an external side effect, and neither is the same as ending the run.

\begin{description}
  \item[Scheduling.] Decide which eligible work should receive available resources next. Scheduling need not interrupt anything that is already running.
  \item[\gls{preemption}.] Temporarily yield a lower-priority activity at an agreed safe boundary so another activity can use a constrained resource. A preempted activity should retain a validated checkpoint and its unspent budget when it is resumable.
  \item[\gls{interruption}.] An authoritative event asks the current activity to stop at the earliest defined safe boundary because risk, policy, priority, or new evidence has changed. An \gls{interruption} request is a control event, not merely a message appended to a prompt.
  \item[\gls{cancellation}.] Request termination of an executing task or tool operation. The request may be rejected, delayed, or acknowledged after the remote system has already performed a side effect.
  \item[Suspension.] Preserve the task's state while no execution is active, pending a resource, condition, deadline, or human decision.
  \item[Resumption.] Continue from a validated checkpoint only after checking that the assumptions, permissions, uncertainty lineage, budget, and external state are still current.
  \item[Stopping.] Terminally end the current run or decision process with an explicit reason. Verified success, safe refusal, abstention, infeasibility, deadline exhaustion, budget exhaustion, and escalation are different terminal or handoff outcomes.
\end{description}

A \gls{cancellation} acknowledgement must not be interpreted as proof that an external effect did not happen. A network request may have reached a payment provider before the local timeout. The remote operation may complete even if the local worker stops waiting. A service may not support cancellation at all. The runtime should therefore track at least the cancellation request, the tool's reported cancellation status, and the known or unknown side-effect status as separate facts. When the effect is consequential, read-after-write reconciliation or human review may be necessary.

A safe preemption boundary depends on the operation. A local summarizer can often pause between chunks. A read-only query may be allowed to complete. A deployment operation may be non-preemptible after dispatch. The runtime should state the cancellation capability for each tool contract rather than assume every tool behaves like a cooperative local task.

\subsection{A proposed lifecycle state machine}
\label{sec:ch16-state-machine}

\EngineeringLabel{} The following state machine is a proposed control contract for the examples in this chapter. It is not a state machine defined by Fan and Lan. Its purpose is to make transitions, authority and persistence testable.

\begin{description}
  \item[READY.] The activity is eligible to be scheduled but has no active reservation.
  \item[RESERVED.] The scheduler has committed the required resource reservation; dispatch is not yet confirmed.
  \item[RUNNING.] The operation has been dispatched and is executing or awaiting a result.
  \item[PREEMPT\_REQUESTED.] The governor has requested that work yield at a safe boundary; the worker has not yet confirmed the checkpoint.
  \item[INTERRUPT\_REQUESTED.] An authoritative interruption event is pending. The worker must reach the defined boundary or report why it cannot.
  \item[SUSPENDED.] No further execution is active. State and reservation reconciliation are persisted.
  \item[RESUME\_CHECK.] The runtime validates checkpoint identity, current policy, permission, estimate freshness, external state and remaining budget.
  \item[STOPPED.] The run is terminal without verified success, with an explicit reason such as safe refusal, abstention, infeasibility or budget exhaustion.
  \item[COMPLETED.] The authorized objective for this run has been verified as complete under its required postconditions.
\end{description}

\texttt{\detokenize{STOPPED}} and \texttt{\detokenize{COMPLETED}} are intentionally distinct. The fact that a controller has no more work scheduled does not prove that the objective was achieved. A run can stop because the deadline expired, because it lacks authority, because the evidence remains insufficient, or because continuation would create unacceptable risk. A completion record needs verified outcome evidence.

\BookDiagram[0.62\textheight]{Tex/diagrams/generated/part2/chapter16/b6-state-machine.png}{Proposed B6 activity state machine separating reservation, execution, preemption, interruption, suspension, resume validation, unsuccessful stopping, and verified completion.}{fig:ch16-state-machine}

Every transition should record its trigger, evidence, authority, budget effect, uncertainty revision consumed or produced, tool side-effect status, persistence effect, and whether it is terminal. Transitions must be guarded. For example, \texttt{\detokenize{RESUME_CHECK}} cannot return to \texttt{\detokenize{RUNNING}} merely because a checkpoint exists. It must first confirm that the permission remains valid, the resource reservation can be reacquired, the estimate has not expired or moved out of domain, and the external state has not changed in a way that invalidates the plan.

\subsection{Learning and calibration updates}
\label{sec:ch16-learning}

The source record leaves a B6-specific learning algorithm unspecified. Sigma's historical mechanism can update parameters or structures, and the synthesis runtime discusses outcomes as potential inputs to later control. That does not mean every B6 comparator learns a common calibration policy, or that Fan and Lan provide an online update rule. A book-level implementation may experiment with calibration updates, but it should keep that update path separate from the immediate runtime control path.\FanLanCite{\S III-D, Sigma; \S VI-A; Appendix C-Q}

A conservative proposal is to admit only verified outcome labels into calibration feedback. The system should preserve the prediction, target definition, model and calibrator versions, evidence, decision, and resulting outcome. Training/update data must be separated from evaluation data. Calibration should be monitored by domain and action class, because aggregate performance can hide a poorly calibrated high-consequence slice. Every updated calibrator should be versioned, and a rollback or disable mechanism should be available if drift degrades calibration. Most importantly, the runtime transition should remain testable with learning disabled. This supports the source-required ablation that removes learning while retaining the immediate control edge.

\section{Java-like contracts, persistence and concurrency}
\label{sec:ch16-contracts}

The previous section described responsibilities. This section makes their boundaries concrete with Java-like contracts. The types are proposed for this book; they are not copied from Fan and Lan or from a particular framework. The examples use ordinary Java records and enums so that the data contract is visible without requiring a specific persistence or orchestration library.

\subsection{Core records}
\label{sec:ch16-core-records}

The uncertainty record must say what it predicts and which calibration context supports it. The resource vector must preserve units. The action proposal must not carry authority merely because a model generated it.

\begin{lstlisting}[language=Java,caption={Proposed records for an uncertainty lineage and a typed resource vector}]
record UncertaintyEstimate(
    UUID estimateId,
    String subjectRef,
    String targetDefinition,
    Prediction prediction,
    String modelVersion,
    String calibratorVersion,
    String calibrationDomain,
    String calibrationEvidenceRef,
    List<String> evidenceRefs,
    UUID parentEstimateId,
    Instant observedAt,
    Instant expiresAt,
    CalibrationStatus status
) {}

sealed interface Prediction permits Probability, Interval, PredictionSet {}
record Probability(double value) implements Prediction {}
record Interval(double lower, double upper) implements Prediction {}
record PredictionSet(Set<String> outcomes) implements Prediction {}

enum CalibrationStatus {
    CALIBRATED, OUT_OF_DOMAIN, STALE, UNVALIDATED, UNKNOWN
}

record ResourceAmount(String resource, BigDecimal amount, String unit) {}
record ResourceVector(List<ResourceAmount> amounts) {}

record BudgetState(
    UUID runId,
    ResourceVector limit,
    ResourceVector reserved,
    ResourceVector consumed,
    Instant deadline,
    long version
) {}
\end{lstlisting}

The \texttt{\detokenize{Prediction}} type is deliberately not a \texttt{\detokenize{double}}. A probability is one possible representation, but an interval or prediction set has different semantics. The \texttt{\detokenize{targetDefinition}} prevents a value from being consumed for a question it was not designed to answer. The calibration status records whether the estimate is applicable now. A stale record is not silently treated as fresh simply because it is still present in a database.

The resource vector stores units alongside values. A production implementation would normally replace the free-form \texttt{\detokenize{resource}} and \texttt{\detokenize{unit}} strings with enums or validated unit types. The abbreviated form keeps the example readable. It must still reject incompatible units and duplicate resource dimensions before a reservation is attempted.

\subsection{Proposal and authorization contracts}

\begin{lstlisting}[language=Java,caption={Proposal and authorization are separate records with separate owners}]
enum Consequence { LOW, MODERATE, HIGH, CRITICAL }
enum Reversibility { REVERSIBLE, COMPENSATABLE, IRREVERSIBLE, UNKNOWN }

enum AuthorizationStatus {
    ALLOWED, DENIED, NEEDS_HUMAN, EXPIRED
}

record ActionProposal(
    UUID proposalId,
    UUID intentionId,
    String toolName,
    String operation,
    String targetScope,
    String payloadDigest,
    Consequence consequence,
    Reversibility reversibility,
    ResourceVector estimatedDemand,
    UUID uncertaintyEstimateId
) {}

record AuthorizationDecision(
    UUID decisionId,
    UUID proposalId,
    String principalId,
    String policyVersion,
    AuthorizationStatus status,
    String reason,
    String requiredApprovalRef,
    Instant decidedAt,
    Instant expiresAt
) {}
\end{lstlisting}

The authorization decision is bound to the proposal identity and a digest of the payload. If the model changes the operation or the target, it needs a new proposal and a new decision. The policy version is recorded so that the audit can reconstruct which rule was applied. The record also captures a required human approval. It does not store a vague Boolean named \texttt{\detokenize{safe}} that could be copied from one action to another.

\subsection{Execution, interruption and outcome contracts}

\begin{lstlisting}[language=Java,caption={Execution records distinguish interruption requests from external cancellation results}]
enum ActivityState {
    READY, RESERVED, RUNNING, PREEMPT_REQUESTED,
    INTERRUPT_REQUESTED, SUSPENDED, RESUME_CHECK,
    STOPPED, COMPLETED
}

enum CancellationStatus {
    NOT_REQUESTED, REQUESTED, ACKNOWLEDGED,
    REJECTED, TIMED_OUT, UNSUPPORTED, UNKNOWN
}

enum SideEffectStatus {
    NONE, NOT_STARTED, IN_PROGRESS, CONFIRMED,
    COMPENSATED, UNKNOWN
}

record InterruptionRequest(
    UUID eventId,
    UUID activityId,
    String cause,
    int priority,
    Instant deadline,
    String safeBoundary
) {}

record ExecutionState(
    UUID activityId,
    ActivityState state,
    long generation,
    UUID checkpointId,
    UUID uncertaintyEstimateId,
    UUID budgetReservationId,
    CancellationStatus cancellationStatus,
    SideEffectStatus sideEffectStatus,
    Instant updatedAt,
    long version
) {}

record StopRecord(
    UUID runId,
    String reason,
    String evidenceRef,
    UUID uncertaintyEstimateId,
    UUID budgetSnapshotRef,
    Instant stoppedAt
) {}
\end{lstlisting}

\texttt{\detokenize{CancellationStatus}} and \texttt{\detokenize{SideEffectStatus}} are separate because cancellation is an attempted control over an operation, while the side-effect status describes what is known about the external world. An acknowledgement can coexist with \texttt{\detokenize{UNKNOWN}} side-effect status if the remote system cannot guarantee whether the request was already processed. That ambiguity must remain visible until reconciliation resolves it.

\subsection{Contract invariants}
\label{sec:ch16-contract-invariants}

The proposed subsystem should enforce at least the following invariants:

\begin{enumerate}
  \item Every scheduler, authorization, monitor, and stop decision records the exact uncertainty estimate revision it consumed.
  \item An estimate can be consumed only if its target, domain, freshness, and status satisfy the component's policy. Unknown or out-of-domain uncertainty is not silently upgraded to calibrated uncertainty.
  \item No external action is dispatched without a valid authorization decision and a resource reservation for the required dimensions.
  \item A model proposal cannot enlarge a budget, approve its own tool call, erase a policy denial, or override an authoritative interruption.
  \item Concurrent reservations cannot drive any resource dimension below its limit.
  \item A reservation is committed or reconciled exactly once, even when duplicate events or retries occur.
  \item Interruption and cancellation requests are idempotent, but acknowledgement is not treated as proof of side-effect reversal.
  \item A resume operation validates the checkpoint, current goal scope, authorization, uncertainty freshness, policy version, deadline, and remaining resources.
  \item A terminal stop has an explicit reason and evidence reference. \texttt{\detokenize{COMPLETED}} requires verified postconditions; \texttt{\detokenize{STOPPED}} does not imply success.
  \item Only verified outcome labels enter calibration feedback, and the prediction/calibrator version used for the original decision remains immutable in the audit record.
\end{enumerate}

These invariants are more important than the exact class names. An implementation in another language or framework should preserve the same boundaries.

\subsection{Persistence and concurrent reservations}
\label{sec:ch16-persistence}

A scheduler often operates under concurrency. Two workers may inspect the same budget at nearly the same time. Both may see one remaining tool call and decide that they can start. Without an atomic reservation, the system overspends even though each worker's local calculation looked valid. A reservation therefore has to be committed against the authoritative budget version.

A relational implementation can use a transaction that locks the relevant budget row or performs a compare-and-swap update against an expected version. The transaction validates every requested resource dimension, creates a unique reservation record, updates the reserved totals, and appends a ledger event. An alternative store can provide equivalent conditional-write semantics. The choice of database is secondary; the invariant is that no two concurrent decisions can both claim the same last unit of capacity.

The external dispatch boundary is more difficult. A database transaction cannot atomically commit with an arbitrary remote tool call. The runtime can use a transactional outbox: commit the reservation and a uniquely identified dispatch event together, then let a dispatcher deliver the event. Delivery may be repeated after a crash, so the tool adapter should use an idempotency key when the remote contract supports it. If it does not, the runtime must reconcile before retrying a possibly non-idempotent operation. It must not claim universal exactly-once execution across a network boundary.

Consider six failure windows:

\begin{enumerate}
  \item \textbf{Reserved but never dispatched.} A worker crashes after reserving capacity. The reservation remains visible and can be released only after a timeout or reconciliation proves that no dispatch occurred.
  \item \textbf{Dispatched but not recorded.} The remote tool receives the request, but the worker crashes before storing the response. The ledger must show an uncertain in-flight operation, not an automatic failure that permits an unsafe duplicate.
  \item \textbf{Interruption races with dispatch.} The governor commits an interruption while a worker is preparing a tool call. A generation or fencing token can prevent an old worker from dispatching after it has lost authority, but an already accepted remote request may still complete.
  \item \textbf{Concurrent use of the final resource unit.} Two schedulers try to reserve the same final tool call. The conditional transaction admits one and rejects or defers the other.
  \item \textbf{Late cancellation response.} A cancellation reply arrives after the external system may have performed the operation. The side-effect status remains unknown until authoritative reconciliation.
  \item \textbf{Stale resume.} A suspended task is resumed after its calibration policy, permission, external world state, or deadline has changed. Resume validation must reject or replan rather than blindly continuing.
\end{enumerate}

\begin{warningbox}
\textbf{Do not promise exactly-once side effects because the local code is idempotent.} Idempotency inside the agent database does not make a remote payment, deployment, or administrative API idempotent. Use the external service's idempotency contract where available; otherwise reconcile the actual state before repeating a possibly consequential operation.
\end{warningbox}

\subsection{Interruption and resume protocol}
\label{sec:ch16-resume-protocol}

A request to interrupt an activity should be an authoritative, versioned event. The worker receives the event, reaches its declared safe boundary, persists its checkpoint, and returns the result of the interruption attempt. If the worker cannot stop safely, it reports the in-flight operation and the reason. The governor then decides whether to wait, escalate, or invoke a tool-specific cancellation path.

Resumption should begin in \texttt{\detokenize{RESUME_CHECK}}. The runtime verifies that the checkpoint belongs to the current activity generation; the goal and scope remain valid; the permission is still in force; the uncertainty lineage is fresh and applicable; the required resource reservation can be reacquired; the deadline has not expired; and no unresolved external side effect makes continuation unsafe. Only after these checks may the state transition back to \texttt{\detokenize{READY}} or \texttt{\detokenize{RUNNING}}.

A resumed task must not simply replay every step from its previous prompt. It should use the execution ledger to identify what was completed, what was only requested, what remains unknown, and which side effects need reconciliation. This is the difference between recovering a durable runtime process and repeating a textual plan.

\subsection{Vertical-slice implementation sequence}
\label{sec:ch16-implementation-sequence}

A team should build the mechanism in small vertical slices so that each new control edge is observable before the next one is added:

\begin{enumerate}
  \item Define target identity, uncertainty lineage, versioning, and evidence references.
  \item Define typed resource dimensions and a budget ledger with transactional reservations.
  \item Add a single calibrated-estimate ingestion path with explicit stale and out-of-domain states.
  \item Propagate the estimate through one read-only diagnostic path and record the revision consumed by planning.
  \item Add a scheduler decision between two competing diagnostics with different resource vectors.
  \item Add authorization for one tool, ensuring that the policy gate—not the model—owns the decision.
  \item Add monitoring and one safe interruption boundary for a cooperative local operation.
  \item Add a non-cooperative external-tool case and preserve unknown side-effect status after cancellation timeout.
  \item Add suspension/resumption with checkpoint and policy revalidation.
  \item Add terminal stop reasons, outcome linking, and a calibration evaluation job that is independent of live dispatch.
\end{enumerate}

Do not start with a sophisticated learned scheduler. First establish the authoritative state, typed resources, and reliable events. A simple deterministic baseline is easier to test and provides a control against which a later adaptive policy can be compared.

\section{Product X: investigate without overspending or overstepping}
\label{sec:ch16-productx}

\BookExampleLabel{} The following scenario is a teaching example created for this book. Its numeric values and event sequence are hypothetical; they are not measurements from EPIC, Sigma, Fan and Lan, or the modern systems discussed earlier.

\subsection{Initial conditions}

Product X's payment service is returning intermittent HTTP 500 responses. The principal authorizes read-only diagnostics, evidence collection, and reporting. A restart, rollback, configuration write, or payment mutation is outside the agent's current authority and would require a separate approval. The agent's resource budget contains limits for inference tokens, model calls, two distinct tool quotas, wall-clock time, and verification capacity. The budget reserves some time and tool capacity for a final status check.

The first evidence package contains a service-health snapshot, a bounded log excerpt, and the deployment timestamp. The current hypothesis is that the recent deployment caused the errors. The estimate is marked calibrated only for the defined event ``the deployment change is associated with the observed error window'' and for the domain represented by the calibration evidence. It is not a probability that rolling back the deployment will be safe or successful.

\subsection{Choosing between investigation paths}

The scheduler receives three proposals: inspect the deployment diff, query a read-only dependency-health endpoint, and run a more expensive trace analysis. Each proposal carries an estimated resource vector and an expected diagnostic contribution. The dependency-health query is cheap in tokens but consumes the only remaining request to a rate-limited dependency service. The trace analysis uses more compute and wall-clock time but no external quota. The deployment diff is inexpensive and may rule out a class of configuration changes.

The scheduler does not rank the options by one arbitrary scalar cost. It checks hard resource limits, deadline, and action authorization, then compares the expected contribution of feasible candidates under the task policy. It chooses the deployment diff first and retains enough capacity for one verification step. This choice is illustrative; the proposed scheduler has no source-prescribed formula.

\subsection{New evidence changes allocation}

The diff shows no relevant payment-path changes. The uncertainty lineage for the deployment-cause proposition is updated with a new linked revision, and the estimate is no longer the same as the first estimate. The dependency-health query becomes more valuable relative to the remaining options. The scheduler reserves the separate dependency quota and a tool call, then dispatches the read-only query.

The response reports elevated latency in the payment provider's status path. That is evidence for a new proposition, not automatic proof that the provider caused the HTTP 500 responses. The system creates or updates the appropriate estimate with the new evidence references. It does not multiply the two estimates as though the log excerpt and status endpoint were independent confirmations.

\subsection{A higher-priority event causes \texorpdfstring{\gls{preemption}}{preemption} of work}

While a low-priority incident-summary task is running, a second alert indicates that the payment failure rate is spreading. The governor emits a preemption request for the summarizer. The summarizer is a cooperative local process and can checkpoint at a safe boundary. Its unconsumed reservation is reconciled, its current uncertainty references are preserved, and its state moves to \texttt{\detokenize{SUSPENDED}}. The newly urgent diagnostic is then scheduled.

This event is not equivalent to cancelling an already dispatched remote write. The example's summarizer is pausable; the external provider request is not assumed to be cancellable. The distinction determines what can safely be promised about the result.

\subsection{A restart proposal is denied}

The model now proposes restarting the payment service because it believes the provider path is stuck. The estimate is not a calibrated prediction of restart safety, and the action is outside the principal's read-only authorization. The authorization gate denies the request and records the proposal, target, evidence, policy version, and reason. The uncertainty estimate may explain why the model proposed the action; it cannot grant the missing permission.

\subsection{Stop with a verified result or escalate with uncertainty}

Suppose the read-only evidence eventually establishes that the service's payment-provider dependency is returning delayed responses and that the application is surfacing those timeouts as HTTP 500 errors. The agent can stop the investigation as complete only if the evidence meets the authorized reporting condition. It records the verified finding, supporting observations, unresolved alternatives, the uncertainty revision used in the final decision, and the resource ledger.

In a different run, the deadline may expire while the cause remains uncertain. The agent then stops with a reason such as \texttt{\detokenize{BUDGET_EXHAUSTED_INSUFFICIENT_EVIDENCE}} and reports what was and was not verified. It does not call the incident resolved, claim that the provider was the cause, or infer that the task was impossible. The result is a safe, explicit stopping decision rather than a false success.

\BookDiagram[0.64\textheight]{Tex/diagrams/generated/part2/chapter16/productx-b6-incident-sequence.png}{Product X teaching sequence: evidence updates a proposal's uncertainty lineage, resource allocation changes, a lower-priority task is preempted, a write action is denied, and the investigation stops with a verified or explicitly insufficient result.}{fig:ch16-productx}

\subsection{What the example demonstrates}

The example uses one uncertainty lineage for each consequential proposition and makes its relationship to each control decision visible. New evidence can change the allocation. The scheduler can protect verification resources. A higher-priority event can preempt cooperative work. A tool proposal can be denied for lack of authority even when the model recommends it. The investigation can complete with evidence or stop without pretending that an unknown result is known. These are the control connections that the B6 hypothesis asks a system to make testable.

\section{Evaluation: can the coupling survive strong baselines?}
\label{sec:ch16-evaluation}

A proposed runtime is not validated because its records look sensible or its diagrams form a complete architecture. The control edge must be instrumented and compared against systems that already have resource scheduling and uncertainty-conditioned action. Fan and Lan's Appendix D specifies the trace, outcomes, safety/efficiency measures, and rejection condition for B6. It also requires at least three distinct ablations. These are experimental requirements, not evidence that the proposed system has already improved performance.\FanLanCite{\S VI-C, B6; Table VI; Appendix D, Table X}

\subsection{Required state trace}
\label{sec:ch16-required-trace}

The experiment must record uncertainty provenance at each stage, allocated budget, interruption events, and stop events. To evaluate the engineering proposal, the trace should also include:

\begin{itemize}
  \item estimate identity, target definition, calibration domain/version, status, and revision consumed by each decision;
  \item evidence references and the relationship between predecessor and successor estimates;
  \item resource limits, reservations, actual usage, releases, and reconciliation events by dimension;
  \item candidate actions, feasibility checks, expected diagnostic contribution, and scheduler decision;
  \item authorization request, policy version, result, reason, and human approval where applicable;
  \item preemption/interruption request, safe-boundary acknowledgement, checkpoint identity, and resume validation;
  \item cancellation request and result, kept separate from the external side-effect status; and
  \item terminal stop reason, supporting evidence, and verified outcome.
\end{itemize}

A task-level success score alone cannot establish that the uncertainty propagation caused the difference. The study needs internal-state instrumentation so that an evaluator can see whether the estimate reached the scheduler or authorization gate and whether it changed the decision. This matches the paper's requirement to hold the base model, prompt, tools, environment, and task budget fixed while recording internal transitions.\FanLanCite{Appendix D}

\subsection{Primary outcomes and safety/efficiency measures}
\label{sec:ch16-outcomes}

Preserve the source's primary B6 outcomes:

\begin{itemize}
  \item \textbf{\gls{selectiverisk}:} the risk among decisions the system elects to make rather than \gls{abstention} from or defer;\index{abstention}
  \item \textbf{Calibration:} agreement between predicted uncertainty and observed outcomes for the specified prediction target and evaluation population; and
  \item \textbf{Irreversible-action errors:} cases in which the system takes an irreversible or consequential action incorrectly or without the required authority.
\end{itemize}

The source also names three safety/efficiency outcomes: help precision, \gls{verificationcost}, and interruption latency.\index{verification cost} These should remain visible in the report. The implementation study may add resource use by dimension, correct stopping rate, false interruption rate, missed interruption rate, task completion after preemption, stale-estimate use, budget-accounting correctness, and side-effect reconciliation accuracy. Those additional measures are engineering extensions; they do not replace the required source outcomes.\FanLanCite{Appendix D, Table X, B6 row}

\subsection{Closest baselines}
\label{sec:ch16-baselines}

The comparison should include strong, well-described baselines rather than intentionally weakened alternatives:

\begin{enumerate}
  \item \textbf{Uncertainty-conditioned action baseline.} Use the strongest applicable baseline for the target task, such as calibrated help seeking or cost-aware exploration. State which actions its uncertainty estimate governs.
  \item \textbf{Resource scheduler baseline.} Use a resource manager such as AgentRM or the resource controls identified in AIOS, without adding the proposed cross-stage uncertainty lineage.
  \item \textbf{Budget-aware search baseline.} Retain a budget-sensitive search policy where the task uses search, so the new scheduler must demonstrate value beyond budget-aware node selection.
  \item \textbf{Fixed-budget/fixed-threshold baseline.} Use the same hard limits and a documented fixed action threshold, providing a simple reproducible control.
  \item \textbf{Heuristic cost-sensitive orchestration baseline.} Include Utility-Guided Agent Orchestration or a faithful reproduction of its cost/uncertainty decision shape, while preserving the source's note that its uncertainty signal is heuristic.
  \item \textbf{Full B6 proposal.} Enable the versioned uncertainty lineage, resource-vector allocation, authorization integration, monitoring, interruption, and stopping policy.
\end{enumerate}

For each baseline, report which state it consumes, which control decision it owns, whether it is calibrated, what resources it controls, and whether it can interrupt or resume work. A baseline that already supports one of the required edges should receive credit for that edge. The comparison is about the remaining coupling, not about rediscovering individual capabilities.

\subsection{Three separate ablations}
\label{sec:ch16-ablations}

Fan and Lan require at least three ablations. They must not be collapsed into one generic ``remove uncertainty'' condition because they isolate different causal contributions.\FanLanCite{Appendix D, text following Table X}

\begin{enumerate}
  \item \textbf{Remove the new state object.} Keep the baseline controls but remove the propagated uncertainty-lineage object. This tests whether the added representation and continuity contribute beyond local signals already available to the baseline.
  \item \textbf{Remove the authoritative control edge while retaining logging.} Continue to record uncertainty and budget events, but prevent the governor from using them to change scheduling, authorization, or interruption. This tests whether the uncertainty state changes execution rather than merely producing a more detailed audit log.
  \item \textbf{Remove the learning update while retaining the runtime transition.} Keep the immediate control path but freeze calibration or adaptive-policy updates. This tests the adaptive value of learning separately from the value of the runtime transition. Because the B6 record does not prescribe a particular online learning algorithm, this ablation applies to the book's proposed learning extension and must not be misrepresented as a source-defined B6 update law.
\end{enumerate}

\BookDiagram[0.56\textheight]{Tex/diagrams/generated/part2/chapter16/b6-evaluation-ablation.png}{Evaluation design for B6: preserve the required trace and outcomes, then isolate the state object, authoritative control edge, and optional learning update in three distinct ablations.}{fig:ch16-evaluation}

\subsection{Task suite and experimental controls}
\label{sec:ch16-task-suite}

The task suite should vary uncertainty quality and drift, consequence severity, reversibility, resource scarcity, task duration, parallel contention, verification cost, human-help availability, interruption timing, safe cancellation capability, and deadline pressure. It should include cases in which additional verification reduces risk, cases in which it adds cost without reducing risk, cases where help seeking is appropriate, and cases where excessive abstention prevents a correct completion. It should include both cooperative local work that can be paused at a safe boundary and external operations that may not be cancellable after dispatch.

The base model, prompt, tool set, environment, and task budget should be held fixed when comparing a controller with its baseline. Use matched seeds and controlled event injection where feasible. Before execution, preregister the task population, experimental unit, model and dependency versions, stochastic replications, sample-size rule, metric formulas, equivalence or non-inferiority margin, uncertainty intervals or statistical tests, exclusion criteria, and stopping rule. A simpler baseline should count as matching only if it is non-inferior under the preregistered primary margin while using no more resources.\FanLanCite{Appendix D}

\subsection{Falsification criteria}
\label{sec:ch16-falsification}

Reject or narrow the B6 proposal if any of the following occurs:

\begin{itemize}
  \item uncertainty cannot be traced across stages or becomes miscalibrated as its target/context changes;
  \item the propagated estimate does not change resource allocation or an authoritative action decision;
  \item the full coupling does not improve selective risk, calibration, or irreversible-action outcomes over strong simpler baselines;
  \item verification cost increases without a corresponding reduction in relevant risk;
  \item help seeking becomes so frequent that the system abstains from tasks it could correctly complete, or help precision becomes poor;
  \item interruption or preemption prevents correct completion, corrupts state, or loses resource-accounting information;
  \item a simpler scheduler or uncertainty policy is non-inferior while using no more resources; or
  \item the mechanism adds operational complexity without measurable benefit.
\end{itemize}

The four source failure conditions remain central: sequential miscalibration, excessive abstention, verification cost without risk reduction, and interruption that prevents correct completion. A result in which the system merely logs more uncertainty without changing control is not evidence for B6. A result in which it changes control but makes more irreversible-action errors is a failure, not a successful implementation.\FanLanCite{\S VI-C, B6; Table VI; Appendix D, Table X}

\section{Boundaries, dependencies and recap}
\label{sec:ch16-boundaries-recap}

B6 is not an independent plug-in that replaces the other mechanisms. Fan and Lan's dependency analysis assigns different responsibilities to the bundles. B4 supplies the commitment boundary that scopes the goal and constrains whether recovery or method changes are acceptable. B2 supplies typed failure and bounded recovery state. B1 supplies memory candidates whose accessibility and value may affect what is considered. B3 controls admission to a bounded workspace. The migrated B5 skill-governance case supplies a route/verify/repair/fallback baseline, but Chapter 17 is reserved for the separate explanation of why B5 is closed by GraSP. B6 governs risk, resource use, verification, tool authorization, interruption, and stopping across these controls.\FanLanCite{\S VI-D; Figs. 8--9}

The current chapter has focused on B6 itself. It has not attempted to compose every residual mechanism into one complete architecture; that is the work reserved for Chapter 18 after the mechanism chapters are complete. Chapter 17 will examine a different kind of result: B5 is not a remaining residual hypothesis in the paper's frozen corpus, because closest-baseline screening identifies GraSP as an implemented E3/D4 convergence case. The distinction is methodological. A candidate should be closed when a reviewed baseline already implements the claimed chain, not retained simply because its historical origin is interesting.

\begin{chapterrecap}
B6 connects calibrated uncertainty to resource allocation, interruption, and stopping. EPIC contributes explicit timing, parallel pathways, and bottleneck control. Sigma contributes uncertainty-bearing state that participates in inference and decision. Modern systems already implement important partial controls: resource managers schedule and admit work, uncertainty-aware systems can guide tools or help seeking, and cost-sensitive orchestrators can choose whether to retrieve, verify, or stop. Fan and Lan's residual claim is the coupling that remains: one calibrated uncertainty lineage should accompany consequential proposals across memory, planning, authorization, execution monitoring, and stopping, while influencing resource allocation and cognitive interruption.

The seven-field analysis preserves explicit unknowns in persistence and learning. The engineering proposal therefore supplies concrete records, resource vectors, authorization contracts, state transitions, persistence rules, and tests without attributing them to the source paper.

A useful B6 implementation must distinguish uncertainty from permission, budget exhaustion from goal failure, preemption from cancellation, cancellation acknowledgement from side-effect reversal, and stopping from verified completion. Its value must be established through a traceable internal state, strong baselines, the three separate ablations, and the source's falsification conditions. If calibrated uncertainty does not change resource/action control or merely increases cost without reducing risk, the proposed coupling should be rejected or narrowed.
\end{chapterrecap}

\begin{reviewquestions}{Review questions and exercises}
\begin{enumerate}
  \item Why is a local confidence value not enough to establish the B6 mechanism?
  \item What does EPIC contribute that a list of parallel workers or a thread pool does not?
  \item What does Sigma contribute that a confidence score attached to one tool response does not?
  \item Why are resource scheduling and uncertainty-conditioned tool use not sufficient separately?
  \item Which fields of the B6 seven-field record are supported by Appendix C-Q, and which persistence or learning details remain unspecified?
  \item Why should a scheduler use a resource vector rather than one scalar cost when tool quotas, tokens, compute and deadlines are scarce?
  \item Give an example where a high-confidence proposal must still be denied authorization.
  \item Distinguish preemption, interruption, cancellation, suspension, resumption and stopping.
  \item Why does a cancellation acknowledgement not prove that a remote side effect was reversed?
  \item What causal contribution does each of the three required ablations isolate?
  \item Describe one result that would falsify the claim that cross-stage uncertainty propagation adds value.
  \item How does B6 complement B4 without taking authority to change or abandon the goal?
\end{enumerate}
\end{reviewquestions}

\section*{Source and example note}

The theoretical account in this chapter is grounded in [Fan \& Lan], principally their historical EPIC/Sigma discussion, modern resource and uncertainty mappings, B6 residual hypothesis, and experimental protocol. EPIC, Sigma, AIOS, AgentRM, budget-aware value search, UALA, KnowNo, Calibrate-Then-Act, Utility-Guided Agent Orchestration, and the runtime interfaces are discussed only in the roles in which Fan and Lan cite and analyze them. The chapter does not introduce a third formal theoretical foundation. The Java-like contracts, resource-vector schema, scheduler responsibilities, authorization record, interruption state machine, persistence protocol, and implementation sequence are engineering proposals created for this book. Product X is a book-created teaching example. All numeric values used in the example are hypothetical and are not empirical findings.

\section*{Looking ahead}

Chapter 17 examines the closed B5 candidate and the closest-baseline result that removes it from the residual list: GraSP. It will explain why closing a proposed gap is a substantive research result, and why the claim remains bounded by the paper's frozen corpus and the evaluated settings. Only after Chapters 9--17 have been completed and reviewed will Chapter 18 synthesize the residual mechanisms and their dependencies into a compositional research agenda.

\chapter{The Closed Candidate B5 and GraSP}
\begin{notebox}
Chapter placeholder. No prose has been drafted. This chapter will be outlined and approved before writing.
\end{notebox}

\chapter{How the Mechanisms Work Together}
\label{ch:mechanism-composition}

\section*{Chapter at a glance}
\phantomsection
\addcontentsline{toc}{section}{Chapter at a glance}

This chapter synthesizes the five residual mechanism hypotheses and the closed B5/GraSP case. It examines cross-bundle dependencies, authority conflicts, implementation order, typed runtime contracts, a composed Product X trace, and the experiments needed to determine whether composition adds value.

\begin{learningbox}
This final chapter of Part II explains how the five residual control bundles identified by Fan and Lan can be composed around an existing skill-governance capability. You will learn to distinguish shared runtime state from shared authority, reconstruct the dependencies among B1, B2, B3, B4, B6, and GraSP, design typed boundaries between them, and evaluate the composition without assuming that connected components automatically cooperate. The proposed Java-oriented contracts and Product X trace are engineering examples, not mechanisms claimed to have been implemented or validated by the source paper.
\end{learningbox}

\section{From mechanism studies to composition}\index{mechanism composition}
\label{sec:ch18-from-mechanisms-to-composition}

Part II has moved from broad architecture labels to the control semantics of individual mechanisms. This chapter asks how those mechanisms can affect one another across state, authority, evidence, failure, and learning boundaries. Connecting components is not enough: the composition must preserve each mechanism's ownership and justify every cross-bundle transition.

\subsection{Five residual bundles and one closed case}\index{residual control bundles}
\label{subsec:ch18-five-and-one}

The table summarizes the five residual control hypotheses and the distinct role of B5/GraSP. The bundle descriptions are compact here; their source mechanisms and engineering bridges remain in the preceding chapters.

\begin{table}[htbp]
\centering
\small
\caption{The five residual bundles and the closed B5 capability in the synthesis.}
\label{tab:ch18-five-and-one}
\begin{tabularx}{\linewidth}{p{1.4cm}p{3.2cm}X}
\toprule
\textbf{Bundle} & \textbf{Control concern} & \textbf{Role in the composition} \\
\midrule
B1 & Memory accessibility, latency, downstream value, and maintenance cost & Supplies memory candidates with provenance and value/cost information. \\
B2 & Typed failure, protected recovery, typed return, and safe compilation & Solves a bounded in-episode failure while preserving the parent task. \\
B3 & Typed-content competition, bounded workspace, broadcast, and admission learning & Decides which heterogeneous evidence becomes available to the next inference cycle. \\
B4 & Persistent intention, reconsideration, and method authority & Preserves the authorized goal and governs whether the current method should continue or change. \\
B6 & Propagated uncertainty, resources, authorization, interruption, and stopping & Governs risk and resource consequences across stages. \\
B5 / GraSP & Calibrated skill routing, typed DAG execution, verification, repair, and fallback & Supplies an existing skill-governance path and mandatory comparison baseline; it is not a residual bundle. \\
\bottomrule
\end{tabularx}
\end{table}

\BookDiagram[0.58\textheight]{Tex/diagrams/generated/part2/chapter18/five-residuals-one-closed-case.png}{The five residual control bundles remain research hypotheses. GraSP is represented separately as the existing capability that closes the B5 candidate and participates in the proposed composition.}{fig:ch18-five-and-one}

\subsection{Connected mechanisms can still conflict}
\label{subsec:ch18-conflicts}

A system can contain all of these components and still behave incoherently. Consider five conflicts. First, memory retrieval may prioritize a highly activated record that is too costly or stale for the current task; workspace admission must still decide whether the item belongs in the next reasoning cycle. Second, a recovery generator may propose a repair that changes the parent objective; the recovery component cannot infer permission to change the intention from the fact that its subtask failed. Third, workspace broadcast can make a proposed tool call visible to the model, but visibility is not permission to execute it. Fourth, a skill graph may pass its local preconditions while its invocation violates a global deadline, risk threshold, or external-action policy. Fifth, a local subsystem may report completion while the parent intention remains unresolved.

There is also a learning conflict. If a composed runtime records only that the task succeeded, it may update the memory selector, workspace admission policy, skill library, or method policy without knowing which admitted evidence or authorized action caused the result. A correct system needs evidence of the path taken, not merely a final scalar reward. The question for the remainder of the chapter is therefore not how to connect every box to every other box. It is how to preserve the necessary dependencies while limiting authority, state mutation, and learning to the boundary where each is justified.

\section{A common runtime state, not a common authority}\index{shared runtime state}
\label{sec:ch18-common-runtime-state}

\subsection{The seven-field lens remains analytical}
\label{subsec:ch18-seven-field-lens}

Chapter~10 introduced Fan and Lan's seven-field representation: operative state substrate, control locus, transition trigger, persistence boundary, failure semantics, learning operator, and resource/governance policy. The representation helps us compare mechanisms. It is not a sequence of seven software stages, and it does not require that every runtime expose every field in a fully populated record. Unknown or undocumented fields must remain unknown rather than being supplied from an architecture's name.

At the composition boundary, the seven fields become a checklist. A message that carries state but not the authority that produced it is incomplete for a control decision. A transition that does not identify its trigger cannot be distinguished from a model suggestion. A recovery result without a parent version cannot safely be applied after the parent has changed. An outcome without attribution cannot justify a learning update. This is the purpose of \gls{outcomeattribution}. A budget value without its scope and accounting unit cannot be compared with another subsystem's limit.

\subsection{Fan and Lan's shared state representation}
\label{subsec:ch18-shared-state-representation}

The source expresses its bounded synthesis runtime using a durable state tuple:

\[
 x_t = \langle i_t, w_t, m_t, s_t, u_t, b_t, q_t \rangle,
\]

where \(i_t\) is the active intention, \(w_t\) is bounded workspace content, \(m_t\) is memory state, \(s_t\) is the compiled-skill state, \(u_t\) is uncertainty, \(b_t\) is the resource budget, and \(q_t\) is the active failure or recovery substate. The tuple identifies shared state that a synthesis runtime must make available across control decisions. It is not a production database schema, does not define a mandatory persistence technology, and does not imply that each mechanism owns or may mutate every element.\FanLanCite{\S VI-A; Fig. 9}

This distinction matters. A shared reference to the active intention does not transfer goal authority to the memory selector. A workspace snapshot can contain a recovery result without owning the recovery branch. The scheduler can inspect the method selected by B4 without changing the committed goal. A skill executor can consume an authorized invocation without obtaining permission to change the principal's objective. The runtime needs common identity and version references, but each subsystem must retain its own authoritative object and decision boundary.

\subsection{Proposal, authorization, transition, and learning}
\label{subsec:ch18-proposal-authorization-transition-learning}

Fan and Lan distinguish a model-generated proposal from the governor's authorized transition. Their compact description can be written as:

\[
 p_t = \operatorname{LLM}(w_t,i_t,q_t), \qquad
 a_t = G(x_t,p_t),
\]
\[
 x_{t+1}=T(x_t,a_t,o_{t+1}), \qquad
 \theta_{t+1}=L(\theta_t,x_t,a_t,o_{t+1}).
\]

The language model proposes or interprets using the context supplied to it. An external governor evaluates the proposal against current state and authority. The runtime then commits an allowed transition using an observation. A later learning operation changes a policy only when the outcome is actually used to update a governor, admission policy, memory policy, or skill-invocation policy. Merely storing a transcript or reflection does not establish that future control has changed.\FanLanCite{\S VI-A}

The source equation is a compact analytical model, not a complete API signature. In a real implementation, the proposal may be produced by a tool, a deterministic subsystem, a human operator, or the model; the relevant authority is determined by the requested transition, not by the identity of the proposer. The key invariant remains: a proposal is data; an authorization is a decision; a transition is a committed state change; and a learning update is a separate operation with its own evidence and version.

\begin{warningbox}
A composite runtime must not treat \texttt{proposal}, \texttt{selected}, \texttt{authorized}, \texttt{dispatched}, \texttt{completed}, and \texttt{learned} as synonyms. A tool call can be authorized but fail during dispatch. A skill node can execute but fail verification. A task can stop without satisfying its goal. A successful episode can be recorded without changing any future policy. Each state needs an explicit receipt.
\end{warningbox}

\subsection{A typed handoff across the seven fields}
\label{subsec:ch18-typed-handoff}

The book-level contracts developed later in this chapter use a common envelope containing a task identifier, event identifier, producer, intended consumer, schema version, causal parent, observed state version, and evidence references. The payload remains specific to the mechanism. A B1 memory candidate carries provenance, validity, activation or value estimates, predicted retrieval latency, and maintenance cost. A B3 workspace decision identifies the admitted cycle snapshot, rejected alternatives, capacity use, and recipients. A B2 failure event records a failure type, parent-state version, and recovery scope. A B2 resolution carries its result type, evidence, dependencies, postconditions, and side-effect history. A B4 decision identifies the goal and method versions, reconsideration trigger, and authority basis. A B6 resource decision contains reservations and uncertainty references. A GraSP execution receipt identifies the skill graph version, node checks, repair history, and fallback route.

These examples do not come from a shared schema in Fan and Lan. They are an engineering proposal for ensuring that the seven analytical fields remain inspectable when one subsystem hands information to another. A receiver should be able to answer: what is this object, who created it, what state was it based on, what is it permitted to influence, and what evidence would make it stale?

\BookDiagram[0.58\textheight]{Tex/diagrams/generated/part2/chapter18/shared-state-and-authority-boundaries.png}{A shared runtime state is exposed through stable references and snapshots, while each mechanism retains its own authoritative decision object. The model proposes; the relevant governor authorizes; the runtime commits; and a separate learning operation may update future policy.}{fig:ch18-shared-state-authority}

\subsection{Why D3 and D4 do not compose automatically}\index{migration depth!D3 and D4}
\label{subsec:ch18-d3-d4-composition}

The source's E and D axes remain independent and case-specific. E describes the evidence relation between a modern mechanism and a historical mechanism. D describes implementation depth for the particular preserved mechanism edge. A D4 memory policy does not make B1 complete, because the residual is the activation--latency--downstream-action-utility coupling. A D4 learning result in one system does not give the composite runtime a D4 code. Nor does GraSP's E3/D4 result establish documented ICARUS lineage or validate the other residual bundles.

At the system level, the right question is which specific edge is preserved, which remains absent, and whether the composite experiment demonstrates a new causal relationship. The book will not assign the composite runtime a single E/D label. The code of one component cannot be promoted to a system-level property merely because the component is present.\FanLanCite{\S II-D; § VI-A; Appendix B; Appendix F}

\section{Responsibilities and authority boundaries}\index{authority boundaries}
\label{sec:ch18-authority-boundaries}

\subsection{Composition matrix}
\label{subsec:ch18-composition-matrix}

The following matrix is a book-level synthesis of the responsibilities described in the preceding mechanism chapters. It is deliberately conservative: each subsystem owns a narrow decision, and the runtime coordinator is not an unrestricted meta-controller.

\begingroup
\footnotesize
\setlength{\tabcolsep}{2pt}
\renewcommand{\arraystretch}{1.25}
\begin{longtable}{P{1.8cm}P{2.5cm}P{3.2cm}P{4.75cm}}
\caption{Proposed responsibilities and authority boundaries in the composite runtime.}\label{tab:ch18-authority-matrix}\\
\toprule
\textbf{Owner} & \textbf{Authoritative object} & \textbf{Decision it owns} & \textbf{Important non-authority} \\
\midrule
\endfirsthead
\toprule
\textbf{Owner} & \textbf{Authoritative object} & \textbf{Decision it owns} & \textbf{Important non-authority} \\
\midrule
\endhead
\midrule
\multicolumn{4}{r}{\small Continued on the next page}\\
\endfoot
\bottomrule
\endlastfoot
B1 & Memory record and memory-policy version & Which memory candidates are accessible or retained under the memory objective and cost model & Does not decide workspace broadcast or tool execution \\
B2 & Failure record and recovery branch & Whether a typed failure opens a bounded child recovery process and whether its result is valid to return & Cannot change the parent goal or promote every trace into a skill \\
B3 & Workspace cycle and admission policy & Which typed content enters the bounded inference workspace and who receives the broadcast & Does not authorize external actions or own memory persistence \\
B4 & Intention and method version & Whether to continue, reconsider, delegate, suspend, abandon, or change the method within the authorized goal & Cannot silently amend the principal's goal or override B6 permissions \\
B6 & Uncertainty lineage and resource/authorization ledger & Risk-sensitive budget allocation, tool authorization, verification, preemption, interruption, and stopping constraints & Does not choose the task goal or replace local mechanism semantics \\
GraSP & Skill library, graph, and node execution records & Whether a retrieved skill graph is applicable, how nodes execute, verify, repair, and escalate & Does not own the goal, global resource budget, or principal authorization \\
Coord. & Event identity, version checks, transaction and routing metadata & Whether a proposed transition is structurally admissible and can be committed under the owning governors' decisions & Does not replace the policy decisions of B1--B6 or GraSP \\
\end{longtable}
\endgroup

\subsection{B1 accessibility is not B3 admission}
\label{subsec:ch18-b1-versus-b3}

B1 makes memory items accessible according to its own retrieval and memory-value logic. B3 operates later or at a different boundary: it receives heterogeneous candidates and decides which information should occupy a limited shared workspace for a particular inference cycle. The candidate types may include memories, tool observations, agent proposals, monitors, partial plans, and verified recovery results. A memory item can be accessible yet lose workspace competition. A workspace can admit an observation that is not a persistent memory. The selector for memory accessibility and the selector for shared content capacity therefore have different domains, triggers, state, and falsifiers.\FanLanCite{Table V; § VI-C, B1 and B3}

\subsection{B2 recovery owns a child problem, not the parent goal}
\label{subsec:ch18-b2-parent-boundary}

B2's recovery branch is scoped by the failure and the parent state. Its job is to investigate a typed impasse and return a resolution whose assumptions and dependencies remain valid. B4 supplies the intention boundary that states what the parent task is still trying to achieve. The recovery process may recommend a different local procedure, but it cannot infer that the principal's goal should change. If the goal itself becomes impossible or irrelevant, that is an intention-level decision, not a side effect of a failed recovery attempt.

\subsection{B3 broadcast is not execution permission}
\label{subsec:ch18-b3-not-permission}

Broadcast makes selected content available to downstream consumers for a bounded cycle. It does not authorize the action described by that content. A workspace may contain a proposed refund, a payment retry, or a destructive database operation because the content is relevant to reasoning. The action still needs the appropriate authorization and resource checks. This distinction prevents a shared workspace from becoming an ambient capability system.

\subsection{B4 owns the goal boundary; B6 owns risk and resource constraints}
\label{subsec:ch18-b4-b6-boundary}

B4 decides whether the active intention remains valid and whether the method should continue, change, delegate, suspend, or terminate under its continuation and reconsideration conditions. B6 evaluates the uncertainty and resource consequences of proposed actions and controls permission, reservations, preemption, interruption, verification allocation, and stopping constraints. B6 may veto an action that violates a risk or resource policy, but it does not silently rewrite the goal. B4 may choose a different method, but it cannot compel an unauthorized external action or exceed a hard resource limit.

This separation is important when a task is expensive but still valuable. B4 may continue to regard the goal as valid; B6 may conclude that the remaining deadline or risk budget does not permit the current route. The resulting decision can be to suspend, ask for human input, delegate, or report an unresolved state. It need not pretend that the goal has become impossible merely because one resource reservation was denied.

\subsection{GraSP is a governed execution path, not a global governor}
\label{subsec:ch18-grasp-boundary}

GraSP provides skill retrieval, typed graph compilation, node-level precondition and postcondition verification, bounded local repair, and global planner or ReAct fallback. Its controls govern the skill path described in Chapter~\ref{ch:closed-b5-grasp}. In a composite runtime, B4 supplies the goal and method scope, while B6 constrains the resource and action permission envelope. B2 can contribute an eligible recovery resolution to a future skill-library candidate, but that proposal requires its own validation and promotion process. The fact that a skill is verified for its local contract does not mean it is safe for every goal, every user, every environment, or every external-action permission.

\subsection{A coordinator preserves boundaries rather than replacing governors}
\label{subsec:ch18-coordinator}

A coordinator is useful for routing events, checking identifiers and versions, deduplicating messages, and committing a transaction after the relevant authorities have approved it. It should not become a single object that decides all memory, workspace, recovery, goal, method, tool, skill, and stopping policy. A coordinator can reject a stale or structurally invalid transition; it should not overrule a B4 goal decision or reinterpret a B6 risk threshold. The safest design is a small protocol coordinator surrounded by specialized governors whose decisions remain attributable.

\BookDiagram[0.60\textheight]{Tex/diagrams/generated/part2/chapter18/authority-responsibility-map.png}{The responsibility map assigns one narrow decision to each mechanism and shows the coordinator as a transaction and routing boundary, not as a universal policy owner.}{fig:ch18-authority-map}

\section{Cross-bundle dependencies}\index{cross-bundle dependencies}
\label{sec:ch18-cross-bundle-dependencies}

Fan and Lan explicitly state that the residual bundles are not independent plug-ins. Their Figure~8 identifies dependencies among the five residuals and the migrated B5 skill governor. The synthesis must preserve those particular relationships instead of replacing them with a generic fully connected diagram.\FanLanCite{\S VI-D; Fig. 8}

\subsection{B4 scopes B2 recovery}
\label{subsec:ch18-b4-scopes-b2}

When an execution step fails, B2 needs a parent objective and a stable parent-state reference. Without that boundary, the recovery process may solve a different problem from the one the principal authorized. B4 provides the commitment boundary: the goal remains stable while the local method or recovery procedure is investigated. B2 then controls how the failure is classified, how the child branch is opened, and what typed result may return.

A useful test compares the same typed recovery mechanism with and without the B4 intention boundary. The expected difference is not simply more successful retries. The experiment should measure whether the recovered parent still pursues the same authorized goal, whether its state remains intact, whether the recovery resolves the diagnosed failure, and what additional latency or resource cost the intention check introduces. If B4 does not alter the recovery scope or protect the objective, the dependency may not justify its extra state.

\subsection{B2 may supply verified resolutions to a GraSP-like governor}
\label{subsec:ch18-b2-to-grasp}

A resolution that solves one recovery episode is not automatically a reusable skill. Three levels must remain distinct. A local resolution completes a bounded episode. A reusable recovery rule is compiled only when the failure type, dependency conditions, result type, and postconditions define where the resolution applies. A generally invocable skill then needs the governance provided by a skill library and graph runtime, including retrieval, typed DAG validation, node verification, bounded repair, and fallback.

The proposed trace-admission experiment compares three policies: admit no B2 traces; admit every apparently successful trace; or admit only typed, independently verified resolutions with explicit applicability conditions. The middle policy is intentionally unsafe as a baseline: a positive local outcome may depend on hidden parent state, transient environment conditions, or a lucky execution path. The third policy tests whether explicit recovery scope and verification improve reuse precision. The experiment is a \gls{compositionalablation} because it isolates the admission edge while leaving the local recovery procedure intact. This is an extension experiment around GraSP, not a claim that the GraSP paper itself admitted B2-derived traces.\FanLanCite{\S VI-D; Appendix D}

\subsection{B1 and migrated skills supply candidates to B3}
\label{subsec:ch18-supply-b3}

B1 supplies memory candidates, and a skill library can supply candidate procedures or skill results. These are instances of \gls{crossbundledependency}, not evidence that the two selectors share one policy. B3's candidate domain is broader: it can include tool observations, monitor alerts, partial plans, agents, memory evidence, and verified recovery results. B3 decides what becomes part of the bounded shared workspace. A skill invocation is not the same as a broadcast decision, just as memory accessibility is not the same as workspace admission.

The dependency is useful because the workspace sees a richer candidate set than either B1 or a skill library can govern alone. It is also risky: a candidate that is repeatedly useful may crowd out rare but decisive evidence. The B3 policy must retain candidate type, provenance, confidence, urgency, interference, and downstream use where those features are relevant to the task. The system should record both admitted and rejected candidates so that suppression can be measured.

\subsection{Workspace evidence can trigger B4 reconsideration}
\label{subsec:ch18-workspace-triggers-reconsideration}

A broadcast cycle may reveal that the active method is failing, that an environmental constraint has changed, that the goal may already be satisfied, or that a previous assumption no longer holds. Such evidence can trigger reconsideration. B3 does not own the intention and does not select the replacement method. It publishes an evidence packet; B4 evaluates it under the intention's continuation and reconsideration conditions and authorizes any change in method or commitment status.

This is a feedback relationship, not a one-way pipeline. A B4 method change can create new evidence requirements, which create new candidates for B3. The runtime must associate each workspace cycle with the intention and method version to which it belongs; otherwise, stale evidence can be admitted into a new plan without an explicit relevance decision.

\subsection{B6 governs risk and resources across the other controls}
\label{subsec:ch18-b6-crosscutting}

B6 spans the other mechanisms because memory retrieval, workspace admission, recovery attempts, skill invocation, verification, and external tool use all have resource and risk consequences. This does not mean that every local limit should be merged into one scalar budget. A workspace has a capacity limit; B2 has a recovery-attempt or recursion bound; GraSP has a local repair budget; a tool may have a rate limit; and the whole task may have a deadline or monetary/API budget. These constraints interact, but they have different units, owners, replenishment rules, and failure semantics.

A proposed runtime-wide ledger can use a \gls{resourcereservation} to reserve global resources while respecting local ceilings. It should not subtract a local attempt counter twice from the global budget unless the accounting rule explicitly says that both represent distinct consumed resources. Nor should a workspace token ceiling be interpreted as a tool-call quota. Risk is also not just another additive currency: a risk policy may impose a hard threshold, require verification, or forbid an action even when resources remain. B6 is the cross-stage control relationship among uncertainty, resources, and action authority, not a license to collapse every policy into one number.\FanLanCite{\S VI-C, B6; \S VI-D; Appendix D}

\subsection{A dependency graph is not an execution sequence}
\label{subsec:ch18-dependency-not-sequence}

Figure~8 describes which controls depend on the state, authority, or evidence supplied by another control. It does not prescribe that every task must execute B1, then B3, then B4, then B2, then GraSP, and finally B6. A short task may not need persistent intention. A high-risk operation may require B6 authorization before any skill dispatch. A new observation may reopen workspace competition after a method switch. A recovery result may be considered for skill promotion after the current episode has ended, not during its execution.

The system is iterative. The order in which the runtime evaluates a dependency can change with the event and the active task. What must remain fixed is the authority rule: each transition is committed by the relevant owner, using the versions and evidence that were valid when the decision was made.

\BookDiagram[0.58\textheight]{Tex/diagrams/generated/part2/chapter18/cross-bundle-dependencies.png}{The dependency graph reconstructs the five relationships identified in Fan and Lan's Figure 8: intention scopes recovery, verified recovery can feed skill governance, memory and skills supply workspace candidates, workspace evidence can trigger reconsideration, and B6 governs risk and resources across the controls.}{fig:ch18-dependencies}

\section{From dependencies to a safe implementation order}\index{implementation order}
\label{sec:ch18-implementation-order}

A dependency graph describes the relationships a running system must preserve. An implementation sequence describes which parts should be built first so later behavior can be observed and audited. Fan and Lan recommend externalizing intention, uncertainty, budgets, and recovery state before adding more adaptive behavior; typed failure and bounded content admission follow; adaptive memory and compiled-skill policies can then be added when outcome attribution is available. They specifically warn that automatic skill compilation before state and failure boundaries are explicit can create shortcuts whose conditions cannot be audited.\FanLanCite{\S VI-D}

\subsection{First, externalize authoritative state}
\label{subsec:ch18-externalize-state}

The first milestone is not an elaborate scheduler or learning policy. It is a durable record of the state that must survive beyond one prompt or model call: active intention and goal/method versions; resource budgets and reservations; uncertainty provenance and calibration version; parent-state and recovery identifiers; action and evidence history; and versioned skill references. The purpose is to make authority and causal history observable before an adaptive policy is allowed to change future decisions.

The minimum acceptance tests should include a restart between proposal and authorization, a concurrent update to the intention, a budget reservation racing with another request, and a recovery result arriving after its parent has changed. In each case the system should either commit against the expected version or reject the stale decision without corrupting the parent task.

\subsection{Next, add typed failure and bounded content admission}
\label{subsec:ch18-typed-failure-and-workspace}

Once durable state and authority references exist, implement B2 typed failure and bounded recovery, then B3 workspace capacity, admission, and broadcast. These controls create observable transitions: a failure is classified and a child state opened; candidates compete; a workspace snapshot is created; and downstream consumers acknowledge what they received. A recovery substate and workspace cycle are different scopes. The first protects parent problem solving; the second governs shared inference context.

Acceptance tests should verify parent-state integrity, typed resolution, bounded recursion, workspace capacity, consumer receipts, rare-evidence retention, and the behavior of late events. A workspace admission policy must not silently increase the number of external actions that a skill or tool can execute.

\subsection{Then add adaptive memory and compiled-skill policies}
\label{subsec:ch18-adaptive-memory-and-skills}

B1 learning and skill-library adaptation should follow after the system can attribute outcomes to evidence, authorized actions, recovery decisions, and verification results. Otherwise, the runtime cannot tell whether a successful task was caused by a retrieved memory, a workspace admission, a recovery action, a skill invocation, or an external event. The first learned policies should therefore be narrow and reversible, with explicit provenance, versioning, evaluation windows, and rollback conditions.

GraSP should be integrated as an already capable baseline. Rebuilding its chain under a new name would neither improve the composition nor test one of the remaining residuals. The experiment should begin by invoking and instrumenting the existing skill-governance path, then add only those interactions justified by the dependency graph.

\subsection{Finally, integrate B6 across the lifecycle}
\label{subsec:ch18-integrate-b6}

The uncertainty and resource controller must consume evidence across stages rather than receiving one score at the final tool boundary. A candidate memory may need a cheap retrieval before expensive verification. A workspace may need to retain a rare high-risk signal despite lower relevance. A recovery branch may receive a bounded retry budget but no permission to repeat a charge. A skill graph may need a reservation for an irreversible operation and an additional verification budget. A deadline may require interruption of low-value work while preserving enough state to resume or report an unresolved result.

The interruption protocol must coordinate with B2's child state, B4's intention lifecycle, and GraSP's node execution. A cancelled task must not appear as a verified skill failure if cancellation was caused by a parent-level budget decision. A resumed task must not reuse a stale authorization token or commit a result derived from an old goal or environment version.

\subsection{Stage gates and dependency-aware roadmap}
\label{subsec:ch18-stage-gates}

\begin{table}[htbp]
\centering
\small
\caption{A proposed implementation order derived from the source's dependencies.}
\label{tab:ch18-stage-gates}
\begin{tabularx}{\linewidth}{p{1.0cm}p{3.2cm}X}
\toprule
\textbf{Stage} & \textbf{Primary work} & \textbf{Gate before proceeding} \\
\midrule
1 & Externalize intention, uncertainty, budgets, recovery identity, evidence, action history, and skill versions & Restart, concurrency, stale-version, and authority-boundary tests pass \\
2 & Add typed recovery and bounded workspace/broadcast & Parent integrity, typed return, capacity, candidate provenance, and consumer-receipt tests pass \\
3 & Integrate existing GraSP and add eligibility rules for reusable B2 resolutions & Skill graph verification, bounded repair/fallback, safe promotion, and rollback tests pass \\
4 & Add adaptive memory and outcome-attributed policy updates & Counterfactual and delayed-outcome tests show updates are attributed to the correct policy \\
5 & Integrate cross-stage uncertainty, resource reservations, interruption, and stopping & Risk, budget, cancellation, resume, and correct-stop tests pass across subsystem boundaries \\
6 & Evaluate staged and full composition & Preregistered comparisons show whether each dependency adds value over simpler baselines \\
\bottomrule
\end{tabularx}
\end{table}

This sequence is a proposed engineering roadmap. It elaborates Fan and Lan's ordering but is not a build plan prescribed by their paper. A project may iterate between stages; the acceptance gates are intended to prevent hidden state and unverified learning from being introduced faster than the runtime can audit them.

\BookDiagram[0.58\textheight]{Tex/diagrams/generated/part2/chapter18/implementation-order-and-stage-gates.png}{The proposed implementation order externalizes authoritative state first, adds typed recovery and bounded content admission next, and introduces adaptive memory, skill policies, and cross-stage resource control only after their evidence can be attributed.}{fig:ch18-implementation-order}

\section{The \gls{boundedsynthesisruntime}}\index{bounded synthesis runtime}
\label{sec:ch18-bounded-synthesis-runtime}

\subsection{Reconstructing the source model}
\label{subsec:ch18-reconstructing-fig9}

Fan and Lan's Figure~9 describes a bounded synthesis runtime in which the language model has proposal and interpretation authority while durable control state remains outside one model call. Memory and skills propose candidates. Workspace admission limits what enters the next inference context. A scheduler governs external actions. Typed failure creates a recovery state. Learning can update admission, memory, skills, and governance evidence. The source presents this as a synthesis runtime for testing the residual bundles around an existing GraSP-like skill governor; it does not claim that the full composition has been empirically validated.\FanLanCite{\S VI-E; Fig. 9}

The figure should be read as a control architecture, not as a set of independent boxes wired together. The durable state carries references and constraints. Proposals flow from models and candidate-producing subsystems. Governors decide whether a transition is permitted. Execution emits observations and receipts. Learning consumes attributable outcomes and changes future selection. Solid paths in the source figure represent execution and learning; dashed paths represent control authority. That distinction is more important than the visual placement of a component.

\BookDiagram[0.60\textheight]{Tex/diagrams/generated/part2/chapter18/bounded-synthesis-runtime.png}{A source-faithful view of the bounded synthesis runtime: durable state remains outside model calls; memory and skills propose candidates; workspace admission bounds inference context; the scheduler governs external actions; typed failure creates recovery state; and attributed outcomes may update later policies.}{fig:ch18-bounded-runtime}

\subsection{A four-plane teaching view}
\label{subsec:ch18-four-planes}

For implementation teaching, the source figure can be reorganized into four explanatory planes. This is a book-level diagramming choice, not a four-layer architecture proposed by Fan and Lan.

The \textbf{durable intent and evidence plane} contains intention references, memory records, uncertainty lineage, resource reservations, recovery identity, action history, and versioned skill metadata. These objects survive individual model calls and support audit. The \textbf{candidate and workspace plane} receives memory, skill, tool, monitor, plan, and agent candidates. It performs typed competition and constructs a bounded workspace snapshot for a specific inference cycle. The \textbf{execution and recovery plane} carries out authorized skill graphs and tool actions, verifies outcomes, opens bounded recovery branches, and returns typed resolutions. The \textbf{governance plane} contains specialized decision authorities: B4 for intention and method decisions, B6 for uncertainty/resource and action permission, plus the local authorities of B1, B2, B3, and GraSP. A small coordinator handles routing, version checks, and transaction sequencing without replacing those authorities.

The planes are explanatory boundaries, not mandatory deployment units. A small local runtime may deploy them in one process. A larger runtime may separate them into services or actors. The architecture is defined by the control contracts and state boundaries, not by the number of executables.

\subsection{Avoid one shared mutable object}
\label{subsec:ch18-no-shared-mutable-object}

A tempting implementation is to create one large mutable `AgentState` object and let every component update whichever fields it needs. This is simple at first and difficult to reason about under concurrency. It gives every component ambient mutation authority; makes it unclear which version a proposal observed; and permits a recovery process to update the parent intention or a memory learner to alter workspace state without an explicit transition.

A safer design uses one stable runtime identity and a set of independently versioned authoritative aggregates. The intention store owns the intention version. The memory store owns memory records and its policy version. The workspace manager owns the current cycle snapshot. The recovery manager owns child branch state. The skill library owns graph and skill versions. The resource ledger owns reservations and accounting. The action ledger records dispatch and outcome evidence. Components exchange references, immutable snapshots, and typed events. A commit includes the expected version of every aggregate it intends to change.

This is an engineering proposal. The source calls for explicit shared state and governance; it does not prescribe a particular database, event bus, aggregate boundary, or consistency model.\FanLanCite{\S VI-A; \S VI-E}

\subsection{The proposal-and-commit protocol}
\label{subsec:ch18-proposal-commit}

A cross-bundle operation should follow a recognizable protocol. A proposer emits an intent-bearing request that states what it wants to change and why. The relevant governor checks its own authority, current evidence, expected versions, constraints, and resource requirements. If several authorities are involved, each returns a scoped decision rather than one ambiguous approval. The coordinator then attempts an atomic commit of the permitted transition. Only after a committed operation is dispatched does the runtime record an external effect. An observation is attached to the action and verified against the expected postcondition. A later learning update is submitted only with the outcome and the policy version that made the decision.

\begin{enumerate}
  \item \textbf{Proposed:} a component requests a specific transition and supplies its evidence references.
  \item \textbf{Authorized:} the owning governor accepts, denies, or defers the request within its scope.
  \item \textbf{Committed:} the runtime updates the relevant versioned state using optimistic concurrency or a transaction.
  \item \textbf{Dispatched:} an external operation is sent only after required permissions and reservations are current.
  \item \textbf{Observed:} the runtime records the response, timeout, cancellation, or ambiguous outcome.
  \item \textbf{Verified:} a verifier establishes the expected result or returns an unresolved/failed state.
  \item \textbf{Learned:} an attributable outcome may update a policy under a separate versioned operation.
\end{enumerate}

A denial is not an exception to be silently retried. A deferred decision may request more evidence or human input. A committed authorization does not guarantee that an external service completed the action. A timeout after dispatch creates an ambiguous external outcome, not proof of failure. The runtime must preserve that distinction to prevent duplicate actions.

\section{Cross-bundle contracts and a Java-oriented architecture}\index{typed event envelope}
\label{sec:ch18-cross-bundle-contracts}

The implementation bridge is intentionally about composition. The previous chapters have already described the internal contracts of memory, recovery, workspace, intention, uncertainty, and GraSP. This section defines the narrower information and authority boundaries needed for them to cooperate. Every type and interface below is a book-level engineering proposal.

\subsection{Runtime identity and version references}
\label{subsec:ch18-runtime-identity}

An event envelope should identify the task, episode, intention, workspace cycle, and causal parent where applicable. It should also record the producer, target consumer, event type, schema version, correlation identifier, idempotency key, creation time, expected aggregate versions, and evidence references. The envelope should be immutable after publication. Payload types should be versioned independently so that a schema update cannot silently reinterpret older recovery or authorization events.

The runtime-wide identity answers which task and causal sequence the event belongs to. Aggregate versions answer which authoritative state the event was based on. They are different. A task may stay the same while its intention changes; an intention may remain active while its method changes; a workspace cycle may advance without changing the parent goal; and a recovery branch may close while the episode continues.

\subsection{Illustrative Java-like contracts}
\label{subsec:ch18-java-contracts}

The following types show how the boundaries could be represented in Java. They are not copied from Fan and Lan or from a single runtime framework. The important feature is not the exact class names. It is that every decision is typed, scoped, versioned, and attributable. The listings are compositional interface and record fragments, not a complete compilable application; framework adapters, imports, and the internal implementations of domain value types are intentionally outside this synthesis chapter. The contracts needed to reason about authority, idempotency, versions, and evidence are shown explicitly, while Part III provides the end-to-end Java implementation.

\begin{lstlisting}[language=Java,caption={A common event envelope and version reference}]
record VersionRef(String aggregateId, long version) {}

record EvidenceRef(
    String evidenceId,
    String kind,
    String contentHash,
    String sourceId
) {}

record EventEnvelope<T>(
    String taskId,
    String episodeId,
    String eventId,
    String correlationId,
    String idempotencyKey,
    String producer,
    String consumer,
    int schemaVersion,
    Instant occurredAt,
    List<VersionRef> observedVersions,
    List<EvidenceRef> evidence,
    T payload
) {}
\end{lstlisting}

A proposal should say what it wants to do; it should not be represented as an already authorized command. A scoped decision records the responsible authority, decision kind, reason, constraints, and state version. The transition receipt is emitted only after the runtime commits the state change.

\begin{lstlisting}[language=Java,caption={Proposal, scoped governor decision, and committed receipt}]
sealed interface Decision<T> permits Permit, Deny, Defer {}

record Permit<T>(
    T authorizedValue,
    String authorityId,
    String policyVersion,
    Instant expiresAt,
    List<String> conditions
) implements Decision<T> {}

record Deny<T>(String authorityId, String reason)
    implements Decision<T> {}

record Defer<T>(String authorityId, String missingEvidence)
    implements Decision<T> {}

record TransitionReceipt(
    String transitionId,
    String taskId,
    String transitionType,
    List<VersionRef> before,
    List<VersionRef> after,
    String authorityDecisionId,
    String outcomeStatus
) {}
\end{lstlisting}

The proposal may be model-generated, but the `Permit` must be created by the governor that owns the decision. The `expiresAt` field prevents a long-lived approval from being reused after its assumptions become stale. The runtime should also validate that the policy version and observed aggregate versions still match at commit time.

\subsection{Separate decision ports}
\label{subsec:ch18-separate-decision-ports}

An implementation can expose a small interface for each authority rather than routing everything through one general-purpose agent controller:

\begin{lstlisting}[language=Java,caption={Illustrative composition ports}]
interface IntentionGovernor {
    Decision<IntentionAction> evaluate(
        IntentionSnapshot intention,
        MethodProposal proposal,
        EvidencePacket evidence);
}

interface RecoveryGovernor {
    Decision<RecoveryCommand> evaluate(
        FailureRecord failure,
        ParentSnapshot parent,
        RecoveryProposal proposal);
}

interface WorkspaceGovernor {
    Decision<WorkspaceSnapshot> admit(
        WorkspaceCycle cycle,
        List<ContentCandidate> candidates);
}

interface ResourceGovernor {
    Decision<ResourceReservation> reserve(
        ResourceRequest request,
        UncertaintySnapshot uncertainty,
        ActionDescriptor action);
}

interface SkillGovernor {
    Decision<SkillInvocation> prepare(
        SkillGraphRef graph,
        TaskContract task,
        ResourceReservation reservation);
}
\end{lstlisting}

The interfaces are deliberately not symmetric. The workspace governor admits content, not actions. The intention governor decides whether a goal/method transition is allowed, not whether a tool operation is safe. The resource governor authorizes a resource and action envelope, not a new goal. The skill governor verifies applicability and execution state within a task contract, not the principal's broader intention. The recovery governor controls a child branch, not the parent task objective.

\subsection{Contract invariants}
\label{subsec:ch18-contract-invariants}

A compositional runtime should enforce invariants that span more than one component:

\begin{enumerate}
  \item \textbf{Goal identity:} a recovery branch and skill invocation inherit the authorized goal reference; only the intention authority can commit a goal amendment.
  \item \textbf{Method versioning:} a workspace snapshot identifies the method version for which its evidence was assembled; a method change invalidates or re-evaluates dependent candidates.
  \item \textbf{Parent integrity:} a B2 result can return only against the parent version and dependency conditions from which it was derived.
  \item \textbf{Action permission:} workspace visibility, memory retrieval, and skill applicability do not substitute for B6 authorization of an external action.
  \item \textbf{Resource reservation:} every dispatch that consumes a governed resource references an active reservation or an explicit exemption.
  \item \textbf{Verification separation:} execution completion is not verified completion; a skill node cannot advance its downstream state edge until its required postcondition is established.
  \item \textbf{Learning attribution:} a policy update identifies the evidence, action, outcome, and policy version it uses; unverified or confounded outcomes are not silently promoted.
  \item \textbf{Idempotent recovery:} retries and resume operations preserve logical operation identity so that an uncertain external result is not duplicated.
\end{enumerate}

\subsection{Event sequence and transaction boundaries}
\label{subsec:ch18-event-sequence}

The event sequence should be durable enough to reconstruct the causal path. A useful high-level stream is `CandidateObserved`, `WorkspaceAdmitted`, `ProposalCreated`, `IntentionDecisionRecorded`, `ResourceReserved`, `ActionAuthorized`, `SkillNodeStarted`, `ActionDispatched`, `ActionObserved`, `PostconditionVerified`, `RecoveryOpened`, `RecoveryResolved`, `IntentionReconsidered`, `TaskStopped`, and `PolicyUpdateCommitted`. These names are illustrative. The exact event vocabulary should be derived from the application's needs and kept stable across versions.

One transaction should not attempt to atomically include an entire distributed episode. Instead, each authoritative transition is committed with its own receipt and causal link. For an external side effect, use an outbox or equivalent reliable-dispatch pattern: commit the authorized intent and an outbox record together, dispatch asynchronously, then record the external observation. A timeout after dispatch remains ambiguous until reconciled. The event stream lets the runtime recover without pretending that a database transaction can roll back an external payment provider.

\subsection{Persistence, idempotency, and concurrency}\index{transactional state transitions}\index{idempotency}
\label{subsec:ch18-persistence-concurrency}

Versioned aggregates and compare-and-swap checks can reject stale proposals. A method-change transaction may invalidate the current workspace snapshot; a recovery-result transaction may be rejected because the parent state changed; a resource reservation may fail because another action consumed the remaining budget. These are normal control outcomes, not exceptional infrastructure failures. The runtime must return a typed conflict or defer result so the correct governor can reconsider.

An intention's events should be ordered by intention version or a comparable logical sequence. Workspace cycles should have their own sequence numbers, and recovery branches should refer to immutable parent snapshots. Resource reservations should have explicit expiry and release rules. External action dispatch should use an idempotency key when the target service supports it. Cancellation should record whether it occurred before dispatch, during dispatch, or after an ambiguous result; these states have different recovery requirements.

Learning updates require a separate concurrency boundary. If two episodes attempt to update the same memory or skill policy, the runtime should preserve the observations and policy version used by each update. It may serialize updates, apply a versioned merge, or use a dedicated learning queue. It should not lose one update or apply it twice merely because an episode was resumed.

\subsection{Graceful degradation and partial composition}\index{partial composition}
\label{subsec:ch18-graceful-degradation}

The source does not require every runtime to use every bundle. A composition should therefore support \gls{partialcomposition} through explicit capability profiles. A low-risk, short-lived request may use native retrieval and a simple tool call without B4 persistence or B3 broadcast. A multi-step task may enable B4 and B2 but not adaptive memory learning. A high-risk task may require B6 verification and a human approval path while leaving the B3 workspace at a simple deterministic selector. The presence of a capability does not mean it must be active for every task.

Degradation must be fail-closed where authority or safety is concerned. If the resource governor is unavailable, the runtime must not interpret that absence as permission. If uncertainty metadata is stale, an action requiring calibrated risk control should be deferred or routed to a conservative policy. If B3 is unavailable, the runtime may fall back to a simpler retrieval path only if the task's risk and context constraints allow it. If B2 cannot preserve a parent version, it must not return a result as applicable. If a skill verifier cannot establish its required postcondition, the skill node remains unverified and escalates or stops according to its contract.

\BookDiagram[0.60\textheight]{Tex/diagrams/generated/part2/chapter18/cross-bundle-state-handoffs.png}{Versioned state references and typed events cross subsystem boundaries. Each event carries provenance, observed versions, and evidence, while the receiving subsystem retains its own decision authority.}{fig:ch18-state-handoffs}

\section{Product X: one incident, several cooperating controls}\index{Product X!composed runtime}
\label{sec:ch18-productx-composition}

The Product X example has appeared throughout Part II. The incident is familiar: the payment service times out after a payment request. The composition problem is not simply to retry until the service responds. The runtime must determine whether the payment was applied, preserve the user's authorized goal, avoid a duplicate charge, and stop honestly when the evidence or budget is insufficient. The trace below is a book-created teaching example, not an empirical result or a sequence prescribed by the source paper.

\subsection{The authorized goal and initial state}
\label{subsec:ch18-productx-goal}

Assume the user has authorized the agent to determine the status of a particular payment and, if appropriate, complete the intended purchase under the previously approved payment terms. The agent has not been authorized to submit a second charge merely because the first request timed out. The intention record therefore separates the desired outcome from the current method. It records the payment identity, approved amount and currency, merchant, idempotency reference, continuation condition, and the boundary that prevents a duplicate charge.

The active runtime state contains the intention version, current workspace cycle, available memories, skill-library version, uncertainty lineage, global deadline and API budget, and no recovery branch yet. The language model may interpret the incident and propose next steps, but it cannot change the goal or dispatch an external payment operation on its own.

\subsection{B1 retrieves candidate memories}
\label{subsec:ch18-productx-b1}

B1 retrieves relevant records: a previous incident involving delayed payment-provider callbacks, the documented status-query procedure, and a service-specific warning that a client timeout does not establish transaction failure. Each candidate carries provenance, validity, and value/cost metadata. B1 may make these items accessible to the runtime, but it does not decide that all should enter the next model context. The retrieval event and downstream-use event are recorded separately so later evaluation can distinguish a retrieved item from one that influenced a decision.

\subsection{B3 forms a bounded workspace cycle}
\label{subsec:ch18-productx-b3}

The candidate set now includes the memory records, the current timeout response, a payment-provider status observation, a monitor alert about elevated latency, and the model's proposed reconciliation plan. B3 evaluates the typed candidates under the workspace capacity. It may admit the provider's ambiguous timeout observation, the idempotency warning, and the status-query procedure while excluding redundant history. The resulting workspace snapshot is tagged with the intention and method versions and broadcast to the designated inference and monitoring consumers.

The broadcast makes the warning available; it does not authorize a second payment request. If later evidence shows that the status endpoint is inconsistent, a new workspace cycle can include that conflict and the relevant monitor observation.

\subsection{B4 preserves the intention}
\label{subsec:ch18-productx-b4}

The model proposes to reconcile the payment status before considering another action. B4 confirms that the goal remains valid and that read-only reconciliation is an acceptable method under the committed intention. The method version advances, but the goal version does not. The decision records why the change occurred and which evidence supports it. If the user had authorized only a status report, the intention would not include permission to complete the purchase; the composition would stop at that narrower goal.

\subsection{GraSP provides an existing skill path}
\label{subsec:ch18-productx-grasp}

The skill library contains a verified `ReadPaymentStatus` procedure with a typed input contract, an authenticated read-only endpoint, and a postcondition that requires a provider result linked to the same payment identity. GraSP retrieves the skill, compiles its typed graph, checks node preconditions, and prepares the invocation. The graph is applicable to the current intention, but that does not itself authorize the external call. The invocation still needs a resource reservation and B6 permission.

\subsection{B6 allocates and authorizes the work}
\label{subsec:ch18-productx-b6}

The resource governor evaluates the uncertainty provenance, endpoint reliability, remaining deadline, API budget, and the risk of acting on an ambiguous payment. It reserves one read-only status request and a bounded verification allocation. It denies any proposed write operation that could submit a duplicate charge without new evidence and explicit authorization. This decision does not alter the goal. It restricts which methods may execute under the present risk and resource state.

\subsection{B2 handles a typed failure}
\label{subsec:ch18-productx-b2}

Suppose the read-only status request returns a typed conflict: the provider's first endpoint cannot determine whether settlement occurred, while the event ledger has not yet received a final callback. The runtime records a typed failure and opens a bounded B2 recovery substate against the current parent version. The recovery process may query a second authorized read-only status endpoint or inspect the notification ledger, subject to its own scope and B6 resource reservation. It returns a typed result such as \texttt{STATUS\_VERIFIED}, \texttt{STATUS\_CONFLICT}, or \texttt{STATUS\_UNRESOLVED}, including evidence, dependencies, and postconditions.

The parent intention remains unchanged. The recovery branch cannot submit a charge or rewrite the goal. If it returns a verified result, the result may be used for the current episode. Promotion into a reusable recovery rule or skill candidate is a separate operation requiring applicability conditions and verification; no trace is automatically compiled merely because the local episode appears successful.

\subsection{New evidence triggers reconsideration}
\label{subsec:ch18-productx-reconsideration}

A verified settlement result or a persistent conflict enters a new B3 workspace cycle. B4 receives the evidence and decides whether the current method remains appropriate. If the payment is settled, the goal may be satisfied and the intention can terminate as complete. If the status remains conflicting but the deadline allows another read-only check, B4 may continue under a new method version and B6 may reserve the additional resources. If the deadline is nearly exhausted, B6 can interrupt low-value work and constrain further verification. The transitions can iterate; there is no requirement that each bundle run only once.

\subsection{Stop with verified evidence or an explicit unresolved result}
\label{subsec:ch18-productx-stop}

If evidence proves that settlement occurred, the agent reports the payment as settled and stops. If evidence establishes that no charge was applied and the user-authorized intention allows a completion action, the runtime must still obtain the required authorization and resource reservation before dispatch. If the deadline expires while the status remains unknown, the agent reports that the result is unresolved and escalates. It does not infer failure from a timeout, claim that the provider caused the incident without evidence, or submit a second charge to force a definitive response.

\begin{table}[htbp]
\centering
\small
\caption{Product X: mechanism responsibility along the incident trace.}
\label{tab:ch18-productx-trace}
\begin{tabularx}{\linewidth}{p{1.35cm}p{3.0cm}X}
\toprule
\textbf{Step} & \textbf{Mechanism} & \textbf{State or authority consequence} \\
\midrule
1 & B4 & Establishes the authorized payment-status goal and duplicate-charge boundary. \\
2 & B1 & Retrieves candidate memories with provenance and value/cost evidence. \\
3 & B3 & Admits a bounded set of heterogeneous candidates for the inference cycle. \\
4 & B4 & Accepts read-only reconciliation as a method change without changing the goal. \\
5 & GraSP & Retrieves and verifies the applicable read-only skill graph. \\
6 & B6 & Reserves resources and authorizes the read-only external request. \\
7 & B2 & Opens a bounded recovery branch for a typed status conflict and returns a verified or unresolved result. \\
8 & B3 and B4 & Broadcast new evidence; reconsider the method if the evidence or constraints changed. \\
9 & B6 and B4 & Stop with verified completion or an explicit unresolved result under the remaining budget. \\
\bottomrule
\end{tabularx}
\end{table}

\BookDiagram[0.62\textheight]{Tex/diagrams/generated/part2/chapter18/productx-composed-run.png}{The Product X trace shows memory candidates, workspace admission, intention control, a GraSP skill path, risk and resource authorization, typed recovery, reconsideration, and evidence-based stopping.}{fig:ch18-productx-composed}

The trace demonstrates the proposed authority boundaries. It does not demonstrate that the full composition improves task success, that the resource policy is calibrated, or that the B2-to-GraSP admission extension is safe. Those are empirical questions for the evaluation section.

\section{Failure modes created by composition}
\label{sec:ch18-compositional-failures}

A component can pass its local tests while the composition fails. Cross-bundle tests should therefore target the transitions where state, authority, or evidence crosses a boundary.

\subsection{Authority collision}
\label{subsec:ch18-authority-collision}

Two governors may issue incompatible decisions based on different state versions. For example, B4 approves a method switch while B6 rejects the required tool action. The correct result is not to let the later component overwrite the earlier decision. The runtime should preserve both scoped decisions, identify the conflicting precondition, and return a typed defer or denial to the owner that can propose a different method. A coordinator should resolve ordering and stale-version conflicts, not policy disagreements outside its authority.

\subsection{Goal drift through recovery or workspace content}
\label{subsec:ch18-goal-drift}

A recovery branch may propose a locally convenient objective, or a persuasive workspace candidate may redefine the task. The parent intention must remain the reference. The system should reject a recovery result that lacks the correct parent reference, and B4 should require an explicit authority event before a goal change. A workspace broadcast can trigger reconsideration but cannot amend the goal by repetition or salience.

\subsection{Duplicate and competing admission policies}
\label{subsec:ch18-admission-duplication}

B1 and B3 can accidentally duplicate one another if each is treated as a generic ranking model. This can lead to double counting, repeated filtering, or a policy that rewards candidates merely because another selector admitted them. The contracts should preserve the candidate domain and decision objective of each stage. The evaluation must record whether a candidate was retrieved, admitted, broadcast, consumed, and used.

\subsection{Unverified recovery becomes a reusable skill}
\label{subsec:ch18-unsafe-skill-promotion}

A B2 branch may succeed under transient conditions. If every apparent success is promoted into the skill library, the runtime may later invoke a procedure without the parent-state assumptions that made it safe. Eligibility must include a typed failure, explicit dependencies, verified postconditions, side-effect classification, and a tested applicability envelope. Skill promotion is separate from recovery completion and must support revocation or rollback.

\subsection{Skill execution bypasses global risk control}
\label{subsec:ch18-skill-bypasses-risk}

A skill may be locally valid yet violate a global deadline, permission, or risk constraint. GraSP verification establishes the skill's local contract; it does not automatically authorize the external action in the wider task. The execution adapter must obtain B6 permission at the correct boundary and re-check time-sensitive authorization before dispatch.

\subsection{Stale uncertainty and resource double counting}
\label{subsec:ch18-stale-uncertainty}

Uncertainty can become stale when new evidence arrives, the model or tool version changes, or the method is revised. A decision based on an old calibration version should not silently remain valid. Resource accounting can fail in the opposite way: a component may count a local limit as already deducted from the global budget, while the ledger subtracts it again. The runtime needs explicit unit, scope, reservation owner, and consumption receipts.

\subsection{Interrupt and resume break an invariant}
\label{subsec:ch18-interrupt-resume}

An interruption may occur while a skill node is running, a recovery branch is open, or an external request is in flight. Resume must revalidate the intention version, parent state, resource reservation, and action idempotency. If an external outcome is ambiguous, the runtime must reconcile before repeating the operation. A cancellation cannot be recorded as a normal skill failure if it was a parent-level resource decision.

\subsection{Incorrect attribution poisons learning}
\label{subsec:ch18-learning-attribution-failure}

If a task succeeds after three candidate memories, two workspace cycles, a recovery branch, and a skill invocation, the final success does not reveal which component made the decisive contribution. Updating every involved policy equally creates confounding feedback. The action ledger should identify the evidence and transitions that are causally relevant to the measured outcome. Where attribution remains uncertain, the learning update should be delayed, down-weighted, or excluded according to a preregistered rule.

\subsection{Emergent loops and the monolith trap}
\label{subsec:ch18-loops-monolith}

A new workspace observation may trigger B4 reconsideration; a method change may refresh memory retrieval; new candidates may reopen workspace competition; a resource update may interrupt recovery; and the interruption may produce new evidence. Such iteration is expected, but it must be bounded. The runtime should impose cycle limits, deadlines, reconsideration thresholds, and explicit terminal outcomes. It should also avoid solving every conflict through a new global controller. A monolithic governor may hide rather than resolve authority boundaries.

\begin{warningbox}
More coordination is not automatically better. A full composition that adds latency, resource consumption, and more failure modes without reducing risk or improving verified outcomes should be rejected in favor of the simpler runtime. The five residual bundles are hypotheses to test, not a checklist that every agent must satisfy.
\end{warningbox}

\section{Compositional evaluation: test the edges, then the system}\index{compositional evaluation}
\label{sec:ch18-compositional-evaluation}

\subsection{The research question}
\label{subsec:ch18-evaluation-question}

The principal question is whether composing the residual mechanisms through their proposed dependencies improves task success, correctness, safety, or efficiency beyond the strongest individual mechanisms and simpler combinations under matched conditions. Fan and Lan do not report a fully composed five-bundle runtime or causal results for the complete composition. Their bounded synthesis runtime is an agenda for experiment, not a result.\FanLanCite{\S VI-E; \S VII-C}

A test must therefore measure internal transitions as well as final outcomes. A higher task score does not show that B4 scoped recovery correctly or that B6's uncertainty signal changed a decision. Conversely, a mechanism may improve an internal invariant without changing the primary task metric in a short trial. Both kinds of result should be reported.

\subsection{Source-prescribed experimental discipline}
\label{subsec:ch18-preregistration}

Fan and Lan's Appendix D requires matched comparisons and preregistration. Each study should hold the base model, prompt, tools, environment, and task budget fixed within the comparison. It should identify the task population, experimental unit, model and dependency versions, stochastic replications, minimum sample size, metric formulae, equivalence or non-inferiority margin, uncertainty intervals or statistical tests, exclusion criteria, and stopping rule. Internal-state instrumentation is required because outcome scores alone do not establish that the proposed control edge caused the difference.\FanLanCite{Appendix D; Table X}

A simpler baseline counts as matching only when it is non-inferior under the preregistered primary margin while using no more resources. The exact margin and sample size are task-specific decisions for the future preregistration. They must not be invented in the chapter as if Fan and Lan had already measured them. Negative outcomes and traces should be retained; the purpose is to find out whether each dependency adds value, not to produce a diagram whose arrows all appear successful.

\subsection{Three minimum ablation types}\index{compositional ablation}
\label{subsec:ch18-three-ablation-types}

For the residual interventions, the source requires at least three distinct forms of ablation: remove the new state object; remove the authoritative control edge while retaining logging; and remove the learning update while retaining the runtime transition. These ask different questions. Removing the state object tests whether explicit representation contributes. Removing the control edge tests whether the authority has an immediate behavioral consequence rather than merely recording what happened. Removing the learning update tests whether experience-dependent adaptation adds value after the runtime transition is already present. They must not be collapsed into one generic condition called \enquote{without memory} or \enquote{without the agent}.\FanLanCite{Appendix D; Table X}

\subsection{Dependency experiment A: B2 with and without B4}
\label{subsec:ch18-ablation-b2-b4}

The first comparison runs typed B2 recovery without a persistent intention boundary and then with B4 scope, using the same typed failures and recovery budget. The primary observables are recovery success by failure type, parent-state integrity, goal drift, unauthorized scope expansion, verified completion, and resource overhead. The expected benefit is not merely fewer retries; it is a recovery branch that resolves the local failure while preserving the authorized parent objective. If the intention boundary does not change the relevant transitions or if a simpler parent-state guard achieves the same outcomes at lower cost, the added B4 dependency should be narrowed.

\subsection{Dependency experiment B: B3 with native retrieval, B1, and B1 plus B6}
\label{subsec:ch18-ablation-b1-b3-b6}

The second comparison uses three candidate sources: native/simple retrieval; B1 memory candidates followed by B3 competition and broadcast; and B1 plus B3 with B6 uncertainty and risk/resource signals. It asks whether memory-item value, heterogeneous workspace competition, and risk-aware admission contribute separately. Measures should include relevant-evidence density, decisive-evidence recall, rare-evidence suppression, workspace cost, downstream use, task success, verification burden, and risk/resource use. A selector that increases average relevance but suppresses a rare decisive observation may be worse for the task.

\subsection{Dependency experiment C: admission of B2 resolutions to GraSP}
\label{subsec:ch18-ablation-recovery-trace-admission}

The third comparison evaluates no-trace admission, admission of every apparently successful recovery trace, and admission only of typed, independently verified B2 resolutions with explicit applicability conditions. Measure repeated-failure reduction, skill invocation precision, unsafe or stale invocation, repair/fallback frequency, and regression under changed parent or environment conditions. This experiment tests whether the B2-to-GraSP dependency is useful; it does not reinterpret GraSP's original experiment or reopen B5.

\subsection{Additional interaction: B4 with B6}
\label{subsec:ch18-ablation-b4-b6}

A further comparison tests whether intention-governed method switching and cross-stage risk/resource control cooperate. Measure unauthorized action rate, correct and premature stopping, method-switch quality, resource consumption, completion after interruption, and unresolved-but-honest outcomes. B6 must be able to veto or interrupt a method without taking ownership of the goal. B4 must not override B6 permission or exceed an enforced budget. This interaction is especially important because a task can remain valuable after one method becomes too risky or expensive.

\subsection{Full composition and staged comparisons}
\label{subsec:ch18-full-composition}

Only after the dependency experiments should the full composition be evaluated. The configuration includes B1, B2, B3, B4, B6, and the existing GraSP skill-governance capability. Compare it with a strong simpler runtime, the strongest applicable individual baselines, and staged compositions that add one dependency at a time. Where relevant, include a fair skill baseline with the GraSP capability disabled or replaced. Do not attempt an indiscriminate full factorial over every possible toggle unless the sample size and task design justify it. Comparisons should be preregistered around concrete dependency claims.

\BookDiagram[0.60\textheight]{Tex/diagrams/generated/part2/chapter18/compositional-ablation-matrix.png}{The evaluation separates the source-required state, authority, and learning ablations from dependency experiments B2/B4, B1/B3/B6, B2-to-GraSP trace admission, B4/B6 interaction, and staged full-composition comparisons.}{fig:ch18-ablation-matrix}

\subsection{Required internal trace}\index{outcome attribution}
\label{subsec:ch18-required-trace}

The trace should link each decision to the state and evidence that supported it. For B1, record candidate memories, provenance, activation/value/cost estimates, retrieval latency, selection, workspace inclusion, and downstream use. For B3, record candidate types, scores or selection rationale, capacity, winner, broadcast recipients, consumer receipts, and later admission updates. For B2, record failure class, parent hash and version, child-state events, returned resolution, postconditions, and compilation eligibility. For B4, record intention, goal and method versions, reconsideration trigger, diagnosis, decision, authority, and terminal outcome. For B6, record uncertainty target, provenance and calibration version, resource budgets, reservations, tool authorization, preemption/interruption, resume, and stop reason. For GraSP, record skill and graph version, route confidence, pre/postcondition receipts, repair operations, verification, and fallback. The shared trace also includes task outcome, latency, resource ledger, and the evidence used for causal attribution.

\subsection{Outcomes and metrics}
\label{subsec:ch18-outcomes-metrics}

No single aggregate score should hide a safety regression or change in stopping behavior. Report task success and verified completion; correct stopping, unresolved-but-honest stopping, and premature stopping; memory recall and use precision; relevant-evidence density and rare decisive-evidence recall; recovery by failure type, parent integrity, repeated-failure rate, and safe transfer; intention drift, method-switch quality, oscillation, and authority violations; calibration, selective risk, irreversible-action errors, help precision, verification cost, and interruption latency; skill invocation precision, graph verification failures, repair/fallback frequency, and safe reuse; and resource use by dimension, including token/API cost, latency, deadline use, interruption overhead, and failure/restart cost.

Metrics should be selected for the task and preregistered. Some tasks do not contain irreversible actions; some do not require persistent memory; and some have no meaningful transfer setting. The report should say when a metric is not applicable rather than filling the table with zeros.

\subsection{Falsification and narrowing criteria}
\label{subsec:ch18-falsification}

Reject or narrow the integrated hypothesis if the full runtime is not non-inferior to a simpler baseline under the preregistered margin, or if it consumes more resources without a commensurate gain. Reject a dependency if removing its control edge does not change the decision it was meant to govern. Reject or redesign B4+B2 if recovery changes the goal or corrupts parent state. Narrow B1+B3 if workspace admission suppresses rare decisive evidence or fails to outperform simpler retrieval. Reject the B6 contribution if it adds verification and interruption cost without reducing risk or improving correct stopping. Reject unsafe B2-to-GraSP promotion if it increases stale or unsafe skill invocation. Reject learning claims if outcomes cannot be causally attributed. A composition that increases oscillation, deadlock, stale transitions, or external-action errors is not rescued by a higher average task score.

A negative result does not invalidate the individual mechanism studies. It narrows the conditions under which the composition is useful. The experiment should also identify which simpler baseline matched the result and which additional state or authority failed to earn its cost.\FanLanCite{Appendix D; Table X; \S VII-C}

\section{B5 as a closed case inside an open composition agenda}\index{B5!closed candidate}\index{GraSP}

GraSP belongs in the synthesis because B2-derived resolutions may need a skill-governance path, because reusable procedures can supply candidates to B3, and because skill execution must remain subject to B4 and B6 boundaries. B5 does not belong in the residual count. Its closure is methodologically important: an apparent research gap should disappear when a sufficiently strong closest baseline already composes the required control chain.\FanLanCite{\S II-E; \S VI-B; \S VI-C, B5; Table V; Appendix C-O}

Three claims must remain separate. The first is Fan and Lan's closure of B5 within their bounded and provisionally coded corpus. The second is GraSP's reported results in its own evaluated settings, which Chapter~\ref{ch:closed-b5-grasp} reviews and which originate in its primary study \cite{xia2026grasp}. The third is any future Java reconstruction or extension of the skill-governance path inside the proposed composition. None of these claims proves the other two. A local reproduction can reveal an implementation difference without disproving the source's mapping; a successful GraSP reproduction does not prove B1--B4 and B6 compose; and a successful full-composition experiment would not establish historical lineage.

The residual agenda remains open because the new causal relationships have not yet been tested as a complete runtime. The correct question is whether B1, B2, B3, B4, and B6 add compositional value around the existing capability, not whether a new B5 controller should be invented. Any future B5-adjacent proposal would need a materially new control relationship that remains absent after GraSP is treated as a mandatory baseline.

\section{What the synthesis establishes, and what remains unknown}\index{compositional limits}

\subsection{Theoretical synthesis}
\label{subsec:ch18-theoretical-synthesis}

The historical architectures contribute precise control invariants rather than a demand to reproduce their entire internal representation. ACT-R helps articulate memory accessibility, latency, and action utility. Soar makes typed failure and scoped recovery consequential before a recovery response is generated. LIDA makes bounded competition consequential through broadcast. BDI and MIDCA distinguish persistent goal commitment from method adaptation. EPIC and Sigma expose the relationship between resource contention, timing, uncertainty, and decision control. ICARUS supplies a language for grounded skills and hierarchical execution. Fan and Lan use these mechanisms to define the five residual bundles after modern baselines are examined.\FanLanCite{\S III; \S VI-B--VI-E}

Modern systems show that many individual control edges are already operational. The open agenda concerns their coupling and the preservation of authority and state across boundaries. This is why a high-depth component does not automatically close a bundle and why a closed candidate such as B5 should not be kept alive merely to make the taxonomy symmetrical.

\subsection{Engineering synthesis}
\label{subsec:ch18-engineering-synthesis}

A runtime can make the relationships testable through \gls{compositionalcontrol}: externalizing authoritative state, separating decision governors, connecting them through typed and versioned events, retaining evidence references, and ensuring that learning follows attributable outcomes. It can implement only the bundles justified by its task, risk, and resource profile. It can preserve partial composition and graceful degradation rather than requiring one monolithic cognitive architecture. These are engineering proposals in this book, not implementation details prescribed by Fan and Lan.

\subsection{Scientific synthesis}
\label{subsec:ch18-scientific-synthesis}

The five residual bundles should be evaluated against strong, current, and task-matched baselines. Each dependency should be instrumented and ablated. The full system should be compared with simpler compositions and rejected when its additional state, control, or resource cost does not earn a measurable benefit. The B5 case demonstrates that the research agenda itself must be revisable when a strong closest baseline closes a proposed gap.

\subsection{Validity limits}
\label{subsec:ch18-validity-limits}

The conclusions remain bounded by the source paper's corpus and methods: the reviewed corpus was frozen on 26 July 2026; the review does not provide a systematic-search recall estimate; the E/D coding is provisional and was not independently recoded in the submitted manuscript; the five-bundle partition depends on the authors' analytical rules and may change with different evidence or codes; modern systems use heterogeneous tasks, models, and experimental protocols; and no experiment in the source establishes causal benefit for a fully composed five-bundle runtime. GraSP closes B5 architecturally within the reviewed corpus, but universal transfer remains an external-validity question.\FanLanCite{\S II-B--II-E; \S VII-C; Appendix C; Appendix F-C}

The chapter's contracts and Product X trace do not remove these limits. They make a proposed implementation inspectable and testable. A future Java system would still require implementation, independent review of the source mapping, controlled experiments, and explicit reports of negative or inconclusive results. The composite architecture may add latency, costs, race conditions, and new failure modes. More mechanisms are not automatically a better agent.

\section{Chapter recap}\index{compositional limits}
\label{sec:ch18-recap}

\begin{chapterrecap}
Five conclusions close Part II:
\begin{enumerate}
  \item Mechanisms compose through state, authority, transition, persistence, failure, learning, and resource boundaries, not merely because their components can be connected.
  \item B1 and existing skills can supply candidates to B3, but memory accessibility, skill invocation, and workspace admission remain different decisions.
  \item B4 scopes the parent intention and recovery boundary; B2 handles typed in-episode recovery; B6 governs risk, resources, authorization, interruption, and stopping across the runtime.
  \item GraSP closes B5 and can participate as an existing skill-governance dependency, but its closure does not validate the remaining composition.
  \item Only controlled compositional experiments can establish whether the combined runtime is worth its additional state, control, learning, and resource costs.
\end{enumerate}
\end{chapterrecap}

\section{Review questions and practical exercise}
\label{sec:ch18-review-questions}

\begin{reviewquestions}{Review questions}
\begin{enumerate}
  \item Why is a graph of connected components insufficient evidence that a runtime has coherent control semantics?
  \item What does each element of Fan and Lan's shared runtime state tuple represent, and why is it not a production schema?
  \item Why must a proposal, authorization, committed transition, observation, verification result, and learning update be recorded separately?
  \item How does B4 scope B2 recovery without taking over the local failure procedure?
  \item Under what conditions may a B2 resolution become a candidate for GraSP-like skill governance?
  \item Why are B1 memory accessibility and B3 workspace admission different decisions?
  \item Why does B3 broadcast not imply permission to execute a tool or skill?
  \item How does B6 constrain B4 and GraSP without owning the goal or the skill semantics?
  \item Why is GraSP included in the synthesis while B5 is excluded from the residual count?
  \item What does E3/D4 say about the GraSP mapping, and what does it not say?
  \item Which experiment tests whether B4 prevents recovery-induced goal drift?
  \item How do native retrieval, B1 candidates, and B1 plus B6 signals isolate different B3 contributions?
  \item Why should apparently successful B2 traces not automatically be compiled into reusable skills?
  \item What evidence is necessary before an outcome can update a future policy?
  \item Give an example in which the full composition should be rejected in favor of a simpler runtime.
\end{enumerate}
\end{reviewquestions}

\paragraph{Practical exercise.} Using the Product X timeout trace, draw the state and authority boundary for each mechanism. Label every proposal, decision, authorization, transition, verification result, and learning update. Then remove one dependency at a time and predict which invariant or outcome should change. Identify the strongest simpler baseline and state a result that would falsify the proposed composition.

\section*{Source and example note}
The theoretical synthesis in this chapter is grounded in Fan and Lan's 2026 review, especially §§ VI-A--VI-E and Appendix D. The Product X sequence, Java-like contracts, four-plane teaching view, coordinator protocol, persistence model, and implementation roadmap are book-level engineering proposals and teaching examples. The source paper does not claim that this exact runtime has been built or experimentally validated. The prior chapters remain the detailed treatments of B1, B2, B3, B4, B6, and the closed B5/GraSP case.\FanLanCite{\S VI-A--VI-E; Appendix D}

\section*{Looking ahead}
Part II is complete. Part III applies the mechanism-level insights to one request-driven enterprise supply-chain disruption management system in Java. That system must preserve the distinctions developed here: state is not authority, visibility is not permission, recovery is not goal replacement, a verified skill is not a universal guarantee, and learning must be tied to attributable outcomes.


\part{Engineering a Production Supply-Chain Agentic System}
\label{part:engineering}
\chapter{Engineering Foundations for Request-Driven Enterprise Agentic Systems}
\begin{notebox}
Chapter placeholder. This chapter establishes the engineering foundations for the single request-driven enterprise supply-chain disruption management system developed in Chapter 20. The chapter outline must be approved before prose is drafted.
\end{notebox}

\chapter[Production Incident Diagnosis and Recovery]{Production Incident Diagnosis\\and Recovery}
\begin{notebox}
Chapter placeholder. No prose has been drafted. This chapter will be outlined and approved before writing.
\end{notebox}



\appendix
\chapter{Application Repository and Architecture}
\label{app:architecture}

This appendix is a navigation aid for the Java application used in Part III. It describes the Maven modules and the main responsibility boundaries visible in the supplied repository. It is not a substitute for the source code, and the mapping from CoALA or Fan and Lan concepts to Java packages is intentionally not one-to-one.

\section{Repository orientation}

The companion project is a Maven multi-module repository. The root \path{pom.xml} declares eleven modules. Each module has its own \path{pom.xml}, production sources under \path{src/main/java}, and, where applicable, tests under \path{src/test/java}. The modules are separated by responsibility so that domain rules, memory, mechanisms, planning, orchestration, and infrastructure can be inspected without treating the entire application as one undifferentiated agent class.

\begin{longtable}{P{0.26\textwidth} P{0.62\textwidth}}
\caption{Maven modules and their primary responsibilities.}\label{tab:appendix-a-modules}\\
\toprule
\textbf{Module} & \textbf{Primary responsibility} \\
\midrule
\endfirsthead
\toprule
\textbf{Module} & \textbf{Primary responsibility (continued)} \\
\midrule
\endhead
\path{supply-chain-agent-domain} & Framework-independent supply-chain entities, value objects, and domain invariants. This is where business facts and valid domain states are represented. \\
\path{supply-chain-agent-core} & Framework-neutral application and domain contracts used across the system. \\
\path{supply-chain-agent-memory} & CoALA-oriented memory models, lifecycle services, retrieval contracts, and repository ports. \\
\path{supply-chain-agent-actions} & Typed action contracts, capability enforcement, memory actions, perception, and simulator-backed grounding. \\
\path{supply-chain-agent-mechanisms} & Framework-neutral, testable engineering realizations of the selected B1, B2, B3, B4, and B6 responsibilities. \\
\path{supply-chain-agent-planning} & Separate planning strategies and their shared contracts: CoT-style sequential planning, bounded ToT-inspired search, and domain-grounded RAP-inspired MCTS. \\
\path{supply-chain-agent-simulator} & Deterministic domain transitions used for scenario execution and planning evaluation. \\
\path{supply-chain-agent-runtime} & Cognitive-cycle orchestration, graph state, decision and lifecycle services, and framework-facing adapters. \\
\path{supply-chain-agent-infrastructure-postgres} & PostgreSQL migrations and repository adapters for memory, planner sessions, checkpoints, and selected runtime records. \\
\path{supply-chain-agent-app} & Application composition and host entry points that connect domain ports, services, runtime, and persistence. \\
\path{supply-chain-agent-evaluation} & Versioned scenario fixtures and deterministic, proposal-oriented planning and workflow evaluation. \\
\bottomrule
\end{longtable}

\section{A responsibility map}

The following map is a guide to where to begin reading. A concept may cross several modules because an architectural role, a mechanism, and its persistence adapter are different responsibilities.

\begin{longtable}{P{0.23\textwidth} P{0.32\textwidth} P{0.33\textwidth}}
\caption{Conceptual responsibilities and useful code locations.}\label{tab:appendix-a-crosswalk}\\
\toprule
\textbf{Concern} & \textbf{Main module or package} & \textbf{Related boundary} \\
\midrule
\endfirsthead
\toprule
\textbf{Concern} & \textbf{Main module or package} & \textbf{Related boundary} \\
\midrule
\endhead
Business entities and invariants & \path{supply-chain-agent-domain} & Domain validation is not delegated to an LLM or graph node. \\
World transitions & \path{supply-chain-agent-simulator} & A command produces a deterministic simulated transition; the world adapter is supplied to the host. \\
CoALA memory & \path{supply-chain-agent-memory} & Repository ports are implemented by PostgreSQL adapters in the infrastructure module. \\
Internal and external actions & \path{supply-chain-agent-actions} & Typed action contracts are registered and dispatched through application/runtime services. \\
B1, B2, B3, B4, and B6 responsibilities & \path{supply-chain-agent-mechanisms} & Mechanism services remain framework-neutral and are invoked by runtime decision or lifecycle services. \\
Planning algorithms & \path{supply-chain-agent-planning} & Planner strategies are selected through a shared planning contract; they are not implemented by LangGraph4j. \\
Cognitive-cycle orchestration & \path{supply-chain-agent-runtime} & LangGraph4j provides graph execution and checkpoint integration. \\
Model-facing integration & \path{supply-chain-agent-runtime} (LangChain4j adapters) & Model access is behind an application-facing adapter; domain authority remains outside the model adapter. \\
Durable adapters & \path{supply-chain-agent-infrastructure-postgres} & PostgreSQL implements selected repository, ledger, and checkpoint boundaries. \\
Composition and configured ports & \path{supply-chain-agent-app} & The host supplies domain, world, authority, and policy adapters. \\
Scenarios and evaluation & \path{supply-chain-agent-evaluation} & Test fixtures and deterministic evaluation are not live supply-chain integrations. \\
\bottomrule
\end{longtable}

\section{The main composition path}

A useful reading path through the code is:

\begin{enumerate}
  \item Start with the domain model and simulator to understand the business facts and valid transitions.
  \item Read the memory, action, mechanism, and planning modules to understand the capabilities available to the cognitive cycle.
  \item Follow the runtime decision and lifecycle services to see when those capabilities are invoked and which state is carried between graph nodes.
  \item Inspect the PostgreSQL infrastructure module to identify which contracts have durable adapters.
  \item Finish with the application host to see how the services, repositories, graph, world adapter, and authority ports are assembled.
\end{enumerate}

The dependency structure is deliberately not a claim that each CoALA component is a separately deployed service. Several architectural responsibilities are implemented by cooperating Java classes within one application. The application uses module boundaries to make those responsibilities testable and traceable.

\section{Mechanism labels, milestone labels, and frameworks}

B1, B2, B3, B4, B5, and B6 are mechanism identifiers used by Fan and Lan. The reference application has engineering implementations for selected B1, B2, B3, B4, and B6 responsibilities. This does not establish full equivalence to every source-described mechanism. B5/GraSP is not presented as an implemented application mechanism.

M-labels such as M0--M8 appear in repository change records, architecture checks, and development documentation. They identify project milestones or workstreams, not additional mechanisms in Fan and Lan's taxonomy. The labels should not be substituted for one another.

Finally, keep planning algorithms separate from runtime frameworks. The sequential, bounded-search, and MCTS-inspired planners are separate strategies behind planning contracts. LangGraph4j coordinates graph state and control flow; LangChain4j provides the model-facing adapter boundary. Neither framework is the source of the application's domain rules, authorization policy, or planning algorithm.

\begin{notebox}
The code references in Chapters 20--24 and their QR links provide the precise source locations used in the worked explanations. This appendix gives a module-level map so that readers can navigate beyond those selected excerpts without assuming that a single module contains the whole implementation of a concept.
\end{notebox}

\chapter{Build, Configuration, and Verification}
\label{app:build}

This appendix records the toolchain and the principal build and verification paths in the supplied application. Version values are taken from the repository's root Maven configuration and wrapper properties. They describe the reference project snapshot, not a recommendation to use the same versions in unrelated projects.

\section{Required toolchain}

The root \path{pom.xml} sets the Java compiler release to 21 and its Maven Enforcer rules require Java 21 and Maven 3.9.16 or later. The repository includes a custom Maven Wrapper bootstrapper that downloads Apache Maven 3.9.16 from the configured Apache distribution endpoints and verifies the archive against a pinned SHA-512 checksum. Use the wrapper from the repository root so the project-defined Maven version is selected.

\begin{itemize}
  \item \textbf{Java:} JDK 21.
  \item \textbf{Maven:} the repository wrapper, pinned to Apache Maven 3.9.16.
  \item \textbf{PostgreSQL:} required for the PostgreSQL-backed integration-test path. The GitHub Actions workflow uses PostgreSQL 16.4.
  \item \textbf{Operating system:} the CI workflow runs on Ubuntu 24.04. The Maven project itself is not limited to that CI operating system, but local shell commands below use POSIX syntax.
\end{itemize}

Selected dependency and build-plugin pins are summarized below.

\begin{center}
\small
\begin{tabularx}{\linewidth}{P{0.43\textwidth} P{0.20\textwidth} X}
\toprule
\textbf{Component} & \textbf{Version} & \textbf{Use in the project} \\
\midrule
LangGraph4j & 1.8.20 & Graph execution and PostgreSQL checkpoint integration. \\
LangChain4j & 1.23.0 & Model-facing Java adapter API. \\
PostgreSQL JDBC & 42.7.8 & JDBC connectivity. \\
HikariCP & 7.1.0 & Bounded JDBC connection pooling. \\
SLF4J API & 2.0.20 & Logging API alignment. \\
JUnit Jupiter & 5.12.2 & Java test framework. \\
Maven Compiler Plugin & 3.14.0 & Java 21 compilation. \\
Maven Surefire Plugin & 3.5.3 & Test execution. \\
\bottomrule
\end{tabularx}
\end{center}

\section{Build and test commands}

From the repository root, the ordinary reactor build is:

\begin{lstlisting}[basicstyle=\ttfamily\small]
./mvnw -B -ntp clean install
\end{lstlisting}

This builds the Maven modules and runs tests whose conditions are satisfied in the local environment. PostgreSQL-backed tests are guarded by the system property \texttt{phase0.postgres=true}; without that property, those tests are not enabled merely because the project can compile.

To run the PostgreSQL integration-test path, first make a PostgreSQL database available and set connection variables in the shell or the local environment. The repository's \path{.env.example} documents the supported names. A typical local environment uses the following values, with the password supplied securely rather than committed to a file:

\begin{lstlisting}[basicstyle=\ttfamily\small]
export PGHOST=localhost
export PGPORT=5432
export PGDATABASE=supply_chain_agent
export PGUSER=postgres
export PGPASSWORD="<read from your local secret store>"
export PGSSLMODE=prefer

./mvnw -B -ntp -Dphase0.postgres=true clean install
\end{lstlisting}

The configuration reader also supports \texttt{PGCONNECT\_TIMEOUT}, \texttt{PGSOCKET\_TIMEOUT}, and \texttt{PGAPPNAME}. Its defaults are 10 seconds, 30 seconds, and \texttt{supply-chain-agent}, respectively. The default host, port, database, user, and SSL mode are \texttt{localhost}, 5432, \texttt{supply\_chain\_agent}, \texttt{postgres}, and \texttt{prefer}. The accepted SSL modes are \texttt{disable}, \texttt{prefer}, \texttt{require}, \texttt{verify-ca}, and \texttt{verify-full}. The password default is empty. The configuration object's string representation redacts the password; this does not remove the caller's responsibility to protect environment values and secrets.

\section{What the CI workflow verifies}

The authoritative GitHub Actions workflow is \path{.github/workflows/ci.yml}. It starts a PostgreSQL 16.4 service, sets Java 21, and invokes the Maven Wrapper with \texttt{-Dphase0.postgres=true clean install}. Before the Maven build, it runs repository architecture and evidence guard scripts. The workflow also runs additional security and evaluation-artifact steps. Those workflow steps are specific to this repository and should not be confused with a single JUnit test command.

A local successful build is useful, but it is not identical to a successful hosted CI run. The CI environment supplies a known PostgreSQL service and runs the repository's scripted guards and other workflow steps. When reporting verification, identify the exact command and environment used rather than writing only that the project ``passes CI''.

\section{The local intake smoke entry point}

The application module configures \texttt{dev.supplychain.agent.app.Phase0Application} as its Maven Exec Plugin main class. This entry point is a deliberately small deterministic smoke example. It creates a fixed-response chat model that returns \texttt{READY}, compiles the supply-chain intake graph with LangGraph4j's in-memory \texttt{MemorySaver}, submits a sample supplier-capacity request, and prints the resulting outcome, intent, and summary. It requires no provider credentials or PostgreSQL connection.

From the repository root, after building the reactor, the configured smoke entry point can be run with:

\begin{lstlisting}[basicstyle=\ttfamily\small]
./mvnw -B -ntp clean install
./mvnw -B -ntp -f supply-chain-agent-app/pom.xml exec:java
\end{lstlisting}

The first command builds and installs the reactor modules into the local Maven repository. The second invokes the app module's configured main class. It is an illustrative local path; a reader should first use the Java and Maven versions required by the project and should inspect the output rather than interpreting a successful smoke run as a full application acceptance test.

It is important not to mistake this smoke entry point for the complete durable cognitive-cycle application. It exercises the intake graph only. The durable application host is assembled through \texttt{DurableCognitiveApplicationHost.open} or \texttt{assemble}, and requires an application-supplied \texttt{Ports} configuration for the domain, lifecycle, authority, planning, perception, simulated world, and policy collaborators. The supplied repository does not turn the intake smoke main class into a universal launcher for a live supply-chain service.

\section{Configuration and operational caution}

The reference application reads PostgreSQL connection properties from environment values. Use a secret manager or local secret injection for passwords; do not place real credentials in source control, sample configuration, screenshots, or test output. The default connection-pool settings in the code are bounded values for this reference composition, not a capacity recommendation for an arbitrary deployment. Database sizing, TLS policy, secret rotation, backups, observability, and operational ownership must be assessed for a real environment.

For module responsibilities, see Appendix~\ref{app:architecture}. For the records and state boundaries that PostgreSQL does and does not own, see Appendix~\ref{app:persistence}.

\chapter{Memory, Runtime State, and Persistence}
\label{app:persistence}

Part III uses several kinds of state that are easy to conflate when reading a graph-based agent. This appendix separates CoALA memory, planning artifacts, graph checkpoints, mechanism-specific records, and the simulated external world. They may share a PostgreSQL data source in the durable host, but they serve different purposes and have different recovery semantics.

\section{Start with ownership, not the database}

A useful first question is not ``Is this stored in PostgreSQL?'' but ``Which component owns this information, and what operation is allowed to change it?'' A graph checkpoint preserves enough runtime state to resume a graph execution. A CoALA memory repository stores information intended to be retrieved or reused as memory. A mechanism-specific repository stores records needed by a mechanism's contract. An action ledger records execution identity and outcomes. None of these automatically becomes the authoritative representation of the simulated supply-chain world.

\section{Persistence map}

The durable host composes a set of PostgreSQL adapters through \texttt{PostgresAgentPersistence} and \texttt{PostgresRuntimePersistence}. The table below summarizes the principal ownership boundaries visible in the application source.

\begin{longtable}{P{0.19\textwidth} P{0.39\textwidth} P{0.32\textwidth}}
\caption{Selected state and persistence responsibilities in the durable host.}\label{tab:appendix-c-persistence}\\
\toprule
\textbf{State or record} & \textbf{Owner and adapter} & \textbf{Important boundary} \\
\midrule
\endfirsthead
\toprule
\textbf{State or record} & \textbf{Owner and adapter} & \textbf{Important boundary} \\
\midrule
\endhead
CoALA working memory & Working-memory service and PostgreSQL working-memory repository. & Active, scoped memory with a lifecycle/expiry policy; it is not the complete graph checkpoint. \\
Episodic memory & Episodic-memory service and PostgreSQL episode repository. & Stores experience records with review, idempotency, and retention behavior; an episode is not automatically verified semantic truth. \\
Semantic memory & Semantic-memory service and PostgreSQL semantic-fact repository. & Stores facts with provenance and admission policy. \\
Procedural memory & Procedural-memory service and PostgreSQL procedural-memory repository. & Stores governed procedure artifacts; activation and use are separate from merely having a stored procedure. \\
Planner sessions and artifacts & PostgreSQL planner-session and planning-artifact stores. & Preserve planner-session and planning-artifact state; a planner session is not a persistent intention. \\
Graph checkpoint & LangGraph4j PostgreSQL saver, created by the checkpoint-saver factory. & Supports graph resume by a stable thread identity derived from the memory scope and execution ID. It is not the repository for every other state type. \\
Runtime lifecycle records & Adapters created by \texttt{PostgresRuntimePersistence}, including observation, approval, action-batch, action-ledger, impasse, child-investigation, and messaging stores. & Records are scoped and have their own contracts; they should not be collapsed into an undifferentiated ``agent state'' table. \\
B1 learned features & PostgreSQL B1 memory-feature repository. & Persists B1 feature/learning evidence separately from the four base CoALA memory repositories. \\
B2 investigation and recovery & PostgreSQL impasse and child-artifact stores, child-step receipts, recovery-procedure draft repository, and immutable artifact store. & A child case disposition, a recovery draft, and parent-objective satisfaction are distinct facts. \\
B3 downstream evidence & PostgreSQL B3 outcome-evidence repository and service. & Broadcast receipts and separately verified downstream outcomes are durable; B3 workspace snapshots are not. \\
B4 recovery intention & PostgreSQL recovery-intention repository. & Versioned intention updates use compare-and-set semantics; intention persistence does not grant action authority. \\
B6 resource governance & PostgreSQL resource-budget ledger, used by the shared execution guard. & Budget reservations and leases are distinct from uncertainty assessments and from the graph's checkpoint state. \\
Optional B6 calibration profiles & Optional PostgreSQL uncertainty-calibration profile repository, assembled when the host receives the corresponding codec. & The profile repository is not a general store for every current uncertainty assessment. \\
\bottomrule
\end{longtable}

\section{Graph checkpoint is not CoALA memory}

The durable host creates the LangGraph4j PostgreSQL saver with the same managed data source used by the PostgreSQL infrastructure. The saver stores graph execution checkpoints. A deterministic thread identifier is derived from the memory-scope identifier and execution identifier, so a resumed run can address the same graph thread while a separate execution receives a different thread identity.

This checkpoint is a runtime recovery mechanism. It does not replace the four CoALA memory repositories. Working, episodic, semantic, and procedural memory have their own models, lifecycle services, retrieval rules, and repositories. A graph state may carry an identifier or a projection of information from one of those stores, but the checkpoint does not thereby become the memory store itself. Conversely, a memory record does not contain everything required to resume a graph at a particular node.

\section{Mechanism-specific boundaries}

Several mechanism boundaries deserve particular attention:

\begin{itemize}
  \item \textbf{B3 workspace.} In the durable host, workspace snapshots are intentionally stored in \texttt{InMemoryWorkspace\allowbreak{}Repository} and are per-cycle/ephemeral. The host persists broadcast receipts and downstream outcomes separately in PostgreSQL. Durable evidence about a broadcast must not be described as durable storage of the workspace snapshot itself.
  \item \textbf{B6 uncertainty.} The lifecycle port owns creation of an immutable uncertainty-assessment revision, exact-revision lookup, and the durable persistence contract for a calibrated assessment. The host can optionally assemble a calibration service backed by a PostgreSQL calibration-profile repository. That profile repository does not, by itself, persist the complete set of assessment revisions. The resource-budget ledger is a separate store with a different responsibility.
  \item \textbf{Simulated world.} The durable host receives the simulation-world adapter through its application-supplied ports. The PostgreSQL persistence composition described here stores graph, memory, planning, approval, and execution records; it does not define a general PostgreSQL repository for the entire \texttt{SupplyChainSnapshot}. A simulator or test-world adapter may maintain its world state separately. Do not infer cross-process durability of the simulated world from a durable graph checkpoint or action ledger.
  \item \textbf{Action reconciliation.} The action ledger, batch store, approvals, and observation records support idempotency, authority revalidation, and recovery from ambiguous outcomes. They do not make an external side effect and a PostgreSQL transaction one atomic operation. The application must reconcile the external-world outcome through the relevant host and domain contracts.
\end{itemize}

\section{How to reason about a restart}

When reading a restart test or designing a recovery path, identify the exact boundary being restarted. Reconstructing the application host, replacing a database connection pool, resuming a graph checkpoint, recovering an action-ledger lease, and restarting a separate JVM are not interchangeable experiments. Each test should name the state it reloads, the operation it retries or reconciles, and the evidence it uses to decide whether an external effect already occurred.

A safe summary is therefore specific: a given test may demonstrate that an intention, approval, budget reservation, checkpoint, or action record can be reloaded and used in a stated scenario. It does not establish that all application state is durable, or that the simulator's world is automatically restored by PostgreSQL. For the corresponding test categories and their limits, see Appendix~\ref{app:testing} and Chapter~24.

\chapter{Testing, Evidence, and Source Traceability}
\label{app:testing}

The application uses several test layers because no single test style can establish every property discussed in Part III. Unit tests, simulator tests, graph-workflow tests, PostgreSQL integration tests, and repository-level CI guards answer different questions. This appendix provides a route through those forms of evidence and explains their limits.

\section{Read a test by its claim}

Before treating a passing test as evidence, identify the claim it checks. A test name is a useful navigation aid, but the assertions, fixtures, test doubles, environment conditions, and setup determine what it actually establishes.

\begin{longtable}{P{0.20\textwidth} P{0.36\textwidth} P{0.34\textwidth}}
\caption{Test layers and the claims they can support.}\label{tab:appendix-d-test-layers}\\
\toprule
\textbf{Test layer} & \textbf{Representative evidence} & \textbf{What it does not establish by itself} \\
\midrule
\endfirsthead
\toprule
\textbf{Test layer} & \textbf{Representative evidence} & \textbf{What it does not establish by itself} \\
\midrule
\endhead
Domain invariants & Supply-chain snapshot invariant tests reject invalid quantities, duplicate identities, and inconsistent world states. & Correct graph orchestration or behavior against a live supply-chain system. \\
Deterministic simulator & Supply-chain simulator tests check typed commands, valid transitions, and domain effects. & External ERP, warehouse, supplier, or carrier behavior; the simulator is a modeled world. \\
Memory, planning, and mechanism units & Service and selector tests exercise local contracts and decisions for memory, B1, B2, B3, B4, B6, and planner strategies. & Correctness of every path when all components are composed into one graph. \\
Graph/workflow tests & Full-graph workflow tests exercise selected paths, including authority and lifecycle boundaries. & A universal proof for all fixtures, all failure conditions, or every research mechanism. \\
PostgreSQL integration tests & Memory-persistence, graph-checkpoint, and restart/reconciliation integration tests exercise database-backed stores and recovery paths. & Durability of state that is deliberately held outside those repositories, including the complete simulated world. \\
Durable host end-to-end test & The durable-host REQ-09 PostgreSQL end-to-end test exercises a composed host with approval, action execution, reconciliation, restart, and memory-reuse assertions. & Coverage of every mechanism branch. B2 child investigation has separate graph/workflow evidence and must not be inferred from this single REQ-09 path. \\
\bottomrule
\end{longtable}

\section{Useful starting points in the test tree}

The paths below are representative starting points. They are not a complete inventory of the repository's tests.

\begin{itemize}
  \item \textbf{Domain and simulator:} \path{supply-chain-agent-domain/src/test/java/dev/supplychain/agent/domain/world/SupplyChainSnapshotInvariantTest.java} and \path{supply-chain-agent-simulator/src/test/java/dev/supplychain/agent/simulator/SupplyChainSimulatorTest.java}.
  \item \textbf{Memory and planning:} tests under \path{supply-chain-agent-memory/src/test/java} and \path{supply-chain-agent-planning/src/test/java}, including planner-specific tests for sequential reasoning, bounded thought search, and MCTS.
  \item \textbf{Mechanisms:} tests under \path{supply-chain-agent-mechanisms/src/test/java/dev/supplychain/agent/mechanisms}, grouped by B1, B2, B3, B4, B6, and feedback responsibilities.
  \item \textbf{Runtime graph:} \path{supply-chain-agent-runtime/src/test/java/dev/supplychain/agent/runtime/cognitive/CognitiveCycleGraphTest.java} and related lifecycle, decision, checkpoint, and framework-compatibility tests.
  \item \textbf{PostgreSQL persistence and recovery:} tests under the PostgreSQL infrastructure module, including memory, checkpoint, action-ledger, budget, intention, and restart/reconciliation integration tests.
  \item \textbf{Application composition and workflow:} \path{supply-chain-agent-app/src/test/java} and \path{supply-chain-agent-evaluation/src/test/java}, especially the durable-host and full-graph workflow tests.
\end{itemize}

\section{PostgreSQL-gated tests}

Several integration-test classes or methods are annotated with JUnit's \texttt{@EnabledIfSystemProperty} condition for \texttt{phase0.postgres=true}. They require a reachable PostgreSQL database and the corresponding environment configuration. The GitHub Actions workflow supplies PostgreSQL 16.4 and runs the Maven build with that property enabled. A local build without the property can compile and run ordinary unit tests while leaving those PostgreSQL-backed tests disabled.

The end-to-end REQ-09 test is especially useful, but it must be read at its actual scope. Its successful path demonstrates the tested approval, simulated action, uncertain acknowledgement, reconciliation, restart, and cross-cycle memory behavior configured by its fixtures. It is not the only evidence for the complete application, and it does not traverse every optional branch. In particular, separate tests are needed to support claims about B2 child investigation and its recovery artifacts. Chapter~24 identifies the principal restart scenarios and their individual purposes.

\section{CI evidence and its limits}

The repository's \path{.github/workflows/ci.yml} runs architecture and composition guard scripts before executing the Maven reactor with PostgreSQL integration tests enabled. It also runs additional security checks and generates selected evaluation artifacts. The guard scripts validate repository-specific constraints such as module composition, dependency boundaries, evidence requirements, and test inventories. They are useful checks, but they do not replace behavioral tests.

A green CI run means that the configured workflow completed for the checked revision and environment. It does not imply that every possible input, concurrency schedule, external integration, or deployment environment has been tested. Similarly, deterministic evaluation fixtures can support repeatable comparisons within their defined protocol, but the existence of a test or generated artifact is not evidence of empirical superiority unless the study design, data, metrics, and results support that conclusion.

The following distinctions should remain explicit:
\begin{itemize}
  \item A \textbf{unit test} supports a local contract or decision under its inputs.
  \item An \textbf{integration test} supports the exercised interaction among components or infrastructure.
  \item A \textbf{full-graph test} supports the workflow path and branches that its fixture actually traverses.
  \item A \textbf{CI run} supports the configured checks for one revision and environment.
  \item An \textbf{empirical research claim} requires an appropriate experimental design and data; it cannot be inferred from ordinary software tests alone.
\end{itemize}

\section{Following source references}

The source links and QR codes in Part III point to specific files in the companion repository. The existing code URLs in Chapters 20--24 are pinned to a single Git revision, \texttt{fa8cd5ec0742d701\allowbreak{}0dbba2c3646852e950\allowbreak{}62a419}. Use those links when you need to inspect the exact source location associated with a printed explanation. A path alone is not a version identifier: a file may change on the repository's default branch after the book's cited revision.

For a useful source-to-test reading sequence, start with the production class cited beside a chapter's claim, then open the named test or its enclosing test class, and finally inspect the collaborating interface or repository adapter. This order helps distinguish what the component promises from what the test fixture supplies. For the broader module map, use Appendix~\ref{app:architecture}; for state ownership and recovery boundaries, use Appendix~\ref{app:persistence}.

\begin{notebox}
The companion project is a reference implementation against a deterministic simulated supply-chain world. Passing tests provide evidence for the behaviors actually exercised. They do not establish live external integration, production deployment readiness, or the complete reproduction of a research mechanism unless those specific claims have their own evidence.
\end{notebox}



\backmatter
\nocite{sumers2024coala,fanlan2026mechanisms}
\printglossaries
\printindex
\printbibliography[heading=bibintoc,title={References}]

\end{document}
